% =====================================================================
% Cybersecurity Laws and Regulations -- Manuale del corso
% Politecnico di Torino -- A.A. 2026/2027
% =====================================================================
\documentclass[11pt,a4paper]{report}
% =====================================================================
% Preambolo -- Manuale CSLR
% =====================================================================
\usepackage[T1]{fontenc}
\usepackage[utf8]{inputenc}
\usepackage[italian]{babel}
\usepackage{lmodern}
\usepackage[a4paper,margin=2.5cm]{geometry}
\usepackage{microtype}
\usepackage{xcolor}
\usepackage{enumitem}
\usepackage{booktabs}
\usepackage{array}
\usepackage{tabularx}
\usepackage{longtable}
\usepackage{amsmath}
\usepackage[most]{tcolorbox}
\usepackage{tikz}
\usetikzlibrary{positioning,arrows.meta}
\usepackage{hyperref}
\hypersetup{
  colorlinks=true,
  linkcolor=blue!50!black,
  urlcolor=blue!60!black,
  pdftitle={Cybersecurity Laws and Regulations -- Manuale del corso}
}

\setlist{itemsep=2pt,topsep=4pt}
\renewcommand{\arraystretch}{1.25}

% ---------------------------------------------------------------------
% Riquadri
% ---------------------------------------------------------------------
% Punto che il docente ha indicato come rilevante per l'esame
\newtcolorbox{esame}{breakable,colback=red!4,colframe=red!60!black,
  fonttitle=\bfseries,title=Per l'esame}

% Nota generica (titolo opzionale)
\newtcolorbox{nota}[1][Nota]{breakable,colback=gray!6,colframe=gray!60!black,
  fonttitle=\bfseries,title=#1}

% Informazione non confermata in aula o calcolata da noi: da verificare
\newtcolorbox{verifica}{breakable,colback=orange!6,colframe=orange!80!black,
  fonttitle=\bfseries,title=Da verificare}

% Soluzione di quiz / simulazioni (titolo opzionale)
\newtcolorbox{soluzione}[1][Soluzione]{breakable,colback=green!4,colframe=green!45!black,
  fonttitle=\bfseries,title=#1}

% Segnaposto per parti ancora da scrivere
\newcommand{\dacompletare}[1]{%
  \begin{tcolorbox}[colback=yellow!8,colframe=yellow!60!black,
    fonttitle=\bfseries,title=Da completare]#1\end{tcolorbox}}

% ---------------------------------------------------------------------
% Intestazione di una lezione
%   \lezione{data}{titolo}{docenti}{argomenti}
% ---------------------------------------------------------------------
\newcommand{\lezione}[4]{%
  \chapter[#1 -- #2]{#2}%
  \begin{tcolorbox}[colback=blue!3,colframe=blue!50!black,
    fonttitle=\bfseries,title=Scheda lezione]
  \textbf{Data:} #1\\
  \textbf{Docenti:} #3\\
  \textbf{Argomenti:} #4
  \end{tcolorbox}}

% ---------------------------------------------------------------------
% Quiz: domanda numerata + opzioni (a), (b), ...
% ---------------------------------------------------------------------
\newcounter{domanda}
\newenvironment{domanda}[1][]{%
  \refstepcounter{domanda}\par\medskip\noindent
  \textbf{Domanda \thedomanda}\ifx&#1&\else\ \textit{(#1)}\fi\textbf{.}~}%
  {\par}
\newenvironment{opzioni}{\begin{enumerate}[label=(\alph*),nosep,leftmargin=2.2em]}{\end{enumerate}}


\title{Cybersecurity Laws and Regulations\\[4pt]
  \large Manuale del corso -- Politecnico di Torino, A.A.~2026/2027}
\author{Marco Pedron e colleghi del corso}
\date{Ultimo aggiornamento: \today}

\begin{document}
\maketitle
\tableofcontents
% =====================================================================
% Premessa: come leggere il manuale (legenda dei riquadri e fonti)
% =====================================================================
\chapter*{Come leggere questo manuale}
\addcontentsline{toc}{chapter}{Come leggere questo manuale}

Questo manuale è scritto da studenti, per studenti: \textbf{non è materiale ufficiale del corso} e non
sostituisce le slide, i testi normativi e le indicazioni del docente. Nasce dalle registrazioni delle
lezioni (trascritte automaticamente e poi revisionate), integrate con slide e materiale dei laboratori.

\section*{Le parti}
\begin{description}[leftmargin=2em,style=nextline]
  \item[Diario delle lezioni] una voce per lezione, in ordine cronologico: sintesi, avvisi pratici,
    contenuti in breve e punti segnalati per l'esame. Serve a seguire l'avanzamento del corso.
  \item[Manuale del corso] un capitolo per norma (GDPR, data breach, ISO/IEC~27001, CRA, NIS2, DORA,
    AI Act, CSA), che unisce lezioni e slide in un'esposizione ordinata. Rimanda alle lezioni del diario.
  \item[Laboratori] i casi pratici del venerdì: contesto, consegna, metodo e soluzione ragionata.
  \item[Schemi di ripasso] tabella delle scadenze, mappe concettuali, schede riassuntive.
  \item[Quiz ed esame] formato dell'esame, tema precedente, simulazioni e quiz per argomento con
    soluzioni e rimando alla lezione.
  \item[Appendice] glossario e acronimi.
\end{description}

\section*{Legenda dei riquadri}

\begin{esame}
Punto che il docente ha indicato come rilevante per l'esame.
\end{esame}

\begin{integrazione}
Contenuto aggiunto dai redattori: \textbf{non detto a lezione e non presente nelle slide}
(spiegazioni, esempi, riferimenti normativi in più). Utile, ma da trattare con cautela all'esame.
\end{integrazione}

\begin{verifica}
Informazione non confermata: passaggio poco chiaro nella registrazione, dato ricostruito o calcolato.
\end{verifica}

\begin{nota}
Nota generica (fonti usate, avvertenze, collegamenti).
\end{nota}

\begin{soluzione}
Soluzione di quiz, simulazioni e laboratori.
\end{soluzione}

\dacompletare{Parte ancora da scrivere (per esempio una norma non ancora trattata a lezione).}

\section*{Fonti}
Registrazioni delle lezioni (Prof.~Giuseppe Vaciago; laboratori con Gianluca Gilardi), slide e materiale
pubblicati sulla piattaforma del corso, testi normativi ufficiali (EUR-Lex). Le slide e le
registrazioni non sono incluse nel repository.


% =====================================================================
\part{Diario delle lezioni}
% Una lezione per file, in ordine cronologico (nome file: AAAA-MM-GG.tex).
% Per aggiungere una lezione: copia lezioni/_template.tex e aggiungi qui l'\input.
% =====================================================================
% =====================================================================
\lezione{21/09/2026}{Lezione introduttiva: il filo rosso del corso}
  {Prof.~Giuseppe Vaciago}
  {organizzazione del corso, laboratori ed esame; un incidente, cinque regimi (GDPR, NIS2, CRA, DORA,
   AI Act) e il diritto penale; le quattro domande del corso; ISO/IEC~27001 e approccio basato sul
   rischio; dati personali e non personali, anonimizzazione, indirizzo IP; regolamenti e direttive;
   livelli di rischio dell'AI Act; Convenzione di Budapest; cinque ``orologi'' di notifica; compliance
   come problema di prova; supply chain e istituzioni}
\label{lez:2026-09-21}
% =====================================================================

\begin{nota}[Fonte]
Trascrizione automatica della registrazione, revisionata. All'inizio il proiettore non funzionava: le
slide sono state mostrate solo in streaming e parte della lezione è stata detta a voce. Una parte
centrale della registrazione (pausa e conversazioni in italiano) è inintelligibile ed è stata
ignorata. Numeri, articoli e date sono ripresi dalle slide \emph{01.~Cybersecurity Laws and
Regulations -- I -- Introductory class: the red thread}.
\end{nota}

\section{Sintesi}
Lezione di presentazione. Il docente spiega organizzazione, laboratori ed esame, poi introduce il
``filo rosso'' del corso con uno scenario: un ransomware colpisce alle 8:14 di un lunedì la piattaforma
di prenotazione di una rete ospedaliera. Lo stesso fatto attiva \textbf{cinque regimi} (GDPR, NIS2,
CRA, DORA, AI Act) più il diritto penale, ciascuno con il proprio bene protetto, la propria autorità e
il proprio \textbf{orologio}. Il messaggio centrale: tutte queste norme si basano su \textbf{risk
assessment} e \textbf{accountability}, e la competenza più utile è capire \emph{quale domanda} pone
ciascuna norma. Seguono la mappa del corso, la distinzione tra dati personali e non personali, la
differenza tra regolamento e direttiva e una panoramica di AI Act, Convenzione di Budapest, scadenze di
notifica, certificazioni, supply chain e autorità.

% ---------------------------------------------------------------------
\section{Organizzazione}
\begin{itemize}
  \item \textbf{Docenti}: Prof.~Vaciago (avvocato, dottorato in digital forensics, viene da Milano il
    lunedì e il venerdì), con Gianluca (esperto di ISO/IEC~27001 e risk assessment) e un terzo docente,
    Michele, che gli studenti incontreranno nei mesi successivi.
  \item \textbf{Teoria e laboratorio} in proporzione circa 50/50. Nei laboratori si analizza un caso
    in gruppo (indicativamente 4--5 persone, comunque non più di 10), con circa un'ora e mezza di lavoro
    dopo la presentazione del caso. Il caso è per circa il 70\% legale e il 30\% tecnico. Le linee guida
    (settore, argomenti) vengono date a fine lezione del lunedì; il venerdì si svolge il lab.
  \item I gruppi \textbf{presentano oralmente} il lavoro (tempi rigidi, es.\ 5 minuti): per il docente
    è un esercizio importante di parlare in pubblico.
  \item Con oltre 240 iscritti (circa 144 in aula e 22 online quel giorno) l'aula non basta: la
    modalità \textbf{ibrida non funziona} per i lab, quindi si pensa a lab interamente online (Zoom con
    stanze separate) o interamente in presenza. Decisione rinviata.
  \item L'uso dell'AI è ammesso; in alcuni lab sarà anzi suggerito.
  \item \textbf{Calendario}: la lezione di venerdì 25/09 sarà \textbf{da remoto} (il docente è a un
    convegno); il primo laboratorio è previsto \textbf{venerdì 2 ottobre}.
  \item Domande sul corso: meglio nella pausa che alle 19:00 (il docente deve prendere il treno).
\end{itemize}

\subsection{L'esame}
\begin{itemize}
  \item Durata \textbf{45 minuti}.
  \item \textbf{4 domande a risposta multipla} con una sola risposta corretta, \textbf{5 punti}
    ciascuna (20 punti). Secondo il docente circa l'80\% degli studenti risponde correttamente.
  \item \textbf{1 domanda aperta} per i punti restanti fino al 30 e lode: è ciò che fa la vera
    differenza. Va scritta in modo \textbf{sintetico e preciso}, mostrando un approccio di compliance
    e la comprensione di \textbf{risk assessment} e \textbf{compliance}.
  \item I laboratori danno \textbf{punti aggiuntivi} (anche frazioni di punto); l'entità sarà decisa
    quando sarà chiara l'organizzazione (vedi lezione del 02/10).
\end{itemize}

\begin{esame}
\begin{itemize}
  \item Obiettivo principale del corso: capire il \textbf{risk assessment} e la \textbf{compliance},
    su cui si basano tutte le norme studiate. Non serve imparare gli articoli a memoria.
  \item Saper fare bene un \textbf{severity assessment} di un data breach vale, a detta del docente,
    ``circa il 20\% dell'esame''.
  \item Tipica domanda aperta: capire \textbf{quale soggetto} deve notificare e \textbf{a quale norma}
    è soggetta l'azienda (es.\ è soggetta a DORA o no?).
\end{itemize}
\end{esame}

% ---------------------------------------------------------------------
\section{Perché studiare le norme europee}
Il corso tratta cinque regolamenti/direttive principali: \textbf{GDPR} (il più importante),
\textbf{NIS2} (sicurezza di soggetti che erogano servizi essenziali, recepita in Italia con legge
nazionale), \textbf{Cyber Resilience Act} (prodotti con elementi digitali, IoT; obblighi di
segnalazione in vigore dall'11 settembre 2026), \textbf{DORA} (settore bancario e assicurativo) e
\textbf{AI Act}; si aggiunge il \textbf{Cybersecurity Act}.

Anche per gli studenti non europei: la regolazione UE fa scuola nel mondo (il GDPR è stato ripreso,
tra gli altri, da Brasile, Svizzera e Singapore; alcuni Stati USA, come la California, si stanno
avvicinando all'AI Act). Soprattutto, l'approccio europeo basato su \textbf{risk assessment e
accountability} è ormai applicato ovunque.

\paragraph{Esperienze dal lavoro del docente.} Due data breach gestiti nel weekend:
\begin{itemize}
  \item un'agenzia di prenotazioni con circa 1.400 interessati in 27 Paesi: 25 minuti per spiegare al
    CEO che quella che lui riteneva la \emph{root cause} non lo era. All'autorità va data la causa
    corretta; se non la si conosce entro le 72 ore, si dichiara che servono ulteriori analisi;
  \item circa 25.000 indirizzi email: severità bassa (la gravità dipende dal \emph{tipo} di dati), ma
    il docente consigliava comunque di avvisare gli interessati; il CEO si opponeva per motivi
    reputazionali, rischiando che l'autorità imponesse poi la comunicazione.
\end{itemize}
Il tecnico che conosce un po' di compliance può gestire da sé queste decisioni, invece di doverle
spiegare a chi non ha competenze IT.

% ---------------------------------------------------------------------
\section{Lo scenario: lunedì, ore 8:14}
Un ransomware colpisce la società che gestisce la piattaforma di prenotazione online di una rete
ospedaliera regionale. Alle 8:20 è un problema tecnico, alle 9:00 un problema legale. Tre elementi da
considerare in ogni incidente:
\begin{itemize}
  \item \textbf{Dati}: 200.000 cartelle cifrate e una copia esfiltrata; sono \textbf{dati sanitari},
    categorie particolari (art.~9 GDPR). Per un ingegnere la differenza tra dato personale e non
    personale è piccola; per il diritto è enorme.
  \item \textbf{Servizio}: prenotazioni ferme per quattro giorni, si torna alla carta. Dieci anni fa il
    problema era solo il dato; oggi conta anche la continuità (``the show must go on''): un paziente
    che non riesce a prenotare può subire danni, con contenzioso conseguente.
  \item \textbf{Supply chain}: l'ingresso è avvenuto tramite un componente vulnerabile del software
    di un fornitore. ``Non è un problema mio'' è falso: in ogni ordinamento chi sceglie il fornitore ne
    risponde, salvo poi rivalersi.
\end{itemize}
In più, lo strumento di triage AI addestrato su quel database inizia a dare risultati inaffidabili e
nessuno se ne accorge per due giorni. Domanda del corso: \emph{quanti obblighi legali sono partiti
alle 8:14?}

\subsection{Un evento, cinque regimi}
\begin{center}
\small
\begin{tabularx}{\linewidth}{>{\raggedright\arraybackslash}p{2.6cm} X X >{\raggedright\arraybackslash}p{2.5cm}}
\toprule
\textbf{Regime} & \textbf{Che cosa protegge} & \textbf{A chi si risponde} & \textbf{Prima scadenza}\\
\midrule
GDPR (Reg.~UE 2016/679) & dati personali dei pazienti & Garante / DPA capofila & 72 h (art.~33)\\
NIS2 (Dir.~UE 2022/2555; D.lgs.~138/2024) & continuità di un servizio essenziale & CSIRT Italia / ACN
  & 24 h early warning (art.~23)\\
CRA (Reg.~UE 2024/2847) & il prodotto con elementi digitali & CSIRT + piattaforma ENISA & 24 h\\
DORA (Reg.~UE 2022/2554) & resilienza di una funzione finanziaria & autorità competente / ESA
  & 4 h dalla classificazione\\
AI Act (Reg.~UE 2024/1689) & il sistema di IA e chi vi è esposto & autorità di vigilanza del mercato
  & 15 giorni (art.~73)\\
\bottomrule
\end{tabularx}
\end{center}
In parallelo c'è il \textbf{diritto penale}: gli stessi fatti sono reati (Convenzione di Budapest,
artt.~615-\emph{ter} ss.\ c.p., accesso abusivo a sistema informatico). Stesso fatto, procedura e
standard probatorio diversi.

\paragraph{GDPR.} L'autorità è il Garante per la protezione dei dati personali, cioè la \emph{data
protection authority}. Entro 72 ore si decide se notificare. La notifica rende pubblico l'incidente
e porta l'autorità a verificare se le misure di sicurezza erano adeguate; se non lo erano, la
sanzione può arrivare a 10 milioni di euro o al 2\% del fatturato. Il CEO tende a chiedere di
``abbassare'' il rischio, ma un incidente ad alto rischio (es.\ dati sanitari) non si può cambiare a
posteriori: l'unico modo sarebbe mentire all'autorità, e chi firma una dichiarazione falsa rischia la
reclusione fino a tre anni. Il DPO o il legale stanno in mezzo tra tecnici e CEO.

\begin{integrazione}[sanzioni e false dichiarazioni]
Per le violazioni degli obblighi di sicurezza e di notifica (artt.~32--34) l'art.~83(4) GDPR prevede
fino a 10 milioni di euro o, per le imprese, fino al 2\% del fatturato mondiale annuo
dell'esercizio precedente, \emph{se superiore}: non è una scelta discrezionale tra i due tetti. Le
false dichiarazioni al Garante sono punite in Italia dall'art.~168 del Codice privacy
(D.lgs.~196/2003), con la reclusione da sei mesi a tre anni.
\end{integrazione}

\paragraph{NIS2.} Si applica perché è in gioco la continuità di un servizio essenziale. Nel contesto
geopolitico attuale proteggere infrastrutture e servizi essenziali (ferrovie, aerei, ospedali,
energia) è cruciale, anche per i rischi che l'AI introduce. Secondo il docente la NIS2 è ``il vero
futuro'': un livello di cybersicurezza molto più spinto dell'art.~32 GDPR, che ha dominato la
generazione 2016--2026. L'ACN può arrivare a fermare un fornitore ritenuto pericoloso per
un'infrastruttura critica, anche se il fornitore non è direttamente soggetto a NIS2. L'early warning
va fatto entro 24 ore, anche di sabato o domenica.

\paragraph{CRA.} Riguarda IoT e prodotti con elementi digitali (smart TV, frigoriferi, smartwatch,
alcuni componenti software); segnalazioni attive dall'11 settembre 2026.

\paragraph{ENISA e la centralizzazione.} ENISA, finora con un ruolo di coordinamento, diventerà
centrale: con il pacchetto \emph{Digital Omnibus} si andrà verso un \textbf{punto unico di
segnalazione} degli incidenti presso ENISA, che poi informerà le autorità competenti. Per il docente
sarà ``la rivoluzione più importante'', perché oggi le autorità nazionali non si parlano abbastanza
mentre lo stesso malware colpisce in tutta Europa. ENISA dovrà assumere molti esperti: un'opportunità
professionale.

\paragraph{DORA e AI Act.} DORA riguarda solo banche, finanza e assicurazioni (non un ospedale).
Per l'AI Act il docente cita in Italia l'\textbf{AgID}, agenzia governativa (non indipendente come il
Garante); il quadro delle autorità sarà ripreso nelle lezioni dedicate.

\begin{esame}
\begin{itemize}
  \item \textbf{Attenzione all'acronimo DPA}: \emph{Data Protection Authority} (l'autorità) ma anche
    \emph{Data Processing Agreement} (il contratto con il fornitore/responsabile). Per il docente
    ``DPA'' è il contratto; per l'autorità dice \emph{data protection authority}. All'esame la
    confusione capita: è un punto insidioso.
  \item Per l'esame vanno studiate \textbf{tutte le autorità} e lo \textbf{schema delle notifiche}
    vigente: il punto unico ENISA del Digital Omnibus non è ancora operativo.
  \item ``Sapere quale domanda pone una norma è la competenza più utile che il corso può dare''.
\end{itemize}
\end{esame}

% ---------------------------------------------------------------------
\section{Le quattro domande del corso}
Ogni lezione risponde a una di queste domande:
\begin{enumerate}
  \item \textbf{Che cosa proteggiamo?} Dati personali e non personali, sistemi informativi e reti,
    prodotti con elementi digitali, modelli di IA.
  \item \textbf{Come dobbiamo proteggerlo?} Approccio basato sul rischio, misure di sicurezza, sistemi
    di gestione, certificazione e valutazione di conformità.
  \item \textbf{Che cosa succede se la protezione fallisce?} Obblighi di notifica e scadenze,
    responsabilità, sanzioni amministrative, reati.
  \item \textbf{Chi decide?} Autorità di controllo, CSIRT, ENISA, agenzie nazionali, organismi
    notificati, tribunali.
\end{enumerate}

\paragraph{IA ad alto rischio $\neq$ IA generativa.} I sistemi ad alto rischio dell'AI Act non sono
gli LLM: sono in gran parte sistemi legati a infrastrutture critiche e macchinari (Allegato~III). L'IA
generativa diventa rilevante per l'uso: ad esempio in ambito HR (selezione dei curriculum) e nella
cybersicurezza, anche per produrre attacchi.

\subsection{ISO/IEC 27001 e approccio basato sul rischio}
La \textbf{ISO/IEC~27001} (seguita da Gianluca) è lo standard ISO sul comportamento che un'azienda
deve tenere per proteggere dati e infrastruttura IT: autenticazione, autorizzazione, resilienza
(backup, disaster recovery, business continuity). Con la NIS2 molte aziende ne hanno bisogno per
stare sul mercato, perché i grandi clienti scelgono i fornitori anche in base alla certificazione. Il
legame con il risk assessment: l'art.~32 GDPR e le altre norme dicono che bisogna essere sicuri, ma
\textbf{non quali misure adottare}. È il principio di \textbf{accountability}: in base a struttura e
dimensione dell'azienda devi fare il meglio possibile, e dimostrarlo. La 27001 colma questo vuoto ed è
il linguaggio comune nel business in Europa; negli Stati Uniti si usano soprattutto gli standard NIST.
Il docente non è un ``fan'' della ISO, ma la considera indispensabile per comunicare con le aziende.

% ---------------------------------------------------------------------
\section{Laboratori: modalità proposta (dopo la pausa)}
Dopo un sondaggio in aula (la maggioranza ha il venerdì poco carico), la proposta è di svolgere i
laboratori \textbf{online}: il docente presenta il caso, i gruppi lavorano in stanze Zoom separate e
possono chiamarlo per domande. Si proverà questa modalità e la si rivedrà se non funziona. Le
presentazioni orali saranno svolte da un sottoinsieme di gruppi (ad esempio 10 per laboratorio), anche
dividendo l'esposizione tra più membri.

% ---------------------------------------------------------------------
\section{La mappa del corso}
\begin{itemize}
  \item \textbf{Parte I -- Dati e privacy}: quadro internazionale (Consiglio d'Europa, OCSE), GDPR,
    applicazioni settoriali e Reg.~(UE) 2018/1807 sui dati non personali, AI Act e approccio basato sul
    rischio.
  \item \textbf{Parte II -- Cybersicurezza e cybercrime}: Convenzione di Budapest (23 novembre 2001),
    data breach e GDPR, NIS2 e Cybersecurity Act, AI Act e cybersicurezza, ISO/IEC~27001 e risk
    assessment, organismi UE e casi di studio.
\end{itemize}
L'oggetto della protezione si sposta nel tempo: dal \emph{record} (Convenzione 108, 1981) al
\emph{processo} (GDPR, 2016), al \emph{servizio} (NIS2 e DORA, 2022), al \emph{prodotto} e al
\emph{modello} (CRA e AI Act, 2024). Cambia il bene, il metodo resta lo stesso.

\paragraph{Cybersecurity Act.} Lezione breve (un'ora, un'ora e mezza): rafforza ENISA e fissa i
principi della standardizzazione e certificazione delle misure di sicurezza; è strettamente legato
alla 27001.

\subsection{Dati non personali e anonimizzazione}
I dati non personali (Reg.~2018/1807) hanno molti meno vincoli. Il docente cita le società che lavorano
sui \textbf{dati sintetici} (spin-off universitari a Torino e Trieste): un settore con molto futuro.
Anonimizzare \textbf{non} significa togliere nome e cognome: giuridicamente nessun elemento del
documento deve poter essere ricondotto, \textbf{direttamente o indirettamente}, alla persona. Esempi:
un riferimento a fatti noti che identifica immediatamente un politico famoso anche senza nominarlo;
``l'unica persona in città con questa malattia''. Nei casi reali le aziende ritengono di aver
anonimizzato, ma in audit emergono molte possibilità di re-identificazione indiretta. Con gli LLM la
re-identificazione è diventata molto più facile: esperimenti svolti dal docente hanno dato buoni
risultati anche con modelli non particolarmente potenti.

\subsection{Privacy come diritto fondamentale, tecnologia no}
La privacy è un diritto fondamentale (artt.~7--8 della Carta), riconosciuto da tutti: per questo è
regolabile in modo stabile (il docente fa risalire la nascita del concetto a un articolo di due
autori statunitensi di fine Ottocento). La tecnologia non è un diritto fondamentale e cambia continuamente.
L'AI Act ne è l'esempio: nato nel 2020, la prima bozza non prevedeva nulla sull'IA generativa, che dal
novembre 2022 ha cambiato il mondo; la parte sui modelli generativi è stata aggiunta in corsa. Quando
la legge insegue la tecnologia ``fallisce sempre'', perché la tecnologia corre più veloce.

\begin{integrazione}[l'origine del concetto]
L'articolo a cui si riferisce il docente è con ogni probabilità \emph{The Right to Privacy} di Samuel
Warren e Louis Brandeis (Harvard Law Review, 1890), due giuristi, scritto in reazione alla diffusione
della fotografia istantanea e della stampa scandalistica; vi si teorizza il \emph{right to be let
alone}.
\end{integrazione}

% ---------------------------------------------------------------------
\section{Tre livelli: internazionale, UE, nazionale}
I dati ignorano i confini, l'\emph{enforcement} no.
\begin{itemize}
  \item \textbf{Internazionale}: Convenzione 108 del Consiglio d'Europa (1981), primo trattato
    vincolante, e Convenzione 108+ (aperta alla firma nel 2018, non ancora in vigore); Linee guida OCSE
    (1980, riviste nel 2013), soft law; Convenzione di Budapest (2001) e Secondo Protocollo (2022) sulle
    prove elettroniche. Cina e Russia non l'hanno firmata.
  \item \textbf{Unione europea}: i \textbf{regolamenti} (GDPR, CSA, DORA, AI Act, CRA) si applicano
    direttamente e in modo identico in tutti gli Stati; le \textbf{direttive} (NIS2) vanno recepite e
    lasciano 27 perimetri nazionali diversi per ambito, sanzioni e vigilanza.
  \item \textbf{Nazionale}: D.lgs.~196/2003 (Codice privacy), D.lgs.~138/2024 (NIS2), L.~90/2024;
    Garante, ACN, CSIRT Italia; reati informatici del codice penale.
\end{itemize}
Oggi circa il 90\% degli atti UE in materia sono regolamenti: un vantaggio rispetto agli USA, dove
federazione e singoli Stati legiferano separatamente (California, New York, Washington hanno leggi
proprie sull'IA che un ordine esecutivo presidenziale non può cancellare).

\begin{nota}[Regolamento e direttiva]
Il docente lo indica come principio generale: probabilmente non sarà una domanda d'esame, ma va
capito. La \textbf{direttiva} obbliga gli Stati a raggiungere un risultato ma lascia libertà su come
recepirla (l'Italia può applicarla diversamente da Francia o Spagna); il \textbf{regolamento} è
immediatamente applicabile e uguale ovunque (il GDPR è lo stesso in Germania e in Italia).
\end{nota}

% ---------------------------------------------------------------------
\section{Dati personali e non personali}
\begin{center}
\small
\begin{tabularx}{\linewidth}{X X}
\toprule
\textbf{Dati personali (GDPR)} & \textbf{Dati non personali (Reg.~2018/1807)}\\ \midrule
Qualsiasi informazione relativa a una persona fisica identificata o \textbf{identificabile} &
  Tutto il resto: dati macchina, telemetria, dati aggregati e anonimizzati\\
Bene protetto: diritto fondamentale & Bene protetto: mercato interno\\
Obblighi legati a finalità, liceità e rischio per la persona & Divieto di obblighi di localizzazione
  dei dati; portabilità per gli utenti professionali\\
La sicurezza è un principio (art.~5(1)(f)) & Gli obblighi di sicurezza arrivano da NIS2, DORA, CRA,
  contratti\\ \bottomrule
\end{tabularx}
\end{center}
\textbf{Dataset misti}: se dati personali e non personali sono inestricabilmente collegati, il GDPR si
applica all'intero insieme (art.~2(2) Reg.~2018/1807). È il caso normale.

\paragraph{L'indirizzo IP è un dato personale?} Sì, perché è \emph{identificabile}, anche se dinamico.
Il percorso per risalire alla persona: si chiede al provider del servizio l'IP usato in un dato
momento, poi al provider di rete/telco chi pagava la connessione. Il risultato è solo un indizio (chi
paga la bolletta non è necessariamente chi ha agito) e in Europa la richiesta passa per l'autorità
giudiziaria (serve una denuncia; il privato non può chiedere direttamente a Google). La conservazione
dei dati di traffico varia molto: in Italia è lunga, mentre in altri Stati (es.\ Germania) la
conservazione generalizzata è stata ritenuta incostituzionale e si può solo ``congelare'' il traffico
da quando nasce il sospetto. Nella pratica investigativa l'IP serve davvero quasi solo nei casi di
diffamazione sui social. Anche un'email come ``nomignolo74@gmail.com'' è dato personale, perché
Google può ricondurla a una persona.

% ---------------------------------------------------------------------
\section{Il GDPR come prima legge sulla cybersicurezza}
\begin{itemize}
  \item \textbf{Art.~5(1)(f)}: integrità e riservatezza come principio del trattamento.
  \item \textbf{Artt.~24--25}: accountability (dimostrare la conformità), protezione dei dati fin dalla
    progettazione e per impostazione predefinita.
  \item \textbf{Art.~32}: misure adeguate al rischio, tenendo conto dello stato dell'arte e dei costi.
    È ``il cardine del corso'': sicurezza relativa al rischio, misurata su uno standard che si muove,
    e qualcuno deve poterla dimostrare.
  \item \textbf{Artt.~33--34 e 35}: notifica in 72 ore, comunicazione agli interessati se il rischio è
    elevato, DPIA per i trattamenti ad alto rischio.
\end{itemize}
Nessun articolo dice quale controllo adottare: è uno spazio voluto, riempito da ISO/IEC~27001 e dalle
metodologie ENISA. Le cinque norme condividono lo stesso metodo (contesto, identificazione, analisi,
valutazione, trattamento, monitoraggio): GDPR artt.~32 e 35, NIS2 art.~21, DORA artt.~5--16, AI Act
artt.~9 e 15, ISO/IEC~27001 cl.~6.1 e Allegato~A. L'approccio basato sul rischio non è una fuga dalle
regole: gli obblighi crescono con ciò che si può rompere, e ``abbiamo seguito la checklist'' non è mai
stata una difesa.

% ---------------------------------------------------------------------
\section{AI Act: quattro livelli di rischio}
\begin{enumerate}
  \item \textbf{Pratiche vietate} (art.~5), es.\ \emph{social scoring}: in Europa vietato in ogni caso.
  \item \textbf{Sistemi ad alto rischio} (art.~6, Allegati I e III): molti legati a infrastrutture
    critiche, ma anche strumenti HR.
  \item \textbf{Obblighi di trasparenza} (art.~50), in vigore dal 2 agosto 2026: ad esempio
    l'indicazione che un contenuto è generato con IA. L'obbligo riguarda gli usi commerciali, ma il
    confine con l'uso personale è sfumato (chi diventa influencer monetizza): il docente consiglia di
    indicarlo sempre, anche per una questione etica e culturale.
  \item \textbf{Rischio minimo}: nessun obbligo specifico (es.\ un chatbot elementare).
\end{enumerate}
Dalle slide: in vigore dal 1° agosto 2024; divieti e alfabetizzazione IA dal 2 febbraio 2025; GPAI dal
2 agosto 2025; trasparenza dal 2 agosto 2026. L'\textbf{AI Omnibus} (in vigore dal 27 luglio 2026)
ha rinviato l'alto rischio: Allegato~III al 2 dicembre 2027, Allegato~I al 2 agosto 2028. Per il
docente il rinvio nasce anche dalle pressioni politiche statunitensi. Per la prima volta un
regolamento UE nomina attacchi di machine learning (data e model poisoning, esempi avversari,
evasione del modello).

% ---------------------------------------------------------------------
\section{Parte II: Budapest, orologi, prova, supply chain, istituzioni}

\paragraph{Convenzione di Budapest.} Primo trattato internazionale sul cybercrime: reati armonizzati
(accesso illegale, intercettazione, danneggiamento di dati e sistemi, frode informatica), poteri
procedurali (conservazione rapida, ordini di produzione, perquisizione e sequestro di dati) e
cooperazione internazionale. Non avrà una lezione dedicata (è materia del corso di Computer Forensics
del docente), ma conta per la \textbf{giurisdizione}: anche ciò che sta nel cloud ha una
giurisdizione, e il sistema di trasferimenti UE--USA è delicato.

\subsection{Un incidente, cinque orologi}
\begin{center}
\small
\begin{tabularx}{\linewidth}{>{\raggedright\arraybackslash}p{2.1cm} X >{\raggedright\arraybackslash}p{2.1cm} >{\raggedright\arraybackslash}p{1.9cm} X}
\toprule
\textbf{Regime} & \textbf{L'orologio parte quando\ldots} & \textbf{Early} & \textbf{Completa} & \textbf{Poi}\\
\midrule
GDPR artt.~33--34 & consapevolezza di una violazione di dati personali & -- & 72 h al Garante
  & interessati senza ingiustificato ritardo, se rischio elevato\\
NIS2 art.~23 & consapevolezza di un incidente significativo & 24 h & 72 h & relazione finale entro 1 mese\\
DORA art.~19 & classificazione dell'incidente come grave & 4 h (max 24 h dalla consapevolezza)
  & 72 h intermedia & relazione finale entro 1 mese\\
CRA art.~14 & vulnerabilità attivamente sfruttata / incidente grave & 24 h & 72 h
  & vulnerabilità: 14 giorni dalla disponibilità di una correzione; incidente grave: 1 mese\\
AI Act art.~73 & incidente grave, stabilito il nesso causale & -- & 15 giorni
  & 2 giorni se infrastruttura critica o violazione diffusa; 10 in caso di morte\\
\bottomrule
\end{tabularx}
\end{center}
Ogni regime conta da un evento diverso: danno alle persone (GDPR), interruzione del servizio (NIS2),
impatto su una funzione (DORA), sfruttabilità del prodotto (CRA), gravità dell'esito (AI Act). Il
docente sottolinea le sole \textbf{4 ore} di DORA.

\paragraph{La compliance è un problema di prova.} Due sensi: preservare le prove digitali e,
soprattutto, poter \textbf{dimostrare} di aver fatto il possibile per prevenire l'incidente
(accountability). Gli attacchi non si possono evitare del tutto (citazione di Eugene Spafford): la
domanda non è \emph{se} ma \emph{quando}, quindi bisogna essere pronti a provare la propria diligenza.
Sistemi di gestione (ISO/IEC~27001), certificazione dei prodotti (Cybersecurity Act, schema EUCC) e
valutazione di conformità (CRA, AI Act) servono a questo; un certificato è un elemento di prova, mai
un'immunità.

\paragraph{Supply chain.} Si risponde anche per chi non si impiega: GDPR art.~28, NIS2
art.~21(2)(d), DORA artt.~28--30, CRA artt.~13--14, AI Act artt.~22--25 regolano il rapporto con i
fornitori lungo la catena.

\paragraph{Istituzioni.} Garante, EDPB ed EDPS per i dati; ENISA, NIS Cooperation Group ed
\textbf{EU-CyCLONe} (task force per le crisi su larga scala) per la cybersicurezza; CSIRT Italia
(dentro l'ACN) per la risposta operativa; ACN come autorità NIS2; EBA, ESMA, EIOPA per la finanza; AI
Office e AI Board per l'IA. Bisogna sapere chi chiamare alle 8:14. Anche qui il docente segnala
opportunità di lavoro presso queste autorità, che cercano nuovi esperti.

% ---------------------------------------------------------------------
\section{Domande aperte e punti da verificare}
\begin{verifica}
\begin{itemize}
  \item \textbf{Punteggio dell'esame}: 4 domande da 5 punti; la domanda aperta vale ``gli altri 11''
    secondo la trascrizione (fino al 30 e lode?). Da confermare con le regole ufficiali.
  \item \textbf{Digital Omnibus}: il docente dice sia che è stato ``applicato a luglio'' sia che il
    punto unico di segnalazione presso ENISA ``non è ancora in vigore'' e non si sa quando lo sarà. Le
    slide citano solo l'\emph{AI Omnibus} in vigore dal 27 luglio 2026. Verificare lo stato delle due
    proposte.
  \item \textbf{Convenzione di Budapest}: il docente corregge a voce un numero di Stati aderenti
    (``non più di 100, credo 68''); le slide non riportano il dato. Verificare.
  \item \textbf{California}: secondo la trascrizione una legge californiana ``applica integralmente
    l'art.~55 dell'AI Act'' (obblighi per i modelli GPAI a rischio sistemico?). Riferimento da
    verificare.
  \item \textbf{AgID e AI Act}: ruolo delle autorità italiane per l'AI Act (AgID, ACN) da verificare.
  \item \textbf{Conservazione dei dati di traffico in Italia}: la durata citata (``cinque anni'') non
    è chiara nella trascrizione.
  \item \textbf{Origine del diritto alla privacy}: la trascrizione parla di un articolo ``del 1898''
    scritto da due ``fotografi'' statunitensi (vedi integrazione nella sezione sulla privacy).
  \item Nomi dei docenti: Gianluca ``Gilardi'' e Michele (cognome non comprensibile) sono ricavati
    dalla trascrizione.
\end{itemize}
\end{verifica}

% =====================================================================
\lezione{25/09/2026}{GDPR (I): ruoli, bilanciamento, principi, basi giuridiche}
  {Prof.~Giuseppe Vaciago}
  {titolare, responsabile e contitolari; bilanciamento degli interessi e caso Google Spain; principi
   dell'art.~5; definizione di trattamento ed esenzione domestica (caso Lindqvist); i quattro pilastri;
   categorie di dati; basi giuridiche; consenso e cookie paywall}
\label{lez:2026-09-25}
% =====================================================================

\begin{nota}[Fonte]
Trascrizione automatica della registrazione, revisionata. Lezione \textbf{da remoto} (circa 60--90
collegati). Slide di riferimento: \emph{02.~Cybersecurity Laws and Regulations -- II Privacy}, dalla
prima slide fino al consenso (il legittimo interesse è stato rinviato al 28/09). Le slide sono state
caricate ruotate: il docente invita a segnalarlo se serve una nuova versione.
\end{nota}

\section{Sintesi}
Prima lezione sul GDPR, ``la madre'' delle altre norme del corso. Il docente insiste soprattutto sui
\textbf{ruoli} (titolare, responsabile, contitolari), che definisce ``la Bibbia del GDPR'', e sul
\textbf{bilanciamento degli interessi}, illustrato con il caso \emph{Google Spain} e con esempi dalla
sua pratica sul diritto all'oblio. Seguono i principi dell'art.~5, la definizione (larghissima) di
trattamento, l'esenzione per le attività domestiche (caso \emph{Lindqvist}), i quattro pilastri del
GDPR, le tre categorie di dati, le basi giuridiche e le condizioni del consenso, con il caso dei
\emph{cookie paywall}.

% ---------------------------------------------------------------------
\section{Organizzazione}
\begin{itemize}
  \item Lunedì 28/09 si conclude la parte generale sul GDPR, \textbf{in presenza}.
  \item Il venerdì successivo (2/10) dovrebbe partire il primo laboratorio; per ora l'idea è farlo
    online, la decisione definitiva sarà presa lunedì in base ai presenti.
  \item Per ragioni logistiche e di aule, alcune lezioni in più si svolgeranno da remoto.
  \item Tutte le lezioni sono \textbf{trasmesse in streaming e registrate}; le registrazioni sono
    sulla piattaforma del corso.
\end{itemize}

\begin{esame}
Questa lezione e quella del 28/09 sono ``molto importanti per l'esame'': contengono molte
\textbf{definizioni} che possono essere usate nelle domande a risposta multipla. Non vanno imparate a
memoria, ma bisogna capirne esattamente il \textbf{significato giuridico}.
\end{esame}

% ---------------------------------------------------------------------
\section{I ruoli: titolare, responsabile, contitolari}
\begin{itemize}
  \item \textbf{Titolare} (\emph{data controller}): persona fisica o giuridica che, da sola o insieme
    ad altri, \textbf{determina le finalità e i mezzi} del trattamento. Le due parole chiave da
    ricordare ``per sempre'': \emph{purposes and means}.
  \item \textbf{Responsabile} (\emph{data processor}): tratta i dati \textbf{per conto} del titolare.
  \item \textbf{Contitolari} (\emph{joint controllers}): due o più titolari che determinano
    \textbf{congiuntamente} finalità e mezzi.
\end{itemize}

\paragraph{Esempio 1 -- La piattaforma di e-commerce.} Un'azienda di prodotti per l'igiene vuole
vendere online tramite una piattaforma SaaS specializzata. La piattaforma chiede di avere anch'essa i
dati dei clienti e propone un accordo di contitolarità. Ma i clienti comprano i prodotti dell'azienda:
è l'azienda che decide le finalità, quindi è il \textbf{titolare}; la piattaforma tratta moltissimi
dati, ma \textbf{per conto} del titolare, quindi è \textbf{responsabile}.

\paragraph{Esempio 2 -- La maratona di beneficenza.} Una catena di negozi di scarpe e una fondazione
per la ricerca sul cancro organizzano insieme una maratona. Nessuna delle due agisce per conto
dell'altra: decidono insieme che cosa fare con i dati degli iscritti, quindi sono
\textbf{contitolari}.

\begin{esame}
\begin{itemize}
  \item Ciò che conta è la \textbf{situazione di fatto}, non ciò che le parti scrivono nell'accordo:
    non si diventa contitolari solo perché lo si dichiara in un contratto. Il docente lo indica come
    possibile spunto per la domanda aperta.
  \item ``Non fate confusione sui ruoli'': è probabilmente l'unico punto su cui il docente dice che
    sarà davvero severo all'esame. Tutto il GDPR, anche la parte su data breach e art.~32, si basa su
    questi ruoli.
\end{itemize}
\end{esame}

% ---------------------------------------------------------------------
\section{Il bilanciamento degli interessi}
La privacy è un diritto fondamentale, ma non è ``scritto nella pietra'': a differenza del diritto alla
vita, spesso va bilanciato con altri interessi, e c'è molta zona grigia. Anche la cybersicurezza è un
bilanciamento: per prevenire incidenti bisogna loggare, monitorare, controllare, il che va in
direzione opposta al diritto di essere lasciati in pace (\emph{right to be let alone}). Oggi la quasi
totalità dei problemi di protezione dei dati riguarda la vita digitale.

\subsection{Il caso Google Spain (CGUE, C-131/12, 2014)}
\begin{itemize}
  \item Il sig.\ Costeja González si rivolge all'autorità spagnola contro il quotidiano \emph{La
    Vanguardia} e contro Google: cercando il suo nome comparivano due pagine del 1998 con l'annuncio di
    un'asta immobiliare legata a debiti previdenziali, vicenda chiusa da anni.
  \item L'autorità spagnola (decisione del 30 luglio 2010) \textbf{respinge} il reclamo contro il
    giornale, la cui pubblicazione era legittima (doveva dare la massima pubblicità all'asta), ma lo
    \textbf{accoglie} contro Google: il motore di ricerca svolge un trattamento di cui è \textbf{titolare} (nelle slide ``responsible'' è usato in senso generico).
  \item La \textbf{Corte di giustizia}: alla luce degli artt.~7 e 8 della Carta, i diritti
    dell'interessato prevalgono di regola sull'interesse economico del motore di ricerca e
    sull'interesse del pubblico a trovare l'informazione, \textbf{salvo} ragioni particolari come il
    ruolo dell'interessato nella vita pubblica.
\end{itemize}
Il punto per il docente non è il caso in sé, ma il \textbf{diritto all'oblio} come esempio di
bilanciamento (oggi ancora più delicato con l'IA). Dopo la sentenza Google, non avendo strumenti per
bilanciare caso per caso, istituì una sorta di comitato di saggi.

\paragraph{Un caso della pratica del docente.} Una persona coinvolta in un processo di grande rilievo
mediatico chiedeva la rimozione di contenuti. Rimuovere la vicenda processuale (``un pezzo di storia
italiana'') non era possibile; rimuovere le sole \textbf{immagini} inopportune che la ritraevano sì,
perché non manipolava la storia. Google ha accolto la richiesta, anche se tecnicamente non si riesce a
eliminare tutte le copie diffuse online.

\begin{nota}[Una regola di vita]
``Ogni volta che risolvi un tuo problema, probabilmente ne crei uno a qualcun altro.'' Anche un piano
di remediation dopo un incidente va valutato per le sue conseguenze, ad esempio sui diritti
fondamentali: è l'approccio del giurista, complementare a quello dell'ingegnere che punta a risolvere
il problema.
\end{nota}

% ---------------------------------------------------------------------
\section{I principi (art.~5)}
\begin{itemize}
  \item \textbf{Liceità, correttezza e trasparenza}. Nessuno legge le informative (leggere tutte
    quelle dei propri dispositivi richiederebbe ore ogni mattina), perciò il titolare deve informare in
    modo chiaro. Molte sanzioni colpiscono trattamenti nascosti (es.\ ``per migliorare i servizi''
    si cedono i dati a centinaia di aziende) o \emph{dark pattern} grafici (un solo clic che spunta
    tutti i consensi).
  \item \textbf{Limitazione della finalità}: finalità determinate, esplicite e legittime; l'ulteriore
    trattamento a fini di archiviazione, ricerca o statistica non è incompatibile.
  \item \textbf{Minimizzazione}: meno dati hai, meno danni puoi subire.
  \item \textbf{Esattezza}: dati inesatti o non aggiornati creano danni (una biografia online con
    l'università sbagliata; una condanna in primo grado ancora online dopo l'assoluzione in appello).
  \item \textbf{Limitazione della conservazione}: legata alla minimizzazione; serve una
    \textbf{retention policy} che fissi per quanto tempo conservare ogni tipo di dato.
  \item \textbf{Integrità e riservatezza}: sicurezza adeguata con misure tecniche e organizzative.
    Con la disponibilità formano la triade RID/CIA.
\end{itemize}

\begin{esame}
Ricordare i \textbf{sei principi} dell'art.~5 (``i sei pilastri dei principi'').
\end{esame}

\paragraph{Un po' di storia.} Il GDPR è il Reg.~(UE) 2016/679, applicato dal 2018; prima c'era la
Direttiva 95/46/CE. Il merito del GDPR è aver dato a tutti gli Stati membri regole uniformi. Il
docente sottolinea che alla sua stesura hanno contribuito anche ingegneri ed esperti IT: la
combinazione giuristi--tecnici è quella giusta.

% ---------------------------------------------------------------------
\section{Trattamento ed esenzione domestica}
\textbf{Trattamento} (art.~4(2)): qualsiasi operazione su dati personali, automatizzata o no:
raccolta, registrazione, organizzazione, conservazione, modifica, estrazione, \textbf{consultazione},
uso, comunicazione, diffusione, cancellazione, distruzione\ldots{} Non va imparato a memoria, ma
bisogna ricordare che quasi ogni attività è un trattamento: anche solo \emph{guardare} i dati.

\paragraph{Esenzione domestica.} Il trattamento per attività esclusivamente personali o domestiche
non rientra nel GDPR, ma l'esenzione si interpreta \textbf{in senso restrittivo}. Il GDPR si rivolge a
chi agisce come titolare o responsabile (tipicamente persone giuridiche), non al privato nella sua
vita personale.

\subsection{Il caso Lindqvist (CGUE, C-101/01)}
Una catechista svedese, durante un corso di informatica, crea da casa un sito per i parrocchiani con
informazioni sui colleghi catechisti, che non ne erano informati. Accusata penalmente di trattamento
illecito, invoca l'esenzione domestica. La Corte la esclude: l'esenzione riguarda la vita privata o
familiare, \textbf{mai la pubblicazione su Internet} di dati accessibili a un numero indefinito di
persone. Nominare persone su una pagina web, con telefono, condizioni di lavoro o hobby, è
trattamento automatizzato (allora ai sensi della Direttiva 95/46).

\paragraph{Ricadute pratiche.} A chi vuole fare un sito ``per la famiglia'' il docente chiede: è
aperto? Con password e senza indicizzazione si resta nella sfera privata; aprendolo al pubblico
l'esenzione diventa problematica. Sui social i termini di servizio parlano di ``community'', ma una
community di miliardi di persone difficilmente è privata: dipende dal bilanciamento e dal giudice.
Esempio della zona grigia: una foto di Street View usata come prova in una causa di separazione a
Napoli. Ogni foto pubblicata può avere conseguenze imprevedibili.

% ---------------------------------------------------------------------
\section{I quattro pilastri del GDPR}
\begin{enumerate}
  \item \textbf{Registro delle attività di trattamento} tenuto dal titolare (o dal suo
    rappresentante).
  \item \textbf{Responsabili del trattamento}: il titolare deve usare solo responsabili che offrano
    garanzie sufficienti su misure tecniche e organizzative adeguate.
  \item \textbf{DPIA}: valutazione d'impatto, prima del trattamento, quando questo può presentare un
    rischio elevato per i diritti e le libertà delle persone.
  \item \textbf{Data breach}: notifica all'autorità di controllo senza ingiustificato ritardo e, ove
    possibile, entro \textbf{72 ore} dalla conoscenza.
\end{enumerate}

\section{Le tre categorie di dati}
Il docente semplifica così:
\begin{itemize}
  \item \textbf{dati comuni}: nome, cognome, email, indirizzo;
  \item \textbf{categorie particolari} (art.~9): opinioni politiche, religione, orientamento sessuale,
    appartenenza sindacale, salute\ldots{} da trattare con maggiore protezione;
  \item \textbf{dati giudiziari} (art.~10): condanne penali e reati.
\end{itemize}
In Europa i dati giudiziari sono tra i più protetti per il principio della \textbf{funzione rieducativa
della pena} (risalente a Cesare Beccaria): chi ha scontato la pena deve poter ricominciare. Negli USA
l'approccio è opposto (esistono app che collegano una foto ai registri penali).

\begin{esame}
Ricordare le \textbf{tre categorie}: comuni, particolari, giudiziari. A ciascuna corrispondono basi
giuridiche diverse.
\end{esame}

\section{Le basi giuridiche}
\begin{itemize}
  \item \textbf{Dati comuni} (art.~6): consenso, contratto, obbligo legale, interessi vitali, compito
    di interesse pubblico, \textbf{legittimo interesse} del titolare (purché non prevalgano i diritti
    fondamentali dell'interessato).
  \item \textbf{Categorie particolari} (art.~9): consenso, sicurezza e protezione sociale, interessi
    vitali, attività di fondazioni e associazioni senza scopo di lucro, dati resi manifestamente
    pubblici, difesa in giudizio, interesse pubblico rilevante, medicina preventiva o del lavoro.
  \item \textbf{Dati giudiziari} (art.~10): solo sotto il controllo dell'autorità pubblica o se
    autorizzati dal diritto dell'Unione o degli Stati membri.
\end{itemize}
Non serve sapere a memoria tutte le basi giuridiche, ma sapere che \textbf{dipendono dalla categoria}
di dati trattati.

\subsection{Il consenso (art.~7)}
\begin{itemize}
  \item Il titolare deve poter \textbf{dimostrare} che l'interessato ha acconsentito.
  \item La richiesta deve essere chiaramente distinguibile dal resto, comprensibile, in linguaggio
    semplice.
  \item Non è libero se l'interessato non ha una scelta reale o non può rifiutare o revocare senza
    subire un pregiudizio (considerando 42), né quando c'è un \textbf{evidente squilibrio} tra
    interessato e titolare (considerando 43).
  \item Si può revocare in ogni momento, ma la revoca non rende illecito il trattamento precedente.
\end{itemize}
\textbf{Squilibrio}: tipico del rapporto di lavoro (un datore che chiede informazioni sulla salute
come condizione per l'assunzione). Su domanda di uno studente il docente chiarisce: il titolare è
l'organizzazione, l'interessato la persona fisica; c'è squilibrio quando chi chiede il consenso ha un
potere su di te (anche l'università rispetto allo studente), molto meno quando sei un cliente libero
di comprare altrove.

\paragraph{Cookie wall e cookie paywall.} Il \emph{cookie wall} (accesso solo accettando tutti i
cookie) è vietato secondo l'EDPB. Il \emph{paywall} offre un'alternativa: abbonarsi o accettare i
cookie di profilazione. Nell'ottobre 2022 i siti del gruppo GEDI lo hanno introdotto e il Garante ha
aperto un'istruttoria. Le autorità austriaca e francese lo ritengono valido se l'abbonamento ha un
costo moderato; la decisione italiana, secondo le slide, è ancora attesa. Domande d'esame sui
paywall non sono previste: il caso serve solo a illustrare il requisito del consenso libero.

% ---------------------------------------------------------------------
\section{Domande aperte e punti da verificare}
\begin{verifica}
\begin{itemize}
  \item I nomi dei soggetti negli esempi di contitolarità (catena di scarpe e fondazione) sono
    incomprensibili nella trascrizione; l'esempio vale comunque in astratto.
  \item Il docente cita un caso di truffa con voce clonata da IA ai danni di Fideuram, avvenuto ``oggi'';
    importi e dettagli non sono affidabili nella trascrizione (ripreso il 28/09).
  \item Il docente dice che la vicenda Lindqvist ``è del 2001'': è il numero di ruolo della causa
    (C-101/01); la data della sentenza va verificata.
\end{itemize}
\end{verifica}

% =====================================================================
\lezione{28/09/2026}{GDPR (II): legittimo interesse, diritti, responsabili, accountability, art.~32 e DPIA}
  {Prof.~Giuseppe Vaciago}
  {interessato; legittimo interesse e videosorveglianza; diritti dell'interessato; sanzioni; data
   retention; DPO; responsabile e sub-responsabile, Data Privacy Framework; accountability e stato
   dell'arte; art.~32 e misure tecniche e organizzative; risk assessment, privacy by design e by default,
   DPIA e consultazione preventiva}
\label{lez:2026-09-28}
% =====================================================================

\begin{nota}[Fonte]
Trascrizione automatica della registrazione, revisionata. Durante la pausa la registrazione contiene
conversazioni in italiano non pertinenti, ignorate. Slide di riferimento: \emph{02.~Cybersecurity Laws
and Regulations -- II Privacy}, dal legittimo interesse fino al \emph{risk assessment} (la parte sul
data breach, moduli M4--M5, è rinviata). In coda, uno scambio in italiano con uno studente collegato
da remoto sul significato di \emph{formal assessment}.
\end{nota}

\section{Sintesi}
Seconda parte del GDPR. Dopo un riepilogo e la definizione di \textbf{interessato}, il docente chiarisce
che il \textbf{legittimo interesse} è del titolare, non dell'interessato, e passa ai diritti, alle
\textbf{sanzioni} (con la distinzione 20M/4\% e 10M/2\%), alla conservazione dei dati e alla figura del
DPO. Ampio spazio al \textbf{responsabile} e al \textbf{sub-responsabile} (con un caso reale di
misurazione del piede con IA e il Data Privacy Framework UE--USA) e all'\textbf{accountability},
legata allo \textbf{stato dell'arte}. Chiudono l'art.~32, le misure tecniche e organizzative e una
prima introduzione al risk assessment (probabilità e gravità), alla privacy by design e by default e
alla DPIA con consultazione preventiva (esempio dell'app Immuni).

% ---------------------------------------------------------------------
\section{Organizzazione: i laboratori}
\begin{itemize}
  \item Funzionamento: il docente presenta un caso (il primo sarà un \textbf{data breach}), gli studenti
    lavorano in piccoli gruppi e poi presentano. I migliori gruppi (3--5, a seconda dei numeri)
    ottengono \textbf{punti extra} per l'esame.
  \item Due obiettivi: imparare un \textbf{ragionamento giuridico} semplice, diverso da quello tecnico,
    e \textbf{parlare in pubblico}.
  \item La grande maggioranza dei presenti è interessata. Decisione annunciata dopo la pausa: i
    \textbf{prossimi due laboratori} si fanno \textbf{da remoto}, poi \textbf{dal 16 ottobre in
    presenza}. Niente modalità ibrida: chi è collegato resta escluso dal lavoro di gruppo.
\end{itemize}

% ---------------------------------------------------------------------
\section{Riepilogo e definizione di interessato}
Riepilogo della lezione del 25/09: ruoli, bilanciamento, Google Spain (``importante l'esempio, non il
caso''), principi, trattamento (``ogni attività può essere un trattamento''), esenzione domestica,
quattro pilastri, tre categorie di dati, consenso.

L'\textbf{interessato} (\emph{data subject}) è la \textbf{persona fisica} a cui si riferiscono i dati,
mai una persona giuridica. Il GDPR va rispettato da chi tratta i dati (titolare, responsabile); la
persona fisica è chi il GDPR protegge. Siamo tutti interessati.

% ---------------------------------------------------------------------
\section{Il legittimo interesse}
\begin{esame}
Il legittimo interesse è \textbf{del titolare}, non dell'interessato: l'interessato ha
\emph{diritti fondamentali}, non ``legittimi interessi''. È un errore frequente e una possibile domanda
d'esame.
\end{esame}

\paragraph{Videosorveglianza.} Chiedere il consenso a ogni persona ripresa è impossibile (e in alcuni
luoghi, come stadi e concerti, le telecamere sono imposte dalla legge). La base giuridica è il
legittimo interesse, frutto di un bilanciamento tra la privacy di chi passa e l'interesse a
ricostruire, ad esempio, un incidente stradale. Le slide ricordano che il bilanciamento resta
obbligatorio anche nei casi tipici (banche, gioiellerie, distributori di benzina).

\paragraph{Marketing ai clienti.} Un negozio di abbigliamento manda offerte a chi ha già comprato: da un
lato il suo legittimo interesse verso un cliente, dall'altro il diritto di non essere disturbati. Serve
il \textbf{balancing test}, caso per caso: è ragionevole proporre prodotti simili a quelli acquistati,
non qualsiasi cosa (altrimenti è spam). Quando la base è il legittimo interesse esiste \textbf{sempre}
il \textbf{diritto di opposizione} (il link ``unsubscribe'').

\paragraph{Il legittimo interesse non vale per le categorie particolari.} Su domanda di uno studente:
per i dati dell'art.~9 non si può usare il legittimo interesse (si veda l'elenco delle basi dell'art.~9
nelle slide). Un'immagine di videosorveglianza che mostra una caratteristica visibile non è per questo
un dato sanitario; una patologia non visibile sì.

\subsection{Europa e Stati Uniti: due culture}
Negli USA si può fare quasi tutto con i dati dei propri clienti, tranne venderli; si ha invece paura
dello Stato (``Big Brother''). In Europa è il contrario: si accetta il trattamento pubblico e si teme
quello privato. Per il docente entrambe le posizioni sono sbagliate: grandi aziende tecnologiche (e ora
quelle dell'IA) hanno un peso economico superiore a molti Stati, e l'unica strada è la
\textbf{cooperazione pubblico--privato} e il bilanciamento. Esempio: i dati sanitari dei wearable
potrebbero alimentare la prevenzione, ma il quadro privacy lo rende difficilissimo, mentre molti
dispositivi a basso costo cedono comunque i dati a Paesi non democratici.

% ---------------------------------------------------------------------
\section{I diritti dell'interessato}
\begin{itemize}
  \item \textbf{Accesso}: riguarda solo i \emph{propri} dati. Esempio del docente: lo studente può
    accedere ai dati che ha fornito all'esame, ma la valutazione del docente non è un suo dato.
  \item \textbf{Rettifica} dei dati inesatti.
  \item \textbf{Cancellazione} (diritto all'oblio).
  \item \textbf{Limitazione del trattamento}: poco usata e difficile da applicare.
  \item \textbf{Portabilità}: ricevere i dati in formato strutturato. Fondamentale in teoria, molto
    complessa in pratica: alla prima richiesta di un attivista austriaco, Facebook consegnò i dati su
    migliaia di pagine stampate. Per il docente la portabilità richiede interoperabilità e open
    source.
  \item \textbf{Opposizione}: chiedere che cessi un trattamento, salvo motivi prevalenti del titolare.
\end{itemize}

% ---------------------------------------------------------------------
\section{Sanzioni}
In generale fino a \textbf{20 milioni di euro o il 4\%} del fatturato mondiale annuo; l'autorità
usa il criterio che dà l'importo più alto (per un gruppo come Google, il 4\%).

\begin{esame}
\textbf{Domanda insidiosa}: per le violazioni delle \textbf{misure di sicurezza}, quindi anche per i
\textbf{data breach}, il tetto è la metà: \textbf{10 milioni di euro o il 2\%} del fatturato. ``Avrete
probabilmente una domanda sul data breach all'esame.''
\end{esame}

% ---------------------------------------------------------------------
\section{Compliance in quattro regole e data retention}
Le quattro regole operative delle slide (non fondamentali per l'esame, ma utili): analizzare quali dati
si raccolgono e dove (\textbf{registro dei trattamenti}, seguendo il flusso dei dati); verificare se il
tempo di conservazione è giustificato (\textbf{policy di conservazione e cancellazione}); informare i
clienti su come modificare o cancellare i dati (informativa); monitorare chi accede ai dati.

\paragraph{Per quanto tempo si conserva un CV?} Di norma \textbf{da uno a due anni} per un'azienda;
circa \textbf{sei mesi} per le agenzie per il lavoro, che ne ricevono moltissimi. Perché cancellare?
\begin{itemize}
  \item \textbf{sicurezza}: in caso di fuga di dati, ciò che non serviva non doveva esserci;
  \item \textbf{esattezza}: un CV vecchio non riflette le competenze attuali e può danneggiare la
    persona.
\end{itemize}
Ogni categoria di dati ha un tempo diverso (in Italia, ad esempio, gli attestati di formazione possono
servire al lavoratore per decenni a fini previdenziali). E la retention va coordinata con i
\textbf{backup}: un backup che conserva dati oltre i termini viola la norma. Per il docente la
retention è ``il vero caos del GDPR'' ed è un problema tecnico, quindi degli ingegneri.

% ---------------------------------------------------------------------
\section{I soggetti del GDPR e il DPO}
Interessato, titolare, responsabile, \textbf{DPO} (\emph{Data Protection Officer}), autorità di
controllo. Il DPO (art.~37) può essere interno o esterno ed è \textbf{obbligatorio} per:
\begin{itemize}
  \item autorità e organismi pubblici;
  \item titolari o responsabili le cui attività principali richiedono il \textbf{monitoraggio regolare
    e sistematico} degli interessati su larga scala;
  \item titolari o responsabili le cui attività principali consistono nel trattamento \textbf{su larga
    scala} di categorie particolari o di dati giudiziari.
\end{itemize}
Per il docente, nel 2026, con NIS2, DORA, Cybersecurity Act, AI Act e la \emph{shadow AI}, il DPO
dovrebbe esserci in ogni azienda. Servono competenze giuridiche ma anche \textbf{tecniche} (reti,
storage, flussi di dati): oggi il 90\% dei DPO sono avvocati, ma il ruolo sarebbe adatto agli
ingegneri. Il DPO deve essere \textbf{indipendente}: risorse adeguate, nessuna istruzione dal titolare
sui propri compiti, nessuna penalizzazione per averli svolti, nessun conflitto di interessi; non è
personalmente responsabile della non conformità. Nella pratica subisce le pressioni del CEO, che vede
compliance e sicurezza solo come costi. Il docente cita una truffa da decine di milioni di euro ai
danni di Fideuram con la voce di un avvocato clonata dall'IA: episodi così stanno cambiando la
percezione. La decisione finale resta comunque del titolare.

% ---------------------------------------------------------------------
\section{Responsabile e sub-responsabile}
\paragraph{Responsabile (art.~28).} Il criterio per distinguerlo dal titolare è il \textbf{potere
decisionale} sul trattamento. Il responsabile deve offrire \textbf{garanzie sufficienti} su misure
tecniche e organizzative adeguate (art.~28(1)). Il contratto (il \emph{Data Processing Agreement},
art.~28(3)) deve indicare: oggetto e durata, natura e finalità, tipo di dati e categorie di
interessati, obblighi e diritti del titolare.

\subsection{Caso: la misura del piede con l'IA}
Un'azienda francese offre ai siti di calzature un servizio: si fotografa il piede con lo smartphone e
l'IA indica la taglia giusta (riducendo i resi). L'azienda sostiene di essere \textbf{titolare}, perché
vuole usare le foto per migliorare il servizio e addestrare l'IA.
\begin{itemize}
  \item Ragionamento del docente: bisogna guardare all'\textbf{interessato}. Il cliente vuole comprare
    un paio di scarpe, non misurarsi il piede; la finalità è del marchio di calzature, che è il
    titolare. Il fornitore è \textbf{responsabile}.
  \item La \textbf{contitolarità} è possibile, ma nel bilanciamento non è nell'interesse del
    consumatore cedere a un terzo un database di milioni di immagini del corpo (i termini prevedono
    l'uso per l'addestramento dell'IA; esistono anche rischi legati a usi feticistici). Con un semplice
    rapporto titolare--responsabile, il titolare può vietare usi ulteriori e fissare i tempi di
    conservazione.
  \item Una decisione giuridica sui ruoli può cambiare la sorte dei dati di milioni di persone.
\end{itemize}

\paragraph{Sub-responsabile.} Il responsabile può ricorrere a un altro responsabile solo con
l'autorizzazione scritta, specifica o generale, del titolare (art.~28(2)), imponendogli gli stessi
obblighi; il responsabile iniziale resta pienamente responsabile verso il titolare (art.~28(4)).
Esempi: il fornitore francese che usa un provider di IA; un'agenzia di selezione che fa il backup dei
CV su Microsoft. Il titolare deve sapere quali sub-responsabili ci sono e \textbf{dove risiedono i
dati}: la catena (responsabile, sub-responsabile, sub-sub-responsabile\ldots) va seguita, e secondo il
docente moltissimi incidenti nascono da sub-responsabili di cui il titolare non era consapevole.

\paragraph{Data Privacy Framework.} Elenco delle aziende statunitensi aderenti al quadro UE--USA per il
trasferimento dei dati. Se il fornitore è nell'elenco con stato attivo (es.\ Microsoft, OpenAI), il
trasferimento è più semplice; se non c'è, servono accordi specifici come le \textbf{clausole
contrattuali standard}. Poiché la grande maggioranza dei fornitori di storage e trattamento è
statunitense, la verifica conta molto. Non è una domanda d'esame.

% ---------------------------------------------------------------------
\section{Accountability e stato dell'arte}
\textbf{Accountability} significa saper sempre spiegare, in modo legale e adeguato, \textbf{perché} si
è presa una decisione. Nella zona grigia bisogna poter dimostrare di stare dalla parte bianca.
L'ingegnere tende a risolvere il problema; il giurista guarda anche al processo e alle conseguenze
(esempio del docente: una ristrutturazione ``efficiente'' delle ferrovie con prepensionamenti di massa,
che decenni dopo pesa sul sistema pensionistico). Con l'IA, che decide in modo probabilistico e poco
spiegabile, tracciare il processo decisionale è ancora più importante.

\begin{esame}
Ciò che il docente vuole vedere dagli studenti: in un incidente si decide in pochi minuti e la
decisione può rivelarsi sbagliata, ma bisogna \textbf{documentarla e renderla tracciabile}, scrivendo
perché la si prende. ``Non diventate avvocati: serve risolvere il problema, ma documentate le
decisioni.''
\end{esame}

\paragraph{Stato dell'arte.} Con un budget illimitato si proteggerebbe tutto: le misure vanno valutate
con un budget limitato e devono essere \textbf{economicamente sostenibili}. Lo stato dell'arte cambia:
quindici anni fa l'\textbf{autenticazione a più fattori} era un'idea innovativa, oggi è di fatto un
obbligo per il titolare, e senza MFA si rischia la sanzione. L'interessato, come privato, è libero di
non usarla (anche se è fortemente consigliata).

\begin{esame}
Due parole da ricordare: \textbf{accountability} e \textbf{stato dell'arte} (spiegati bene
nell'art.~32).
\end{esame}

% ---------------------------------------------------------------------
\section{Art.~32: sicurezza del trattamento}
Tenendo conto dello stato dell'arte, dei costi di attuazione, della natura, dell'ambito, del contesto e
delle finalità del trattamento e dei rischi per i diritti e le libertà delle persone, titolare e
responsabile mettono in atto \textbf{misure tecniche e organizzative adeguate al rischio}. L'elenco è
\textbf{aperto} (pseudonimizzazione, cifratura, riservatezza e integrità, disponibilità e resilienza,
test periodici): prima c'era una lista chiusa di misure, ma con nuove vulnerabilità ogni settimana una
lista non regge. Bisogna decidere ogni volta e saper spiegare perché la soluzione scelta è adeguata.

\begin{integrazione}[la vecchia lista di misure]
La ``lista'' a cui il docente si riferisce è l'Allegato~B al D.lgs.~196/2003 nella versione originaria
(\emph{Disciplinare tecnico in materia di misure minime di sicurezza}), superato con l'arrivo del GDPR e
del suo approccio basato sul rischio.
\end{integrazione}

\begin{itemize}
  \item Va verificata la sicurezza del \textbf{responsabile}, che materialmente detiene i dati: la scelta
    tra provider (grandi hyperscaler statunitensi o provider europei) è una decisione caso per caso,
    con vantaggi e svantaggi di compliance e sicurezza.
  \item Chiunque agisca sotto l'autorità di titolare o responsabile tratta i dati solo su istruzione
    (art.~32(4)): servono \textbf{istruzioni e formazione}.
  \item Per ruoli e responsabilità in caso di incidente il docente richiama la matrice \textbf{RACI}
    (chi è responsabile, chi decide, chi è consultato, chi è informato): come in un'evacuazione
    antincendio, ognuno deve sapere cosa fare.
\end{itemize}

\paragraph{Misure tecniche e organizzative (report ENISA 2021).} Tecniche: cultura della
cybersicurezza, formazione, rapporti con terze parti (anche con ENISA), procedura di data breach,
accesso sicuro ai sistemi, sicurezza dei dispositivi, della rete, fisica, backup, cloud, siti web,
condivisione delle informazioni. Organizzative: modello organizzativo, informative, policy (retention,
data breach, valutazione dei rischi), privacy by design, consensi, documentazione (registro dei
trattamenti, DPA). Le misure organizzative sono ``pezzi di carta'', ma concreti.

\begin{esame}
Conoscere la \textbf{differenza tra misure tecniche e misure organizzative}. Gli elenchi delle slide
sono solo esempi e non vanno imparati a memoria.
\end{esame}

% ---------------------------------------------------------------------
\section{Risk assessment: prima introduzione}
Il docente presenta una visione ``più giuridica'' (\emph{anticipatory compliance}); Gianluca ne darà
una più ibrida, tecnico-giuridica, e la \textbf{definizione di riferimento di risk assessment} arriverà
con lui tra qualche settimana.
\begin{itemize}
  \item \textbf{Analisi preliminare}: mappare sistemi \emph{e} dati. Il tecnico guarda ai sistemi, ma
    la qualità dei dati cambia tutto.
  \item \textbf{Strategia sui dati} e \textbf{piano di remediation}: portare il rischio, ad esempio,
    da 8 a 3 su 10.
  \item \textbf{Minimizzazione e approccio by design}.
\end{itemize}

\paragraph{Dati finanziari: ENISA e GDPR.} Per il GDPR i dati finanziari sono \textbf{dati comuni};
nella metodologia ENISA di severity assessment (che si vedrà nel laboratorio) pesano più dei dati
semplici, perché nasce da una logica di cybersicurezza vicina al settore finanziario: l'obiettivo
finale del criminale è quasi sempre il denaro.

\paragraph{By design e by default.} \emph{By design}: valutare la privacy fin dal primo momento della
progettazione, non a prodotto finito (``a quel punto è troppo tardi''). \emph{By default}: per
impostazione predefinita si trattano solo i dati necessari; esempio dell'\textbf{opt-in} europeo
(serve il tuo sì) contro l'\textbf{opt-out} statunitense (si tratta finché non ti opponi).

\paragraph{Il modello in tre fasi delle slide.} Valutazione generale (artt.~24 e 32) e approccio by
design/by default (art.~25); valutazione formale (art.~35, rischio elevato); consultazione preventiva
(art.~36); compliance ed enforcement (art.~83). Il docente avverte che questa slide è
``discutibile'' e di non prenderla come definizione definitiva.

\paragraph{DPIA e consultazione preventiva: l'esempio Immuni.} L'app di contact tracing del periodo
COVID fu considerata ad alto rischio: dopo la DPIA il rischio residuo restava elevato, quindi serviva la
\textbf{consultazione preventiva} con il Garante, che indicò ulteriori misure di mitigazione prima del
via libera. Dopo l'approvazione il lavoro non finisce: la compliance va \textbf{monitorata di continuo},
perché ogni applicazione può generare nuovi rischi.

\paragraph{Probabilità e gravità.} Ogni valutazione del rischio si basa su due parole:
\textbf{probabilità} (\emph{likelihood}) e \textbf{gravità} (\emph{severity}). Le misure tecniche e
organizzative servono a mitigare il rischio così stimato; l'approccio è circolare (valutazione,
misure, test).

\subsection{Dopo la lezione: valutazione generale e specifica}
Nello scambio in italiano con uno studente collegato si chiarisce:
\begin{itemize}
  \item la \textbf{valutazione generale} (artt.~24 e 32) riguarda tutto l'ambiente IT e ogni misura di
    sicurezza, in modo continuo; la valutazione dell'\textbf{art.~35} (DPIA) riguarda un
    \textbf{trattamento specifico} (es.\ un nuovo smartwatch) ed è più dettagliata. Più che
    ``formale'', sarebbe meglio chiamarla ``specifica'';
  \item probabilità $\times$ gravità dà il \textbf{rischio inerente}; applicando le misure di mitigazione
    si ottiene il \textbf{rischio residuo}, cioè quello che resta dopo le misure. Se resta alto, si va dall'autorità, che può indicare altre
    misure o vietare il trattamento;
  \item con l'AI Act, invece, il livello di rischio dei sistemi ad alto rischio non si autodefinisce: è
    la norma a classificarli (es.\ IA sul posto di lavoro).
\end{itemize}

% ---------------------------------------------------------------------
\section{Domande aperte e punti da verificare}
\begin{verifica}
\begin{itemize}
  \item \textbf{Calendario dei laboratori}: dalla trascrizione, i due lab da remoto dovrebbero essere il
    2 e il 9 ottobre, poi in presenza dal 16 ottobre. Il 2/10 si è poi svolta solo una simulazione (vedi
    lezione del 02/10).
  \item \textbf{Sanzione a Google per Street View}: il docente cita una sanzione di 400 milioni di euro
    ``di ieri'' a Google Ireland. Notizia non verificata.
  \item \textbf{Data Privacy Framework}: secondo il docente un noto fornitore di IA non compare
    nell'elenco. Da verificare sul sito ufficiale, perché la situazione cambia nel tempo.
  \item \textbf{Fideuram}: l'importo della truffa con voce clonata non è affidabile nella trascrizione.
  \item \textbf{Modello in tre fasi}: il docente ha dichiarato la slide ``discutibile''; la definizione
    di risk assessment da usare all'esame è quella che darà Gianluca.
  \item \textbf{Riferimenti dell'art.~28}: paragrafi (1), (2), (3), (4) ripresi dalle slide; il docente
    non li ha letti in aula.
\end{itemize}
\end{verifica}

% =====================================================================
\lezione{02/10/2026}{Simulazione di laboratorio: severity assessment di un ransomware}
  {Prof.~Giuseppe Vaciago}
  {organizzazione dei laboratori e punti extra; caso Alpha S.r.l.\ (ransomware su piattaforma
   e-commerce B2B); ruoli e obblighi del responsabile; autovalutazione sul sito del Garante; modulo di
   notifica; metodologia ENISA di severity assessment ($SE = DPC \times EI + CB$); decisione sulla notifica}
\label{lez:2026-10-02}
% =====================================================================

\begin{nota}[Fonte]
Trascrizione automatica della registrazione, revisionata. Lezione \textbf{online}, svolta come
\textbf{simulazione} di laboratorio, con interazione via chat. Scale e soglie della metodologia sono
riprese dal documento \emph{03.~Data breach severity methodology} (ENISA, \emph{Recommendations for a
methodology of the assessment of severity of personal data breaches}, v1.0, dicembre 2013). In aula
si è usato anche lo strumento di autovalutazione sul sito del Garante (tradotto con il traduttore del
browser). I nomi delle aziende citate nel caso sono esempi di fantasia.
\end{nota}

\section{Sintesi}
Invece del laboratorio vero e proprio, il docente svolge una simulazione guidata su un caso
volutamente semplice: un ransomware blocca per 24 ore una piattaforma di e-commerce B2B che lavora per
conto di grandi marchi. Si ragiona su ruoli (la piattaforma è \textbf{responsabile}, i marchi sono
\textbf{titolari}), sul modulo di notifica al Garante e soprattutto sulla \textbf{metodologia ENISA di
severity assessment}, applicata passo per passo. Il risultato ($SE = 2$, rischio medio) porta alla
notifica al Garante; il docente consiglia di informare comunque gli utenti. Messaggio centrale: ogni
scelta nella formula pesa molto e va \textbf{motivata} (accountability); il CEO chiederà quasi sempre
di abbassare il punteggio.

% ---------------------------------------------------------------------
\section{Organizzazione}
\begin{itemize}
  \item \textbf{Calendario}: lunedì 5/10 lezione (in presenza) che chiude la parte di protezione dei
    dati con riferimento alla cybersicurezza; a fine lezione viene presentato il caso del primo
    \textbf{laboratorio ufficiale}, che si svolge \textbf{venerdì 9/10} online.
  \item \textbf{Punti extra}: \textbf{0,25 punti} sul voto d'esame per ogni laboratorio, per le
    elaborazioni migliori; non è garantito che tutti li ricevano. Le regole definitive saranno pubblicate
    dopo il primo laboratorio ufficiale, su una pagina dedicata.
  \item \textbf{Gruppi}: massimo 5 persone (6 tollerato, non 8--10), formati dagli studenti; il docente
    può aiutare chi non trova un gruppo. Si può lavorare anche da soli, ma il docente consiglia
    fortemente il gruppo.
  \item \textbf{Modalità}: il lab si svolge nell'ora e mezza del venerdì. I gruppi possono usare una
    propria call (Zoom, Meet) o trovarsi al Politecnico; per le domande il docente apre una stanza
    separata con il gruppo. Più tempo solo in casi eccezionali.
  \item Il documento ENISA viene caricato sulla piattaforma: \textbf{non è materia d'esame}, ma è ``la
    Bibbia'' per il laboratorio.
\end{itemize}

% ---------------------------------------------------------------------
\section{Il caso: Alpha S.r.l.}
\begin{itemize}
  \item \textbf{Alpha S.r.l.} gestisce in modalità SaaS l'e-commerce di circa \textbf{20 grandi
    clienti} (marchi di moda, bricolage\ldots): chi compra sul sito del marchio viene in realtà servito
    dalla piattaforma di Alpha, basata a sua volta su Shopify (un \textbf{sub-responsabile}). Alpha
    gestisce anche la logistica delle consegne e trattiene una percentuale sulle vendite.
  \item Un \textbf{ransomware} blocca tutte le transazioni per \textbf{24 ore}. Il team di risposta
    (DPO, legale, IT, CISO) verifica che i backup sono integri e, con i piani di \textbf{business
    continuity} e \textbf{disaster recovery}, ripristina tutto in 24 ore.
  \item Domanda: che cosa deve fare il DPO dal punto di vista della protezione dei dati? Per decidere se
    la notifica al Garante è obbligatoria serve un \textbf{severity assessment}.
\end{itemize}

\subsection{Primi passi}
\paragraph{Autovalutazione sul sito del Garante.} Lo strumento chiede: c'è stato un incidente di
sicurezza con perdita di riservatezza, integrità o disponibilità? Sì: almeno la
\textbf{disponibilità} è certa. Riguarda dati personali? Sì: i clienti di Alpha sono persone giuridiche,
ma ci sono i dati dei loro dipendenti e soprattutto quelli degli \textbf{utenti finali} che comprano
online. È quindi una violazione di dati personali. La pagina del Garante rimanda alla metodologia ENISA.

\paragraph{Chi notifica.} Alpha è \textbf{responsabile}: deve informare \textbf{senza ingiustificato
ritardo} ogni titolare, cioè ciascuno dei suoi clienti. Il paradosso tipico degli incidenti: notifica il
titolare, ma tutte le informazioni le ha il responsabile, che deve collaborare.
\begin{itemize}
  \item Con 20 clienti, \textbf{ogni titolare} fa la propria notifica (e l'autorità dovrà poi riunire
    20 pratiche sullo stesso caso).
  \item Eccezione pratica raccontata dal docente: in un breach di un'app di pagamento usata da moltissimi
    servizi, il Garante chiese al responsabile un'unica notifica per conto di tutti i titolari, per non
    riceverne un numero enorme. È una soluzione eccezionale, non la regola.
\end{itemize}

\begin{esame}
Ricordare ruoli e flusso: il \textbf{responsabile} informa il \textbf{titolare}; il \textbf{titolare}
notifica all'autorità. Nel caso, il titolare è il marchio, non la piattaforma.
\end{esame}

\subsection{Il modulo di notifica}
\begin{itemize}
  \item \textbf{Titolare}: il marchio cliente (si ipotizza che Alpha gli abbia girato tutte le
    informazioni).
  \item \textbf{Date}: data della violazione e data in cui il titolare ne ha avuto ragionevole certezza.
    È decisiva per le 72 ore. Nel caso: attacco il 30 settembre, consapevolezza il 1° ottobre. Spesso
    l'esfiltrazione precede di molto il ransomware ed è difficile datarla.
  \item \textbf{Natura}: disponibilità certa; la \textbf{riservatezza} è il vero nodo. Con la doppia
    estorsione l'attaccante chiede un riscatto e pubblica un campione dei dati esfiltrati sul dark web:
    a quel punto c'è una presunzione di perdita di riservatezza. A volte la causa si chiarisce solo dopo
    lunghe analisi.
  \item \textbf{Causa}: se non è nota, si dichiara che al momento non è nota.
  \item \textbf{Numero di interessati}: qualche centinaio di dipendenti dei clienti e, soprattutto, gli
    utenti: per la simulazione \textbf{50.000}.
  \item \textbf{Contatti del DPO}.
\end{itemize}

% ---------------------------------------------------------------------
\section{La metodologia ENISA}
\[
SE = DPC \times EI + CB
\]
\begin{itemize}
  \item \textbf{DPC} (\emph{Data Processing Context}): tipo di dati e contesto del trattamento.
  \item \textbf{EI} (\emph{Ease of Identification}): quanto è facile risalire alla persona.
  \item \textbf{CB} (\emph{Circumstances of the Breach}): perdita di riservatezza, integrità,
    disponibilità e intento malevolo; i punteggi si \textbf{sommano}.
\end{itemize}

\begin{center}
\small
\begin{tabularx}{\linewidth}{l X}
\toprule
\textbf{Fattore} & \textbf{Valori (punteggio di base)}\\ \midrule
DPC & dati semplici 1; comportamentali 2; finanziari 3; sensibili (artt.~9--10) 4\\
EI  & trascurabile 0,25; limitata 0,5; significativa 0,75; massima 1\\
CB -- riservatezza & 0 nessuna evidenza di trattamento illecito; +0,25 destinatari noti;
  +0,5 numero indeterminato di destinatari\\
CB -- integrità & 0; +0,25 alterazione recuperabile; +0,5 non recuperabile\\
CB -- disponibilità & 0 recupero senza difficoltà; +0,25 indisponibilità temporanea;
  +0,5 indisponibilità totale\\
CB -- intento malevolo & +0,5\\ \bottomrule
\end{tabularx}
\end{center}

\begin{center}
\small
\begin{tabular}{ll}
\toprule
\textbf{Punteggio} & \textbf{Severità}\\ \midrule
$SE < 2$ & bassa\\
$2 \le SE < 3$ & media\\
$3 \le SE < 4$ & alta\\
$SE \ge 4$ & molto alta\\ \bottomrule
\end{tabular}
\end{center}

\begin{esame}
La formula $SE = DPC \times EI + CB$ è ``la regola d'oro'': va \textbf{ricordata a memoria} e descritta
con chiarezza all'esame. È una ``formula stupida'' ma con conseguenze molto concrete.
\end{esame}

\subsection{Applicazione al caso}
\paragraph{DPC: quali dati?} Discussione con la classe:
\begin{itemize}
  \item \textbf{semplici} (nome, cognome, email, codice fiscale, telefono): base 1;
  \item \textbf{comportamentali}? Lo storico degli acquisti su più marchi permette di ricostruire un
    profilo (gusti, capacità di spesa). Il CEO direbbe di no, il DPO potrebbe dire di sì: è un punto su
    cui ``si può litigare'' a lungo;
  \item \textbf{finanziari}? Probabilmente no: i pagamenti passano spesso da servizi esterni (es.\
    Stripe) e il CEO afferma che i dati delle carte non sono conservati. Attenzione: per il GDPR i dati
    finanziari sono dati \emph{comuni}; è ENISA, con la sua logica di cybersicurezza, a pesarli 3;
  \item \textbf{sensibili}? Il data center potrebbe contenere dati sanitari dei dipendenti di Alpha, ma
    il titolare qui è il marchio, che non tratta dati di quel tipo.
\end{itemize}
Scelta del docente: \textbf{DPC = 1} (dati semplici), pur senza escludere i comportamentali, purché la
scelta sia motivata. ``Non ho una risposta: è la vostra scelta, il vostro livello di accountability.''

\paragraph{EI.} Nome, cognome, data di nascita, telefono insieme: \textbf{EI = 1}. Esempio opposto: un
solo indirizzo email composto da un numero di matricola renderebbe l'identificazione difficile
(0,25), mentre \texttt{nome.cognome@polito.it} no.

\paragraph{CB.}
\begin{itemize}
  \item \textbf{Riservatezza +0,25}: per ora i dati sono nelle mani dell'attaccante e forse sul dark
    web, ma non c'è certezza di una pubblicazione aperta. Per il docente il ``numero indeterminato di
    destinatari'' (+0,5) corrisponde a una pubblicazione accessibile a tutti; si potrebbe sostenere
    anche +0,5.
  \item \textbf{Integrità 0}: nessuna alterazione.
  \item \textbf{Disponibilità +0,25}: indisponibilità temporanea (24 ore).
  \item \textbf{Intento malevolo +0,5}: certamente sì.
\end{itemize}
\[
SE = 1 \times 1 + (0{,}25 + 0 + 0{,}25 + 0{,}5) = 2 \quad\Rightarrow\quad \text{media (``medio-bassa'')}
\]
Nella sezione sulle misure adottate conviene descrivere il ripristino in 24 ore: aiuta a giustificare
un punteggio contenuto.

\subsection{Dal punteggio alla decisione}
Lettura data in aula:
\begin{itemize}
  \item rischio \textbf{basso}: nessuna notifica al Garante;
  \item rischio \textbf{medio}: \textbf{notifica al Garante};
  \item rischio \textbf{alto o molto alto}: notifica al Garante e \textbf{comunicazione agli
    interessati}.
\end{itemize}
Con $SE = 2$ si è al confine: in teoria si potrebbe discutere se evitare la notifica, ma con tante
ipotesi incerte la strada sicura è notificare. Il docente consiglia anche di \textbf{avvisare gli
utenti}, pur non essendo obbligatorio, per creare un rapporto di fiducia e perché i dati potrebbero
essere usati per frodi. Se fossero stati coinvolti dati delle carte (DPC 3), il risultato sarebbe stato
alto o molto alto e la comunicazione agli interessati immediata.

\begin{integrazione}[le soglie del GDPR]
Nel GDPR la notifica al Garante (art.~33(1)) è dovuta salvo che sia \emph{improbabile} che la violazione
presenti un rischio per i diritti e le libertà delle persone; la comunicazione agli interessati
(art.~34(1)) è dovuta quando il rischio è \emph{elevato}. Le fasce ENISA sono uno strumento per
motivare queste due valutazioni, non la norma stessa.
\end{integrazione}

\subsection{Il dilemma del consulente}
Uno studente osserva che nessuno ha interesse a dare una valutazione peggiore del necessario, e che
quindi il punteggio verrà sempre abbassato. Il docente conferma: nel 95\% dei casi il CEO chiederà di
``ammorbidire'' la valutazione. Abbassare i punteggi è legittimo quando ha senso ed è motivato; farlo in
modo malevolo, nascondendo o falsando informazioni, espone alla sanzione (10 milioni o 2\% del
fatturato) e al rischio penale per false dichiarazioni al Garante (reclusione da sei mesi a tre anni),
che colpisce chi firma la dichiarazione. Sulla riservatezza: dati finiti sul dark web sono in generale
da considerare \textbf{persi}, anche se poi rimossi, ma la risposta va data caso per caso.

% ---------------------------------------------------------------------
\section{Domande aperte e punti da verificare}
\begin{verifica}
\begin{itemize}
  \item \textbf{Soglia di notifica}: nella parte finale la trascrizione è confusa (``obbligati con 2,5''?
    ``con 3''?). Qui si è usata la lettura esplicita data poco prima: basso $\rightarrow$ nessuna
    notifica; medio $\rightarrow$ Garante; alto/molto alto $\rightarrow$ anche interessati. Confrontare
    con la spiegazione teorica della lezione del 05/10.
  \item \textbf{Riservatezza +0,25 o +0,5}: nella metodologia ENISA +0,25 indica dati esposti a
    destinatari \emph{noti}; un attaccante che potrebbe rivendere i dati è difficilmente un destinatario
    noto. La scelta del docente è dichiaratamente discutibile.
  \item \textbf{Date del caso} (attacco il 30/09, consapevolezza il 1/10) e numero di utenti (50.000):
    ipotesi della simulazione, dette a voce.
  \item Il docente ha anche detto che la metodologia ENISA risale al 2014 (lezione del 28/09); il
    documento è datato dicembre 2013.
\end{itemize}
\end{verifica}

% =====================================================================
\lezione{05/10/2026}{DPIA e data breach; presentazione del caso Arvenza}
  {Prof.~Giuseppe Vaciago}
  {regole del laboratorio; compliance anticipatoria; DPIA (contenuto, rischio inerente e residuo,
   consultazione preventiva); privacy by design e by default; data breach: definizione, artt.~33 e 34,
   consapevolezza del titolare, fattori di rischio, contenuto della notifica, responsabile del trattamento,
   remediation plan, registro dei data breach; presentazione del Lab~1 (Arvenza)}
\label{lez:2026-10-05}
% =====================================================================

\begin{nota}[Fonte]
Trascrizione automatica della registrazione, revisionata (pausa tra il minuto 94 e il 112 circa, con
formazione dei gruppi: non riportata). Slide di riferimento: \emph{02.~Cybersecurity Laws and
Regulations -- II Privacy} (moduli M4 \emph{DPIA} e M5 \emph{Data Breach}) e
\emph{03.~Data breach severity methodology~1.0} (ENISA, 2013). Il caso del laboratorio è nel PDF
\emph{Arvenza Breach Case Study} (vedi Capitolo~\ref{lab:01}).
\end{nota}

\section{Sintesi}
Lezione in due parti. La prima, più breve, è sulla \textbf{valutazione d'impatto (DPIA)}: il
docente la inquadra nel passaggio da una valutazione del rischio ``statica'' (basata sugli incidenti
già accaduti) a una \textbf{compliance anticipatoria}, che deve prevedere l'evoluzione del rischio
digitale. La seconda, centrale, è sul \textbf{data breach}: definizione, notifica al Garante entro 72
ore, comunicazione agli interessati, chi è ``il titolare'' in concreto, perché servono un piano prima
e un remediation plan dopo, e perché il \textbf{registro dei data breach} è la prova
dell'\emph{accountability}. In chiusura viene presentato il caso Arvenza per il laboratorio di
venerdì~9/10.

% ---------------------------------------------------------------------
\section{Organizzazione}
\begin{itemize}
  \item \textbf{Laboratori non obbligatori}, ma fanno parte del programma d'esame: durante il lab il
    docente dà indicazioni utili per le \textbf{domande aperte}. Chi non partecipa deve comunque
    ascoltare la registrazione (anche in differita; il docente suggerisce di usare un'IA per trovare
    i punti in cui parla, dato che lunghi tratti sono muti).
  \item \textbf{Punti bonus}: ogni laboratorio consegnato con esito sufficiente vale un piccolo
    punteggio aggiuntivo (in aula si parla di 0,25 punti a lab, su un totale previsto di 8--10 lab).
    Una parte dei gruppi (indicativamente il 25--30\%, forse di più) potrà presentare il proprio
    lavoro il lunedì successivo.
  \item \textbf{Formazione dei gruppi}: si è dedicata parte della lezione a creare i gruppi (in aula
    e online) per partire subito venerdì. Si può lavorare anche da soli, ma il docente lo sconsiglia.
  \item \textbf{Tempi del lab}: un'ora e mezza in aula, consegna via email al docente entro la fine
    del lab (indicate le 13:15); nessun lavoro richiesto prima o dopo. Solo per questo primo lab è
    consentito prepararsi in anticipo. Le consegne in ritardo non sono considerate, salvo situazioni
    eccezionali.
  \item \textbf{Formato libero}: il docente non fornisce un modello di severity assessment; ogni
    gruppo sceglie il formato che ritiene più efficace per presentare il rischio al Garante.
  \item Se pochi gruppi parteciperanno, il docente si riserva di cambiare formato ai laboratori.
  \item Il docente è disponibile per \textbf{tesi} (meno per tirocini, perché come avvocato ha meno
    contatti con aziende che cercano profili IT).
\end{itemize}

% ---------------------------------------------------------------------
\section{Dalla valutazione statica alla compliance anticipatoria}

Il settore in cui la valutazione del rischio è nata è la \textbf{salute e sicurezza sul lavoro}:
nell'industria (anche a Torino, nell'automotive) per decenni si moriva sul lavoro per mancanza di
prevenzione. In Italia il percorso arriva a compimento con il \textbf{D.Lgs.~81/2008}. In quel
contesto la valutazione può essere relativamente \emph{statica}: le macchine fanno sempre la stessa
cosa, l'errore nasce soprattutto dall'interazione uomo--macchina e la prevenzione passa in gran parte
dalla formazione.

Il ``corpo digitale'' è diverso: è ovunque e cambia a una velocità enorme (si pensi all'IA
generativa dal 2022 a oggi). Ne derivano tre conseguenze:
\begin{itemize}
  \item non si può risolvere il problema in modo \textbf{tattico} (``mettiamo l'antivirus''): serve una
    \textbf{strategia}, anche perché una PMI non ha il budget per prevenire tutto;
  \item il punto di rottura resta il \textbf{comportamento umano}, che va anticipato;
  \item la valutazione va scritta pensando a come verrà letta tra uno, due, tre anni:
    \textbf{compliance anticipatoria}, basata su elementi di \emph{foresight}. L'IA è insieme
    soluzione e fonte di nuovi rischi.
\end{itemize}

\begin{nota}[Indicazione per i laboratori]
Ogni valutazione fatta ``oggi'' deve reggere quando sarà riletta in futuro: chiedersi sempre come
apparirà quel documento tra qualche anno.
\end{nota}

% ---------------------------------------------------------------------
\section{La valutazione d'impatto (DPIA)}

\subsection{L'approccio del GDPR: fiducia nel titolare}
Il docente illustra la logica dell'\emph{accountability} con l'esempio di uno studente
(``Simone''): il GDPR si fida che il titolare faccia i compiti, senza che nessuno li controlli
preventivamente. La DPIA la redige il titolare e, se il rischio risulta alto, ha due strade: chiedere
la consultazione preventiva all'autorità o rinunciare al trattamento. Nella prassi, osserva il
docente, la tentazione è valutare il rischio come ``medio'' per evitare entrambe. Nell'AI Act invece,
secondo il docente, il sistema cambia: per i sistemi ad alto rischio non basta l'autovalutazione e
serve un passaggio con le autorità prima dell'uso (tema rinviato alla parte sull'AI Act).

\subsection{Quando serve una DPIA}
La DPIA va fatta quando un trattamento è \textbf{probabilmente ad alto rischio}, ma stabilirlo è
spesso una \textbf{zona grigia} lasciata alla valutazione del titolare (a differenza di norme come
NIS2, dove si entra o non si entra in una categoria). Esempi citati:
\begin{itemize}
  \item \textbf{videosorveglianza (CCTV)}: riconosciuta come trattamento che incide sui diritti
    fondamentali, quindi soggetta a DPIA;
  \item \textbf{whistleblowing}: chi lavora con un ente (dipendenti, collaboratori, fornitori) può
    segnalare violazioni in forma anonima; le piattaforme di whistleblowing sono soggette a DPIA.
\end{itemize}
Di conseguenza può capitare che per lo stesso strumento un'azienda faccia la DPIA e un'altra no,
perché i rispettivi DPO valutano il rischio in modo diverso.

\subsection{Contenuto e ciclo}
Dalle slide, la DPIA contiene: descrizione sistematica del trattamento e delle finalità;
valutazione di \textbf{necessità e proporzionalità}; valutazione dei rischi per i diritti e le libertà
degli interessati; misure previste per affrontarli e per dimostrare la conformità. Seguono
documentazione e monitoraggio.

Il rischio va valutato da \textbf{due punti di vista}: \textbf{giuridico} (impatto sui diritti
fondamentali) e \textbf{tecnico} (sicurezza del trattamento).

\paragraph{Necessità e proporzionalità.} Esempio del docente: chiedere a tutti gli studenti l'ultimo
certificato medico per dividerli in classi in base alle patologie può avere una logica sanitaria, ma
non è proporzionato, perché dà accesso a una quantità di informazioni sensibili che potrebbero essere
usate contro di loro.

\paragraph{Rischio inerente e residuo.}
\[
\text{rischio inerente} = \text{severità} \times \text{probabilità}
\qquad
\text{rischio residuo} = \text{rischio inerente} - \text{mitigazione}
\]
\begin{itemize}
  \item residuo \textbf{alto}: consultazione preventiva (art.~36) oppure stop del trattamento;
  \item residuo \textbf{medio o basso}: si prosegue, ma il docente raccomanda di ripetere la
    valutazione \textbf{ogni sei mesi}, perché nel digitale il rischio cambia in fretta (problema
    tipico: DPIA scritte nel 2020 e mai riviste).
\end{itemize}

\begin{esame}
Non fermarsi a ``rischio = severità $\times$ probabilità'': quello è solo il \textbf{rischio
inerente}. La valutazione si chiude sul \textbf{rischio residuo} (inerente meno mitigazioni), ed è il
residuo che decide se proseguire, consultare l'autorità o fermarsi. Il docente ha detto che lo
ripeterà anche in altri contesti.
\end{esame}

\paragraph{Natura e approccio (slide).} Valutazione \emph{a domande} (checklist, che il docente
preferisce perché evita dimenticanze), basata sul contesto; team interdisciplinare; analisi fin dalla
fase di progettazione; verifiche periodiche; approccio by design / by default. Una sola DPIA può
coprire trattamenti simili per natura, ambito, contesto, finalità e rischi. La pubblicazione non è
obbligatoria, ma l'EDPB suggerisce di pubblicarne almeno una sintesi o le conclusioni. Esiti possibili
della consultazione preventiva: ulteriori misure oppure stop del trattamento per mancanza di misure
di mitigazione disponibili.

\subsection{Privacy by design e by default (art.~25)}
\begin{itemize}
  \item \textbf{By design}: fin dall'inizio parte legale e parte tecnica lavorano insieme per
    integrare le garanzie nel trattamento, tenendo conto di stato dell'arte, costi di attuazione,
    natura, ambito, contesto, finalità e rischi (gli stessi criteri dell'\emph{accountability}).
  \item \textbf{By default}: per impostazione predefinita si trattano solo i dati necessari alla
    finalità (quantità, portata, conservazione, accessibilità). Contrasto con la cultura
    statunitense dell'\emph{opt-out}: in Europa vale l'\emph{opt-in}.
\end{itemize}

\begin{esame}
La DPIA non è, secondo il docente, una parte consistente dell'esame e una domanda specifica è
``rara'', ma \textbf{non si può escludere}; molte nozioni (proporzionalità, rischio residuo, by
design/by default) servono comunque per altre domande. Il ciclo della DPIA non va imparato a memoria.
\end{esame}

% ---------------------------------------------------------------------
\section{Data breach}

\begin{esame}
Sul data breach ci sarà \textbf{quasi sempre almeno una domanda} all'esame: è la parte più rilevante,
anche perché tutte le norme successive (DORA, AI Act, CRA, NIS2) prevedono notifiche di incidenti
costruite sullo stesso principio.
\end{esame}

\subsection{Perché il GDPR è stato il primo}
Il GDPR (deciso nel 2016, applicabile dal 2018) è stato la prima norma a imporre la notifica a
\textbf{tutti} i titolari, senza distinzione di settore: anche il pub che conserva le preferenze
alimentari dei clienti. Le norme successive sono invece segmentate (DORA solo per il settore
finanziario, NIS2 per alcune categorie). Anche dati apparentemente banali ``possono aprire molte
porte''.

\subsection{Definizione e tipologie}
L'\textbf{art.~4(12)} definisce la violazione dei dati personali come la violazione di sicurezza che
comporta accidentalmente o in modo illecito la \textbf{distruzione, perdita, modifica, divulgazione
non autorizzata o accesso} a dati personali trasmessi, conservati o comunque trattati.
Collegamento con la triade CIA:
\begin{itemize}
  \item distruzione e perdita $\rightarrow$ \textbf{disponibilità};
  \item modifica $\rightarrow$ \textbf{integrità};
  \item divulgazione o accesso non autorizzati $\rightarrow$ \textbf{riservatezza}.
\end{itemize}
Esempi dalle slide: smarrimento di un dispositivo non cifrato; ransomware (indisponibilità, ma spesso
anche riservatezza).

\begin{esame}
La definizione dell'art.~4(12) va \textbf{memorizzata bene}, insieme al collegamento con riservatezza,
integrità e disponibilità.
\end{esame}

\subsection{Gli obblighi: artt.~33 e 34}
\begin{itemize}
  \item \textbf{Art.~33}: notifica all'autorità di controllo \textbf{senza ingiustificato ritardo e,
    ove possibile, entro 72 ore}. Il ``ove possibile'' conta: se era materialmente impossibile (esempio
    del docente: incendio presso un fornitore) il ritardo può essere giustificato; se era possibile e
    non è stato fatto, è un problema.
  \item \textbf{Art.~34}: comunicazione agli \textbf{interessati} quando la violazione è suscettibile
    di presentare un \textbf{rischio elevato} per i loro diritti e libertà. Le aziende la temono per
    il danno reputazionale.
\end{itemize}

\subsection{Chi è ``il titolare'' quando parte l'orologio}
Le 72 ore decorrono da quando il \textbf{titolare è venuto a conoscenza} della violazione. Il titolare
è la persona giuridica, ma un'azienda ``è consapevole'' quando lo è chi la rappresenta: il
\textbf{legale rappresentante}.
\begin{itemize}
  \item Il legale rappresentante è nominato dai soci; nel 90\% dei casi coincide con il CEO, ma va
    verificato (in Italia con la \textbf{visura camerale}). Nelle multinazionali possono esistere
    deleghe specifiche, ad esempio in materia di trattamento dei dati.
  \item Per qualsiasi comunicazione legale a un'azienda conviene rivolgersi al legale
    rappresentante, che ha il dovere di rispondere.
  \item \textbf{Verità legale e verità reale}: se il legale rappresentante è informato per telefono
    martedì ma l'unica prova è un'email di giovedì, in giudizio conta ciò che si può provare.
    Da qui l'importanza delle \textbf{evidenze} della data di consapevolezza.
  \item Se il legale rappresentante non è raggiungibile (ferie, volo) e non c'è una procedura di
    delega, l'autorità lo leggerà come \textbf{carenza organizzativa}: secondo il docente in quel caso
    la sanzione è molto probabile.
\end{itemize}

\begin{esame}
Distinguere \textbf{errore umano} e \textbf{carenza organizzativa}. L'errore del singolo (il click sul
link di phishing) è umano e può essere compreso; la mancanza di procedure (chi avvisa chi, chi
decide quando il CEO è in volo) è un difetto prevedibile dell'organizzazione ed è ciò che le autorità
sanzionano. Strumento utile: la \textbf{matrice RACI} (\emph{Responsible, Accountable, Consulted,
Informed}), che sarà ripresa nelle esercitazioni su procedure e policy. Distinzione da usare nei lab.
\end{esame}

\subsection{Fattori di rischio per i diritti e le libertà}
Dalle slide, sei elementi da considerare come questione di metodo:
\begin{enumerate}[label=\alph*)]
  \item tipo di violazione;
  \item natura, numero e grado di sensibilità dei dati;
  \item facilità di associare i dati a una persona (es.\ dati pseudonimizzati: milioni di record
    ridotti a codici);
  \item gravità delle conseguenze (dati utilizzabili subito contro l'interessato: frode, sostituzione
    di persona);
  \item numero di interessati coinvolti;
  \item caratteristiche del titolare (un ospedale non è una casa vacanze).
\end{enumerate}
Il metodo operativo è la \textbf{metodologia ENISA} di severity assessment (vedi
Capitolo~\ref{lab:01} e la scheda al Capitolo~\ref{cap:schede}). Tutto è relativo: un milione di
interessati è molto per l'Italia, poco per l'India. Non esistono soluzioni bianco/nero.

\subsection{Contenuto della notifica (art.~33(3))}
Natura della violazione (categorie e numero approssimativo di interessati e di registrazioni);
contatti del \textbf{DPO}; probabili conseguenze; misure adottate o proposte per porre rimedio e
attenuarne gli effetti.

\begin{nota}[Attenzione alla sigla DPA]
``DPA'' indica sia la \emph{Data Protection Authority} sia il \emph{Data Processing Agreement}
(contratto con il responsabile ex art.~28). Il docente invita a chiedere chiarimenti, anche all'esame,
se la sigla è ambigua.
\end{nota}

Oltre ai contenuti, il titolare deve \textbf{motivare le proprie decisioni}, in particolare quando
decide di non notificare, ritarda la notifica o non comunica la violazione agli interessati. La
motivazione parte dal severity assessment: è l'applicazione diretta dell'\emph{accountability}.

\subsection{Lo scopo della notifica e la capacità di reazione}
In origine l'idea era lo \textbf{scambio di informazioni in tempo reale} per proteggere altre
organizzazioni da attacchi simili (il docente cita un lavoro del 2013, precedente al GDPR, e il ruolo
di centralizzazione di ENISA). Oggi la giurisprudenza delle autorità europee chiede di
\textbf{dimostrare la capacità di reazione}:
\begin{enumerate}
  \item un \textbf{piano prima}: procedure già in vigore (vedi il caso del legale rappresentante
    irraggiungibile);
  \item un \textbf{remediation plan dopo}. Alcuni esperti temono che un piano di rimedio equivalga
    ad ammettere una colpa, ma la violazione c'è già stata: più il piano è solido, più si abbassa il
    rischio. In alcuni casi il docente, in base all'accountability, ha deciso di non notificare
    perché il remediation plan riduceva il rischio in modo convincente.
\end{enumerate}

\begin{nota}[Indicazione per il laboratorio]
Il remediation plan può incidere sulla valutazione del rischio: considerarlo sempre, non limitarsi a
calcolare il numero. La slide sul remediation plan è indicata come ``molto importante'' per il lab.
\end{nota}

\subsection{Mappatura severità -- obblighi ed esenzioni dell'art.~34(3)}
Mappatura indicata in aula (e ripresa nel caso Arvenza):
\begin{center}
\begin{tabular}{lll}
\toprule
\textbf{Severità} & \textbf{Esito} & \textbf{Base} \\ \midrule
Bassa ($SE<2$) & solo registro & art.~33(5) \\
Media ($2 \le SE < 3$) & notifica all'autorità & art.~33 \\
Alta / molto alta ($SE \ge 3$) & autorità e interessati & art.~34, salvo 34(3) \\
\bottomrule
\end{tabular}
\end{center}
Il docente ammette che, da avvocato, si tende a ``tirare'' un po' questi livelli, ma vale ciò che è
scritto nella metodologia.

Esenzioni dalla comunicazione agli interessati (art.~34(3)):
\begin{enumerate}[label=\alph*)]
  \item misure tecniche e organizzative adeguate applicate ai dati violati (es.\ cifratura);
  \item misure successive che scongiurano il rischio elevato: difficile da sostenere, perché non si
    può garantire che una carta d'identità o di credito sottratta non venga usata;
  \item sforzo sproporzionato: si sostituisce con una \textbf{comunicazione pubblica}. Esempio vissuto
    dal docente: un'azienda di trasporto pubblico milanese con molti utenti senza email e non più
    rintracciabili, per cui si è scelto un avviso sui giornali.
\end{enumerate}

\subsection{Gestione del processo e costi}
\begin{itemize}
  \item \textbf{Preparazione}: piano di risposta agli incidenti, ruoli e responsabilità (incluso il
    CEO), simulazioni, modalità di interazione con le autorità. \textbf{Prevenzione}: formazione e
    misure di sicurezza.
  \item \textbf{Costi} di un data breach (slide): danno diretto, abbandono anomalo dei clienti, danno
    reputazionale, attività di help desk, indagini interne, spese legali, ripristino dei sistemi. Per il
    docente il più importante oggi è il \textbf{danno reputazionale}.
  \item \textbf{Aspetti pratici}: verificare quali notifiche sono obbligatorie; \textbf{analizzare con
    cura le ipotesi}, perché ogni assunzione può essere letta in due modi; usare i modelli e i canali
    ufficiali. Il modulo di notifica conta circa 55--60 domande: si è annunciata una simulazione di
    compilazione.
\end{itemize}

\subsection{Il responsabile del trattamento}
Il fornitore è oggi ``la base di tutti i problemi''. L'art.~28(3) rimanda al contratto; l'art.~33(2)
obbliga il responsabile a informare il titolare \textbf{senza ingiustificato ritardo} (nei contratti
si trovano tipicamente 24--36 ore). Nella pratica, soprattutto con i grandi fornitori, la
comunicazione arriva anche un mese dopo.

\begin{esame}
Se il responsabile avvisa con un mese di ritardo, il titolare non viola le 72 ore
\textbf{purché dimostri} che la violazione era imprevedibile e a lui ignota: l'orologio parte dalla
comunicazione del responsabile. Il titolare può poi chiedere i danni al responsabile per
inadempimento contrattuale. Tema ripreso nel Lab~2 (Capitolo~\ref{lab:02}).
\end{esame}

\subsection{Remediation plan (slide)}
Correggere la vulnerabilità che ha permesso l'incidente; \textbf{lesson learned}; condividere
informazioni con gli stakeholder per rafforzare il perimetro; se la causa è l'errore umano, rafforzare
la formazione. Sono esempi, non un elenco completo.

\subsection{Il registro dei data breach (art.~33(5))}
Strumenti: \textbf{incident log} (sintetico ma completo e aggiornato) e \textbf{report interni}
(modelli standard, decisioni documentate). L'art.~33(5) impone di documentare ogni violazione, le
circostanze, le conseguenze e i provvedimenti adottati, così che l'autorità possa verificare il
rispetto dell'articolo.

\begin{esame}
\textbf{Il registro è la prova dell'accountability.} Il controllo dell'autorità avviene su due livelli:
\begin{enumerate}
  \item \textbf{esistenza del registro}: se manca, la sanzione è automatica (e, se l'autorità trova i
    dati sul dark web, la massima possibile);
  \item \textbf{merito della valutazione}: se il rischio è stato dichiarato basso ma era evidentemente
    alto, si viene sanzionati comunque; la buona fede aiuta, ``me l'ha chiesto il CEO'' no.
\end{enumerate}
Il compito tipico dell'esperto di sicurezza sarà ridurre il rischio in modo che il titolare possa
legittimamente non notificare, ma sempre lasciando traccia del ragionamento.
\end{esame}

\paragraph{Firmare digitalmente il registro?} Su domanda di uno studente: un documento firmato
digitalmente (o con marca temporale, o inviato via PEC, o accompagnato da un hash) non può essere
modificato senza lasciare traccia e ``ferma l'orologio''. In Italia solo una minoranza di aziende
(5--10\% secondo il docente) lo fa; le altre tengono file Excel o PDF modificabili. Esistono anche
strumenti dedicati (es.\ OneTrust).

\paragraph{Alterare il registro è reato.} Le false dichiarazioni al Garante sono punite
dall'\textbf{art.~168 del Codice privacy} (reclusione da sei mesi a tre anni). Ne risponde non solo
il legale rappresentante ma anche chi ha materialmente alterato il documento: obbedire a un ordine
non scusa (\emph{ignorantia legis non excusat}). Se si deve correggere un file, meglio dirlo
all'autorità spiegando il motivo. Il docente racconta di un cliente che, dopo un'ispezione, gli
chiese di modificare il registro: rifiutò.

% ---------------------------------------------------------------------
\section{Presentazione del Lab~1: il caso Arvenza}
Il docente presenta il caso (dettagli e soluzione nel Capitolo~\ref{lab:01}): mailbox paghe
compromessa via phishing, sei gruppi di dati, tre posizioni in sala (CEO, IT Manager, HR Director).
Gli studenti fanno il \textbf{DPO}, che il docente descrive come funzione insieme legale e IT
(``nel mio sogno''), anche se in pratica è quasi sempre affidata a giuristi.
\begin{itemize}
  \item Formula ENISA: $SE = DPC \times EI + CB$.
  \item Va prodotto un \textbf{severity assessment} con un \textbf{parere} che concili o respinga le
    tre posizioni: ``non c'è una sola soluzione''.
  \item Non basta il numero: la mitigazione conta, e bisogna proporre una strategia (anche di
    prevenzione futura).
  \item Uso dell'IA: ammesso, ma se il lavoro è solo dell'IA sarà difficile presentarlo a voce.
\end{itemize}

% ---------------------------------------------------------------------
\section{Domande aperte e punti da verificare}
\begin{verifica}
\begin{itemize}
  \item \textbf{AI Act e autorizzazione preventiva}: la trascrizione fa dire al docente che per i
    sistemi ad alto rischio serve un'autorizzazione dell'autorità nazionale prima dell'uso. Il passaggio
    è confuso (si parla anche di FRIA) e va verificato sulle slide dell'AI Act.
  \item \textbf{Revisione semestrale della DPIA}: è una raccomandazione del docente, non un termine
    fissato dal GDPR (che chiede un riesame quando cambia il rischio, art.~35(11)). Confermare.
  \item Nella trascrizione il docente cita l'``art.~32.2'' per l'obbligo del responsabile di informare
    il titolare; le slide (e il testo del GDPR) indicano correttamente l'\textbf{art.~33(2)}.
  \item \textbf{Punti bonus}: il valore (0,25 a lab) e il numero di lab sono ricostruiti dal parlato,
    poco nitido; il totale massimo citato non è chiaro.
  \item Il ``paper del 2013'' sullo scambio di informazioni citato dal docente: autore e titolo non
    ricavabili con certezza dalla trascrizione.
  \item Concorso nel reato: il docente parla di ``associazione criminale'' tra legale rappresentante,
    IT manager ed esecutore; in termini tecnici potrebbe trattarsi di concorso di persone nel reato.
    Da non usare all'esame senza verifica.
\end{itemize}
\end{verifica}

% =====================================================================
\lezione{09/10/2026}{Laboratorio 1 -- Severity assessment del caso Arvenza}
  {Prof.~Giuseppe Vaciago; Gianluca Gilardi (supporto)}
  {organizzazione dei gruppi; svolgimento del Lab~1 (data breach della mailbox paghe, metodologia
   ENISA); chiarimenti sul caso durante il lavoro}
\label{lez:2026-10-09}
% =====================================================================

\begin{nota}[Fonte]
Trascrizione automatica della registrazione: la sessione è in gran parte silenziosa (gruppi al
lavoro). Sono riportate solo le indicazioni udibili del docente. Il caso, la metodologia e la
soluzione sono trattati nel Capitolo~\ref{lab:01}.
\end{nota}

\section{Sintesi}
Primo laboratorio, gestito su Zoom (il docente era collegato da remoto). Dopo la registrazione dei gruppi (alla fine circa 25), gli
studenti, nel ruolo di DPO di Arvenza Industrie, preparano in circa un'ora e mezza un
\textbf{severity assessment} con metodologia ENISA e un parere che concili le posizioni di CEO,
IT Manager e HR Director. Il docente e Gianluca rispondono alle domande.

\section{Organizzazione}
\begin{itemize}
  \item Ogni gruppo comunica in chat nome del gruppo e componenti.
  \item \textbf{Consegna}: un solo componente invia l'email al docente, con gli altri in copia; nel
    documento nomi e \textbf{numero di matricola} di tutti. Domande ammesse fino alle 12:45--12:50,
    consegna entro le \textbf{13:15}; oltre quell'ora il lavoro non è considerato.
  \item Il docente sperimenta una sala secondaria su Zoom per le domande, ma conclude che è meglio
    avere tutte le domande nella sala principale, perché le risposte servono a tutti.
  \item La valutazione arriva entro il weekend; i gruppi migliori presentano lunedì~12/10 (cinque
    minuti circa ciascuno, in tutto una quarantina di minuti).
\end{itemize}

\section{Indicazioni del docente durante il lab}
\begin{itemize}
  \item L'obiettivo è il severity assessment ``come presentato lunedì'', combinando le esigenze dei
    tre punti di vista. \textbf{Non c'è una regola fissa}: la scelta giusta è quella coerente con il
    principio di \emph{accountability}.
  \item Chiarimenti forniti a tutti sui dati del caso:
    \begin{itemize}
      \item i dati bancari (anche quelli dei tre dipendenti impersonati) \textbf{non erano cifrati};
      \item le buste paga 2023 vanno lette con attenzione perché la loro situazione è diversa da
        quella delle 2024--26 (le password di gennaio--aprile erano nella stessa mailbox);
      \item i documenti sanitari e i CV \textbf{non erano cifrati}.
    \end{itemize}
\end{itemize}

\begin{esame}
Il ragionamento conta più dell'opzione scelta: ogni punteggio deve poggiare su un fatto, e ogni
scostamento dal valore di base va motivato. Lo svolgimento completo è nel Capitolo~\ref{lab:01}.
\end{esame}

\section{Domande aperte e punti da verificare}
\begin{verifica}
\begin{itemize}
  \item Il testo del caso annuncia una \textbf{Parte~B} (aggiornamento dell'8 ottobre) dopo il
    debrief: non risulta né dalla registrazione né dal materiale nella cartella del lab.
  \item Orari esatti di chiusura delle domande e di consegna: ricostruiti dal parlato (12:45--12:50 e
    13:15).
\end{itemize}
\end{verifica}

% =====================================================================
\lezione{12/10/2026}{Debrief del Lab~1; introduzione a ISO/IEC~27001; presentazione del caso Sylvane}
  {Prof.~Giuseppe Vaciago; Gianluca Gilardi}
  {presentazioni dei gruppi sul caso Arvenza e commenti; triade CIA; esiti della gestione del rischio
   (mitigare, accettare, trasferire); ISMS e sue componenti; famiglia ISO 27k; presentazione del
   Lab~2 (Sylvane Retraite)}
\label{lez:2026-10-12}
% =====================================================================

\begin{nota}[Fonte]
Trascrizione automatica della registrazione, revisionata (pausa tra il minuto 66 e l'87 circa). La
parte delle presentazioni è in larga parte inintelligibile: si riportano solo i commenti dei docenti
e i contenuti ricorrenti. Slide di riferimento: \emph{04.~polito\_iso\_2026} (Fundamentals and
ISO/IEC Standards). Il caso del Lab~2 è nel Capitolo~\ref{lab:02}.
\end{nota}

\section{Sintesi}
La lezione ha tre parti. Nei primi 50 minuti alcuni gruppi presentano il proprio lavoro sul caso
Arvenza (Lab~1) e i docenti commentano. Segue la prima lezione di \textbf{Gianluca Gilardi}:
introduzione all'\textbf{ISO/IEC~27001} come sistema di gestione basato sul rischio, con triade CIA,
tre esiti della gestione del rischio e componenti dell'ISMS. Infine Gianluca presenta il caso
\textbf{Sylvane Retraite} per il laboratorio di venerdì~16/10.

% ---------------------------------------------------------------------
\section{Organizzazione}
\begin{itemize}
  \item \textbf{Esiti del Lab~1}: ogni gruppo ha ricevuto una valutazione via email. Soglie
    comunicate: sotto 20 nessun punto bonus; da 20 a 25 si ottiene il punto bonus; da 25 a 30 si può
    anche presentare il caso in aula. Chi ritiene di avere diritto al bonus ma non è in elenco deve
    segnalarlo al docente.
  \item \textbf{Ritardi}: tollerati i lavori arrivati entro le 14:00 e quelli inviati in tempo ma
    senza allegato (``siete esperti di sicurezza: ricordatevi l'allegato''); non accettati quelli
    arrivati il giorno dopo.
  \item \textbf{Regole formali per i prossimi lab}: nomi di tutti i componenti in prima pagina e
    \textbf{file nominato con il nome del gruppo}. Con oltre 90 studenti e più di 25 gruppi, un lavoro
    non attribuibile viene scartato.
  \item \textbf{Consegna del Lab~2}: caricamento sul repository SharePoint indicato nel caso
    (preferito); in alternativa email al docente o sezione \emph{Elaborati} del portale, ma è meglio un
    solo punto di raccolta.
  \item Su Zoom non si useranno più le sale separate: le domande si fanno nella sala principale (o in
    chat se brevi).
  \item Le slide ISO vengono caricate sul portale la sera stessa. Agli studenti senza gruppo viene
    chiesto di trovarne uno (anche online).
  \item Il materiale del caso Sylvane è online dalla mattina; si può studiare in anticipo, ma il lavoro
    si svolge venerdì.
\end{itemize}

% ---------------------------------------------------------------------
\section{Debrief del Lab~1 (Arvenza)}

\subsection{Che cosa è emerso dalle presentazioni}
Le presentazioni, da cinque--sette minuti, hanno mostrato approcci diversi, tutti basati sulla
formula ENISA:
\begin{itemize}
  \item \textbf{EI = 1} per tutti i gruppi di dati (nome, codice fiscale, IBAN identificano
    direttamente la persona);
  \item \textbf{intento malevolo +0,5} ovunque (l'attaccante voleva deviare gli stipendi);
  \item due \textbf{perimetri}: solo ciò che risulta letto dai log, oppure l'intera mailbox;
  \item più gruppi hanno sottolineato che \textbf{la severità non è il rischio}: alcuni dati hanno
    severità alta ma nessuna prova di accesso (rischio basso), altri severità media ma rischio molto
    alto (IBAN già usati per la frode);
  \item date di inizio delle 72 ore diverse: 20/9 (alert di sicurezza) oppure 28/9 (segnalazione del
    payroll bureau);
  \item decisioni variabili: quasi tutti hanno respinto la posizione del CEO e combinato le altre due.
\end{itemize}

\subsection{Commenti di Gianluca}
\begin{esame}
\begin{itemize}
  \item \textbf{Certificati di malattia}: molti li hanno valutati DPC~2 senza motivare. Gianluca è
    personalmente d'accordo, ma ai sensi dell'\textbf{art.~9 GDPR} ogni dato relativo alla salute ha
    la stessa natura e ricade nel \textbf{DPC~4} del modello ENISA. Abbassarlo è possibile solo
    \textbf{argomentando}.
  \item \textbf{Data di consapevolezza}: un ispettore del Garante contesterebbe quasi certamente una
    data di consapevolezza così tardiva. La violazione era forse attribuibile già all'inizio della
    cronologia (richieste di modifica dell'IBAN respinte e verificate con gli interessati). Chi
    sostiene una data tardiva deve essere pronto a dimostrare perché la conoscenza effettiva è
    arrivata dopo.
  \item In generale: si può sostenere quasi tutto, ma ogni affermazione e decisione deve avere
    \textbf{fondamento ed evidenze}. Nel caso successivo servirà un'analisi ancora più accurata.
\end{itemize}
\end{esame}

\paragraph{Commenti del docente.} Si è apprezzata la divisione dei ruoli nel gruppo; in una
presentazione devono parlare tutti i componenti, non uno solo. Quando si parla in pubblico si guarda
chi \emph{non} sta ascoltando, non il docente. Un gruppo ha riconosciuto di aver sovrastimato il CB
aggiungendo integrità e disponibilità ``per prudenza''.

% ---------------------------------------------------------------------
\section{ISO/IEC 27001: introduzione (Gianluca Gilardi)}

\begin{nota}[Il docente]
Gianluca è stato avvocato per 13 anni, poi \emph{digital forensics} per otto anni (spiegava la
tecnologia ad avvocati e giudici e i requisiti legali ai tecnici); da otto anni è CEO della società di
legal tech fondata dal prof.~Vaciago. Assiste aziende durante data breach e audit.
\end{nota}

\subsection{Perché ISO 27001}
\begin{itemize}
  \item La certificazione ISO/IEC 27001 significa che una \textbf{terza parte fidata} ha verificato
    che l'organizzazione ha misure e controlli appropriati per la sicurezza delle informazioni.
  \item \textbf{Non è richiesta da nessuna legge}: è uno schema volontario. Dal punto di vista del
    mercato, però, conoscerla è quasi un requisito per chi lavora in cybersecurity, anche senza
    certificarsi.
  \item È un sistema \textbf{basato sul rischio}: i controlli si adottano dove il rischio valutato lo
    richiede, per portarlo a un livello accettabile.
  \item Esistono alternative: negli USA i clienti chiedono più spesso il \textbf{SOC~2}; in Europa è
    più accettata la ISO. La differenza è di mercato, non tecnica.
  \item ``Sono troppo piccolo'': nel 2026 nessuna organizzazione è esente, perché tutte raccolgono,
    trattano, conservano e trasmettono informazioni che sono \textbf{asset} esposti a rischio.
    Battuta del settore: non chiederti \emph{se} subirai un incidente, ma \emph{quando}.
\end{itemize}

\subsection{I tre esiti della gestione del rischio}
Secondo Gianluca ogni processo di gestione del rischio si chiude con una di tre decisioni:
\begin{enumerate}
  \item \textbf{mitigare}: adottare controlli che riducano il rischio a un livello accettabile;
  \item \textbf{accettare} il rischio (mitigato o no);
  \item \textbf{trasferire} il rischio a un terzo (il ``Santo Graal'' del risk manager), come con
    l'assicurazione auto. Bisogna però poter dimostrare che la scelta del terzo è solida: non basta
    una firma su un foglio.
\end{enumerate}

\begin{nota}[Dalle slide successive del corso]
Le slide \emph{06.~27001\_2026} elencano anche una quarta opzione, \textbf{evitare} (cessare
l'attività che genera il rischio), che corrisponde allo ``stop del trattamento'' visto per la DPIA
(Capitolo~\ref{lez:2026-10-05}).
\end{nota}

\subsection{La triade CIA}
Dalle slide (definizioni ISO 27000): la sicurezza delle informazioni è la preservazione di
\textbf{riservatezza, integrità e disponibilità}; possono rilevare anche autenticità,
accountability, non ripudio e affidabilità.
\begin{itemize}
  \item \textbf{Riservatezza}: l'informazione non è resa disponibile né divulgata a persone, entità o
    processi non autorizzati. Non tutte le informazioni hanno la stessa sensibilità:
    \textbf{classificare} serve a non sprecare risorse su informazioni irrilevanti.
  \item \textbf{Integrità}: accuratezza e completezza.
  \item \textbf{Disponibilità}: accessibile e utilizzabile su richiesta da un'entità autorizzata.
\end{itemize}

\begin{esame}
\begin{itemize}
  \item Le \textbf{definizioni formali} di riservatezza, integrità e disponibilità sono molto
    importanti ed è probabile che vengano chieste all'esame.
  \item Anche una \textbf{perdita temporanea di disponibilità} di dati personali (poche ore) è un
    data breach. Il severity assessment potrà concludere che il rischio è così basso da non
    richiedere notifiche, ma la violazione va comunque \textbf{annotata nel registro}.
\end{itemize}
\end{esame}

\subsection{Che cos'è un ISMS}
Un \textbf{Information Security Management System} è un approccio sistematico per stabilire,
attuare, gestire, monitorare, riesaminare, mantenere e migliorare la sicurezza delle informazioni
di un'organizzazione. Un \emph{sistema}, non un patchwork: componenti tecniche, organizzative e
documentali devono stare in armonia (approccio olistico).

\paragraph{Monitoraggio e data governance.} Nel Lab~1 era difficile stabilire quando il titolare
fosse venuto a conoscenza della violazione: è il problema di chi non monitora. Il peccato più grave
nella gestione dei dati è non saper provare \emph{cosa}, \emph{dove}, \emph{quando}, \emph{chi} e
\emph{perché}: prima ancora della sicurezza viene la \textbf{data governance}.

\paragraph{Componenti dell'ISMS (slide).}
\begin{enumerate}[label=\alph*)]
  \item \textbf{policy}: la parte più facile da scrivere (spesso quasi un copia-incolla dei
    controlli), più difficile da mantenere aggiornata;
  \item \textbf{persone con responsabilità definite}: senza responsabilità formali (chi monitora, chi
    segnala, chi decide) non c'è un sistema di gestione ma un ``sistema della speranza'';
  \item \textbf{processi di gestione}: definizione delle policy, consapevolezza e competenze
    (gli auditor intervistano anche chi sta alla reception), pianificazione, attuazione, operatività,
    valutazione delle prestazioni, \textbf{riesame della direzione}, miglioramento;
  \item \textbf{informazioni documentate}: niente ``tradizione orale'';
  \item \textbf{valutazione del rischio} per la sicurezza delle informazioni;
  \item \textbf{trattamento del rischio}, con scelta e attuazione dei controlli.
\end{enumerate}
Le componenti a)--d) sono comuni a qualsiasi sistema di gestione; e) e f) sono specifiche
dell'ISMS.

\begin{itemize}
  \item \textbf{Ruolo della direzione}: le decisioni strategiche sulla sicurezza vengono dall'alto.
    Uno dei primi controlli in audit è come il management è coinvolto e informato.
  \item \textbf{Miglioramento continuo}: un ISMS nasce e cresce. Se l'auditor non vede miglioramenti
    solleva un rilievo, che può essere un'\emph{opportunità di miglioramento}, una \emph{non
    conformità minore} o una \emph{non conformità maggiore}.
\end{itemize}

\subsection{Valutare e trattare il rischio}
\begin{itemize}
  \item Rischio = severità (o impatto) $\times$ probabilità. Dalle slide: rischio come ``effetto
    dell'incertezza sugli obiettivi''; il rischio informativo è la possibilità che una minaccia sfrutti
    una vulnerabilità di un asset.
  \item Il risk manager prudente valuta sempre lo \textbf{scenario peggiore}, compresi i
    \textbf{cigni neri}: il COVID-19 è stato il più duro esercizio di gestione del rischio, e il
    mondo non era pronto perché la probabilità era stata giudicata troppo bassa.
  \item Il management chiederà di tagliare i costi: bisogna \textbf{prioritizzare} in base a
    criticità degli asset, impatto, probabilità e \textbf{rapporto costi/benefici} dei controlli.
  \item Esistono \emph{quick win}: i \textbf{backup}. Gianluca cita un cliente bloccato oltre due mesi
    da un ransomware perché non aveva backup dei dati cifrati.
  \item Alcuni controlli sono specifici per rischi ristretti, altri hanno portata più ampia: anche
    questo incide sul costo.
\end{itemize}

\subsection{Benefici e famiglia ISO 27k}
\begin{itemize}
  \item Ridurre il rischio per la sicurezza delle informazioni significa anche \textbf{ridurre il
    rischio di compliance}: è un investimento in conformità (le sanzioni sono rilevanti).
  \item La ISO 27001 è lo standard più usato per attuare e certificare un ISMS. La famiglia 27k conta
    \textbf{oltre 70 standard} (cloud, protezione dei dati personali, ecc.).
  \item Secondo Gianluca le altre certificazioni 27k sono ``figlie'' della 27001: senza la 27001 non
    si possono ottenere.
\end{itemize}

\begin{integrazione}[27k ``figlie'' della 27001]
Era vero per la ISO/IEC~27701:2019 (estensione della 27001). Dall'edizione 2025 la 27701 è uno standard
autonomo, certificabile anche senza 27001 (vedi Capitolo~\ref{cap:iso}).
\end{integrazione}

\begin{esame}
Il docente sottolinea che l'\textbf{ISMS è la base della ISO 27001}: è il primo concetto da fissare.
\end{esame}

% ---------------------------------------------------------------------
\section{Presentazione del Lab~2: Sylvane Retraite (Gianluca)}
Nuovo caso di data breach, più complesso sotto certi aspetti e più semplice sotto altri (dati meno
eterogenei, cronologia più breve). Dettagli e svolgimento nel Capitolo~\ref{lab:02}.
\begin{itemize}
  \item \textbf{Wall clock} 105 minuti (tempo reale del lab) e \textbf{exercise clock} tre giorni (il
    tempo della storia, nei tre giorni dopo la presa di conoscenza).
  \item La violazione arriva da una \textbf{comunicazione del responsabile del trattamento}: questo
    semplifica l'individuazione del momento in cui parte l'orologio per il titolare.
  \item Il responsabile ha già notificato l'incidente a un'autorità di settore \textbf{ai sensi del
    DORA}: diverse norme, diverse notifiche.
  \item Stavolta il \textbf{modulo di severity assessment è obbligatorio}; se si decide di notificare,
    si compila anche il file Excel del modulo di notifica (modello usato da molte autorità europee), e
    per gli interessati va scritto il testo effettivo della comunicazione (art.~34).
  \item Durante il lab arriveranno \textbf{curveball}, nuove informazioni che possono cambiare la
    valutazione: ``la vita è ciò che accade mentre fai l'esercizio''.
  \item \textbf{Honne e tatemae}: il modulo chiede sia la giustificazione da dare all'autorità
    (\emph{tatemae}, la facciata) sia il vero ragionamento interno (\emph{honne}, la verità, anche
    ``temo di essere licenziato'').
\end{itemize}

% ---------------------------------------------------------------------
\section{Domande aperte e punti da verificare}
\begin{verifica}
\begin{itemize}
  \item \textbf{Cognome di Gianluca}: ``Gilardi'' (dalla trascrizione del 09/10 e dall'indirizzo del
    repository SharePoint del Lab~2). La trascrizione oscilla tra ``Gilardi'' e ``Gillardi''.
  \item Soglie del bonus (20 / 25 / 30): dal parlato del docente; confermare con l'email ricevuta.
  \item ``Senza 27001 non si ottengono le altre certificazioni 27k'': affermazione di Gianluca, vera
    per le estensioni come la 27701; da verificare per gli altri standard della famiglia.
  \item Gianluca indica come \emph{honne} il ``paragrafo 10'': nel modulo il paragrafo~9 è la
    giustificazione (tatemae) e il 10 la strategia interna (honne). Coerente, ma verificare.
\end{itemize}
\end{verifica}

% =====================================================================
\lezione{16/10/2026}{Laboratorio 2 -- Severity assessment e notifica del caso Sylvane Retraite}
  {Gianluca Gilardi}
  {svolgimento del Lab~2 (breach presso il fornitore del canale SFTP di un fondo pensione francese);
   worksheet di severity assessment; modulo di notifica alla CNIL; curveball}
\label{lez:2026-10-16}
% =====================================================================

\begin{nota}[Fonte]
Trascrizione automatica della registrazione: quasi tutta silenzio o rumore (all'inizio è stato
captato l'audio di un altro evento, non pertinente). Si riportano solo le indicazioni udibili di
Gianluca. Caso, metodologia e svolgimento ragionato nel Capitolo~\ref{lab:02}.
\end{nota}

\section{Sintesi}
Secondo laboratorio, condotto da Gianluca. I gruppi, nel ruolo di DPO del fondo pensione francese
Sylvane Retraite, compilano il \textbf{worksheet di severity assessment} (obbligatorio) e, se
concludono che la notifica è dovuta, il \textbf{modulo di notifica alla CNIL}. Durante il lavoro
arrivano nuove informazioni (\emph{curveball}) da integrare nella valutazione.

\section{Organizzazione}
\begin{itemize}
  \item Consegna sul repository SharePoint indicato nel caso: il solo worksheet, oppure un unico zip
    con worksheet e modulo di notifica se si notifica. Nomi del gruppo e dei componenti nel worksheet
    e nel modulo di caricamento; file nominati con il nome del gruppo.
  \item Gli ultimi cinque minuti sono stati lasciati per le domande; consegna entro le
    \textbf{13:30}.
  \item Lunedì~19/10: breve \textbf{debrief} sui lavori, poi si prosegue con la ISO.
\end{itemize}

\section{Indicazioni di Gianluca durante il lab}
\begin{esame}
Gianluca chiede una valutazione \textbf{complessiva} della violazione (severity e rischio
complessivi), non un punteggio per ogni singola voce dell'elenco dei dati: si valutano i singoli
elementi ``nella propria testa'' e si arriva a un \textbf{punteggio globale}, perché è questo che
chiede l'autorità di protezione dei dati.
\end{esame}
\begin{itemize}
  \item Le nuove informazioni sono state mostrate a schermo; l'ultima è stata annunciata come
    ``l'ultimo ostacolo'', dopo il quale i gruppi avevano tutti gli elementi disponibili.
\end{itemize}

\section{Domande aperte e punti da verificare}
\begin{verifica}
\begin{itemize}
  \item Nella cartella del lab è conservata solo la \textbf{curveball~1} (stessa comunicazione inviata
    a una società sorella che fa wealth management). Le altre informazioni mostrate a schermo, tra cui
    ``l'ultimo ostacolo'', non sono nel materiale né udibili: da recuperare (chat di Zoom o compagni).
  \item Contenuto del debrief del 19/10: da riportare nella voce di quella lezione.
\end{itemize}
\end{verifica}

% =====================================================================
\lezione{19/10/2026}{ISO/IEC 27001: famiglia 27k, clausole 4--10 e primi controlli tecnologici}
  {Gianluca Gilardi; Prof.~Giuseppe Vaciago}
  {famiglia ISO 27k; perché la 27001 è diffusa; ciclo di implementazione dell'ISMS; clausole 4--10;
   ciclo di certificazione triennale; 93 controlli dell'Allegato A; controlli 8.1--8.8;
   restituzione del laboratorio 2 (Sylvane Retraite); presentazioni dei gruppi}
\label{lez:2026-10-19}
% =====================================================================

\begin{nota}[Fonte]
Trascrizione automatica della registrazione, revisionata (nei primi 14 minuti non c'è lezione; pausa
intorno ai minuti 70--76). Numeri, clausole e controlli sono ripresi dalle slide
\emph{04.~polito\_iso\_2026} (\emph{Fundamentals and ISO/IEC Standards}). Le parti sul laboratorio~2
fanno riferimento al caso \emph{Sylvane Retraite} (cartella \texttt{lab/lab\_02}).
\end{nota}

\section{Sintesi}
Lezione tenuta in gran parte da Gianluca per preparare il laboratorio sul \emph{risk register}. Dopo una
panoramica della famiglia ISO/IEC~27k e delle ragioni per cui la 27001 è lo standard di riferimento in
Europa, si ripercorrono le fasi di implementazione di un ISMS e le clausole 4--10, con molti esempi
dalla pratica professionale (audit, costi, ciclo triennale). Nella parte centrale si restituisce il
laboratorio~2 sulla severità del data breach e un gruppo presenta il proprio lavoro. Si chiude con i
primi otto controlli tecnologici dell'Allegato~A (8.1--8.8). Messaggio di fondo: l'ISMS è un
\textbf{processo basato sul rischio}, non un timbro; contano i processi più dei documenti.

% ---------------------------------------------------------------------
\section{Organizzazione}
\begin{itemize}
  \item \textbf{Esercitazione sul risk register}: Gianluca annuncia un laboratorio in cui compilare un
    registro dei rischi secondo la ISO~27001 per una società di consulenza; i rischi saranno già forniti,
    agli studenti spetterà valutarli e scegliere i controlli tecnologici. (Il laboratorio si è poi tenuto
    venerdì 30/10: vedi Capitolo~\ref{lab:03}.)
  \item \textbf{Norme ISO non distribuite}: 27001 e 27002 sono documenti a pagamento e non possono essere
    caricati sulla piattaforma per ragioni di copyright.
  \item \textbf{Esiti del laboratorio~2}: pubblicati sulla piattaforma la sera precedente, con avviso in
    bacheca.
  \item \textbf{Presentazioni dei laboratori}: i gruppi invitati (quelli con il punteggio più alto)
    devono essere presenti in aula; il professore ha escluso dal punto aggiuntivo i gruppi invitati e
    assenti senza preavviso. Chi non può presentare deve avvisare per email: per il docente è una
    questione di rispetto, non di regolamento.
  \item \textbf{Calendario}: venerdì 23/10 lezione \textbf{online} (Gianluca prosegue con i controlli
    ISO); lunedì 26/10 lezione del professore in presenza su un nuovo argomento; il venerdì successivo
    (30/10) il laboratorio.
\end{itemize}

% ---------------------------------------------------------------------
\section{La famiglia ISO/IEC 27k}

La famiglia conta oltre 70 standard e continua a crescere (Gianluca segnala la recentissima ISO/IEC~27090
sulla sicurezza dei sistemi di IA, approvata pochi giorni prima e quindi assente dalle slide). Quelli
rilevanti in generale sono:
\begin{center}
\small
\begin{tabularx}{\linewidth}{>{\raggedright\arraybackslash}p{2.4cm} X}
\toprule
\textbf{Standard} & \textbf{Ruolo} \\ \midrule
ISO/IEC 27000 & panoramica e vocabolario dell'ISMS (liberamente disponibile) \\
ISO/IEC 27001 & \textbf{requisiti} dell'ISMS: principi di alto livello ed elenco dei controlli; è lo
  standard certificabile \\
ISO/IEC 27002 & guida ai controlli: come attuarli e quali effetti devono produrre \\
ISO/IEC 27003 & guida interpretativa dei requisiti (``should'', ``can'', ``may''): utile per capire
  cosa conta per l'auditor \\
ISO/IEC 27005 & guida alla gestione del rischio di sicurezza delle informazioni \\
ISO/IEC 27701 & estensione privacy (PIMS) di 27001 e 27002 \\ \bottomrule
\end{tabularx}
\end{center}

\begin{itemize}
  \item Gli standard ISO sono \textbf{neutrali rispetto ai fornitori}: indicano l'effetto da ottenere,
    non quali prodotti usare.
  \item Anche la ISO segue un \textbf{approccio basato sul rischio}, come GDPR, DORA e NIS2. Nessuna
    metodologia di risk assessment è imposta: tassonomia dei rischi, scale di probabilità e impatto e
    scelta dei controlli spettano all'organizzazione.
  \item La \textbf{27701} è un'estensione: per ottenerla serve prima la certificazione 27001. Nella
    pratica di Gianluca nessun cliente l'ha mai chiesta; ha senso solo per trattamenti molto sensibili su
    larga scala (es.\ dati sanitari). (Integrazione: vale per l'edizione 2019; con la 27701:2025 lo
    standard è diventato autonomo, vedi Capitolo~\ref{cap:iso}.)
  \item Molte aziende non si certificano ma dichiarano di \textbf{seguire} i principi della 27001, per
    evitare il peso delle verifiche e dei rinnovi.
\end{itemize}

\subsection{Perché la 27001 è così diffusa}
Secondo le slide: è applicabile a qualsiasi organizzazione; è breve (corpo principale di circa 25
pagine); è allineata ad altri sistemi di gestione; è collaudata (deriva dalla BS~7799-1 del 1995); ha
principi semplici ma utili; è certificabile; molti professionisti la usano e quindi ``parlano la stessa
lingua''. Gianluca aggiunge che circa il 75--80\% dei controlli 27001 si ritrova nel framework
statunitense SOC~2, per cui il lavoro fatto è in buona parte riutilizzabile.

\paragraph{Benefici.} Sicurezza complessiva, consapevolezza, riduzione dei rischi di business, fiducia
degli stakeholder, conformità legale e contrattuale, maturità dei processi, trasparenza del budget di
sicurezza e opportunità di marketing: nelle gare la certificazione può decidere l'aggiudicazione. Davanti
a un'autorità dopo un incidente, essere certificati aiuta a dimostrare di aver fatto il possibile, ma
\textbf{non garantisce} che il sistema sia sicuro.

\paragraph{Diffusione (ISO Survey 2024).} Al 31/12/2024 risultano 96.709 certificati ISO/IEC~27001 validi
(contro circa 1,47 milioni della ISO~9001 sulla qualità). Nella classifica per Paese guida la Cina
(33.359); l'Italia è sesta con 3.284. Per il mercato del lavoro, conoscere la 27001 (o essere
\emph{lead auditor}) è un titolo molto apprezzato; un corso completo da lead auditor dura circa 40 ore.

% ---------------------------------------------------------------------
\section{Implementare un ISMS}

\paragraph{Le fasi.} Il ciclo (ISO~27003) è: \textbf{contesto} (bisogni dell'organizzazione e perimetro)
$\rightarrow$ \textbf{rischi} $\rightarrow$ \textbf{controlli} $\rightarrow$ \textbf{monitoraggio}
$\rightarrow$ \textbf{miglioramento}. Il perimetro può coprire tutta l'azienda o un solo prodotto o
servizio. Sicurezza significa governare \textbf{persone, processi e tecnologia}: la tecnologia è la parte
più semplice.

\paragraph{Propensione al rischio (\emph{risk appetite}).} È il livello di rischio che l'azienda è
disposta ad accettare senza ulteriori interventi. Esempio: su una scala 1--9, si accetta tutto ciò che sta
sotto 3; il resto va mitigato, trasferito o evitato. Nell'esperienza di Gianluca, però, nessun rischio
\emph{inerente} che compare nel registro è davvero accettabile così com'è: i rischi davvero minimi non
vengono nemmeno percepiti e quindi non vengono mappati.

\paragraph{Controlli efficaci.} Un controllo inefficace è il caso peggiore: si spendono soldi e tempo e il
rischio si materializza comunque. Il monitoraggio dell'efficacia è quindi centrale; se i controlli
falliscono servono \emph{gap analysis} e piano di miglioramento (aggiornare, trasferire il rischio,
oppure abbandonare il prodotto o il processo).

\subsection{Le clausole 4--10 (attività obbligatorie)}
Le clausole 4--10 non possono essere escluse da chi dichiara la conformità alla norma.
\begin{center}
\small
\begin{tabularx}{\linewidth}{>{\raggedright\arraybackslash}p{3.2cm} X}
\toprule
\textbf{Clausola} & \textbf{Contenuto e osservazioni in aula} \\ \midrule
4 Contesto & fattori interni ed esterni, parti interessate, perimetro dell'ISMS (dettaglio non
  richiesto all'esame) \\
5 Leadership & impegno dell'alta direzione, politica, ruoli. Punto spesso critico: se il vertice non
  conosce rischi e controlli non può decidere budget e priorità; ``tutto parte dall'alto'' \\
6 Pianificazione & processo di valutazione e di trattamento del rischio (6.1), obiettivi (6.2),
  pianificazione delle modifiche (6.3) \\
7 Supporto & risorse, competenze, \textbf{consapevolezza} (molto verificata dagli auditor e richiesta
  anche da GDPR e AI Act), comunicazione, informazioni documentate \\
8 Attività operative & pianificazione e controllo; esecuzione del risk assessment (8.2) e del piano di
  trattamento (8.3), con evidenze documentate \\
9 Valutazione delle prestazioni & monitoraggio e misurazione (KPI: ``non si controlla ciò che non si
  misura''), audit interno, riesame della direzione \\
10 Miglioramento & miglioramento continuo, non conformità e azioni correttive \\ \bottomrule
\end{tabularx}
\end{center}

\paragraph{Valutazione del rischio.} La 27001 richiede un risk assessment \textbf{documentato}. Nella
convenzione usata nel corso (e nel Lab~3) l'impatto si valuta su riservatezza, integrità e disponibilità
e la probabilità è unica; la norma in sé non impone un metodo. Dopo il trattamento si
valuta quanto il controllo riduce il rischio e si calcola il \textbf{rischio residuo}, che a sua volta va
accettato, rifiutato, trasferito o ulteriormente mitigato.

\paragraph{Audit interno e riesame.} L'audit interno non deve essere un timbro: in un audit recente
l'auditor ha contestato all'azienda l'assenza di un \textbf{programma pluriennale} di audit interni con
perimetro definito in anticipo. La direzione deve riesaminare l'ISMS almeno una volta l'anno (l'ideale è
ogni trimestre), lasciando evidenza del proprio coinvolgimento.

\paragraph{Ciclo di certificazione.} La certificazione segue un ciclo \textbf{triennale}: audit di
certificazione il primo anno, due audit di mantenimento negli anni successivi (solo su parti della
norma), poi ricertificazione completa. Più l'organizzazione è certificata da tempo, più gli auditor sono
severi: il primo anno controllano soprattutto i documenti, poi testano i processi e intervistano il
personale.

\paragraph{Costi.} La consulenza per arrivare alla certificazione costa da poche migliaia a decine di
migliaia di euro; il costo vero è far funzionare davvero i processi e controllare i controlli. Dal punto
di vista dell'azienda la compliance è un costo che non produce valore diretto, ma protegge clienti e
pubblico. Implementare un ISMS da zero richiede in genere \textbf{un anno -- un anno e mezzo} per
un'organizzazione di medie dimensioni.

\begin{nota}[Cosa sapere dell'ISMS (slide)]
Requisiti generici e flessibili; attenzione alle parti interessate (anche fornitori, consulenti e, per le
informazioni pubbliche, il pubblico); approccio basato sul rischio; processi e miglioramento continuo;
\textbf{i processi contano più delle policy}; non puntare subito alla certificazione, ma prima
implementare i processi; l'implementazione è un progetto lungo. L'auditor chiede anche le
\textbf{dipendenze} tra processi e unità organizzative.
\end{nota}

% ---------------------------------------------------------------------
\section{Allegato A: i 93 controlli}
Controllo = misura che mantiene o modifica il rischio. La revisione 2022 (che sostituisce quella del
2013) conta \textbf{93 controlli}, di cui 11 nuovi, in quattro temi:
\textbf{organizzativi (37)}, \textbf{persone (8)}, \textbf{fisici (14)}, \textbf{tecnologici (34)}.
I controlli fisici sono spesso trascurati (armadi chiusi a chiave, scrivanie sgombre). La ISO~27002:2022
assegna a ogni controllo \textbf{cinque attributi}: tipo (preventivo, detective, correttivo), proprietà
CIA, concetti di cybersicurezza (\emph{identify, protect, detect, respond, recover}), capacità operative,
domini di sicurezza.

Non tutti i controlli valgono per ogni organizzazione: una piccola società di consulenza che non sviluppa
software non ha codice sorgente da proteggere né ambienti separati di sviluppo e produzione (8.31). Lo
standard evolve: un controllo oggi sufficiente può non bastare alla prossima revisione.

\begin{esame}
\begin{itemize}
  \item Numeri da ricordare: \textbf{93 controlli} in 4 temi (37 / 8 / 14 / 34); clausole \textbf{4--10}
    obbligatorie; ciclo di certificazione \textbf{triennale}.
  \item La certificazione 27001 riguarda l'\textbf{organizzazione} e non garantisce l'assenza di
    incidenti.
  \item Il dettaglio della clausola~4 (contesto) è stato indicato come \emph{non} strettamente richiesto
    all'esame; il focus del corso sono i controlli tecnologici.
\end{itemize}
\end{esame}

% ---------------------------------------------------------------------
\section{Controlli tecnologici 8.1--8.8}
\begin{description}
  \item[8.1 Dispositivi endpoint degli utenti.] Antimalware, aggiornamenti, MDM, cifratura del disco. Il
    livello di protezione segue la classificazione dei dati trattati: il portatile del CEO va protetto più
    di quello della reception (nella realtà sono spesso i dirigenti a chiedere meno vincoli).
  \item[8.2 Diritti di accesso privilegiato.] Principio del \textbf{minimo privilegio}. Problema
    riscontrato ``otto volte su dieci'': i diritti di amministratore locale concessi a tutti per comodità.
    Esempio: un account pubblicitario aziendale compromesso, con decine di migliaia di euro di campagne
    addebitate, partito da un PC dell'agenzia esterna su cui era stato installato software non autorizzato.
  \item[8.3 Restrizione dell'accesso alle informazioni.] Mentre l'8.2 riguarda la configurazione
    dell'utente, l'8.3 riguarda l'informazione: nessuna discrepanza tra classificazione del dato e diritti
    concessi.
  \item[8.4 Accesso al codice sorgente.] Caso di un'azienda ad-tech con sviluppatori esterni (anche extra
    UE) tutti con accesso completo al repository: servono tracciamento degli accessi, accessi limitati ai
    moduli di competenza, eventualmente repository separati.
  \item[8.5 Autenticazione sicura.] Per Gianluca l'MFA va imposta ovunque, anche con fattori biometrici.
    Un attacco \emph{brute force} riuscito nel 2026 rivela password deboli e nessun limite ai tentativi.
    I segreti vanno protetti per tutto il ciclo di vita: mai inviare credenziali in chiaro sull'email
    personale del neoassunto senza cambio password obbligatorio al primo accesso.
  \item[8.6 Gestione della capacità.] Riguarda soprattutto la \textbf{disponibilità}: con servizi cloud la
    capacità è demandata al fornitore; con infrastruttura propria va monitorata per pianificare la
    scalabilità.
  \item[8.7 Protezione dal malware.] Non solo antivirus: protezione a strati (IDS/IPS, rete, collegamento
    verso l'ISP). ``Non chiederti \emph{se} sarai violato, ma \emph{quando} e come ripartirai'': un'azienda
    industriale senza protezione né backup è rimasta ferma un mese e mezzo e ha poi pagato il riscatto.
  \item[8.8 Gestione delle vulnerabilità tecniche.] Inventario degli asset, scansione, prioritizzazione
    in base al rischio, processo di patching. Casi reali: server Windows 2012 mai aggiornati e di cui
    nessuno conosceva l'esistenza; attacco partito da un server interno ignoto al responsabile IT
    (``risposta che nessuna autorità accetterà''); molti incidenti causati da firewall perimetrali non
    aggiornati.
\end{description}

% ---------------------------------------------------------------------
\section{Restituzione del laboratorio 2 (Sylvane Retraite)}

\paragraph{Esito atteso.} L'incidente era un data breach da notificare \textbf{all'autorità di controllo e
agli interessati}. Errore frequente: ``non ci sono prove di esfiltrazione, quindi il rischio è basso e non
comunichiamo agli interessati''.
\begin{itemize}
  \item La \textbf{perdita di riservatezza} si ha con il semplice \textbf{accesso} non autorizzato ai dati,
    anche senza esfiltrazione. Se l'attaccante ha avuto accesso al server, compresa la copia di backup
    ``nascosta'' (il secondo colpo di scena del caso), tutti quei dati sono coinvolti e la severità sale.
  \item Dati finanziari dettagliati sono ideali per campagne di \emph{phishing} mirate: il rischio per gli
    interessati è alto.
  \item Sulla \textbf{facilità di identificazione}: per il primo insieme (nome, cognome, valore del fondo
    pensione) si poteva discutere; con il backup (codice fiscale, indirizzo, email, dati bancari) quella
    difesa non regge.
  \item Il professore precisa che, in pratica, l'impatto per l'interessato è molto diverso se i dati sono
    stati solo visti o anche portati via; ma dal punto di vista formale la violazione di riservatezza c'è
    in entrambi i casi.
\end{itemize}

\paragraph{La presentazione di un gruppo.} Divisione dei dati in due insiemi (dati del canale di scambio e
backup), punteggi calcolati separatamente e insieme, entrambi sopra la soglia. Decisioni: notifica alla
CNIL entro 72 ore senza attendere il quadro completo (i dettagli si integrano dopo); notifica anche del
backup; lettera agli aderenti (``non vi chiederemo mai password o denaro''); segnalazione all'ACPR come
grave incidente DORA; orologio GDPR fatto partire dalla comunicazione del fornitore, purché documentata;
comunicazione coerente con la società sorella del gruppo; ogni passaggio documentato e firmato. Gianluca
osserva che, se il vettore d'attacco è la compromissione dell'infrastruttura del responsabile del
trattamento, la responsabilità di quest'ultimo resta anche se i dati stavano sul server del titolare.

\paragraph{Il punteggio ``esattamente 3''.} Le fasce di rischio della sezione~6 del worksheet del Lab~2 (costruite sulla metodologia ENISA)
sono una \textbf{guida, non una regola matematica}: le righe che delimitano il valore 3 si sovrappongono.
Con un punteggio di confine la decisione spetta a chi valuta, che se ne assume la responsabilità in base al
principio di accountability; l'autorità potrà comunque arrivare a conclusioni diverse.

\begin{esame}
\begin{itemize}
  \item \textbf{Violazione di riservatezza = accesso non autorizzato}, anche senza esfiltrazione né
    lettura provata.
  \item \textbf{Non presumere nulla che non si possa provare}: le dichiarazioni false all'autorità di
    controllo sono un reato, di cui risponde personalmente chi le rende. Nel dubbio, attenersi alle
    evidenze.
  \item Il reato richiede il dolo (\emph{mens rea}), che va dimostrato; ma un pubblico ministero poco
    esperto di aspetti tecnici può partire dall'ipotesi opposta, e l'onere di spiegare ricade su chi ha
    firmato. L'interpretazione delle evidenze va curata con estrema attenzione.
\end{itemize}
\end{esame}

\noindent Il professore chiede infine agli studenti un riscontro sul livello di difficoltà dei
laboratori, percepito come alto.

% ---------------------------------------------------------------------
\section{Domande aperte e punti da verificare}
\begin{verifica}
\begin{itemize}
  \item \textbf{ISO/IEC 27090} (sicurezza dei sistemi di IA): citata a voce come approvata pochi giorni
    prima della lezione; non è nelle slide. Verificare numero e stato di pubblicazione.
  \item \textbf{SOC~2}: la sovrapposizione del 75--80\% con i controlli 27001 è una stima del docente.
  \item \textbf{Origine della 27001}: in aula ``prima versione del 1995''; le slide precisano che nel 1995
    fu pubblicata la BS~7799-1, da cui la norma deriva.
  \item \textbf{Fasce ENISA}: nel worksheet del Lab~2 (sezione~6) le fasce si sovrappongono sul valore 3; nella
    metodologia ENISA (sezione~3.2) sono disgiunte: $2 \le SE < 3$ media, $3 \le SE < 4$ alta. Vedere la
    metodologia (slide \emph{03}) e il Capitolo~\ref{cap:data-breach}.
  \item Il passaggio sulle presentazioni dei gruppi è in parte confuso nella trascrizione: le regole esatte
    (soglia di punteggio per presentare, punto aggiuntivo) vanno confermate dagli avvisi in piattaforma.
\end{itemize}
\end{verifica}

% =====================================================================
\lezione{23/10/2026}{Controlli tecnologici 8.8--8.34 e introduzione al risk management ISO}
  {Gianluca Gilardi (lezione online)}
  {gestione delle vulnerabilità; configurazione; cancellazione sicura; data masking; DLP; backup;
   ridondanza; logging e limite dei 21 giorni; monitoraggio; rete e zero trust; crittografia;
   sviluppo sicuro; change management; dati di test; risk management ISO vs GDPR; asset, minacce,
   vulnerabilità}
\label{lez:2026-10-23}
% =====================================================================

\begin{nota}[Fonte]
Trascrizione automatica della registrazione (lezione online), revisionata. Numeri e titoli dei controlli
dalle slide \emph{04.~polito\_iso\_2026} (slide 39--49); la parte sul rischio dalle slide
\emph{06.~27001\_2026\_2310} (\emph{Risk Management}). Lezione chiusa in anticipo dal docente.
\end{nota}

\section{Sintesi}
Gianluca completa la rassegna dei 34 controlli tecnologici dell'Allegato~A, insistendo sui punti in cui
la sicurezza entra in tensione con il GDPR e il diritto del lavoro (DLP, logging, monitoraggio). Molti
controlli sono ``buon senso'', ma proprio per questo vengono spesso trascurati. Nell'ultima parte avvia il
\textbf{risk management secondo la ISO}: qui impatto e probabilità si misurano sull'\textbf{organizzazione},
non sui diritti e le libertà degli interessati come nella valutazione di severità di un breach.

% ---------------------------------------------------------------------
\section{Organizzazione}
\begin{itemize}
  \item Le slide sul risk management vengono caricate in piattaforma subito dopo la lezione; domande
    sul tema si possono porre lunedì 26/10.
  \item Gianluca chiede al professore un'ora del lunedì per completare il risk management prima del
    laboratorio del venerdì successivo (30/10).
\end{itemize}

% ---------------------------------------------------------------------
\section{Controlli tecnologici 8.8--8.34}

\begin{nota}[Attualità: perché l'8.8 serve davvero]
Nelle 48 ore precedenti la lezione, alcune vulnerabilità molto critiche in librerie PHP hanno costretto
mezzo mondo a spegnere sistemi in attesa delle patch. Avere un inventario degli asset, sapere quali
sistemi sono esposti e un processo di patching già definito è ciò che ha permesso di gestire la crisi.
\end{nota}

\begin{description}
  \item[8.9 Gestione della configurazione.] Si definisce una \emph{baseline} sicura; ogni modifica
    rilevante segue un processo con valutazione dei rischi, piano di rollback e benefici attesi. In ogni
    momento si deve poter dire quanto la configurazione corrente si discosta dalla baseline.
  \item[8.10 Cancellazione delle informazioni.] Cancellazione sicura a fine ciclo di vita, inclusi i
    backup (serve sapere quante copie esistono e dove). Con i grandi fornitori cloud non si ottengono
    evidenze di cancellazione di un singolo bucket: si deve almeno documentare come si è gestito il
    problema, perché l'auditor lo chiederà.
  \item[8.11 Data masking.] Nasce per la conformità GDPR ma vale per ogni informazione riservata;
    soprattutto negli ambienti non di produzione. Alternativa diffusa: i dati sintetici.
  \item[8.12 Data leakage prevention.] Controllo introdotto nell'ultima revisione: rilevare e bloccare
    l'esfiltrazione via email e altri canali significa monitorare le comunicazioni, con tensione rispetto
    al GDPR e all'art.~4 dello Statuto dei lavoratori. Chi esclude 8.8 o 8.12 dalla dichiarazione di
    applicabilità deve avere una motivazione documentata molto solida.
  \item[8.13 Backup.] Un backup vale quanto la possibilità di ripristinarlo: serve un piano di
    \textbf{test di ripristino} periodico (backup ``vuoti'' per uno script fallito equivalgono a nessun
    backup).
  \item[8.14 Ridondanza.] Disaster recovery, eventualmente con sito \emph{hot}; equilibrio tra requisiti
    di sicurezza, operativi e budget. Tutela la disponibilità.
  \item[8.15 Logging.] La ``traccia documentale'' di cosa è successo ai dati e ai sistemi (vedi sotto il
    nodo dei 21 giorni).
  \item[8.16 Monitoraggio.] Guardare i log, in tempo reale o quasi (SIEM, SOC), e gestire gli alert.
    Caso reale: un alert inviato il venerdì sera non è stato letto da nessuno e nel weekend l'attaccante
    ha esfiltrato terabyte di dati sensibili. Raccogliere senza presidiare non serve.
  \item[8.17 Sincronizzazione degli orologi.] Senza orologi allineati (fusi, UTC, derive) non si ricostruisce
    una timeline affidabile degli eventi.
  \item[8.18 Utility privilegiate] e \textbf{8.19 Installazione di software}: niente software che richiede
    privilegi elevati senza intervento della sicurezza; installazione limitata a software vagliato
    (eventuale whitelist tramite MDM).
  \item[8.20--8.22 Rete.] Segregazione (VLAN, mai reti piatte), attenzione all'IoT (telecamere IP, DVR
    esposti in rete e accessibili da terzi), paradigma \textbf{zero trust}: anche un dispositivo collegato
    via cavo nei locali aziendali va autenticato. Il Wi-Fi è di per sé un rischio. La segregazione limita
    il \emph{blast radius}: va separato soprattutto l'ambiente industriale (SCADA) dal back office, perché
    l'attacco arriva di solito dal PC che controlla le macchine.
  \item[8.24 Crittografia.] Cifrare tutto, ma la crittografia è forte quanto il controllo sulle chiavi:
    delegare la gestione delle chiavi al cloud provider è comodo, ma se il fornitore le perde o blocca
    l'accesso i dati diventano inaccessibili. Protegge riservatezza e integrità (hash).
  \item[8.25--8.30 Sviluppo.] Ciclo di vita sicuro (rilevante solo per chi sviluppa, o in settori con
    lunghe pipeline come il farmaceutico); requisiti di sicurezza fin dalla progettazione, anche quando lo
    sviluppo è esternalizzato; principi di architettura sicura (caso di un sistema di sviluppo ``ombra''
    fuori dal perimetro, compromesso e con credenziali in chiaro); \emph{secure coding} (es.\ OWASP); test di
    sicurezza prima del rilascio (VA/PT, oltre ai test funzionali); supervisione del fornitore esterno.
  \item[8.31 Separazione di sviluppo, test e produzione.] Caso di una compagnia assicurativa che usava
    dati di produzione nei test e talvolta testava direttamente in produzione.
  \item[8.32 Change management.] Ogni cambiamento \textbf{rilevante} va valutato sotto il profilo della
    sicurezza delle informazioni, documentato, possibilmente testato, con piano di rollback. Non solo IT:
    anche un trasloco di sede o una riorganizzazione aziendale.
  \item[8.33 Dati di test.] Preferire dati sintetici, ma attenzione: conservano le proprietà statistiche
    dei dati di origine. L'uso di dati di produzione è ammesso solo in casi limite documentati (es.\ un
    ticket di assistenza riproducibile solo con quei dati), cancellandoli subito dopo.
  \item[8.34 Protezione durante audit e test.] Esempio: un cliente che commissiona un penetration test
    sull'istanza SaaS in produzione concorda prima il carico sostenibile, per non compromettere la
    disponibilità.
\end{description}

\begin{nota}[Buon senso, ma obbligatorio]
Molti controlli sembrano ovvi (backup, monitoraggio\ldots). Sono requisiti della ISO, ma anche buone
pratiche che chiunque dovrebbe adottare, certificato o no. Il buon senso, osserva Gianluca, non è
distribuito in modo uniforme.
\end{nota}

\subsection{Logging e privacy: il limite dei 21 giorni}
Secondo un documento del Garante privacy citato in aula, i \textbf{metadati della posta elettronica} dei
dipendenti, generati dai sistemi di logging, possono essere conservati per non più di \textbf{21 giorni},
salvo accordo sindacale o autorizzazione. Il problema operativo:
\begin{itemize}
  \item molti servizi (es.\ Microsoft~365, Google Workspace) hanno retention non personalizzabili e più
    lunghe;
  \item per la sicurezza servono log molto più lunghi: in Italia il tempo medio di rilevazione di un
    incidente è stimato in circa sei mesi;
  \item un documento di accountability che giustifichi una conservazione più lunga per motivi di
    sicurezza non è bastato a evitare una sanzione (circa \texteuro{}50.000) a una Regione.
\end{itemize}
Prima di raccogliere dati di monitoraggio occorre chiedersi se sono dati personali: se sì, si applica
tutto il GDPR.

\subsection{Privacy by design: un caso}
Un'applicazione legacy a cui la sicurezza è stata aggiunta solo a posteriori è costata a una banca
spagnola una sanzione di \texteuro{}500.000 per mancanza di privacy by design: sicurezza e privacy vanno
integrate dalla fase di progettazione.

% ---------------------------------------------------------------------
\section{Risk management secondo la ISO: primi concetti}

\paragraph{ISO e GDPR: due prospettive.} Nel GDPR il rischio si misura sui \textbf{diritti e le libertà
degli interessati}; nella ISO si misura sull'\textbf{organizzazione}. Un'azienda preoccupata per le
conseguenze di un incidente sulla propria operatività non ha, solo per questo, un problema di notifica
GDPR: sono valutazioni distinte.

\paragraph{Impatto e probabilità in chiave ISO.} L'impatto è la conseguenza per l'organizzazione dello
sfruttamento di una vulnerabilità; la probabilità è quella che la vulnerabilità venga sfruttata.
Esempio della settimana: la vulnerabilità PHP (esecuzione di codice da remoto) aveva da subito impatto
altissimo e probabilità bassa; quando è stato pubblicato un \emph{proof of concept} la probabilità è
schizzata in alto e tutti hanno corso a installare le patch.

\paragraph{Definizioni (slide).}
\begin{itemize}
  \item \textbf{Asset}: ciò che ha valore e va protetto. Tipi: informazioni, hardware, software,
    \textbf{persone}, beni intangibili (reputazione, marchio, avviamento). Non tutto va protetto (un
    bilancio destinato al deposito pubblico no). Le persone sono l'asset più dimenticato: chi detiene
    conoscenza non documentata (\emph{key people}) è un asset critico, e il suo abbandono impatta su
    disponibilità e integrità delle informazioni. Da qui il peso dell'ingegneria sociale.
  \item \textbf{Minaccia}: evento o attore che \emph{potrebbe} danneggiare un asset. Categorie: dolose,
    accidentali, ambientali, tecniche.
  \item \textbf{Vulnerabilità}: debolezza sfruttabile (patch mancanti, password deboli o di default,
    scarsa formazione, configurazioni insicure, lacune fisiche come un cassetto senza chiave, applicazioni
    progettate male).
\end{itemize}

\begin{esame}
Nella ISO l'impatto è valutato sull'\textbf{organizzazione} (triade CIA); nella metodologia di severità
ENISA e nel GDPR sui \textbf{diritti e le libertà degli interessati}. Non confondere i due piani.
\end{esame}

% ---------------------------------------------------------------------
\section{Domande aperte e punti da verificare}
\begin{verifica}
\begin{itemize}
  \item \textbf{Limite dei 21 giorni}: ricostruito dal parlato (``linee guida del Garante di fine 2024''
    sui metadati della posta elettronica). Verificare il provvedimento del Garante e le condizioni per
    superare il limite.
  \item \textbf{Sanzione alla Regione Lombardia} (circa \texteuro{}50.000) e \textbf{sanzione di
    \texteuro{}500.000 a una banca spagnola}: importi e soggetti dal parlato, da verificare.
  \item \textbf{8.12 DLP ``nuovo'' nella revisione 2022}: detto in aula; le slide indicano solo che i
    controlli nuovi sono 11. (Integrazione: confermato, l'8.12 è tra gli 11 controlli nuovi del 2022.)
  \item Vulnerabilità PHP ``delle ultime 48 ore'': non identificata con precisione nella trascrizione.
  \item ``Tempo medio di rilevazione di circa sei mesi in Italia'': stima del docente, senza fonte.
\end{itemize}
\end{verifica}

% =====================================================================
\lezione{26/10/2026}{Risk management ISO e introduzione al Cyber Resilience Act}
  {Gianluca Gilardi; Prof.~Giuseppe Vaciago}
  {minacce e threat intelligence; scale di probabilità e impatto; prioritizzazione; metodologia ENISA
   per i dati personali; ISO vs ENISA; risk register; opzioni di trattamento; Statement of
   Applicability; piano di trattamento; casi di digital forensics su dispositivi smart; CRA: ambito,
   esclusioni, requisiti essenziali, prodotti importanti e critici, obblighi della filiera, segnalazione;
   briefing del laboratorio~3}
\label{lez:2026-10-26}
% =====================================================================

\begin{nota}[Fonte]
Trascrizione automatica della registrazione, revisionata (inizio con prove audio; pausa intorno ai
minuti 63--78). Slide di riferimento: \emph{06.~27001\_2026\_2310} (\emph{Risk Management}) per la prima
parte, \emph{05.~PoliTO\_CRA} per la seconda; briefing del laboratorio in
\texttt{lab/lab\_03/materiale}. La parte del CRA prosegue e si conclude nella lezione del 02/11
(Capitolo~\ref{lez:2026-11-02}).
\end{nota}

\section{Sintesi}
Tre blocchi. Gianluca completa il \textbf{risk management secondo la ISO}: scale, prioritizzazione,
confronto con la metodologia ENISA per i dati personali, contenuto del risk register, opzioni di
trattamento, Statement of Applicability e piano di trattamento. Il professore introduce poi il
\textbf{Cyber Resilience Act} partendo da casi in cui dispositivi smart sono diventati prove in giudizio,
per mostrare quanto IoT e dispositivi indossabili siano sottovalutati sul piano della sicurezza; segue la
prima metà delle slide (ambito, esclusioni, requisiti, filiera, segnalazione). Chiude il briefing del
laboratorio~3 sul risk register.

% ---------------------------------------------------------------------
\section{Organizzazione}
\begin{itemize}
  \item \textbf{Laboratorio~3, venerdì 30/10}: compilazione di un risk register ISO~27001. Materiale nella
    cartella della piattaforma ``lab 2026 1030'': briefing (PDF), playbook (Markdown) e una copia di prova
    del file Excel senza scenari, per esercitarsi. Il playbook va letto prima.
  \item Domande sull'esercitazione: per email a Gianluca, o in aula il venerdì.
  \item Il professore resta disponibile durante la pausa per chi vuole discutere il laboratorio~2.
  \item La parte finale del CRA (slide residue e un ``meme'' sulla filiera) è rinviata al lunedì
    successivo.
\end{itemize}

% =====================================================================
\section{Risk management ISO (Gianluca)}
% =====================================================================

\subsection{Minacce, vulnerabilità e threat intelligence}
Le minacce sfruttano vulnerabilità tecniche (individuabili con penetration test e scansioni) o umane
(ingegneria sociale). L'ultima revisione della 27001 chiede di individuare le proprie \textbf{fonti di
threat intelligence}: non si deve conoscere ogni minaccia da soli, ma attingere a fonti esterne, come i
bollettini dell'ACN. La probabilità di sfruttamento di molte vulnerabilità note è già stimata nelle fonti
di threat intelligence.

\subsection{Scale, prioritizzazione, propensione al rischio}
\begin{itemize}
  \item Le matrici sono quasi sempre \textbf{qualitative}: alto, medio, basso (al massimo quattro livelli,
    con ``molto alto''); le scale quantitative sono rare.
  \item Formula di base: \textbf{Rischio = Impatto $\times$ Probabilità}.
  \item Prioritizzare: di solito si parte dal rischio complessivo, ma conta anche il tempo di attuazione
    delle contromisure. Le \emph{quick win} possono passare avanti a rischi più alti che richiedono
    progetti lunghi. Decide il vertice (il board), che però chiede a chi valuta costi, tempi e proposta di
    priorità.
  \item La ISO \textbf{non fissa} alcuna matrice: si deve essere coerenti e saper difendere le scelte. La
    più usata è la \textbf{3$\times$3}.
  \item \textbf{Risk appetite} (slide): disponibilità complessiva ad assumere rischi, decisione strategica
    del board. \textbf{Risk tolerance}: livello accettato per singole attività o processi, livello
    operativo con soglie specifiche.
\end{itemize}

\paragraph{Livelli di impatto sull'organizzazione (indicazioni del docente).} Basso: perdita economica
limitata, nessun fermo operativo, recupero facile. Medio: un ``mal di testa'' misurabile (sanzioni, fermo
limitato nel tempo) ma superabile. Alto (o molto alto): evento che può mettere a rischio
\textbf{l'esistenza stessa} dell'organizzazione, come la perdita totale dei dati senza backup.

\subsection{La metodologia ENISA per i dati personali}
ENISA non ha pubblicato solo la metodologia di severità dei breach, ma anche linee guida per le PMI sulla
valutazione del rischio dei trattamenti. Sono pensate per i dati personali, ma ogni dato personale è
un'informazione: il metodo è un buon punto di partenza per costruirne uno proprio.
\begin{enumerate}
  \item \textbf{Contesto del trattamento}: quali dati, finalità, mezzi, luogo (on premise o cloud, con
    rischi di terze parti), categorie di interessati, destinatari; tutte le fasi: raccolta, conservazione,
    uso, trasferimento, cancellazione. Output: il perimetro della valutazione.
  \item \textbf{Impatto sugli individui}, su quattro livelli (basso, medio, alto, molto alto), valutato
    separatamente per perdita di riservatezza, integrità e disponibilità: \textbf{si prende il più alto}.
  \item \textbf{Probabilità della minaccia}, su quattro aree (risorse di rete e tecniche; processi e
    procedure; persone e terze parti; settore e dimensione), ciascuna con punteggio 1--3. Totale: 4--5
    bassa, 6--8 media, 9--12 alta. Le definizioni sono qualitative (``improbabile'', ``ragionevole
    possibilità'', ``probabile''), senza percentuali.
  \item \textbf{Valutazione del rischio} con la matrice probabilità $\times$ impatto.
\end{enumerate}

\begin{center}
\small
\textbf{Matrice ENISA (slide)}\par\smallskip
\begin{tabular}{llll}
\toprule
\textbf{Probabilità} & \textbf{Impatto basso} & \textbf{Impatto medio} & \textbf{Impatto alto/molto alto} \\ \midrule
Bassa & basso & medio & alto \\
Media & medio & alto & alto \\
Alta  & alto & alto & alto \\ \bottomrule
\end{tabular}
\end{center}

\subsection{ISO 27001 e ENISA a confronto}
\begin{center}
\small
\begin{tabularx}{\linewidth}{l X X}
\toprule
\textbf{Aspetto} & \textbf{ISO/IEC 27001} & \textbf{ENISA (dati personali)} \\ \midrule
Ambito & sicurezza delle informazioni in generale & trattamenti di dati personali \\
Impatto & sull'organizzazione (triade CIA) & sui diritti e le libertà degli individui \\
Calcolo & metodologie libere & approccio strutturato in 4 passi \\
Minacce & categorie generali & questionari per contesto \\
Controlli & 93 controlli (Allegato A) & misure proporzionate al livello di rischio \\
Conformità & certificazione ISO & art.~32 GDPR \\ \bottomrule
\end{tabularx}
\end{center}
I due approcci si usano spesso insieme: ENISA per i trattamenti di dati personali, ISO per l'ISMS, un
unico registro con doppia valutazione. La metodologia ENISA è apprezzata anche nella comunità ISO.

\subsection{Il risk register}
Non esiste un formato obbligatorio (Excel o piattaforme SaaS). Campi che, per esperienza, vanno considerati
irrinunciabili:
\begin{itemize}
  \item \textbf{ID univoco} del rischio, per tracciarlo in tutti i documenti e progetti;
  \item asset coinvolti, minaccia, vulnerabilità;
  \item probabilità e impatto \textbf{inerenti} e, dopo i controlli, \textbf{residui};
  \item livello di rischio calcolato con la matrice;
  \item \textbf{risk owner}: chi decide e firma il trattamento e ne risponde;
  \item trattamento proposto (misure già adottate e da adottare).
\end{itemize}

\subsection{Opzioni di trattamento e controlli}
\begin{itemize}
  \item \textbf{Mitigare} (ridurre): controlli che riducono probabilità o impatto; è l'opzione più comune.
  \item \textbf{Accettare}: il rischio rientra nella propensione al rischio (vale per l'inerente o per il
    residuo); richiede approvazione formale.
  \item \textbf{Evitare}: rinunciare all'attività che genera il rischio.
  \item \textbf{Trasferire}: condividere il rischio con un terzo (assicurazione, outsourcing). Non elimina
    la responsabilità: anche il trasferimento è una decisione di cui si risponde e che deve essere coerente
    con la propensione al rischio. La \emph{shared responsibility matrix} dei cloud provider è una forma di
    trasferimento parziale.
\end{itemize}
La 27001 dà i nomi dei controlli, la 27002 i mezzi per attuarli (es.\ controllo d'accesso attuato con
l'MFA). \textbf{Tipi di controllo}: preventivi (VA, penetration test), detective (IDS/IPS, DLP),
correttivi (disaster recovery, business continuity).

\subsection{Statement of Applicability e piano di trattamento}
\paragraph{SoA.} Documento \textbf{obbligatorio} per la conformità 27001: elenca tutti i controlli
dell'Allegato~A, dichiara se ciascuno è applicabile, ne giustifica inclusione o esclusione e ne indica lo
stato di attuazione. Documenta le decisioni prese nel trattamento del rischio. Per Gianluca è ancora più
importante saper giustificare perché un controllo \textbf{non} è attuato; l'auditor può contestare il
legame tra un controllo e un rischio e imporre di correggere la SoA.

\paragraph{Piano di trattamento.} Per ogni rischio da trattare: opzione scelta, azioni concrete, owner di
ogni azione (accountability), tempi di completamento, risorse (budget e personale: il management chiederà
sempre ``quanto costa?''), giustificazione. La giustificazione discende in gran parte dalla propensione al
rischio: si accetta perché sotto soglia; si mitiga perché l'inerente è sopra soglia e i controlli portano il
residuo entro la soglia; si trasferisce ciò che resta fuori dal proprio controllo; si evita quando nessun
residuo è accettabile. La matrice \textbf{RACI} (responsabile, accountable, consultato, informato) è
fondamentale nella procedura di gestione dei data breach, ma non è richiesta per il piano di trattamento.

\paragraph{Monitoraggio continuo.} I rischi cambiano, nascono e muoiono: il registro non si chiude in un
cassetto. La threat intelligence dimostra all'auditor l'attenzione continua al rischio.

\begin{esame}
\textbf{Messaggi chiave del risk management (slide finali):} è il fondamento della cybersicurezza; servono
asset, minacce, vulnerabilità, impatto e probabilità; si segue un processo strutturato (ISO, NIST, FAIR,
ENISA: non serve reinventare la ruota); il trattamento è \textbf{strategico}, non solo tattico (non ha
senso spendere per mitigare un processo che verrà dismesso a breve); le decisioni si documentano nella
\textbf{SoA}; il monitoraggio è continuo.
\end{esame}

% =====================================================================
\section{Cyber Resilience Act (Prof.~Vaciago)}
% =====================================================================

\subsection{Introduzione: i dispositivi smart come fonte di prova e di rischio}
Fuori programma, per introdurre il tema. Ogni dispositivo digitale tocca due ambiti: \textbf{digital
forensics} e \textbf{cybersecurity}. Le stesse informazioni utili alla polizia sono utili a chi vuole
commettere un reato.
\begin{itemize}
  \item \textbf{Caso Fitbit}: in un'indagine per omicidio i dati dei braccialetti dei coniugi hanno
    smentito la versione del marito, che ha poi confessato.
  \item \textbf{Caso Amazon Echo}: la polizia cercava registrazioni ambientali; Amazon ha chiarito che non
    registra in continuo (solo dopo la parola di attivazione o per attivazioni accidentali). L'aspetto
    curioso: inizialmente ha negato i dati sostenendo che solo il proprietario potesse chiederli.
  \item \textbf{Smart TV}: studi recenti (il professore cita la ricercatrice Anna Maria Mandalari)
    mostrano funzioni nascoste che attivano la telecamera per rilevare le reazioni degli spettatori;
    erano previste nelle condizioni d'uso.
  \item \textbf{Caso Ross Compton}: accusato di aver incendiato la propria casa per frodare
    l'assicurazione, è stato condannato anche grazie ai dati del suo pacemaker (battito incompatibile con
    la sua versione). Dibattito sull'uso processuale dei dati sanitari.
\end{itemize}
Il professore richiama il rischio legato alla quantità di dati sanitari ceduti spontaneamente a
dispositivi e assistenti di IA, spesso in mano a soggetti statunitensi o cinesi, e la superficie
d'attacco offerta da oggetti come termostati o pacemaker. Opinione del docente: il CRA è più importante
di AI Act e NIS2, anche se meno ``di moda''.

\subsection{CRA e NIS2: prodotto contro organizzazione}
\begin{itemize}
  \item \textbf{NIS2} riguarda le reti e le \textbf{aziende}; è una \textbf{direttiva}, attuata in Italia con
    il D.Lgs.~138/2024 (in Italia i soggetti obbligati sono, secondo il docente, circa 20.000).
  \item \textbf{CRA} riguarda il \textbf{prodotto}; è un \textbf{regolamento}, direttamente applicabile in
    tutti gli Stati membri.
\end{itemize}
Obiettivo: tutelare consumatori e imprese che acquistano o usano prodotti o software con una componente
digitale, con requisiti obbligatori per fabbricanti e rivenditori. Problemi affrontati: livello di
sicurezza inadeguato e aggiornamenti insufficienti; impossibilità per gli utenti di capire quali prodotti
siano sicuri. La sicurezza è la prima voce su cui si risparmia per tenere basso il prezzo. Soluzioni:
regole armonizzate per l'immissione sul mercato, requisiti lungo progettazione, sviluppo e manutenzione,
\textbf{dovere di diligenza per l'intero ciclo di vita}. Il considerando~3 spiega perché serve: né il
Cybersecurity Act né la NIS2 coprono direttamente requisiti obbligatori di sicurezza dei prodotti.

\subsection{Ambito di applicazione}
Il CRA si applica ai \textbf{prodotti con elementi digitali messi a disposizione sul mercato} il cui uso
previsto o ragionevolmente prevedibile include una \textbf{connessione dati} diretta o indiretta, logica o
fisica, a un dispositivo o a una rete.
\begin{itemize}
  \item \emph{Prodotto con elementi digitali}: prodotto software o hardware e le relative soluzioni di
    elaborazione dati da remoto, inclusi i componenti immessi sul mercato separatamente.
  \item Connessione \emph{logica} (tramite interfaccia software), \emph{fisica} (cavi, onde radio,
    interfacce elettriche, ottiche, meccaniche), \emph{indiretta} (come parte di un sistema più ampio
    collegabile).
\end{itemize}

\begin{esame}
\textbf{Il software è un prodotto ai fini del CRA.} Errore tipico all'esame: pensare solo a smartwatch e
smart TV ed escludere le applicazioni. Le due parole chiave sono \textbf{mercato} (messa a disposizione
sul mercato, cioè prospettiva commerciale) e \textbf{connessione}.
\end{esame}

\subsection{Esclusioni}
\begin{itemize}
  \item Prodotti già disciplinati da norme settoriali (es.\ regolamento sui \textbf{dispositivi medici}):
    il CRA può essere limitato o escluso se la disciplina settoriale è coerente con il quadro complessivo e
    garantisce un livello di protezione pari o superiore; la Commissione può chiarire con atti delegati.
    Un anello o un orologio smart non è un dispositivo medico; il confine resta discusso (in futuro alcuni
    sistemi di IA generativa potrebbero esservi ricondotti).
  \item \textbf{Pezzi di ricambio} identici, fabbricati secondo le stesse specifiche.
  \item Prodotti sviluppati o modificati esclusivamente per \textbf{sicurezza nazionale o difesa}, o per
    trattare informazioni classificate.
\end{itemize}
Per l'esame basta sapere che esiste un'esclusione generale a favore delle norme settoriali più specifiche.

\subsection{Approccio basato sul rischio e requisiti essenziali}
Il fabbricante deve svolgere un \textbf{cybersecurity risk assessment} del prodotto, che informa tutte le
fasi del ciclo di vita, è documentato e aggiornato durante il periodo di supporto e considera uso previsto e
prevedibile, condizioni d'uso, applicabilità dei requisiti e gestione delle vulnerabilità. Rispetto a un
risk assessment su un processo (es.\ selezione del personale), quello sul dispositivo è
\textbf{profondamente tecnico}: bisogna entrare in ogni processo del dispositivo e nella sua filiera di
fornitori.

\paragraph{Uso prevedibile.} Il fabbricante non può conoscere tutti gli usi futuri; quando emerge un uso non
previsto, una risposta è aggiornare le condizioni d'uso. Il professore cita l'esempio di un assistente di
IA generativa le cui condizioni sono state aggiornate per escludere che si sostituisca a un giudice.

\paragraph{Requisiti per l'immissione sul mercato.} Il prodotto deve rispettare i requisiti essenziali
dell'\textbf{Allegato~I, Parte~I} (se installato, mantenuto e usato correttamente, con gli aggiornamenti
di sicurezza) e i processi del fabbricante quelli dell'\textbf{Allegato~I, Parte~II} (gestione delle
vulnerabilità). Tra i requisiti della Parte~I: nessuna vulnerabilità sfruttabile nota, configurazione
sicura per impostazione predefinita, aggiornamenti di sicurezza, protezione da accessi non autorizzati,
riservatezza e integrità dei dati, \textbf{minimizzazione dei dati}, disponibilità delle funzioni
essenziali, superficie d'attacco ridotta, registrazione dell'attività interna, cancellazione sicura dei dati
da parte dell'utente. Non serve conoscere gli allegati nel dettaglio.

\begin{esame}
\textbf{Spunto per le domande aperte: il rovescio della minimizzazione.} Meno dati significa meno rischio
per l'interessato e per la sicurezza, ma anche meno elementi per il debugging, per l'usabilità e
soprattutto per l'\textbf{analisi forense} dopo un incidente. ``Vince la privacy, muore l'esperto di
cybersecurity'': c'è una tensione da saper argomentare.
\end{esame}

\subsection{Prodotti importanti e critici}
\begin{center}
\small
\begin{tabularx}{\linewidth}{>{\raggedright\arraybackslash}p{3.4cm} X X}
\toprule
& \textbf{Prodotti importanti} & \textbf{Prodotti critici} \\ \midrule
Dove & Allegato~III, classi I e II (es.\ password manager) & Allegato~IV (es.\ dispositivi hardware con
  \emph{security box}), individuati con atti della Commissione \\
Criteri & funzioni critiche per la sicurezza di altri prodotti, reti o servizi, oppure rischio
  significativo di effetti negativi su molti prodotti o sulla salute e sicurezza degli utenti & --- \\
Conseguenza & procedure di valutazione della conformità più rigorose & certificato europeo di
  cybersicurezza con livello di affidabilità almeno \emph{sostanziale} \\ \bottomrule
\end{tabularx}
\end{center}
Il professore ammette che la distinzione pratica non è ancora chiara: per ora l'unico obbligo già operativo
è la \textbf{segnalazione degli incidenti}.

\subsection{Security by design e obblighi lungo la filiera}
\begin{itemize}
  \item \textbf{Fabbricante}: sicurezza integrata in progettazione e sviluppo (come la privacy by design
    del GDPR: la sicurezza è parte della privacy); risk assessment come base di tutto; \textbf{due
    diligence sui componenti di terzi} (supply chain); gestione delle vulnerabilità durante il periodo di
    supporto.
  \item \textbf{Importatore}: verifica che il fabbricante abbia svolto la valutazione di conformità e
    redatto la documentazione tecnica, che ci siano marcatura CE e dichiarazione UE di conformità; azioni
    correttive o ritiro se il prodotto non è conforme; informa fabbricante e autorità di vigilanza del
    mercato delle vulnerabilità scoperte.
  \item \textbf{Distributore}: due diligence, verifica di marcatura CE e documentazione, azioni correttive,
    segnalazione di vulnerabilità e rischi.
\end{itemize}
\paragraph{Il problema pratico.} Il professore cita un grande distributore italiano di hardware e software:
se molti prodotti a basso costo extra UE non rispettano il CRA, dovrà scegliere se toglierli dal catalogo.
Se gli Stati membri applicassero le regole con rigore diverso, i consumatori potrebbero comprare
all'estero ciò che in Italia non si trova più.

\subsection{Obblighi di segnalazione (anticipazione)}
Il fabbricante notifica ogni \textbf{vulnerabilità attivamente sfruttata} al CSIRT tramite la
\textbf{piattaforma unica di segnalazione}, a cui ha accesso anche ENISA: \textbf{early warning entro 24
ore}, \textbf{notifica entro 72 ore}, \textbf{final report entro 14 giorni} dalla disponibilità di una
misura correttiva. Un incidente può far scattare insieme obblighi CRA e di altre norme (24 ore CRA, 72 ore
GDPR\ldots): il tema è ripreso nella lezione del 02/11 (Capitolo~\ref{lez:2026-11-02}) e nella tabella
delle scadenze (Capitolo~\ref{cap:scadenze}).

% =====================================================================
\section{Briefing del laboratorio 3 (Gianluca)}
% =====================================================================
Ruolo: \textbf{responsabile della sicurezza delle informazioni} di un piccolo studio professionale di
consulenza (molto digitale). Il registro contiene già \textbf{18 rischi}: non va inventato nulla.
\begin{itemize}
  \item Per ogni rischio: probabilità e impatto \textbf{inerenti}, separatamente per riservatezza,
    integrità e disponibilità; controlli di trattamento; rischio \textbf{residuo}; decisione finale.
  \item Solo i \textbf{controlli tecnologici} (8.x): niente persone, processi o management.
  \item Matrice \textbf{3$\times$3}; un rischio \textbf{alto} (inerente o residuo) non è accettabile, il
    residuo deve essere medio o basso.
  \item Decisione tra le quattro opzioni; un rischio trasferito resta sotto la responsabilità dell'azienda;
    l'accettazione si usa quando l'inerente è già medio o inferiore.
  \item Dei 34 controlli, 3 non sono applicabili a uno studio di consulenza, 26 sono applicabili e 5 vanno
    valutati caso per caso; ogni esclusione va motivata nelle righe apposite (è ciò che finirebbe nella
    SoA, che però non è oggetto dell'esercizio).
  \item Un foglio di \textbf{coerenza} controlla la completezza formale, non la correttezza nel merito.
  \item Non esiste una risposta unica: \textbf{conta il ragionamento}, non i numeri.
\end{itemize}

\begin{esame}
\textbf{Errore tipico: fare la media} dei punteggi di riservatezza, integrità e disponibilità. Il livello
del rischio è il \textbf{valore massimo} dei tre: un rischio alto solo sulla disponibilità è comunque alto.
\end{esame}

\noindent Svolgimento e metodo completi nel Capitolo~\ref{lab:03}.

% ---------------------------------------------------------------------
\section{Domande aperte e punti da verificare}
\begin{verifica}
\begin{itemize}
  \item \textbf{Threat intelligence nella 27001:2022}: citata come requisito dell'ultima revisione, senza
    numero di controllo nelle slide.
  \item \textbf{Linee guida ENISA per le PMI}: titolo e anno del documento non indicati.
  \item \textbf{Data di applicazione del CRA}: il professore dice ``in vigore da settembre, il 21 se ricordo
    bene''; secondo le slide della lezione del 02/11 gli obblighi di segnalazione si applicano
    dall'\textbf{11 settembre 2026}. Verificare.
  \item \textbf{Circa 20.000 soggetti NIS2 in Italia}: stima a voce del docente.
  \item \textbf{Nuovo decreto legislativo italiano su IA e rapporti di lavoro} (``in vigore da venerdì''),
    citato a proposito dell'uso dell'IA nella selezione del personale: estremi non ricavabili dalla
    trascrizione.
  \item Nomi propri dei casi (Fitbit, Ross Compton, ricercatrice Mandalari) ricostruiti da una trascrizione
    rumorosa: verificare l'ortografia prima di citarli.
  \item Esempi di prodotti importanti e critici: dalle slide (password manager; dispositivi hardware con
    security box); l'elenco completo è negli Allegati III e IV del CRA.
\end{itemize}
\end{verifica}

% =====================================================================
\lezione{30/10/2026}{Laboratorio 3: risk register ISO/IEC 27001}
  {Gianluca Gilardi}
  {compilazione in gruppo di un risk register ISO/IEC~27001:2022 (controlli tecnologici dell'Allegato~A,
   clausola~8) per uno studio professionale}
\label{lez:2026-10-30}
% =====================================================================

\begin{nota}[Fonte]
Trascrizione automatica della registrazione: quasi tutta silenzio (gruppi al lavoro). Contenuti ripresi
dal briefing, dal playbook e dal file Excel del laboratorio (\texttt{lab/lab\_03/materiale}). Il caso,
il metodo e uno svolgimento ragionato sono nel Capitolo~\ref{lab:03}.
\end{nota}

\section{Sintesi}
Laboratorio in gruppo di circa 90 minuti. Ogni gruppo, nel ruolo di responsabile della sicurezza delle
informazioni di uno studio professionale, completa la valutazione e il trattamento di 18 scenari di
rischio già individuati: probabilità e impatto inerenti per riservatezza, integrità e disponibilità,
controlli dell'Allegato~A (solo clausola~8), misure, giustificazione, efficacia, rischio residuo e
decisione di trattamento.

\section{Organizzazione}
\begin{itemize}
  \item Il file dell'esercitazione è stato caricato nella cartella dei materiali della piattaforma poco
    prima dell'inizio; chi aveva già provato la copia di prova e letto il playbook partiva avvantaggiato.
  \item Prima cosa da compilare: nome del gruppo e tutti i partecipanti nel primo foglio, altrimenti il
    lavoro non può essere attribuito né valutato. Nome del file e caricamento secondo le istruzioni del
    file (come nel laboratorio precedente).
  \item \textbf{Consegna entro le 13:15}: i file caricati dopo vengono scartati.
  \item Esiti promessi entro domenica; i gruppi selezionati presentano lunedì 02/11.
\end{itemize}

\section{Contenuti: le indicazioni date in apertura}
Gianluca ripercorre il foglio di lavoro:
\begin{itemize}
  \item i rischi, le minacce e gli asset sono già compilati: si inseriscono probabilità e impatto
    inerenti, mentre il punteggio è calcolato in automatico dal file;
  \item si indicano i controlli applicabili, le misure concrete e il \textbf{perché} di quelle misure,
    e si valuta quanto sono efficaci rispetto a quel rischio specifico;
  \item si stimano probabilità e impatti residui (C, I, A) e si sceglie la decisione di trattamento
    (ridotto, accettato, trasferito, evitato);
  \item la propensione al rischio non ammette nulla sopra il livello medio: un rischio alto non si può
    semplicemente accettare;
  \item il foglio riepilogativo verifica solo la completezza, non la correttezza.
\end{itemize}

\begin{esame}
``Non esiste una risposta corretta'': tutto può essere discusso, purché il ragionamento sia ragionevole e
giustificato. Il playbook suggerisce una soluzione di riferimento, ma non è vincolante. Lo stesso criterio
vale per l'esame: conta la motivazione delle scelte di rischio inerente, controlli e rischio residuo.
\end{esame}

\section{Domande aperte e punti da verificare}
\begin{verifica}
\begin{itemize}
  \item Al termine Gianluca accenna a un intervento del professore, non registrato in modo comprensibile.
  \item Criteri di valutazione e soglia per le presentazioni: vedere gli avvisi in piattaforma.
\end{itemize}
\end{verifica}

% =====================================================================
\lezione{02/11/2026}{Conclusione CRA e simulazione CRA--GDPR (Pulse K)}
  {Prof.~Giuseppe Vaciago; Gianluca (laboratori)}
  {obblighi di segnalazione CRA; CRA e AI Act; impatti del CRA; software nel perimetro CRA;
   simulazione di incidente Pulse K; presentazione del laboratorio di audit}
\label{lez:2026-11-02}
% =====================================================================

\begin{nota}[Fonte]
Trascrizione automatica della registrazione, revisionata. Le date, gli orari e i numeri del caso
Pulse K sono ripresi dalle slide \emph{Case1\_PulseK\_CRA\_GDPR\_Simulation} (la trascrizione era
inaffidabile su questi punti). Slide di riferimento: \emph{05.~PoliTO\_CRA}.
\end{nota}

\section{Sintesi}
Conclusione della parte teorica sul Cyber Resilience Act: obblighi di segnalazione, rapporto con
l'AI Act, impatti e sfide. Segue la prima simulazione di incidente (smartwatch per bambini Pulse K),
che obbliga a lavorare insieme su CRA e GDPR: qualificazione del caso, orologi di notifica,
severità ENISA e piano di remediation. Messaggio centrale: a fine corso si deve saper
\textbf{(1) riconoscere la norma applicabile} e \textbf{(2) seguire il percorso corretto in caso di
incidente}. Infine Gianluca presenta il laboratorio di venerdì, in cui gli studenti faranno gli auditor.

% ---------------------------------------------------------------------
\section{Obblighi di segnalazione del CRA: riepilogo}

Il fabbricante (\emph{manufacturer}) deve notificare ogni \textbf{vulnerabilità attivamente
sfruttata} presente in un prodotto con elementi digitali di cui venga a conoscenza. La notifica
passa dalla \textbf{Single Reporting Platform} gestita da ENISA e raggiunge \textbf{due destinatari}:
\begin{itemize}
  \item il \textbf{CSIRT} designato come coordinatore (per l'Italia il CSIRT Italia presso l'ACN);
  \item l'\textbf{ENISA}, che ha accesso alla piattaforma.
\end{itemize}

\paragraph{Vulnerabilità attivamente sfruttata (art.~14(2)).}
\begin{itemize}
  \item \textbf{Early warning}: entro \textbf{24 ore} dalla conoscenza;
  \item \textbf{Notifica della vulnerabilità}: senza ingiustificato ritardo e comunque entro \textbf{72 ore};
  \item \textbf{Final report}: entro \textbf{14 giorni} dal momento in cui è disponibile una misura
    correttiva o di mitigazione.
\end{itemize}

\paragraph{Incidente grave (art.~14(4)).} Lo schema è parallelo: early warning entro 24 ore,
notifica dell'incidente entro 72 ore, \textbf{final report entro un mese}. Il professore non ha
messo questa scadenza nelle slide per non aggiungere complessità, ma ha anticipato che comparirà
nella tabella riassuntiva preparata per l'esame.

\paragraph{Date di applicazione.} Obblighi di segnalazione dall'\textbf{11 settembre 2026};
applicazione integrale dall'\textbf{11 dicembre 2027}.

\paragraph{Differenza principale con il GDPR.}
Il CRA fissa una scadenza anche per il \emph{final report}, il GDPR no. Nel GDPR entro 72 ore va
fatta la notifica, ma raccogliere tutte le informazioni del modulo in 72 ore è quasi impossibile:
nella prassi la notifica completa arriva dopo almeno un mese e, trascorse due o tre settimane, è
il Garante a sollecitarne il completamento. Non esiste però un termine vincolante.

\begin{esame}
\begin{itemize}
  \item Primo passo: \textbf{qualificare il caso} (CRA, NIS2, DORA, GDPR?).
  \item Secondo passo: ricordare con precisione le \textbf{tempistiche} applicabili
    (vedi tabella al Capitolo~\ref{cap:scadenze}).
  \item Errore tipico: confondere l'\emph{early warning} del CRA con la \emph{notifica preliminare}
    del GDPR.
\end{itemize}
\end{esame}

% ---------------------------------------------------------------------
\section{CRA e AI Act: presunzione di conformità}

I prodotti con elementi digitali che rientrano nel CRA e sono anche \textbf{sistemi di IA ad alto
rischio} ai sensi dell'AI Act si considerano conformi ai requisiti di cybersicurezza dell'AI Act
(art.~15 AI Act) se sono soddisfatte tre condizioni:
\begin{enumerate}
  \item il prodotto soddisfa i requisiti essenziali dell'\textbf{Allegato I, Parte I} del CRA;
  \item i processi del fabbricante rispettano l'\textbf{Allegato I, Parte II} del CRA
    (gestione delle vulnerabilità);
  \item il livello di protezione richiesto dall'AI Act è dimostrato nella \textbf{dichiarazione di
    conformità UE} rilasciata ai sensi del CRA.
\end{enumerate}

\noindent Anticipazione sull'AI Act: esistono \textbf{pratiche vietate} (es.\ \emph{social scoring})
e \textbf{sistemi ad alto rischio}, cioè sistemi che possono causare un danno serio alle persone e
sono soggetti a misure specifiche, anche di cybersicurezza. L'Allegato III dell'AI Act verrà
approfondito più avanti.

\begin{esame}
Da ricordare: esiste un'interazione forte tra l'\textbf{art.~15 dell'AI Act} e il CRA per i sistemi
di IA ad alto rischio. Gli allegati non vanno imparati a memoria.
\end{esame}

% ---------------------------------------------------------------------
\section{Impatti e sfide del CRA}

Il settore \textbf{tech e sviluppo software} sarà probabilmente il più colpito, data la pervasività
degli elementi digitali: potranno servire cambiamenti significativi nelle pratiche di sviluppo e nei
prodotti. Le sfide principali sono:
\begin{itemize}
  \item complessità dei requisiti;
  \item integrazione della \emph{security by design} e \emph{by default} in prodotti e processi;
  \item gestione della \emph{supply chain};
  \item investimenti necessari;
  \item gestione continua del rischio.
\end{itemize}

% ---------------------------------------------------------------------
\section{Perché il CRA è complesso: il quadro d'insieme}

\begin{itemize}
  \item \textbf{NIS2}: riguarda solo alcune categorie (soggetti \emph{essenziali} e \emph{importanti}),
    legate a infrastrutture critiche, con requisiti di sicurezza specifici e numerosi.
  \item \textbf{GDPR}: si applica a chiunque tratti dati personali, senza distinzioni di categoria.
    L'art.~32 discende dal principio di \emph{accountability} e non prescrive regole tecniche puntuali.
  \item \textbf{AI Act}: riguarda in primo luogo i \emph{provider} di sistemi di IA.
  \item \textbf{CRA}: l'esempio classico è lo smartwatch o qualsiasi oggetto IoT, ma il regolamento
    copre anche il \textbf{software}. Un'applicazione può quindi ricadere insieme sotto CRA, NIS2
    (che include il settore digitale) e GDPR.
\end{itemize}

\begin{esame}
\textbf{Le due competenze da portare a casa dal corso:}
\begin{enumerate}
  \item \textbf{ricondurre il rischio alla norma corretta}: capire quale quadro normativo si applica
    a un prodotto, un sistema di IA o una situazione. Non serve diventare giuristi: per i dettagli
    ci si può appoggiare a consulenti esperti della singola norma, ma il perimetro va riconosciuto da soli;
  \item \textbf{in caso di incidente, individuare il percorso corretto} da seguire: a chi notificare,
    cosa e entro quando.
\end{enumerate}
I requisiti tecnici, legali e organizzativi elencati negli allegati non vanno imparati a memoria.
Il problema vero, secondo il professore, è non capire in quale quadro normativo ci si trova.
\end{esame}

% ---------------------------------------------------------------------
\section{Chiarimento: il software nel perimetro del CRA}

Su domanda di uno studente, il professore ha approfondito la nozione di \emph{prodotto con elementi
digitali} applicata al software:
\begin{itemize}
  \item l'elemento chiave è la \textbf{connessione} (logica o fisica) a un dispositivo o a una rete;
  \item un \textbf{software installabile} (un'applicazione) rientra nel CRA; un sito web di per sé no;
  \item rientra anche il software \textbf{connesso a un hardware};
  \item restano esclusi gli ambiti coperti da regimi settoriali (es.\ software sanitario soggetto a
    normativa specifica) e l'ambito militare e della difesa.
\end{itemize}

\begin{esame}
\textbf{Il software rientra nel CRA.} Domande di dettaglio su questa definizione non saranno chieste,
perché la norma non è ancora pienamente applicabile e alcune interpretazioni restano incerte.
\end{esame}

% =====================================================================
\section{Simulazione \#1 -- Pulse K (Volta Wear): CRA e GDPR}
\label{sec:pulsek}
% =====================================================================

\begin{nota}
Le slide della simulazione non sono materia d'esame, ma questo tipo di esercizio (notifica alle
autorità) è molto probabile all'esame. Il caso aggiunge un livello di complessità in più: i dati
riguardano \textbf{minori}. Lo strumento \emph{CRA--GDPR Incident Desk} usato in aula è disponibile
sulla piattaforma del corso per esercitarsi.
\end{nota}

\subsection{Scenario}

\begin{itemize}
  \item \textbf{Volta Wear S.r.l.}, Torino (stabilimento principale in Italia), fondata nel 2019,
    48 dipendenti, fatturato 2025 di \texteuro{}9,4 milioni: \textbf{piccola impresa}.
  \item \textbf{Prodotto}: Pulse~K, smartwatch GPS per bambini di 5--12 anni, sul mercato UE da
    settembre 2024. Vendite tramite web shop proprio e retailer, tra cui ElektroMax GmbH (Germania).
  \item \textbf{182.000 orologi attivi}: Italia 118k, Francia 36k, Germania 17k, Spagna 11k.
\end{itemize}

\paragraph{Tre componenti, un solo prodotto.}
\begin{itemize}
  \item \textbf{Pulse K watch}: GPS, tasto SOS, messaggi vocali con contatti approvati;
  \item \textbf{Pulse Cloud}: progettato e gestito da Volta Wear su server UE in affitto; collega
    orologio e app;
  \item \textbf{Pulse Family app}: i genitori vedono la posizione live, impostano le zone ``Casa''
    e ``Scuola'', inviano messaggi vocali.
\end{itemize}

\paragraph{Dati trattati.}
\begin{center}
\begin{tabularx}{0.9\linewidth}{lXl}
\toprule
\textbf{Soggetto} & \textbf{Dato} & \textbf{Conservazione} \\ \midrule
Bambino  & Nome, anno di nascita, foto (facoltativa) & Durata dell'account \\
Bambino  & Posizione live e storico posizioni        & 30 giorni \\
Bambino  & Zone sicure ``Casa'' e ``Scuola''         & Durata dell'account \\
Bambino  & Messaggi vocali                           & 7 giorni \\
Genitore & Nome, email, cellulare                    & Durata dell'account \\ \bottomrule
\end{tabularx}
\end{center}
Non presenti nel cloud: dati di pagamento (gestiti dal payment provider) e password (salvate solo
come hash).

\subsection{La postura di sicurezza prima dell'incidente}

\begin{itemize}
  \item \textbf{Team}: un CTO e 9 sviluppatori, nessun addetto dedicato alla sicurezza.
  \item \textbf{Segnalazione vulnerabilità}: indirizzo \texttt{security@voltawear.eu} e una policy di
    due righe sul sito.
  \item \textbf{Componenti}: nessun elenco del software di terze parti incluso nel prodotto (nessuna SBOM).
  \item \textbf{Notti e weekend}: un solo sviluppatore reperibile; nessun legale né DPO reperibile;
    nessuna procedura di emergenza.
  \item \textbf{Customer care}: esternalizzato a CallPro S.r.l., con script fissi.
  \item \textbf{Privacy}: DPO esterno; \textbf{DPIA svolta al lancio nel 2024 e mai aggiornata}.
  \item \textbf{Formazione}: un e-learning generico di 45 minuti all'anno, uguale per tutti.
  \item \textbf{Ticket di aprile 2026}: ``aggiornare la libreria di login TokenGate, fix di sicurezza
    disponibile da marzo'', rinviato a dopo la campagna \emph{back to school}.
\end{itemize}

\begin{esame}
\textbf{DPIA non aggiornata.} Alla domanda ``esiste una DPIA?'' la risposta formalmente è sì, ma una
DPIA mai aggiornata può valere quasi quanto non averla. Averla resta meglio che non averla: tutto
dipende da quanto è cambiato il software dal 2024. Dettagli di questo tipo compaiono sia nei casi
reali sia nelle tracce d'esame.
\end{esame}

\subsection{Cronologia dell'incidente (i sei \emph{inject})}

La falla nel cloud era sfruttata da \textbf{mercoledì 4 novembre 2026, ore 21:00}.

\begin{center}
\small
\begin{tabularx}{\linewidth}{c p{2.6cm} X}
\toprule
\textbf{\#} & \textbf{Quando} & \textbf{Cosa succede} \\ \midrule
1 & Gio 5/11, 23:10 & Le richieste di posizione salgono a 40 volte il volume abituale, da centinaia
  di IP diversi. Lo sviluppatore reperibile blocca alcuni range IP e chiude il ticket come
  ``bot scraping, nessun impatto''. Nessun altro viene informato. \\ \midrule
2 & Ven 6/11, 08:30--11:05 & CallPro riceve 23 chiamate in due ore e mezza: gli orologi ricevono
  messaggi vocali da adulti sconosciuti. Gli operatori seguono lo script (riavviare l'orologio,
  cambiare il PIN). Alle 11:05 un team leader scrive al customer-care manager di Volta Wear;
  nessuno avvisa CTO o DPO. \\ \midrule
3 & Ven 6/11, 12:20 & Un ricercatore scrive a \texttt{security@}: cambiando il numero dell'orologio
  nella richiesta dell'app, qualsiasi utente autenticato vede la posizione live di qualsiasi bambino
  e può inviargli messaggi vocali. La falla è già sfruttata: un canale Telegram (``KidRadar'') vende
  la posizione live di qualsiasi Pulse~K a \texteuro{}15 a ricerca. Il ricercatore allega screenshot
  e annuncia la pubblicazione tra 30 giorni. \\ \midrule
4 & Ven 6/11, 16:00 & Riunione di crisi. I log confermano \textbf{41.250 orologi interrogati},
  \textbf{9.870 con posizione esposta} (posizione live, storico di 30 giorni, zone) e
  \textbf{37 bambini contattati} da sconosciuti. Cause: libreria TokenGate obsoleta (login
  falsificabile) e mancato controllo di \emph{ownership} (il cloud non verificava che l'orologio
  richiesto appartenesse all'account). CVSS 9.1 (critico). Il CTO propone di disattivare la
  condivisione della posizione; il CEO rifiuta: ``i genitori ci contano, non allarmiamo nessuno prima
  di lunedì, sistemiamo in silenzio''. \\ \midrule
5 & Sab 7/11 & Il retailer tedesco chiede se deve sospendere le vendite; una giornalista di
  \emph{Le Parisien} chiede un commento entro le 18:00; il marketing pubblica su Instagram:
  ``Pulse K è sicuro. Nessun dato è stato compromesso''. \\ \midrule
6 & Dom 8/11, 20:00 & Fix sul cloud: controllo di ownership aggiunto, TokenGate aggiornato, tutte le
  sessioni chiuse. Lunedì 9/11 alle 09:00 esce l'app 4.2. KidRadar resta online e ora vende gli
  \emph{storici} delle posizioni. \\ \bottomrule
\end{tabularx}
\end{center}

\subsection{Lezioni emerse dai singoli inject}

\paragraph{Inject 1 -- Sicurezza contro privacy.} Il riflesso tipico del tecnico di cybersicurezza è
``c'è un attacco, lo blocco''. Manca la sensibilità privacy: ogni attacco può avere impatto sui dati
personali, soprattutto quando riguarda minori.

\paragraph{Inject 2 -- Informazioni che non risalgono.} Il call center aveva segnali chiari, ma
nessuno ha avvisato DPO o CTO. Domanda aperta del debrief: ciò che sapeva un operatore del call
center esternalizzato vale come conoscenza dell'azienda?

\paragraph{Inject 3 -- Ricercatori e \emph{bug bounty}.} In alcuni casi l'attività è legale ed etica.
Diventa problematica quando la segnalazione si accompagna a richieste di denaro, che di fatto
equivalgono a un riscatto.

\paragraph{Inject 4 -- Bilanciamento di interessi.} Il CEO punta su reputazione e continuità del
servizio, ma non coglie che le conseguenze di un incidente gestito male possono essere molto più
gravi. ``Sistemare in silenzio'' è incompatibile con un orologio normativo da 24 ore.

\paragraph{Inject 5 -- Rischio reputazionale e comunicazione.} Esperienza maturata dal 2018 a oggi:
\begin{itemize}
  \item i clienti ricevono ormai spesso comunicazioni di data breach (es.\ il caso Trenitalia) e,
    soprattutto i più giovani, le accettano come normali;
  \item nel 2026 il rischio reputazionale di \emph{comunicare} un breach esiste ma raramente fa perdere
    clienti in modo significativo; il rischio vero è \textbf{non comunicare correttamente,
    minimizzare o mentire};
  \item i dati finiranno comunque sul dark web: la verità emergerà;
  \item ``chiedere scusa nel modo giusto'' è meglio che occultare;
  \item se i clienti sono già stati informati, con la stampa si è in posizione di forza: si può dire
    di aver già avvisato gli interessati e rifiutare dettagli coperti da indagini in corso;
  \item il post Instagram (``nessun dato compromesso'') contraddice i fatti noti all'azienda.
\end{itemize}
È compito del consulente spiegare all'alta direzione perché la trasparenza conviene.

\paragraph{Inject 6 -- Il fix fa partire un orologio CRA.} Si tratta del termine di 14 giorni per il
\emph{final report} sulla vulnerabilità (art.~14(2)(c)).

\subsection{Task 1 -- Qualificazione del caso ai sensi del CRA}

\begin{itemize}
  \item \textbf{È un prodotto con elementi digitali?} Sì. Pulse Cloud rientra come soluzione di
    elaborazione dati da remoto (\emph{remote data processing}) del prodotto.
  \item \textbf{Si applica un regime di esclusione?} No: non è un dispositivo medico, un veicolo,
    un prodotto dell'aviazione civile o dell'equipaggiamento marittimo, non è sviluppato
    esclusivamente per sicurezza nazionale o difesa, non è un pezzo di ricambio identico.
  \item \textbf{Ruolo}: Volta Wear è il \textbf{fabbricante}.
  \item \textbf{Categoria}: i wearable per bambini sono \textbf{prodotti importanti di Classe~I}
    (Allegato III).
  \item \textbf{Date}: l'evento (novembre 2026) cade dopo l'11/9/2026 (segnalazione già applicabile)
    e prima dell'11/12/2027 (applicazione integrale). Quindi \textbf{l'art.~14 si applica,
    l'Allegato~I non ancora}.
  \item \textbf{Dimensione}: piccola impresa; per le micro e piccole imprese l'art.~64(10) prevede una
    regola sanzionatoria specifica (vedi il riquadro \emph{Da verificare} in fondo alla lezione).
  \item \textbf{Verdetto}: c'è evidenza affidabile che un attore malevolo ha sfruttato la
    vulnerabilità, e l'incidente compromette la capacità del prodotto di proteggere dati sensibili o
    importanti. Sono soddisfatti entrambi i presupposti: \textbf{vulnerabilità attivamente sfruttata}
    e \textbf{incidente grave}, cioè \textbf{due binari di notifica CRA}.
  \item C'è inoltre, senza dubbio, un \textbf{data breach ai sensi del GDPR}.
\end{itemize}

\subsection{Task 2 -- Gli orologi: CRA e GDPR}

\paragraph{Consapevolezza ai fini del CRA.}
\begin{itemize}
  \item \textbf{Vulnerabilità attivamente sfruttata}: in aula (professore e Gianluca concordi) si è
    adottata la lettura prudente, che fa coincidere la consapevolezza con la prima evidenza
    disponibile all'azienda: l'alert di giovedì 5/11 alle 23:10. Lettura alternativa: venerdì 6/11
    alle 12:20, quando l'email del ricercatore fornisce l'evidenza affidabile dello sfruttamento.
  \item \textbf{Incidente grave}: venerdì 6/11 alle 16:00 (riunione di crisi, i log confermano).
  \item \textbf{Mitigazione disponibile}: domenica 8/11 alle 20:00.
\end{itemize}

\paragraph{Consapevolezza ai fini del GDPR.}
Il titolare è consapevole quando ha un ragionevole grado di certezza che \textbf{dati personali}
siano stati compromessi. Qui le opzioni (5 o 6 novembre) distano circa un giorno e la scelta è poco
delicata. Diventa invece delicatissima quando, per esempio, l'attacco avviene il 1° e il CEO viene
informato l'8:
\begin{itemize}
  \item per sostenere una data di consapevolezza posticipata servono \textbf{evidenze documentali} del
    momento in cui il titolare è stato informato: la \textbf{tracciabilità} nella gestione
    dell'incidente è essenziale;
  \item senza evidenze si rischia una \textbf{dichiarazione falsa} al Garante, punita con la
    reclusione da sei mesi a tre anni;
  \item \textbf{posizione di Gianluca}: l'orologio GDPR parte quando si è consapevoli della
    compromissione di dati personali. L'attacco può essere in corso da giorni, settimane o mesi
    (ricognizione, movimento laterale, ecc.): queste fasi fanno parte dell'incidente per CRA, DORA e
    NIS2, ma non rilevano per l'orologio GDPR finché non toccano dati personali;
  \item \textbf{aggiunta del professore}: se l'attacco era evidente e qualcuno (es.\ il CISO) l'ha
    nascosto pensando di gestirlo da solo, è più prudente non far partire l'orologio dalla data in
    cui è stato informato il CEO, perché un ritardo così lungo non sarebbe credibile. Un conto è un
    attacco occulto rimasto ignoto per mesi, altro è un attacco evidente tenuto nascosto.
\end{itemize}

\begin{center}
\textbf{Tabella delle scadenze}\par\smallskip
\small
\begin{tabularx}{\linewidth}{l X l l}
\toprule
\textbf{Binario} & \textbf{Obbligo} & \textbf{Decorrenza} & \textbf{Scadenza} \\ \midrule
CRA  & Early warning vulnerabilità (14(2)(a))     & Gio 5/11 23:10       & Ven 6/11 23:10 \\
CRA  & Notifica vulnerabilità (14(2)(b))          & Gio 5/11 23:10       & \textbf{Dom} 8/11 23:10 \\
CRA  & Final report vulnerabilità (14(2)(c))      & Dom 8/11 20:00 (fix) & \textbf{Dom} 22/11 20:00 \\
CRA  & Early warning incidente grave (14(4)(a))   & Ven 6/11 16:00       & \textbf{Sab} 7/11 16:00 \\
CRA  & Notifica incidente (14(4)(b))              & Ven 6/11 16:00       & Lun 9/11 16:00 \\
CRA  & Final report incidente (14(4)(c))          & Notifica 14(4)(b)    & $\approx$ Mer 9/12 \\
CRA  & Informare gli utenti (14(8))               & Consapevolezza       & Tempestivamente \\
GDPR & Notifica al Garante (art.~33), lettura prudente & Gio 5/11 23:10  & \textbf{Dom} 8/11 23:10 \\
GDPR & Notifica al Garante (art.~33), lettura tardiva  & Ven 6/11 16:00  & Lun 9/11 16:00 \\
GDPR & Comunicazione alle famiglie (art.~34)      & Consapevolezza       & Senza ingiustificato ritardo \\
\bottomrule
\end{tabularx}
\end{center}

\noindent Canale CRA: Single Reporting Platform ENISA, verso il CSIRT dello Stato dello stabilimento
principale (Italia). L'early warning deve indicare gli Stati membri in cui il prodotto è messo a
disposizione (Italia, Francia, Germania, Spagna).

\noindent\textbf{Quanto costa ``aspettare lunedì''}: entro lunedì mattina sono già scaduti l'early
warning sulla vulnerabilità, la notifica della vulnerabilità, l'early warning sull'incidente grave e,
nella lettura prudente, anche le 72 ore GDPR.

\subsection{Task 3 -- Severità del breach (metodologia ENISA 2013)}

\[
SE = DPC \times EI + CB
\]
Scale e valori sono riepilogati nella scheda di ripasso (Capitolo~\ref{cap:schede}).

\paragraph{Gruppo A -- 31.380 famiglie (solo dati di profilo).}
Nome, anno di nascita, foto del bambino; nome, email, cellulare del genitore.
DPC = 1 (dati semplici), EI = 1 (massima), CB = 0,5 (riservatezza, destinatari ignoti) + 0,5
(intento malevolo: vendita su Telegram).
\[
SE_A = 1 \times 1 + 1 = 2 \quad\Rightarrow\quad \text{media}
\]

\paragraph{Gruppo B -- 9.870 famiglie (profilo e posizione).}
Profilo, posizione live, storico di 30 giorni, zone ``Casa'' e ``Scuola''; 37 bambini contattati.
DPC = 2 (dati comportamentali), EI = 1, CB = 0,5 + 0,5.
\[
SE_B = 2 \times 1 + 1 = 3 \quad\Rightarrow\quad \text{alta (parte bassa della fascia)}
\]

\paragraph{Critica alla metodologia.} Il professore la definisce ``folle'' su questo punto: la
metodologia ENISA, con un'impostazione pensata per il settore finanziario, \textbf{non distingue tra
minori e adulti}. Gianluca precisa che, in senso stretto, i dati di profilo sono DPC~1, ma il valore
può essere \textbf{aumentato motivatamente} in base al numero di interessati e alla loro
vulnerabilità (\emph{accountability}).

\paragraph{Decisioni.}
\begin{itemize}
  \item \textbf{Art.~33}: notifica al Garante \textbf{obbligatoria}. Il Garante italiano è l'autorità
    capofila (stabilimento principale a Torino, trattamento transfrontaliero: art.~56).
  \item \textbf{Art.~34}: comunicazione agli interessati del Gruppo~B \textbf{sì}, anche se il punteggio
    (3) consentirebbe di discuterne e il CEO probabilmente spingerebbe per evitarla. Le ragioni:
    \begin{itemize}
      \item trasparenza;
      \item interessati estremamente vulnerabili: i dati possono interessare criminali, inclusi pedofili;
      \item secondo Gianluca, lavoratori, malati e bambini sono categorie particolarmente seguite
        dalle autorità di controllo, quindi il rischio va valutato al livello più alto possibile.
    \end{itemize}
  \item \textbf{A chi si comunica}: ai \textbf{genitori} (titolari della responsabilità genitoriale),
    non ai bambini.
\end{itemize}

\subsection{Task 4 -- Piano di remediation (da svolgere a casa)}

Per mancanza di tempo il task è rimasto da svolgere individualmente con l'Incident Desk.
Il piano va articolato su quattro assi:
\begin{itemize}
  \item \textbf{tecnico}: chiudere le due falle e fare in modo che il prossimo componente obsoleto
    venga individuato in giorni, non in mesi;
  \item \textbf{organizzativo}: chi decide alle 23:10 di un giovedì, e quali template devono già esistere;
  \item \textbf{formazione}: almeno tre azioni, ciascuna con destinatari, contenuto e KPI, collegate a
    un inject;
  \item \textbf{comunicazione}: famiglie, tutti gli utenti, retailer, stampa, ricercatore, rettifica
    del post Instagram.
\end{itemize}

\begin{esame}
Serve un \textbf{piano di risposta agli incidenti integrato} che coordini contemporaneamente CRA,
GDPR e, dove applicabile, NIS2 (tema della prossima lezione).
\end{esame}

% =====================================================================
\section{Prossimo laboratorio (venerdì): il ruolo dell'auditor}
% =====================================================================

Presentato da Gianluca. Cambia il ruolo: non si gestisce più l'incidente dall'interno, si è
\textbf{auditor esterno}.

\paragraph{Scenario.}
\begin{itemize}
  \item \textbf{Sestante Mobility}, \textbf{media impresa} (non micro né piccola: si applicano tutti
    gli obblighi previsti per le medie imprese).
  \item Prodotto: box telematico installato su veicoli di flotte aziendali, collegato a una porta del
    veicolo da cui estrae dati, con piattaforma cloud (SaaS) e app mobile.
  \item Clienti: \emph{fleet operator} in diversi Paesi.
  \item Dati raccolti: viaggi, eventi di guida e un punteggio del conducente calcolato dal dispositivo.
    Il conducente ha nome e cognome, quindi si tratta di \textbf{dati personali}.
\end{itemize}

\paragraph{Eventi (date nel futuro).}
\begin{itemize}
  \item \textbf{Agosto 2027}: un ricercatore segnala che il firmware accetta manomissioni.
  \item \textbf{Settembre 2027}: nei dispositivi ricondizionati (ritirati da ex clienti e riassegnati
    a nuovi clienti) restano dati del cliente precedente.
  \item \textbf{Novembre 2027}: un conducente scopre di essere tracciato 24 ore su 24, anche fuori
    dall'orario di lavoro (i veicoli restano in uso al conducente), e il cliente non riesce a
    disattivare il tracciamento.
\end{itemize}

\paragraph{Domande a cui rispondere.}
\begin{itemize}
  \item Si può mantenere il prodotto sul mercato, in particolare le unità vendute dopo il
    \textbf{12 dicembre 2027}?
  \item Come gestire un cliente che rinnova parte dei dispositivi entro il \textbf{24 aprile 2028}?
    Le unità sono conformi?
  \item Un nuovo dispositivo non ancora sul mercato: è valutabile? Senza evidenze, l'auditor deve
    dichiararlo fuori perimetro o non valutabile.
  \item Che ruolo ha Sestante ai fini GDPR: titolare, responsabile o entrambi?
\end{itemize}

\begin{esame}
\textbf{Regole metodologiche dell'audit:}
\begin{itemize}
  \item le date di cut-off si applicano alla \textbf{singola unità} immessa sul mercato, non al
    modello: unità dello stesso modello possono essere conformi o no a seconda della data di immissione;
  \item \textbf{una fonte} = asserzione, da prendere con riserva; \textbf{due fonti concordanti} =
    fatto accertato (evidenza);
  \item \textbf{un certificato ISO non implica conformità CRA né GDPR}: il GDPR è una norma giuridica,
    la ISO uno standard tecnico di sicurezza delle informazioni;
  \item la scala di severità dell'audit è diversa da quella ISO (spiegata nel playbook);
  \item il perimetro va preso sul serio: \textbf{non inventare né presumere nulla}. Ciò che non è
    nelle evidenze non esiste.
\end{itemize}
\end{esame}

\paragraph{Materiale e consegna.} 16 documenti brevi (poche frasi ciascuno) e le trascrizioni di due
brevi interviste (tra cui il CTO), con un ordine di analisi suggerito: prima disclosure, configurazione
e impostazioni di default del prodotto, poi le informazioni di secondo livello. Briefing e playbook
vengono caricati in anticipo; la piattaforma online è attiva solo il giorno del laboratorio.
Deliverable: un testo breve, da compilare online e caricare come file di testo secondo le istruzioni.

% ---------------------------------------------------------------------
\section{Domande aperte e punti da verificare}

\begin{verifica}
\begin{itemize}
  \item \textbf{Data di consapevolezza CRA}: la trascrizione fa dire alla classe ``4 novembre'', ma il
    4 è il giorno di inizio dello sfruttamento, di cui l'azienda non sapeva nulla. Qui si è usato
    l'alert del 5/11 alle 23:10. Confermare con i propri appunti.
  \item \textbf{Tabella delle scadenze}: in aula il calcolo non è stato completato; le date sono
    calcolate applicando art.~14 CRA e art.~33 GDPR alle decorrenze scelte.
  \item \textbf{Art.~64(10) CRA}: per micro e piccole imprese non si applicano sanzioni per il solo
    mancato rispetto del termine di 24 ore dell'early warning. Verificare sul testo ufficiale.
  \item \textbf{Art.~14(8) CRA}: l'obbligo di informare gli utenti non ha un termine in ore.
    Verificare sul testo ufficiale.
  \item Cognome del co-docente Gianluca (non ricavabile con certezza dalla trascrizione).
\end{itemize}
\end{verifica}

% =====================================================================
\lezione{06/11/2026}{Laboratorio 4: conformity audit CRA--GDPR (Sestante Mobility)}
  {Gianluca Gilardi}
  {lavoro di gruppo sul caso Sestante Mobility; audit esterno di un fabbricante;
   registro dei finding, opinione e management letter; consegna tramite workbench online}
\label{lez:2026-11-06}
% =====================================================================

\begin{nota}[Fonte]
Trascrizione automatica della registrazione (lezione~14): la sessione è quasi interamente silenzio,
perché i gruppi lavorano in autonomia. Si sentono solo l'apertura, due avvisi sul tempo residuo e i
saluti finali. Il contenuto del caso, la consegna e i numeri provengono dal materiale del laboratorio:
\emph{Conformity\_Audit\_Playbook} e slide \emph{Conformity\_Audit\_Exercise\_Briefing}. Lo
svolgimento completo è nel Capitolo~\ref{lab:04}.
\end{nota}

\section{Sintesi}
Quarto laboratorio, presentato da Gianluca al termine della lezione del 02/11
(Capitolo~\ref{lez:2026-11-02}). Gli studenti lavorano da \textbf{auditor esterni} nominati dal
comitato di audit di \textbf{Sestante Mobility S.p.A.}, una media impresa che produce un box
telematico per flotte aziendali. Partendo da una \emph{evidence room} di 16 documenti e da due
interviste, devono stabilire se il fascicolo tecnico consente di tenere il prodotto sul mercato ai
sensi del CRA e quale ruolo ha l'azienda ai fini del GDPR. Il lavoro si chiude con un registro dei
finding, un'opinione di audit e una management letter. Rispetto alla simulazione Pulse K cambia la
prospettiva: non si gestisce un incidente dall'interno, si valuta a posteriori ciò che l'azienda può
dichiarare di sé.

\section{Consegna}
\begin{itemize}
  \item Si lavora sul \textbf{Conformity Audit Workbench}, uno strumento online che ricalca i cinque
    passi del playbook (scope, evidence room, registro, opinione, testo finale). Il deliverable è il
    testo prodotto al passo~5, copiato con l'apposito pulsante in un file di testo semplice, con il
    nome del gruppo e dei membri in cima. Un file per gruppo, caricato entro la fine della sessione.
  \item Il file ha tre blocchi: \textbf{A}~risposte di scope (quali obblighi CRA per quali unità;
    ruolo GDPR per tre flussi di dati), \textbf{B}~registro di almeno quattro finding di livello
    Medium o superiore, con un controllo che funziona e almeno un criterio esaminato e scartato,
    \textbf{C}~opinione e management letter. Bonus: mappatura di ogni finding su un controllo
    dell'Annex~A.8 della ISO/IEC~27001.
\end{itemize}

\section{Indicazioni udibili in aula}
\begin{itemize}
  \item Gianluca dà per scontato che il playbook sia stato studiato prima del laboratorio; il link
    allo strumento è stato condiviso in chat.
  \item \textbf{Nessuno deve completare l'audit}: l'obiettivo è arrivare ad \textbf{almeno quattro
    finding}, che per il docente sono ``più che sufficienti''.
  \item Il risultato può essere giusto o sbagliato e tutto è discutibile; ciò che non si discute è la
    \textbf{logica}. Serve un percorso logico coerente, come già suggerito dopo il laboratorio
    precedente.
  \item Le domande si fanno in chat; il microfono viene aperto su richiesta.
  \item Chiusura prevista verso le 13:15, come nel laboratorio precedente.
  \item La discussione degli elaborati era annunciata per la lezione del lunedì successivo. Il 09/11
    però nessun gruppo ha presentato e le presentazioni sono state rinviate
    (Capitolo~\ref{lez:2026-11-09}).
\end{itemize}

\begin{esame}
Il laboratorio applica le regole di metodo annunciate il 02/11: le date di applicazione del CRA
valgono \textbf{per singola unità} e non per modello; \textbf{una fonte è un'asserzione, due fonti
concordanti sono un fatto}; un certificato ISO/IEC~27001 non prova la conformità né al CRA né al GDPR;
ciò che non è nelle evidenze non si presume. Sono competenze direttamente utili per le domande
d'esame sulla qualificazione di un caso.
\end{esame}

\section{Domande aperte e punti da verificare}
\begin{verifica}
\begin{itemize}
  \item L'orario di chiusura (``1.15'') e i due avvisi sul tempo residuo sono ricavati dalla
    trascrizione; le slide indicano per il laboratorio la fascia 11:30--13:00.
  \item Il debrief in aula (cinque domande indicate nelle slide di briefing) non risulta ancora
    svolto: verificare se e quando verrà ripreso.
  \item Cognome del co-docente Gianluca: l'indirizzo di consegna riportato nel materiale del
    laboratorio fa pensare a ``Gilardi'', ma non è confermato in aula.
\end{itemize}
\end{verifica}

% =====================================================================
\lezione{09/11/2026}{Hard law e soft law; Cybersecurity Act (CSA)}
  {Prof.~Giuseppe Vaciago}
  {le quattro regole dell'esame; hard law vs soft law; caso del Google Form per la cena di Natale;
   persona autorizzata vs responsabile del trattamento; Cybersecurity Act: ENISA, quadro europeo di
   certificazione, schemi, livelli di affidabilità, EUCC; certificazione continua; security by design;
   rischio della supply chain; CSA~2}
\label{lez:2026-11-09}
% =====================================================================

\begin{nota}[Fonte]
Trascrizione automatica della registrazione, revisionata. Le slide sul Cybersecurity Act non sono
ancora nella cartella \emph{material}: date, articoli e nomi di organismi sono ricostruiti dal parlato
e, dove serve, segnalati nei riquadri \emph{verifica}. Non ci sono state presentazioni dei gruppi:
nessuno si è proposto e sono state rinviate.
\end{nota}

\section{Sintesi}
Il docente apre ripetendo le ``regole del gioco'' per l'esame e introduce la distinzione, per lui
fondamentale, tra \textbf{hard law} e \textbf{soft law}, con un caso pratico (dati sanitari raccolti
con un Google Form personale). Poi presenta il \textbf{Cybersecurity Act}: il regolamento che rafforza
\textbf{ENISA} e crea il \textbf{quadro europeo di certificazione} della cybersicurezza per prodotti,
servizi e processi ICT. È la soft law europea che aiuta a dimostrare il rispetto della hard law, in
particolare del CRA. Chiudono la lezione la certificazione come processo continuo, la security by
design e la dipendenza dai fornitori. Il caso di studio sul CSA è rinviato alla lezione successiva.

% ---------------------------------------------------------------------
\section{Organizzazione}
\begin{itemize}
  \item \textbf{Presentazioni del laboratorio}: può presentare chi ha preso più di 25 nella
    valutazione del lab (valutazione fatta da Gianluca); non è obbligatorio, ma chi viene invitato e non
    vuole presentare deve avvisare prima. Il docente le consiglia come esercizio per il lavoro (presentare
    a un pubblico o in una call). Oggi nessuno era pronto; non c'è valutazione sulla presentazione.
  \item Il lunedì successivo si valuterà se proseguire con i laboratori (il docente percepisce interesse
    per i punti del lab, ma difficoltà con le presentazioni).
  \item \textbf{Venerdì 13/11}: lezione frontale su \textbf{NIS2}, \textbf{online} (Gianluca non sarà in
    Italia). Era prevista il 19/10 e va recuperata per arrivare a dicembre senza ritardi.
  \item \textbf{Venerdì 20/11}: nuovo laboratorio.
  \item Prossima lezione: caso di studio sul CSA (punti di vista di un ospedale, di un fornitore e di
    un valutatore tecnico). Le slide si possono già guardare.
\end{itemize}

% ---------------------------------------------------------------------
\section{Le regole del gioco per l'esame}
Il corso ha visto GDPR, risk assessment ISO/IEC~27001 e CRA; oggi il CSA, poi restano NIS2, DORA e
AI Act. La confusione a fine corso è normale, perché non siamo abituati a lavorare sulle norme. Le
regole:
\begin{enumerate}
  \item conoscere \textbf{definizioni e ambito} di ogni norma;
  \item conoscere le \textbf{tempistiche di notifica} degli incidenti: a fine corso il docente pubblicherà
    sulla piattaforma una tabella con tutte le scadenze (vedi Capitolo~\ref{cap:scadenze});
  \item tutte queste norme seguono un \textbf{approccio basato sul rischio}: niente errori su rischio
    inerente e misure di mitigazione (lab di risk assessment con Gianluca);
  \item distinguere \textbf{hard law} e \textbf{soft law}.
\end{enumerate}
Saper gestire un incidente con una metodologia di severity assessment (quella ENISA è la più nota)
non serve solo all'esame: nel contesto attuale, che il docente descrive vicino a una ``guerra
ibrida'', gestire le notifiche è decisivo anche sul piano della compliance.

\begin{esame}
\begin{itemize}
  \item Ambito della norma e timeline di notifica sono il cuore dell'esame.
  \item La differenza hard law / soft law ``non è una domanda d'esame'', ma serve in pratica per
    rispondere alle domande.
\end{itemize}
\end{esame}

% ---------------------------------------------------------------------
\section{Hard law e soft law}
\paragraph{Hard law.} Norma vincolante con conseguenze \textbf{predeterminate}. Esempio: il limite di
velocità in autostrada è fissato dal codice della strada; chi gestisce l'autostrada (ANAS, Autostrade
per l'Italia) non può cambiarlo. Se vado a 160 la sanzione è definita dalla legge (eventualmente con
differenze tra livelli nazionale e regionale): non può essere decisa ``a simpatia''. Anche nella
hard law c'è un margine di discrezionalità (nel penale, una pena da 5 a 10 anni), ma è delimitato.

\paragraph{Soft law.} Regole interne, policy, linee guida, standard. Il gestore dell'autostrada può
decidere nuova segnaletica e nuovi sistemi di sicurezza: è una regola interna. In cybersecurity:
\begin{itemize}
  \item la policy aziendale sull'\textbf{autenticazione a più fattori}: nessuna legge la impone (l'art.~32
    GDPR e il principio di \emph{accountability} chiedono misure adeguate, non l'MFA);
  \item le \textbf{linee guida ENISA} (approccio, indicazioni, policy).
\end{itemize}
La violazione di una soft law può avere conseguenze, ma \textbf{non previste dalla legge}: l'autorità
o l'azienda hanno un ampio margine di discrezionalità nel decidere la sanzione.

\paragraph{Più hard law o più soft law in futuro?} Gli studenti rispondono ``più hard law'' (l'Europa
ne produce molta). Il docente concorda solo in parte: la hard law fissa i principi fondamentali, ma
\textbf{la sostanza della cybersecurity è la soft law}, che ogni organizzazione deve creare; anche
CRA e GDPR non si applicano senza soft law. Una legge richiede uno--tre anni (in Italia anche cinque),
troppo per una realtà che cambia di continuo, soprattutto con l'AI: secondo lui la hard law diventerà
una sorta di \emph{kill switch} (pochi principi invalicabili, la cui violazione impone di ``spegnere''),
il resto sarà soft law.
\begin{nota}[Le policy in azienda]
Un'organizzazione ha policy che vengono da GDPR, NIS2, DORA, AI Act\ldots: non si può gestire solo la
``propria'' (es.\ solo DORA perché si è ICT risk manager). Serve governare l'interconnessione tra le
policy. Battuta del docente: lasciare le policy di sicurezza solo all'ufficio legale o all'HR è il modo
migliore per renderle inapplicabili.
\end{nota}

% ---------------------------------------------------------------------
\section{Caso: il Google Form per la cena di Natale}
\paragraph{Scenario.} Una grande azienda organizza la cena di Natale ufficiale. Un dipendente crea con il
\textbf{proprio account Google personale} un Google Form per raccogliere partecipanti, intolleranze e
allergie. La policy aziendale impone gli strumenti Microsoft~365 (Microsoft Forms) anche per i sondaggi.

\paragraph{Prima domanda: è un data breach?}
\begin{itemize}
  \item Con soli dati comuni (nome, cognome) sarebbe un problema GDPR piccolo: è solo l'elenco dei
    dipendenti.
  \item Le \textbf{intolleranze} sono però \textbf{dati relativi alla salute}, raccolti su un account
    privato fuori dal controllo del titolare.
  \item La risposta della classe (``in buona fede non succede niente'') è ragionevole, ma all'esame è
    rischiosa: va fatto il ragionamento di \textbf{severity assessment}. Con dati sanitari il punteggio
    di base del contesto (DPC) parte più alto; c'è una perdita di riservatezza; l'intento non è malevolo.
    Paradosso: questo elenco può avere una severità più alta di una lista di due milioni di nomi.
  \item Conclusione del docente: deve prevalere il ragionamento razionale, e non si notifica al Garante
    solo perché un dipendente ha violato una policy; ma il rischio va comunque valutato
    (\emph{risk-based approach}).
\end{itemize}

\paragraph{Seconda domanda: quali conseguenze per il dipendente?}
La violazione è di una \textbf{policy} (soft law), quindi le conseguenze sono
\textbf{giuslavoristiche}. In Italia i provvedimenti disciplinari sono di fatto cinque (non da sapere
per l'esame): richiamo verbale, richiamo scritto, sospensione, licenziamento con o senza giusta causa.
Licenziare qui è eccessivo, ma una conseguenza va data: un richiamo registrato, con conseguenze più
gravi in caso di recidiva. Il collegamento con la hard law c'è, ma nessuna legge impone Google Forms o
Microsoft Forms: nella soft law il margine è ampio e va usato in modo razionale. Se il dipendente
contesta il provvedimento, decide un giudice del lavoro.

\paragraph{Chi paga le sanzioni GDPR.}
Il GDPR sanziona il \textbf{titolare} (la persona giuridica), non il dipendente. L'azienda può agire sul
dipendente solo sul piano del diritto del lavoro: anche dopo una sanzione di 2 milioni di euro, di norma
non può chiedergli di rimborsarla, ma può licenziarlo. Per i dirigenti (\emph{leadership team}) la
responsabilità personale è più concreta: per questo a quel livello serve una polizza \textbf{D\&O}
(\emph{Directors and Officers}).

\subsection{Dipendente = persona autorizzata, non responsabile del trattamento}
Domanda di uno studente: chi è il \emph{data processor} nel caso? Il docente corregge la confusione:
\begin{itemize}
  \item il \textbf{titolare} (\emph{data controller}) è l'azienda;
  \item il \textbf{dipendente} è una \textbf{persona autorizzata} al trattamento (art.~29 GDPR): agisce
    sotto l'autorità del titolare, dentro la sua organizzazione;
  \item un consulente esterno (es.\ una Big Four come KPMG, Deloitte, PwC) che svolge la \textbf{stessa}
    attività di un dipendente è invece \textbf{responsabile del trattamento} (\emph{data processor},
    art.~28), perché è un soggetto esterno.
\end{itemize}
\begin{esame}
Non confondere \textbf{responsabile del trattamento} (soggetto esterno che tratta per conto del titolare)
e \textbf{persona autorizzata} (dipendente del titolare). In caso di dubbi, rivedere le slide sul GDPR.
\end{esame}

% ---------------------------------------------------------------------
\section{Cybersecurity Act: perché un quadro europeo}
Il CSA è l'introduzione della \textbf{soft law dell'Unione europea} in materia di cybersicurezza. Ha
\textbf{due pilastri}:
\begin{enumerate}
  \item i \textbf{poteri e il mandato di ENISA};
  \item il \textbf{quadro europeo di certificazione} (schemi di certificazione).
\end{enumerate}
``Ricordate: \emph{scheme} ed ENISA.'' Il docente ammette che finora il successo del CSA è limitato, ma
si aspetta che cresca.

\paragraph{Il problema.} Standard e certificazioni erano nazionali. Una multinazionale deve rispettare
regole diverse in ogni paese (l'esempio di Google Street View in 140 paesi). Per la cybersicurezza i
problemi sono quattro:
\begin{itemize}
  \item \textbf{interdipendenza}: un fornitore può colpire molti clienti, anche in paesi diversi;
  \item \textbf{asimmetria informativa}: il venditore conosce il prodotto meglio dell'acquirente;
  \item \textbf{frammentazione}: schemi nazionali diversi;
  \item \textbf{fiducia verificabile} (\emph{verifiable trust}): le dichiarazioni di sicurezza richiedono
    requisiti ed \textbf{evidenze}. È il punto più importante del CSA.
\end{itemize}

\paragraph{Il parallelo con ISO/IEC~27001.} ISO (organizzazione svizzera) ha creato lo standard 27001,
che certifica la sicurezza dell'\textbf{organizzazione}. Il CSA fa la stessa cosa a livello europeo,
con ENISA e la Commissione, ma per i \textbf{prodotti} (e servizi e processi ICT): uno smartphone, un
firewall. L'azienda si certifica 27001 (o 42001 per l'AI) per dimostrare di avere un'organizzazione
conforme; il fabbricante usa uno schema del CSA per dimostrare la sicurezza del prodotto.

\paragraph{Il collegamento con il CRA.} Prodotti $\Rightarrow$ CRA. Il CRA dice quali regole rispettare,
ma non come; gli schemi del CSA danno requisiti concreti che offrono \textbf{la possibilità, non la
certezza}, di dimostrare la conformità al CRA. Il CSA non è obbligatorio.
\begin{esame}
\begin{itemize}
  \item CSA = \textbf{soft law} europea; CRA = \textbf{hard law} sui prodotti. Il CSA aiuta a rispettare il
    CRA, ma non lo sostituisce.
  \item ISO/IEC~27001 = organizzazione; CSA = prodotti, servizi e processi ICT.
  \item Oggi in Europa \textbf{non esistono certificazioni di cybersicurezza obbligatorie} (resta da capire
    il ruolo della marcatura CE).
\end{itemize}
\end{esame}

\paragraph{Tre concetti.} \textbf{Cybersecurity}: proteggere sistemi, reti e informazioni.
\textbf{Cyber resilience}: mantenere le funzioni essenziali o ripristinarle dopo un incidente.
\textbf{Cyber assurance} (da non confondere con la \emph{cyber insurance}): creare \textbf{evidenze} che
certi requisiti sono soddisfatti; è il vero obiettivo del CSA. Il limite della soft law resta
l'incertezza delle conseguenze.

% ---------------------------------------------------------------------
\section{ENISA e la governance della certificazione}
\textbf{ENISA} è l'Agenzia dell'Unione europea per la cybersicurezza (agenzia indipendente, legata alla
Commissione). Compiti:
\begin{itemize}
  \item \textbf{policy}: linee guida, formazione ed esercitazioni;
  \item \textbf{cooperazione operativa}: aiuta i CSIRT nazionali a lavorare insieme (con il
    \emph{Digital Omnibus} si va verso la centralizzazione delle notifiche presso ENISA);
  \item \textbf{schemi candidati}: prepara le regole tecniche di valutazione (compito centrale nel CSA);
  \item \textbf{conoscenza delle minacce}: raccoglie e condivide analisi e buone pratiche.
\end{itemize}
Opinione del docente: le norme (CRA, Digital Omnibus) danno a ENISA molti compiti, ma ha poco personale.
Il CSA~2 la rafforzerà ancora.

\paragraph{Gli attori.}
\begin{center}
\small
\begin{tabularx}{\linewidth}{>{\raggedright\arraybackslash}p{4.2cm} X}
\toprule
\textbf{Attore} & \textbf{Ruolo}\\
\midrule
ENISA & prepara gli schemi candidati\\
ECCG (\emph{European Cybersecurity Certification Group}) & rappresentanti delle autorità nazionali;
  coopera alla preparazione ed esprime un parere sulla bozza\\
Commissione europea & individua il bisogno di uno schema e lo \textbf{adotta}\\
Autorità nazionali di certificazione & vigilano sulla certificazione; in Italia l'\textbf{ACN} (Agenzia per la
  Cybersicurezza Nazionale)\\
Laboratori e organismi di certificazione & eseguono i test e rilasciano i certificati\\
\bottomrule
\end{tabularx}
\end{center}
Opportunità professionali: esperti esterni per ENISA o la Commissione (bandi sul sito), rappresentanti
nazionali nell'ECCG (non solo da ACN: il docente cita una rappresentante che lavora al Garante privacy).

\paragraph{Il processo di uno schema.} \textbf{Richiesta} (la Commissione individua il bisogno per una
categoria di prodotti e chiede a ENISA uno schema candidato) $\rightarrow$ \textbf{bozza} (ENISA definisce
requisiti e metodi di valutazione e consulta gli stakeholder e gli esperti) $\rightarrow$
\textbf{parere} dell'ECCG $\rightarrow$ \textbf{adozione} da parte della Commissione.

\paragraph{Framework, schema, certificato.}
\begin{itemize}
  \item \textbf{Framework} = il CSA: definisce ruoli e regole comuni. \textbf{Non contiene schemi}: chi cerca
    uno schema nel testo del CSA ``non ha capito''.
  \item \textbf{Schema} = il \emph{rule book} della valutazione, creato da ENISA con il processo sopra:
    ambito, requisiti di sicurezza, test e \textbf{livelli di affidabilità}. Il primo e più importante è
    l'\textbf{EUCC}.
  \item \textbf{Certificato} = attesta che un prodotto rispetta lo schema. È l'obiettivo concreto
    dell'azienda: l'unico modo per ridurre l'incertezza creata dalla hard law (``devi essere sicuro'': come?).
\end{itemize}
Uno schema non è un certificato, e un certificato non copre tutta l'azienda.

\paragraph{Che cosa può coprire uno schema.}
\begin{itemize}
  \item \textbf{Prodotti ICT}: elemento o componente digitale valutabile (firewall, chip, software); non
    solo prodotti di sicurezza;
  \item \textbf{servizi ICT}: servizi che memorizzano, trasmettono o elaborano informazioni (es.\ cloud storage);
  \item \textbf{processi ICT}: attività per progettare o mantenere prodotti e servizi; per ora nessuno schema;
  \item \textbf{servizi di sicurezza gestiti}: svolti per un cliente (incident response, audit, penetration test).
\end{itemize}
Processi e servizi gestiti toccano il terreno della 27001 (organizzazione). Formalmente ISO certifica
l'organizzazione e il CSA prodotti e servizi, ma per il docente la cooperazione tra ISO ed ENISA
(con sede in Grecia) è decisiva: due organizzazioni che si contendono la governance creerebbero un
\emph{single point of failure} proprio dove serve resilienza.

% ---------------------------------------------------------------------
\section{Livelli di affidabilità e autovalutazione}
Il livello di affidabilità (\emph{assurance level}) esprime la fiducia che i requisiti siano soddisfatti,
in proporzione al rischio:
\begin{center}
\small
\begin{tabular}{lll}
\toprule
\textbf{Livello} & \textbf{Rischio considerato} & \textbf{Chi valuta}\\
\midrule
Base (\emph{basic}) & rischi noti di base & anche \textbf{autovalutazione} del fabbricante\\
Sostanziale (\emph{substantial}) & attaccanti con risorse limitate & terza parte indipendente\\
Elevato (\emph{high}) & attaccanti con risorse significative & terza parte indipendente\\
\bottomrule
\end{tabular}
\end{center}
L'autovalutazione è possibile solo per il livello base, ma in pratica oggi non ci sono esempi: la
certificazione è quasi sempre di \textbf{terza parte}. (``È come fare i compiti e dire: passerò l'esame.
Poi il professore non sempre è d'accordo.'')

\paragraph{Come funziona in pratica (esempio 27001).} Una società di consulenza come Bureau Veritas aiuta
l'azienda a soddisfare i requisiti; poi un organismo di certificazione distinto (in Italia il docente
cita CEPAS) verifica e certifica; il tutto fa capo a ISO. Per il CSA le stesse società potranno operare
sugli schemi europei, riferendosi a ENISA; il sistema si costruirà nei prossimi 2--4 anni. Il docente
invita a vederlo come un'opportunità professionale, anche imprenditoriale.

% ---------------------------------------------------------------------
\section{Che cosa fa (e non fa) un certificato}
\begin{itemize}
  \item Dà \textbf{evidenze} su un prodotto o servizio \textbf{identificato con precisione}, su requisiti e
    test specifici, a un livello di affidabilità definito.
  \item \textbf{Non prova che non esistano vulnerabilità}: ho il requisito MFA, non per questo sono al
    sicuro (oggi gli attacchi aggirano anche l'MFA).
  \item Il certificato aiuta a evitare sanzioni per violazione del CRA, ma \textbf{non in automatico}: dopo
    un incidente si verifica, secondo il principio di accountability, se i requisiti certificati erano
    davvero rispettati.
\end{itemize}

\subsection{EUCC}
\textbf{EUCC} (\emph{EU Common Criteria-based scheme}): primo schema adottato sulla base del CSA.
\begin{itemize}
  \item valuta \textbf{prodotti e componenti ICT}: niente servizi, processi o management, quindi non è in
    conflitto con ISO;
  \item si basa sui \textbf{Common Criteria} (criteri internazionali per i requisiti di sicurezza);
  \item applicabile dal \textbf{27 febbraio 2025};
  \item livelli di affidabilità \textbf{sostanziale ed elevato} (un livello base in autovalutazione non
    avrebbe senso).
\end{itemize}
Elementi della valutazione: \textbf{target of evaluation} (TOE, il prodotto esatto da certificare);
\textbf{security target} (documento che descrive il TOE e i requisiti di sicurezza dichiarati);
\textbf{protection profile} (insieme riutilizzabile di requisiti per un tipo di prodotto); metodi, analisi
delle vulnerabilità, ipotesi e limiti operativi. Un certificato per la versione 3 di un modulo
\textbf{non copre} automaticamente la versione 4 in un'altra configurazione.

\paragraph{Il problema dei prodotti.} L'organizzazione cambia persone e strumenti, ma processi e
sistemi di autenticazione e autorizzazione (la base della 27001) restano stabili. I prodotti invece
cambiano a ogni release (iPhone 15, 17, 18\ldots): ogni prodotto va ricertificato. È un'opportunità
economica, ma pesa sulle piccole aziende. Previsione del docente: solo alcune aziende useranno l'EUCC, e
sarà un segno di qualità (come Parmigiano Reggiano e ``parmesan'': conta il processo, che richiede tempo).

\subsection{La certificazione continua dopo il rilascio}
Il certificato si riferisce a un prodotto definito, ma prodotto e minacce cambiano: valutare,
certificare, \textbf{monitorare e mantenere} (controllare l'oggetto, registrare i risultati, osservare i
cambiamenti, rivederne l'impatto). Una patch di sicurezza non invalida automaticamente il certificato,
ma il suo \textbf{impatto va valutato}. È lo stesso monitoraggio continuo di DPIA e FRIA (AI Act).

\subsection{Strumenti diversi per domande diverse}
\begin{itemize}
  \item \textbf{ISO/IEC~27001}: sistema di gestione della sicurezza dell'organizzazione;
  \item \textbf{EUCC}: prodotti;
  \item \textbf{marcatura CE}: il prodotto rispetta le norme UE che la richiedono. Timore del docente:
    EUCC e marcatura CE potrebbero sovrapporsi per i prodotti ICT, imponendo un doppio percorso
    (difficile da chiedere, ad esempio, a un produttore cinese);
  \item \textbf{penetration test}: attacco simulato autorizzato, anch'esso soft law con possibili conseguenze
    legali; probabilmente confluirà in EUCC o nella parte VAPT della 27001.
\end{itemize}
Riepilogo per norma: CSA $\rightarrow$ le dichiarazioni di sicurezza si valutano con un
\textbf{quadro comune di certificazione}; NIS2 $\rightarrow$ reti e sistemi informativi; CRA
$\rightarrow$ obblighi sulle vulnerabilità dei prodotti con elementi digitali; GDPR $\rightarrow$ la norma
più generale; DORA $\rightarrow$ settore finanziario; poi AI Act.

% ---------------------------------------------------------------------
\section{Security by design e supply chain}
\paragraph{Security by design.} Protezione e manutenzione si pianificano quando il prodotto è progettato
(come privacy by design, forensics by design, ``legal by design''): ridurre la superficie d'attacco,
\textbf{sapere cosa contiene il software} (\textbf{SBOM}, \emph{software bill of materials}), gestire le
segnalazioni di vulnerabilità e distribuire gli aggiornamenti. Con le componenti AI spesso non si sa da
cosa è composto un prodotto. Esempio dal lavoro del docente: in un data breach su un macchinario
sanitario si è scoperto un fornitore canadese a 45 minuti dalla scadenza della notifica.

\paragraph{Rischio della supply chain.} Domande da porsi a ogni nuovo fornitore: il componente può
essere attaccato (sicurezza tecnica)? \textbf{Possiamo continuare a operare senza questo fornitore}
(continuità)? Chi controlla servizi e decisioni chiave (governance)?
\begin{itemize}
  \item Esempi: CrowdStrike (2024); un mega breach finanziario causato dall'unico fornitore di servizi
    KYC/AML (lavoro che nessuno vuole fare, quindi monopolio e poca attenzione alla sicurezza); il caso
    Italiaonline con un fornitore israeliano consapevole di essere insostituibile.
  \item Se c'è concorrenza il fornitore si cambia; senza alternative è un problema di resilienza e
    continuità operativa.
  \item \textbf{Governance}: conta chi controlla i dati, non dove risiedono. Dati in Europa ma controllo in
    un altro Stato: se la capogruppo (es.\ negli USA) cambia la governance può chiudere l'entità europea.
    La \emph{data residency} è un'illusione se il controllo è altrove.
\end{itemize}

\paragraph{CSA~2.} È una \textbf{proposta} di revisione del CSA (come da NIS a NIS2): più poteri a ENISA,
semplificazione della certificazione, più attenzione al rischio della supply chain. Per ora basta sapere
che esiste.

% ---------------------------------------------------------------------
\section{Domande aperte e punti da verificare}
\begin{verifica}
\begin{itemize}
  \item \textbf{Severity del caso Google Form}: il docente non ricordava se i dati sanitari partano da DPC 2 o
    3. Nella metodologia ENISA (slide \emph{03}) la categoria \emph{sensitive data} (salute, opinioni
    politiche\ldots) ha punteggio di base \textbf{4}; controllare anche il valore del fattore
    \emph{malicious intent} nella tabella delle circostanze (CB).
  \item \textbf{Articoli GDPR}: persona autorizzata (art.~29) e responsabile del trattamento (art.~28)
    ricostruiti dalla trascrizione (in quel passaggio quasi inintelligibile): ricontrollare.
  \item \textbf{EUCC}: data di applicazione (27/02/2025) e livelli (sostanziale, elevato) dal parlato; fonte
    attesa: regolamento di esecuzione EUCC e slide CSA quando saranno disponibili.
  \item \textbf{Autovalutazione}: ammessa solo per il livello base, come detto in aula; verificare l'articolo
    del CSA sulle slide.
  \item \textbf{CSA~2}: stato e data della proposta da verificare (``probabilmente in vigore il prossimo anno''
    secondo il docente).
  \item Organismo di certificazione italiano citato come ``CEPAS'' (trascrizione incerta).
\end{itemize}
\end{verifica}


% =====================================================================
\part{Manuale del corso}
% Capitoli tematici, uno per norma: slide + diario + integrazioni segnalate.
% =====================================================================
\chapter{Introduzione e quadro normativo europeo}
\label{cap:introduzione}

\begin{nota}[Fonti]
Lezione introduttiva del 21/09/2026 (Capitolo~\ref{lez:2026-09-21}) e slide \emph{01.~Cybersecurity
Laws and Regulations -- I -- Introductory class: the red thread}. Per hard law e soft law: lezione del
09/11/2026 (Capitolo~\ref{lez:2026-11-09}). Questo capitolo dà la mappa del corso: i dettagli di ogni
norma sono nei capitoli successivi.
\end{nota}

% =====================================================================
\section{Perché un diritto della cybersicurezza}
\label{sec:intro-perche}

Il corso studia come il diritto protegge beni \textbf{immateriali} (dati, sistemi, prodotti, modelli di
IA) e chi risponde quando la protezione fallisce. Per un ingegnere la differenza tra un dato personale e
uno non personale, o tra un incidente e una violazione di dati, può sembrare minima; per il diritto è
enorme, perché da quella qualificazione dipendono obblighi, scadenze, autorità e sanzioni.

Tre ragioni per studiare le norme europee, anche per chi lavorerà fuori dall'Unione:
\begin{itemize}
  \item la regolazione UE \textbf{fa scuola}: il GDPR è stato ripreso da molti ordinamenti (Brasile,
    Svizzera, Singapore tra gli esempi citati), e alcuni Stati USA si avvicinano all'AI Act;
  \item l'approccio europeo basato su \textbf{valutazione del rischio} (\emph{risk assessment}) e
    \textbf{responsabilizzazione} (\emph{accountability}) è ormai il linguaggio comune della sicurezza
    in tutto il mondo;
  \item chi ha competenze tecniche \emph{e} un minimo di compliance può gestire da sé le decisioni
    durante un incidente (notificare o no, con quale motivazione), invece di doverle spiegare a chi non
    conosce l'IT. Gli esempi del docente: un CEO convinto di conoscere la \emph{root cause} di un
    breach, che non era quella reale; un altro che rifiutava di avvisare gli interessati per timore del
    danno reputazionale, rischiando che l'autorità imponesse poi la comunicazione.
\end{itemize}

\paragraph{Diritto e tecnologia.} La privacy è un \textbf{diritto fondamentale} (artt.~7 e 8 della Carta
dei diritti fondamentali dell'UE): per questo si può regolare in modo stabile. La tecnologia, invece,
cambia continuamente, e quando la legge la ``insegue'' arriva sempre in ritardo. L'esempio è l'AI Act:
la prima bozza (2020) non prevedeva nulla sull'IA generativa, esplosa dal novembre 2022, e la parte sui
modelli generativi è stata aggiunta in corsa. Le norme che durano sono quelle che fissano un
\textbf{metodo} (valutare il rischio, scegliere misure adeguate, dimostrarlo) invece di un elenco di
tecnologie.

\begin{esame}
\begin{itemize}
  \item L'obiettivo del corso è capire \textbf{risk assessment} e \textbf{compliance}, su cui si
    basano tutte le norme studiate: non serve imparare gli articoli a memoria.
  \item Saper fare un \textbf{severity assessment} di un data breach vale, secondo il docente,
    ``circa il 20\% dell'esame'' (Capitolo~\ref{cap:data-breach}).
  \item Tipica domanda aperta: capire \textbf{quale soggetto} deve notificare e \textbf{a quale norma}
    è soggetta l'organizzazione (ad esempio: è soggetta a DORA o no?).
  \item ``Sapere quale domanda pone una norma è la competenza più utile che il corso può dare.''
\end{itemize}
\end{esame}

% =====================================================================
\section{Il filo rosso: lunedì, ore 8:14}
\label{sec:intro-scenario}

Tutto il corso ruota intorno a uno scenario. Lunedì alle 8:14 un \textbf{ransomware} colpisce la
società che gestisce la piattaforma di prenotazione online di una rete ospedaliera regionale. Alle 8:20
è un problema tecnico; alle 9:00 è un problema legale. In ogni incidente vanno guardati tre elementi:
\begin{itemize}
  \item \textbf{dati}: 200.000 cartelle cifrate e una copia esfiltrata; sono \textbf{dati sanitari},
    cioè categorie particolari (art.~9 GDPR);
  \item \textbf{servizio}: prenotazioni ferme per quattro giorni, si torna alla carta. Dieci anni fa
    contava soprattutto il dato; oggi conta anche la \textbf{continuità}, perché un paziente che non
    riesce a prenotare può subire un danno (e ne nasce un contenzioso);
  \item \textbf{supply chain}: l'ingresso è avvenuto attraverso un componente vulnerabile del software
    di un fornitore. ``Non è un problema mio'' è falso: chi sceglie il fornitore ne risponde, salvo
    rivalersi su di lui.
\end{itemize}
In più, lo strumento di \emph{triage} basato su IA, addestrato su quel database, comincia a dare
risultati inaffidabili e nessuno se ne accorge per due giorni. La domanda del corso è: \emph{quanti
obblighi legali sono partiti alle 8:14?}

\subsection{Un evento, cinque regimi (più il diritto penale)}
\label{sec:intro-regimi}

Nessuna norma, da sola, governa quella mattina. Ne valgono diverse contemporaneamente, ciascuna con il
proprio \textbf{bene protetto}, la propria \textbf{autorità} e il proprio \textbf{orologio}.

\begin{center}
\small
\begin{tabularx}{\linewidth}{>{\raggedright\arraybackslash}p{3.1cm} X >{\raggedright\arraybackslash}p{3.6cm} >{\raggedright\arraybackslash}p{1.8cm}}
\toprule
\textbf{Regime} & \textbf{Che cosa protegge} & \textbf{A chi si risponde} & \textbf{Capitolo}\\
\midrule
GDPR (Reg.~UE 2016/679) & i dati personali dei pazienti & Garante / autorità capofila &
  \ref{cap:gdpr}, \ref{cap:data-breach}\\
NIS2 (Dir.~UE 2022/2555; in Italia D.lgs.~138/2024) & la continuità di un servizio essenziale &
  CSIRT Italia / ACN & \ref{cap:nis2}\\
CRA (Reg.~UE 2024/2847) & il prodotto con elementi digitali & CSIRT + piattaforma unica ENISA &
  \ref{cap:cra}\\
DORA (Reg.~UE 2022/2554) & la resilienza di una funzione finanziaria & autorità competente / ESA &
  \ref{cap:dora}\\
AI Act (Reg.~UE 2024/1689) & il sistema di IA e chi vi è esposto & autorità di vigilanza del mercato &
  \ref{cap:aiact}\\
\bottomrule
\end{tabularx}
\end{center}

Nello scenario ospedaliero non tutti i regimi si applicano davvero: DORA, per esempio, riguarda solo
banche, finanza e assicurazioni. Il punto è un altro: \textbf{stessi fatti, domande diverse}. Il primo
passo di ogni analisi è capire quali norme sono attivate e da quale evento parte ciascun orologio
(Sezione~\ref{sec:intro-orologi}).

\paragraph{Il diritto penale in parallelo.} Gli stessi fatti sono anche \textbf{reati}: la Convenzione
di Budapest sul cybercrime (23 novembre 2001) armonizza le fattispecie (accesso illegale,
intercettazione illecita, danneggiamento di dati e sistemi, abuso di dispositivi, falso e frode
informatica), i poteri procedurali (conservazione rapida dei dati, ordini di esibizione, perquisizione
e sequestro di dati) e la cooperazione internazionale; il Secondo Protocollo addizionale (2022)
riguarda le prove elettroniche. In Italia le fattispecie corrispondono agli artt.~615-\emph{ter},
617-\emph{quater}, 635-\emph{bis} e 640-\emph{ter} del codice penale. Lo stesso evento si divide in due
fascicoli: uno \textbf{amministrativo}, in cui l'organizzazione deve dimostrare di essere stata
diligente, e uno \textbf{penale}, in cui è lo Stato a dover dimostrare che qualcuno non lo è stato.
Procedura e standard probatorio sono diversi. La Convenzione non ha una lezione dedicata, ma conta per
la \textbf{giurisdizione}: anche ciò che sta nel cloud ha una giurisdizione.

\begin{integrazione}[i reati informatici citati]
Art.~615-\emph{ter} c.p.: accesso abusivo a un sistema informatico o telematico.
Art.~617-\emph{quater} c.p.: intercettazione, impedimento o interruzione illecita di comunicazioni
informatiche o telematiche. Art.~635-\emph{bis} c.p.: danneggiamento di informazioni, dati e programmi
informatici. Art.~640-\emph{ter} c.p.: frode informatica. La L.~90/2024 ha inasprito diverse di queste
pene. La Convenzione di Budapest non è stata firmata, tra gli altri, da Russia e Cina (come ricordato
in aula).
\end{integrazione}

% =====================================================================
\section{Le quattro domande e la mappa del corso}
\label{sec:intro-domande}

Ogni lezione risponde a una di quattro domande:
\begin{center}
\small
\begin{tabularx}{\linewidth}{>{\raggedright\arraybackslash}p{3.3cm} X >{\raggedright\arraybackslash}p{3.2cm}}
\toprule
\textbf{Domanda} & \textbf{Contenuti} & \textbf{Dove nel manuale}\\
\midrule
1. \textbf{Che cosa} proteggiamo? & dati personali e non personali, reti e sistemi informativi,
  prodotti con elementi digitali, modelli di IA & Sez.~\ref{sec:intro-dati}; cap.~\ref{cap:gdpr},
  \ref{cap:cra}, \ref{cap:aiact}\\
2. \textbf{Come} dobbiamo proteggerlo? & approccio basato sul rischio, misure di sicurezza, sistemi di
  gestione, certificazione e valutazione di conformità & Sez.~\ref{sec:intro-rischio};
  cap.~\ref{cap:iso}, \ref{cap:csa}\\
3. \textbf{Che cosa succede} se la protezione fallisce? & obblighi di notifica e scadenze,
  responsabilità, sanzioni amministrative, reati & cap.~\ref{cap:data-breach}, \ref{cap:scadenze}\\
4. \textbf{Chi decide}? & autorità di controllo, CSIRT, ENISA, agenzie nazionali, organismi notificati,
  tribunali & Sez.~\ref{sec:intro-autorita}\\
\bottomrule
\end{tabularx}
\end{center}

\paragraph{Le due parti del corso.} La \textbf{Parte~I (dati e privacy)} dà gli oggetti e il linguaggio:
quadro internazionale (Consiglio d'Europa, OCSE), GDPR, applicazioni settoriali e Reg.~(UE) 2018/1807
sui dati non personali, AI Act e approccio basato sul rischio. La \textbf{Parte~II (cybersicurezza e
cybercrime)} li mette sotto pressione: Convenzione di Budapest, data breach, NIS2 e Cybersecurity Act,
AI Act e sicurezza, ISO/IEC~27001 e risk assessment, organismi UE e casi di studio.

\subsection{L'oggetto che si sposta}
Ogni generazione di norme risponde ai limiti della precedente e allarga ciò che il diritto considera
``la cosa da proteggere''. Cambia il bene; il metodo resta lo stesso.
\begin{center}
\small
\begin{tabularx}{\linewidth}{>{\raggedright\arraybackslash}p{1.3cm} >{\raggedright\arraybackslash}p{4.2cm} X}
\toprule
\textbf{Anno} & \textbf{Strumento} & \textbf{Che cosa cambia}\\
\midrule
1980 & Linee guida OCSE & le misure di sicurezza diventano un principio (soft law)\\
1981 & Convenzione 108 (Consiglio d'Europa) & primo trattato vincolante sulla protezione dei dati:
  si protegge un \emph{archivio}\\
2001 & Convenzione di Budapest & il cybercrime diventa una categoria internazionale\\
2016 & GDPR (e prima direttiva NIS) & la sicurezza diventa un obbligo sanzionabile: si protegge un
  \emph{processo}\\
2018 & Reg.~(UE) 2018/1807 & entrano in scena i dati non personali\\
2019 & Cybersecurity Act & mandato permanente di ENISA e certificazione europea\\
2022 & NIS2 e DORA & resilienza, supply chain, responsabilità degli organi di vertice: si protegge un
  \emph{servizio}\\
2024 & AI Act e CRA & entrano nel perimetro il \emph{prodotto} e il \emph{modello}\\
2026 & segnalazioni CRA, misure NIS2, AI Omnibus & la frontiera continua a muoversi\\
\bottomrule
\end{tabularx}
\end{center}

% =====================================================================
\section{Le fonti del diritto}
\label{sec:intro-fonti}

\subsection{Tre livelli}
I dati ignorano i confini; l'\emph{enforcement} no. Per questo la materia si studia su tre livelli
contemporaneamente:
\begin{itemize}
  \item \textbf{internazionale}: Convenzione 108 del Consiglio d'Europa (1981), primo trattato
    vincolante, e Convenzione 108+ (aperta alla firma nel 2018; secondo le slide 34 ratifiche delle 38
    necessarie, quindi non ancora in vigore); Linee guida OCSE sulla privacy (1980, riviste nel 2013),
    soft law che ha influenzato tutto ciò che è venuto dopo; Convenzione di Budapest (2001) e Secondo
    Protocollo (2022);
  \item \textbf{Unione europea}: regolamenti e direttive (sotto); è da qui che proviene ormai quasi tutto
    il diritto vincolante della cybersicurezza;
  \item \textbf{nazionale}: in Italia il D.lgs.~196/2003 (Codice privacy, modificato), il
    D.lgs.~138/2024 (recepimento della NIS2), la L.~90/2024; le istituzioni (Garante per la protezione
    dei dati personali, ACN, CSIRT Italia); i reati informatici del codice penale.
\end{itemize}
Per un ingegnere la conseguenza è concreta: la stessa architettura (una region cloud, un fornitore, un
modello) è un oggetto giuridico diverso a seconda di dove stanno i dati, dove è offerto il servizio e chi
ne subisce gli effetti.

\begin{nota}[Corti europee: attenzione]
La \textbf{Corte di giustizia dell'Unione europea} (CGUE, Lussemburgo) è un'istituzione dell'UE e
interpreta il diritto dell'Unione (per esempio nei casi \emph{Google Spain} e \emph{Lindqvist},
Capitolo~\ref{cap:gdpr}). La \textbf{Corte europea dei diritti dell'uomo} (CEDU, Strasburgo)
appartiene invece al Consiglio d'Europa e \textbf{non fa parte dell'UE}: le slide lo segnalano
esplicitamente come trappola.
\end{nota}

\subsection{Regolamento e direttiva}
\label{sec:intro-reg-dir}
\begin{center}
\small
\begin{tabularx}{\linewidth}{>{\raggedright\arraybackslash}p{3cm} X X}
\toprule
 & \textbf{Regolamento} & \textbf{Direttiva}\\
\midrule
Efficacia & si applica direttamente e in modo identico in tutti gli Stati membri & obbliga gli Stati a
  raggiungere un risultato; va \textbf{recepita} con una legge nazionale\\
Margine nazionale & minimo (salvo le aperture previste dal regolamento stesso) & ampio: ambito, sanzioni e
  vigilanza possono differire da Stato a Stato\\
Effetto pratico & stesse regole in Germania e in Italia & 27 ``perimetri'' nazionali diversi\\
Esempi del corso & GDPR, Cybersecurity Act, DORA, AI Act, CRA & NIS2 (in Italia D.lgs.~138/2024)\\
\bottomrule
\end{tabularx}
\end{center}
Secondo il docente oggi circa il 90\% degli atti UE della materia sono regolamenti: un vantaggio
rispetto agli Stati Uniti, dove federazione e singoli Stati legiferano separatamente (California, New
York e Washington hanno leggi proprie sull'IA). Il docente presenta la distinzione come un
\textbf{principio generale da capire}, più che come probabile domanda d'esame.

\begin{integrazione}[base giuridica della distinzione]
La distinzione è fissata dall'art.~288 del Trattato sul funzionamento dell'Unione europea (TFUE): il
regolamento ha portata generale, è obbligatorio in tutti i suoi elementi ed è direttamente applicabile
in ciascuno Stato membro; la direttiva vincola lo Stato cui è rivolta quanto al risultato da
raggiungere, lasciando agli organi nazionali la scelta della forma e dei mezzi. Anche
un regolamento può però lasciare spazi agli Stati: il GDPR, per esempio, consente norme nazionali su età
del consenso dei minori, dati sanitari, rapporti di lavoro e sanzioni penali (in Italia attuate con il
D.lgs.~101/2018, che ha adeguato il Codice privacy).
\end{integrazione}

\subsection{Hard law e soft law}
\label{sec:intro-hard-soft}
La \textbf{hard law} è la norma vincolante con conseguenze \emph{predeterminate} dalla legge (GDPR, CRA,
NIS2, DORA, AI Act); la \textbf{soft law} comprende policy aziendali, linee guida (ad esempio quelle
ENISA), standard (ISO) e schemi di certificazione, la cui violazione ha conseguenze non fissate dalla
legge e lasciate a un'ampia discrezionalità. Le norme del corso dicono di ``essere sicuri'' ma non come:
la sostanza operativa della cybersicurezza è soft law, che ogni organizzazione deve costruire per
dimostrare il rispetto della hard law. Il tema è sviluppato nella lezione del 09/11
(Capitolo~\ref{lez:2026-11-09}) e nel capitolo sul Cybersecurity Act (Capitolo~\ref{cap:csa}); secondo
il docente non è una domanda d'esame, ma serve per rispondere bene alle altre.

% =====================================================================
\section{Dato personale e dato non personale}
\label{sec:intro-dati}

\begin{center}
\small
\begin{tabularx}{\linewidth}{>{\raggedright\arraybackslash}p{2.6cm} X X}
\toprule
 & \textbf{Dati personali (GDPR)} & \textbf{Dati non personali (Reg.~UE 2018/1807)}\\
\midrule
Definizione & qualsiasi informazione relativa a una persona fisica identificata o
  \textbf{identificabile} & tutto il resto: dati macchina, telemetria, dati aggregati e anonimizzati\\
Bene protetto & un diritto fondamentale (artt.~7--8 della Carta) & il mercato interno, non l'individuo\\
Obblighi & legati a finalità, liceità e rischio per la persona & divieto di obblighi di localizzazione
  dei dati; portabilità per gli utenti professionali\\
Sicurezza & è un principio del trattamento (art.~5(1)(f) GDPR) & gli obblighi arrivano da altre fonti:
  NIS2, DORA, CRA, contratti\\
\bottomrule
\end{tabularx}
\end{center}

\paragraph{I dataset misti sono la regola.} Se dati personali e non personali sono inestricabilmente
collegati, il GDPR si applica all'intero insieme (art.~2(2) Reg.~2018/1807). La maggior parte dei
dataset reali è mista.

\paragraph{Anonimizzare non è togliere il nome.} Giuridicamente un dato è anonimo solo se nessun
elemento consente di ricondurlo alla persona, \textbf{né direttamente né indirettamente}. Esempi
citati: un riferimento a fatti noti che identifica subito un politico famoso anche senza nominarlo;
``l'unica persona in città con questa malattia''. Negli audit le aziende convinte di aver anonimizzato
scoprono spesso molte possibilità di re-identificazione indiretta, e con i modelli linguistici la
re-identificazione è diventata molto più facile. Le slide aggiungono che l'anonimizzazione ``non è una
porta a senso unico'': sono proprio le misure di sicurezza a mantenere anonimi i dati ``anonimi''. I dati
non personali hanno molti meno vincoli, e il docente indica nei \textbf{dati sintetici} un settore in
crescita.

\paragraph{L'indirizzo IP è un dato personale?} Sì, perché rende la persona \emph{identificabile}, anche
quando è dinamico. Il percorso per risalire alla persona passa per il fornitore del servizio (quale IP
era in uso in quel momento) e poi per l'operatore di rete (chi pagava la connessione); in Europa la
richiesta passa dall'autorità giudiziaria, e il risultato è solo un indizio (chi paga la bolletta non è
necessariamente chi ha agito). Anche un indirizzo email come \texttt{nomignolo74@gmail.com} è un dato
personale, perché il fornitore può ricondurlo a una persona.

\begin{integrazione}[pseudonimizzazione e IP dinamico]
\begin{itemize}
  \item Il dato \textbf{pseudonimizzato} (art.~4(5) GDPR: non più attribuibile all'interessato senza
    informazioni aggiuntive conservate separatamente) resta un \textbf{dato personale}; solo il dato
    davvero anonimo esce dal GDPR (considerando 26). La pseudonimizzazione è una misura di sicurezza
    (art.~32(1)(a)), non un modo per uscire dalla norma.
  \item Sull'IP dinamico la CGUE si è pronunciata nella sentenza \emph{Breyer} (C-582/14, 2016): è dato
    personale per il gestore del sito se questi dispone di mezzi giuridici ragionevoli per ottenere
    dal fornitore di accesso le informazioni che consentono l'identificazione.
\end{itemize}
\end{integrazione}

% =====================================================================
\section{L'approccio basato sul rischio}
\label{sec:intro-rischio}

\subsection{Il GDPR come prima legge sulla cybersicurezza}
Molto prima della NIS2, il GDPR ha dato alla sicurezza un nome giuridico, un metodo e un prezzo:
\begin{itemize}
  \item \textbf{art.~5(1)(f)}: integrità e riservatezza come \emph{principio} del trattamento, non come
    allegato o opzione;
  \item \textbf{artt.~24--25}: accountability (dimostrare la conformità), protezione dei dati fin dalla
    progettazione e per impostazione predefinita;
  \item \textbf{art.~32}: misure adeguate al rischio, tenendo conto dello stato dell'arte e dei costi;
  \item \textbf{artt.~33--34 e 35}: notifica in 72 ore, comunicazione agli interessati se il rischio è
    elevato, valutazione d'impatto (DPIA) per i trattamenti ad alto rischio.
\end{itemize}
L'art.~32 è ``il cardine del corso'' perché dice tre cose che tutte le norme successive ripetono: la
sicurezza è \textbf{relativa al rischio}; si misura rispetto a uno standard che si muove (lo
\textbf{stato dell'arte}); qualcuno deve poterla \textbf{dimostrare}. Nessun articolo del GDPR indica
quale controllo adottare: è un silenzio voluto, riempito da ISO/IEC~27001 e dalle metodologie ENISA.

\subsection{Una grammatica comune}
\begin{center}
\small
\begin{tabularx}{\linewidth}{>{\raggedright\arraybackslash}p{2.4cm} >{\raggedright\arraybackslash}p{2.6cm} X}
\toprule
\textbf{Strumento} & \textbf{Dove} & \textbf{Che cosa chiede}\\
\midrule
GDPR & artt.~32, 35 & misure adeguate al rischio per diritti e libertà; DPIA se il rischio è elevato\\
NIS2 & art.~21 & misure di gestione del rischio ``multi-rischio'', proporzionate a esposizione,
  dimensione e impatto\\
DORA & artt.~5--16 & un quadro di gestione del rischio ICT approvato dall'organo di gestione\\
AI Act & artt.~9, 15 & un sistema di gestione del rischio lungo tutto il ciclo di vita; accuratezza,
  robustezza, cybersicurezza\\
ISO/IEC~27001 & cl.~6.1 e Allegato~A & identificare, analizzare, valutare, trattare: in modo
  documentato, verificabile, certificabile\\
\bottomrule
\end{tabularx}
\end{center}
Il metodo è sempre lo stesso: \textbf{contesto} $\rightarrow$ \textbf{identificazione} $\rightarrow$
\textbf{analisi} $\rightarrow$ \textbf{valutazione} $\rightarrow$ \textbf{trattamento} $\rightarrow$
\textbf{monitoraggio}. Il rischio si stima con due grandezze, \textbf{probabilità} e \textbf{gravità};
le misure riducono il rischio \emph{inerente} a un rischio \emph{residuo}, ed è quest'ultimo che decide
(Capitolo~\ref{cap:gdpr}, DPIA; Capitolo~\ref{cap:iso}).

L'approccio basato sul rischio \textbf{non è una fuga dalle regole}: è una regola che dice che gli
obblighi crescono con ciò che si può rompere. Per questo la stessa azienda deve di più a un ospedale che a
una newsletter, e ``abbiamo seguito la checklist'' non è mai stato una difesa.

\subsection{La compliance è un problema di prova}
In due sensi: preservare le \textbf{prove digitali} e, soprattutto, poter \textbf{dimostrare} di aver
fatto il possibile per prevenire l'incidente (accountability). Gli attacchi non si possono evitare del
tutto (la citazione di Eugene Spafford chiude le slide sul GDPR): la domanda non è \emph{se} ma
\emph{quando}. Sistemi di gestione (ISO/IEC~27001), certificazione dei prodotti (Cybersecurity Act,
schema EUCC) e valutazione di conformità (CRA, AI Act) servono a questo. Un certificato è un
\textbf{elemento di prova}, mai un'immunità: GDPR (artt.~42--43), NIS2 (art.~24) e Cybersecurity Act
trattano la certificazione come prova da pesare, non come porto sicuro.

\subsection{Il perimetro è la supply chain}
Tutti i regimi del corso sono arrivati alla stessa conclusione: si risponde anche di chi non si impiega.
\begin{center}
\small
\begin{tabularx}{\linewidth}{>{\raggedright\arraybackslash}p{2cm} >{\raggedright\arraybackslash}p{2.4cm} X}
\toprule
\textbf{Norma} & \textbf{Dove} & \textbf{Effetto sui contratti con i fornitori}\\
\midrule
GDPR & art.~28 & responsabili vincolati da contratto, istruzioni, diritti di audit, consenso ai
  sub-responsabili\\
NIS2 & art.~21(2)(d) & sicurezza della catena di approvvigionamento come misura obbligatoria\\
DORA & artt.~28--30 & contenuti contrattuali obbligatori, registro degli accordi, sorveglianza UE sui
  fornitori ICT critici\\
CRA & artt.~13--14 & obblighi che seguono il prodotto (fabbricante, importatore, distributore): SBOM,
  gestione delle vulnerabilità, aggiornamenti\\
AI Act & artt.~22--25 & ruoli lungo la catena del valore; chi modifica un sistema può diventarne il
  fornitore\\
\bottomrule
\end{tabularx}
\end{center}

% =====================================================================
\section{Le autorità}
\label{sec:intro-autorita}

Ogni regime porta la propria autorità: disegnare correttamente questa mappa è metà della risposta a un
incidente. ``Alle 8:14 bisogna sapere chi chiamare.''
\begin{center}
\small
\begin{tabularx}{\linewidth}{>{\raggedright\arraybackslash}p{3.2cm} X}
\toprule
\textbf{Ambito} & \textbf{Autorità}\\
\midrule
Protezione dei dati & Garante per la protezione dei dati personali (autorità indipendente); Comitato
  europeo per la protezione dei dati (EDPB: coerenza, linee guida); Garante europeo (EDPS) per le
  istituzioni UE\\
Cybersicurezza -- UE & ENISA (certificazione, panorama delle minacce, piattaforma unica di
  segnalazione); NIS Cooperation Group (strategia e indirizzi); EU-CyCLONe (crisi su larga scala)\\
Risposta operativa & CSIRT Italia, all'interno dell'ACN (``il numero da chiamare''); rete europea dei
  CSIRT; CERT di settore\\
Vigilanza NIS2 & ACN come autorità nazionale competente; autorità di settore per soggetti specifici\\
Finanza (DORA) & EBA, ESMA, EIOPA e autorità nazionali; sorveglianza dei fornitori ICT critici\\
Intelligenza artificiale & AI Office e AI Board; autorità nazionali di vigilanza del mercato; autorità
  di notifica e organismi notificati\\
\bottomrule
\end{tabularx}
\end{center}
Due conseguenze, secondo le slide: \textbf{non si sceglie il foro}, e un difetto di coordinamento tra
queste autorità è un rischio dell'organizzazione, non loro. Per l'AI Act il docente cita in Italia
l'\textbf{AgID}, agenzia governativa (non indipendente come il Garante).

\paragraph{Verso un punto unico di segnalazione.} Con il pacchetto \emph{Digital Omnibus} si andrà verso
un \textbf{punto unico di segnalazione} degli incidenti presso ENISA, che poi informerà le autorità
competenti: per il docente sarà ``la rivoluzione più importante'', perché oggi le autorità nazionali non
si parlano abbastanza mentre lo stesso malware colpisce in tutta Europa. Non è però ancora operativo.

\begin{esame}
\begin{itemize}
  \item Vanno studiate \textbf{tutte le autorità} e lo \textbf{schema delle notifiche vigente}: il
    punto unico ENISA del Digital Omnibus non è ancora operativo.
  \item \textbf{Attenzione alla sigla DPA}: \emph{Data Protection Authority} (l'autorità) ma anche
    \emph{Data Processing Agreement} (il contratto con il responsabile, art.~28 GDPR). Il docente usa
    ``DPA'' per il contratto e dice \emph{data protection authority} per l'autorità; nel dubbio, all'esame
    si può chiedere chiarimento.
\end{itemize}
\end{esame}

\begin{integrazione}[autorità italiane per l'AI Act]
La legge italiana sull'intelligenza artificiale (L.~132/2025) designa come autorità nazionali per l'AI
Act l'\textbf{AgID} (autorità di notifica, con compiti di promozione) e l'\textbf{ACN} (vigilanza,
incluse le attività ispettive e sanzionatorie, e profili di cybersicurezza), ferme restando le
competenze di Garante e autorità di settore. Da confrontare con il capitolo sull'AI Act
(Capitolo~\ref{cap:aiact}).
\end{integrazione}

% =====================================================================
\section{Un incidente, cinque orologi}
\label{sec:intro-orologi}

La parte difficile della risposta a un incidente non è tecnica: ogni regime conta il tempo in modo
diverso e \textbf{da un evento diverso}. La tabella mostra solo da che cosa parte ciascun orologio e che
cosa misura; fasi, scadenze intermedie e trappole sono nella tabella comparativa del
Capitolo~\ref{cap:scadenze}.
\begin{center}
\small
\begin{tabularx}{\linewidth}{>{\raggedright\arraybackslash}p{2.3cm} X >{\raggedright\arraybackslash}p{3.4cm} >{\raggedright\arraybackslash}p{2.2cm}}
\toprule
\textbf{Regime} & \textbf{L'orologio parte quando\ldots} & \textbf{Che cosa misura} & \textbf{Prima scadenza}\\
\midrule
GDPR artt.~33--34 & il titolare è consapevole di una violazione di dati personali & il danno alle
  persone & 72 h\\
NIS2 art.~23 & si è consapevoli di un incidente significativo & l'interruzione del servizio & 24 h\\
DORA art.~19 & l'incidente è classificato come grave & l'impatto su una funzione & 4 h\\
CRA art.~14 & si è consapevoli di una vulnerabilità attivamente sfruttata o di un incidente grave &
  la sfruttabilità del prodotto & 24 h\\
AI Act art.~73 & si accerta un incidente grave e il nesso causale con il sistema & la gravità
  dell'esito & 15 giorni\\
\bottomrule
\end{tabularx}
\end{center}
Il docente sottolinea le sole \textbf{4 ore} di DORA. Le notifiche corrono \textbf{in parallelo}: una
notifica fatta a un'autorità non adempie gli obblighi verso un'altra (esempi nei laboratori:
GDPR + DORA nel caso Sylvane, Capitolo~\ref{lab:02}; GDPR + CRA nella simulazione Pulse~K,
Sezione~\ref{sec:pulsek}).

\begin{verifica}
\begin{itemize}
  \item Ruolo delle autorità italiane per l'AI Act (AgID, ACN): il docente ha citato solo l'AgID; il
    quadro va confermato nella parte del corso sull'AI Act.
  \item Stato del \emph{Digital Omnibus}: in aula è stato detto sia che è stato ``applicato a luglio''
    sia che il punto unico di segnalazione non è in vigore; le slide citano solo l'\emph{AI Omnibus}
    (in vigore dal 27 luglio 2026).
  \item Le date e le scadenze delle slide cambiano spesso (``tre scadenze di queste slide sono cambiate
    negli ultimi dodici mesi''): verificare sempre sul testo vigente.
\end{itemize}
\end{verifica}

\chapter{GDPR}
\label{cap:gdpr}

\dacompletare{Fonte: slide \emph{02.~Cybersecurity Laws and Regulations -- II Privacy}.
Scaletta prevista:
\begin{itemize}
  \item storia: le quattro generazioni della protezione dei dati;
  \item definizioni: dato personale, dati particolari (art.~9), anonimizzazione e pseudonimizzazione;
  \item ruoli: titolare, responsabile, sub-responsabile, contitolari, interessato;
  \item principi (art.~5) e basi giuridiche (art.~6);
  \item diritti dell'interessato;
  \item accountability, privacy by design e by default, DPIA, DPO;
  \item governance esterna: autorità di controllo, EDPB, tribunali;
  \item trasferimenti extra UE (Capo V) e TIA;
  \item sanzioni e responsabilità.
\end{itemize}}

\chapter{Data breach e valutazione della severità}
\label{cap:data-breach}

\dacompletare{Fonte: slide \emph{03.~Data breach severity methodology\_1.0}; Simulazione Pulse K
(Sezione~\ref{sec:pulsek}). Scaletta prevista:
\begin{itemize}
  \item definizione di violazione dei dati personali (art.~4(12)): riservatezza, integrità, disponibilità;
  \item incidente di sicurezza e data breach: differenza;
  \item art.~33 (notifica all'autorità) e art.~34 (comunicazione agli interessati);
  \item metodologia ENISA: DPC, EI, CB e scala di severità;
  \item registro dei data breach (art.~33(5));
  \item ruolo del responsabile (art.~33(2));
  \item notifica preliminare e notifica completa; modulo del Garante;
  \item linee guida EDPB 9/2022.
\end{itemize}}

\chapter{ISO/IEC 27001 e gestione del rischio}
\label{cap:iso}

\begin{nota}[Fonti]
Lezioni di Gianluca Gilardi del 12/10, 19/10, 23/10 e 26/10/2026 (Capitoli~\ref{lez:2026-10-12},
\ref{lez:2026-10-19}, \ref{lez:2026-10-23} e~\ref{lez:2026-10-26}); richiami del professore il 21/09, 28/09,
05/10 e 09/11 (Capitoli~\ref{lez:2026-09-21}, \ref{lez:2026-09-28}, \ref{lez:2026-10-05}
e~\ref{lez:2026-11-09}); laboratorio sul risk register (Capitolo~\ref{lab:03}). Slide
\emph{04.~polito\_iso\_2026} (\emph{Fundamentals and ISO/IEC Standards}) e \emph{06.~27001\_2026\_2310}
(\emph{Risk Management}); briefing e playbook del Lab~3. Come fonte secondaria, gli appunti degli studenti
dell'A.A.~2025/26, usati solo dove coerenti con le slide di quest'anno.

Le norme ISO sono a pagamento e protette da copyright (per questo non sono sulla piattaforma del corso): qui
si riportano solo i \textbf{titoli} di clausole e controlli e la loro idea, in parole nostre. Il diario è la
cronaca delle lezioni; questo capitolo riordina e approfondisce.
\end{nota}

% =====================================================================
\section{Standard tecnici e soft law: perché la 27001}
\label{sec:iso-standard}
% =====================================================================

\subsection{Che cos'è uno standard ISO}
L'ISO è un'organizzazione internazionale con sede in Svizzera che pubblica \textbf{standard tecnici}:
documenti \textbf{volontari}, che nessuna legge impone di per sé. Nel linguaggio del corso sono \textbf{soft
law} (Capitolo~\ref{lez:2026-11-09} e Capitolo~\ref{cap:csa}): regole che un'organizzazione sceglie di
seguire, si aggiornano molto più in fretta delle leggi e non hanno sanzioni fissate dal legislatore. La
ISO/IEC~27001 certifica l'\textbf{organizzazione}, mentre gli schemi del Cybersecurity Act certificano
prodotti, servizi e processi ICT.

Tre caratteristiche sottolineate a lezione:
\begin{itemize}
  \item gli standard ISO sono \textbf{neutrali rispetto ai fornitori}: dicono quale effetto ottenere, non quale
    prodotto comprare;
  \item seguono un \textbf{approccio basato sul rischio}, come GDPR, NIS2, DORA e AI Act;
  \item non impongono una metodologia di valutazione del rischio: tassonomia dei rischi, scale e scelta dei
    controlli spettano all'organizzazione, che deve però essere coerente e saperle difendere.
\end{itemize}

\begin{integrazione}[ISO e IEC]
ISO (\emph{International Organization for Standardization}) ha sede a Ginevra; la sigla ``ISO/IEC'' indica
che lo standard è sviluppato insieme all'IEC (\emph{International Electrotechnical Commission}) nel comitato
tecnico congiunto JTC~1, sottocomitato SC~27 (sicurezza delle informazioni, cybersicurezza e privacy).
\end{integrazione}

\subsection{Il vuoto che la 27001 riempie}
Il punto di partenza del corso è l'\textbf{art.~32 GDPR}: titolare e responsabile devono adottare misure
tecniche e organizzative \textbf{adeguate al rischio}, tenendo conto dello stato dell'arte e dei costi. La
norma \textbf{non dice quali misure}: l'elenco è aperto e la scelta va motivata caso per caso
(\emph{accountability}). Lo stesso vale per le altre norme del corso. La 27001 riempie questo spazio con un
metodo e un catalogo di controlli riconosciuti, ed è il \textbf{linguaggio comune} con cui in Europa aziende,
clienti e auditor parlano di sicurezza (negli Stati Uniti prevalgono gli standard NIST e, lato clienti, i
report SOC~2). Il professore, pur ``non essendo un fan'' della ISO, la considera indispensabile per
comunicare con le aziende.

Tutte le norme del corso condividono lo stesso ciclo (contesto, identificazione, analisi, valutazione,
trattamento, monitoraggio):
\begin{center}
\small
\begin{tabularx}{\linewidth}{>{\raggedright\arraybackslash}p{2.6cm} >{\raggedright\arraybackslash}p{3.2cm} X}
\toprule
\textbf{Norma} & \textbf{Dove sta il rischio} & \textbf{Che cosa offre la 27001}\\
\midrule
GDPR & artt.~24, 25, 32, 35 & metodo per scegliere e documentare le misure dell'art.~32; prova di diligenza
  (accountability)\\
NIS2 & art.~21 (misure di gestione del rischio) & processi e controlli già mappati sulle misure richieste;
  molte aziende la adottano perché i grandi clienti scelgono i fornitori anche in base alla certificazione\\
DORA & artt.~5--16 (quadro di gestione del rischio ICT) & struttura dell'ISMS riutilizzabile nel settore
  finanziario\\
CRA & risk assessment del \textbf{prodotto} & nulla di diretto sul prodotto: la 27001 riguarda
  l'organizzazione; può sostenere i processi del fabbricante (es.\ gestione delle vulnerabilità)\\
AI Act & artt.~9 e 15 & stessa logica di sistema di gestione (per l'IA esiste la ISO/IEC~42001)\\
\bottomrule
\end{tabularx}
\end{center}
Le slide chiudono con l'idea che la spinta regolatoria (NIS2, DORA, CRA, AI Act) renda la 27001 una
\textbf{spina dorsale comune} della compliance.

\begin{integrazione}[NIS2, art.~21 e standard]
L'art.~21(2) NIS2 elenca le misure minime di gestione del rischio: politiche di analisi del rischio e di
sicurezza dei sistemi; gestione degli incidenti; continuità operativa (backup, disaster recovery, gestione
delle crisi); sicurezza della catena di approvvigionamento; sicurezza in acquisizione, sviluppo e
manutenzione, inclusa la gestione delle vulnerabilità; valutazione dell'efficacia delle misure; igiene
informatica e formazione; crittografia; sicurezza delle risorse umane, controllo degli accessi e gestione
degli asset; autenticazione a più fattori e comunicazioni protette. Quasi ogni voce ha uno o più controlli
corrispondenti nell'Allegato~A (es.\ 8.13 backup, 8.24 crittografia, 8.5 autenticazione, 8.8 vulnerabilità).
L'art.~25 NIS2 incoraggia l'uso di standard europei e internazionali, ma la certificazione 27001 non è una
presunzione di conformità alla NIS2: resta da dimostrare che le misure coprono davvero quanto richiesto.
Per il CRA la presunzione di conformità passa invece da norme armonizzate di prodotto e dagli schemi del
Cybersecurity Act (Capitolo~\ref{cap:cra}).
\end{integrazione}

\subsection{Che cosa la certificazione dimostra e che cosa no}
\begin{itemize}
  \item Certificazione 27001 significa che una \textbf{terza parte fidata} ha verificato che l'organizzazione
    ha un sistema di gestione con misure e controlli appropriati.
  \item \textbf{Non è richiesta da nessuna legge}; in alcuni Paesi o settori può essere un requisito
    contrattuale o di mercato (gare, grandi clienti).
  \item Davanti a un'autorità, dopo un incidente, aiuta a dimostrare di aver fatto il possibile, ma
    \textbf{non garantisce} che il sistema sia sicuro né che l'organizzazione sia conforme al GDPR, alla NIS2 o
    al CRA. Gli obblighi di legge (es.\ notifica di un data breach) restano tutti.
  \item Molte aziende non si certificano ma dichiarano di \textbf{seguire} i principi della 27001, per evitare
    il peso di audit e rinnovi.
  \item ``Sono troppo piccolo'' non è una risposta: ogni organizzazione raccoglie, tratta, conserva e
    trasmette informazioni, cioè asset esposti a rischio. Non ci si chiede \emph{se} ci sarà un incidente, ma
    \emph{quando}.
\end{itemize}

\begin{esame}
La certificazione ISO/IEC~27001 riguarda l'\textbf{organizzazione}, è \textbf{volontaria} e \textbf{non prova}
la conformità legale né l'assenza di incidenti: sostiene l'accountability, non la sostituisce. Vedi la
domanda~\ref{q:iso} del quiz.
\end{esame}

% =====================================================================
\section{La famiglia ISO/IEC 27000}
\label{sec:iso-famiglia}
% =====================================================================
La famiglia ``27k'' conta \textbf{oltre 70 standard} e continua a crescere. Le slide li dividono in standard
che: (a)~fissano i requisiti dell'ISMS e di chi lo certifica; (b)~danno supporto, guida o interpretazione
per costruirlo e migliorarlo; (c)~danno linee guida per settori specifici; (d)~riguardano la valutazione
della conformità. La 27001 è lo standard più usato per attuare e \textbf{certificare} un ISMS.

\begin{center}
\small
\begin{tabularx}{\linewidth}{>{\raggedright\arraybackslash}p{2.5cm} >{\raggedright\arraybackslash}p{2.6cm} X}
\toprule
\textbf{Standard} & \textbf{Tipo} & \textbf{Ruolo}\\
\midrule
ISO/IEC 27000 & vocabolario & panoramica dell'ISMS e termini comuni (definizioni di CIA, rischio, controllo);
  liberamente scaricabile\\
ISO/IEC 27001 & requisiti & requisiti dell'ISMS (clausole 4--10) e Allegato~A con i 93 controlli;
  \textbf{l'unico certificabile} di questo elenco\\
ISO/IEC 27002 & linea guida & guida ai controlli: come attuarli e quale effetto devono produrre; definisce
  gli attributi dei controlli\\
ISO/IEC 27003 & linea guida & interpretazione dei requisiti della 27001 (``should'', ``can'', ``may''): utile
  per capire che cosa guarderà l'auditor\\
ISO/IEC 27005 & linea guida & gestione del rischio di sicurezza delle informazioni: come svolgere
  valutazione e trattamento richiesti dalla 27001\\
ISO/IEC 27701 & estensione & privacy: sistema di gestione delle informazioni personali (PIMS) per titolari e
  responsabili, come estensione di 27001 e 27002\\
ISO/IEC 27017 & settoriale & sicurezza dei servizi cloud, per fornitori e clienti (appunti 2025/26)\\
ISO/IEC 27018 & settoriale & protezione dei dati personali nel cloud pubblico (appunti 2025/26)\\
\bottomrule
\end{tabularx}
\end{center}

\begin{itemize}
  \item Il rapporto tra 27001 e 27002, detto a lezione: la 27001 dà il \textbf{nome} del controllo, la 27002
    i \textbf{mezzi} (es.\ controllo degli accessi attuato con l'MFA). Ci si certifica sulla 27001; l'auditor
    usa la 27002 come riferimento per giudicare come i controlli sono attuati.
  \item Secondo Gianluca le altre certificazioni 27k sono ``figlie'' della 27001; la 27701, in particolare,
    nella sua pratica professionale non è mai stata chiesta da un cliente: ha senso per trattamenti molto
    sensibili su larga scala (es.\ dati sanitari).
  \item A lezione è stata citata anche la \textbf{ISO/IEC~27090} sulla sicurezza dei sistemi di IA,
    ``appena approvata'' e quindi assente dalle slide.
\end{itemize}

\begin{integrazione}[Requisiti contro linee guida; edizioni recenti]
\begin{itemize}
  \item La 27001 è scritta in forma di requisiti (``\emph{shall}''), verificabili in audit; la 27002 in forma
    di raccomandazioni (``\emph{should}''). Per questo solo la prima è certificabile.
  \item Le slide citano la \textbf{27701:2019}, che è un'estensione e presuppone la 27001. Nell'ottobre 2025 è
    uscita una nuova edizione che rende il PIMS uno standard di sistema di gestione \textbf{autonomo}:
    l'affermazione ``senza 27001 non si ottiene la 27701'' vale per l'edizione 2019.
  \item 27017 e 27018 nascono come codici di pratica che aggiungono controlli e guida alla 27002 per il
    cloud; non sono certificabili da sole, ma di solito vengono verificate insieme alla 27001.
  \item La 27001:2022 è stata modificata nel 2024 (Amd~1) per chiedere di considerare il cambiamento
    climatico tra i fattori di contesto (clausole 4.1 e 4.2), come per tutti i sistemi di gestione ISO.
\end{itemize}
\end{integrazione}

\paragraph{Perché la 27001 è così diffusa (slide).} È applicabile a qualsiasi organizzazione; è breve (corpo
di circa 25 pagine) e concentrata su ISMS e controlli; è allineata agli altri sistemi di gestione (qualità,
privacy, servizi, continuità operativa); è collaudata (deriva dalla BS~7799-1 del 1995); ha principi semplici
(parti interessate, leadership, miglioramento continuo, approccio per processi e basato sul rischio); è
certificabile; molti professionisti la usano e ``parlano la stessa lingua''. Gianluca aggiunge che circa il
75--80\% dei controlli si ritrova nel framework statunitense SOC~2, quindi il lavoro è riutilizzabile.

\paragraph{Benefici (slide).} Sicurezza complessiva (compresi risposta agli incidenti e resilienza),
consapevolezza e cultura, riduzione dei rischi di business, fiducia degli stakeholder, conformità legale,
regolamentare e contrattuale, approccio comune a requisiti diversi, maturità dei processi (IT, HR,
governance), budget di sicurezza più trasparenti e giustificabili, marketing. Ridurre il rischio di
sicurezza significa anche ridurre il \textbf{rischio di compliance}.

\paragraph{Diffusione (ISO Survey 2024).} Al 31/12/2024 i certificati 27001 validi sono 96.709 (contro circa
1,47 milioni della ISO~9001). Guida la Cina (33.359); l'Italia è \textbf{sesta} con 3.284. Sul mercato del
lavoro, conoscere la norma o essere \emph{lead auditor} è un titolo apprezzato.

% =====================================================================
\section{L'ISMS}
\label{sec:iso-isms}
% =====================================================================

\subsection{Sicurezza delle informazioni e triade CIA}
Dalle definizioni della ISO~27000, la sicurezza delle informazioni è la \textbf{preservazione di
riservatezza, integrità e disponibilità} delle informazioni; possono rilevare anche autenticità,
accountability, non ripudio e affidabilità.
\begin{center}
\small
\begin{tabularx}{\linewidth}{>{\raggedright\arraybackslash}p{2.6cm} X X}
\toprule
\textbf{Proprietà} & \textbf{Idea della definizione} & \textbf{Esempio di perdita}\\
\midrule
Riservatezza & l'informazione non è resa disponibile né rivelata a persone, entità o processi non
  autorizzati & email al destinatario sbagliato; accesso non autorizzato a un server, anche senza
  esfiltrazione\\
Integrità & l'informazione è accurata e completa & modifica non autorizzata di un record; IBAN alterato\\
Disponibilità & l'informazione è accessibile e utilizzabile su richiesta da chi è autorizzato & ransomware;
  guasto senza ridondanza; anche poche ore di indisponibilità\\
\bottomrule
\end{tabularx}
\end{center}
\noindent Non tutte le informazioni hanno la stessa sensibilità: \textbf{classificarle} serve a non sprecare
risorse su ciò che non conta (un bilancio destinato al deposito pubblico non va protetto in riservatezza).

\begin{esame}
\begin{itemize}
  \item Le \textbf{definizioni} di riservatezza, integrità e disponibilità sono indicate come probabile
    domanda d'esame.
  \item Anche una \textbf{perdita temporanea di disponibilità} di dati personali è un data breach: la
    valutazione di severità può escludere le notifiche, ma la violazione va annotata nel registro
    (Capitolo~\ref{cap:data-breach}).
\end{itemize}
\end{esame}

\subsection{Definizione e componenti}
Un \textbf{Information Security Management System} è un approccio sistematico per stabilire, attuare,
gestire, monitorare, riesaminare, mantenere e migliorare la sicurezza delle informazioni di
un'organizzazione, al servizio dei suoi obiettivi di business. È un \emph{sistema}, non un insieme di
toppe: componenti tecniche, organizzative e documentali devono stare insieme, e si governano
\textbf{persone, processi e tecnologia} (la tecnologia è la parte più facile).

\begin{center}
\small
\begin{tabularx}{\linewidth}{>{\raggedright\arraybackslash}p{4.4cm} X}
\toprule
\textbf{Componente (slide)} & \textbf{Osservazioni del docente}\\
\midrule
a) Politica & la parte più facile da scrivere, la più difficile da tenere aggiornata\\
b) Persone con responsabilità definite & senza ruoli formali (chi monitora, chi segnala, chi decide) non c'è
  un sistema di gestione ma un ``sistema della speranza''\\
c) Processi di gestione & definizione delle politiche, consapevolezza e competenze, pianificazione,
  attuazione, operatività, valutazione delle prestazioni, riesame della direzione, miglioramento\\
d) Informazioni documentate & niente ``tradizione orale''\\
e) Valutazione del rischio & specifica dell'ISMS\\
f) Trattamento del rischio, con scelta e attuazione dei controlli & specifica dell'ISMS\\
\bottomrule
\end{tabularx}
\end{center}
Le componenti a)--d) sono comuni a qualsiasi sistema di gestione (qualità, ambiente\ldots); e) ed f) sono ciò
che caratterizza l'ISMS: \textbf{è il rischio a definirlo}. Un ISMS si basa sulla valutazione del rischio e
sui livelli di accettazione stabiliti dall'organizzazione; i suoi benefici derivano dalla riduzione della
probabilità e/o dell'impatto degli incidenti.

\paragraph{Data governance prima della sicurezza.} Il caso del Lab~1 (Capitolo~\ref{lab:01}) mostra il
problema di chi non monitora: non si riesce a stabilire quando si è venuti a conoscenza della violazione. Il
peccato più grave è non saper provare \emph{cosa}, \emph{dove}, \emph{quando}, \emph{chi} e \emph{perché}.

\paragraph{Il ruolo della direzione.} Le decisioni strategiche sulla sicurezza vengono dall'alto: se il
vertice non conosce rischi e controlli non può decidere budget e priorità. Uno dei primi punti verificati in
audit è come il management è coinvolto e informato (clausola~5 e riesame della clausola~9.3).

\subsection{Fasi di implementazione e miglioramento continuo}
La ISO~27003 descrive un ciclo in cinque fasi; dopo il miglioramento si riparte dal contesto.
\begin{center}
\begin{tikzpicture}[node distance=4mm, every node/.style={draw, rounded corners, align=center,
    font=\small, minimum height=1.2cm, text width=2.35cm}]
  \node (c) {\textbf{Contesto}\\bisogni e perimetro};
  \node (r) [right=of c] {\textbf{Rischi}\\valutazione};
  \node (k) [right=of r] {\textbf{Controlli}\\processi e misure};
  \node (m) [right=of k] {\textbf{Monitoraggio}\\prestazioni ed efficacia};
  \node (i) [right=of m] {\textbf{Miglioramento}\\continuo};
  \draw[-{Stealth}] (c) -- (r); \draw[-{Stealth}] (r) -- (k);
  \draw[-{Stealth}] (k) -- (m); \draw[-{Stealth}] (m) -- (i);
  \draw[-{Stealth}] (i.south) -- ++(0,-0.45) -| (c.south);
\end{tikzpicture}
\end{center}
\begin{itemize}
  \item Il \textbf{perimetro} può coprire tutta l'azienda o un solo prodotto o servizio.
  \item Un ISMS nasce e cresce: se l'auditor non vede miglioramenti, solleva un rilievo
    (\S\ref{sec:iso-certificazione}).
  \item Un \textbf{controllo inefficace} è il caso peggiore: si spendono soldi e il rischio si materializza
    comunque. Se i controlli falliscono servono \emph{gap analysis} e piano di miglioramento (aggiornare,
    trasferire il rischio, abbandonare il processo).
  \item Implementare un ISMS da zero richiede in genere \textbf{un anno -- un anno e mezzo} per
    un'organizzazione media.
\end{itemize}

\begin{nota}[Che cosa sapere dell'ISMS (slide)]
Requisiti generici e flessibili; centralità delle parti interessate e delle loro aspettative; approccio
basato sul rischio; processi e miglioramento continuo; \textbf{i processi contano più delle politiche e dei
documenti}; prima si attuano i processi, poi ci si certifica (``\emph{implement first, certify later}'');
l'implementazione è un progetto lungo e complesso.
\end{nota}

\begin{integrazione}[Il ciclo PDCA]
Il miglioramento continuo dei sistemi di gestione ISO si descrive tradizionalmente con il ciclo
\textbf{Plan--Do--Check--Act}. La 27001 lo citava esplicitamente fino all'edizione 2005; dal 2013 adotta la
struttura comune a tutti i sistemi di gestione (clausole 4--10), che vi corrisponde così: \textbf{Plan} =
clausole 4--7 (contesto, leadership, pianificazione, supporto), \textbf{Do} = clausola~8 (attività
operative), \textbf{Check} = clausola~9 (valutazione delle prestazioni), \textbf{Act} = clausola~10
(miglioramento). Il ciclo delle cinque fasi della 27003 qui sopra è la stessa idea.
\end{integrazione}

% =====================================================================
\section{Le clausole 4--10}
\label{sec:iso-clausole}
% =====================================================================
Le clausole 4--10 contengono i requisiti dell'ISMS: chi dichiara la conformità alla norma \textbf{non può
escluderne nessuno}. (L'Allegato~A, invece, si applica secondo il rischio, con le esclusioni motivate nella
SoA.) Sintesi dei contenuti e delle osservazioni in aula:

\begin{center}
\small
\begin{tabularx}{\linewidth}{>{\raggedright\arraybackslash}p{2.8cm} X X}
\toprule
\textbf{Clausola} & \textbf{Contenuto} & \textbf{Osservazioni del corso}\\
\midrule
4 Contesto & questioni interne ed esterne (4.1), parti interessate e loro requisiti (4.2), perimetro
  dell'ISMS (4.3), istituzione e miglioramento dell'ISMS (4.4) & dettaglio non richiesto all'esame; le parti
  interessate includono fornitori, consulenti, a volte il pubblico\\
5 Leadership & impegno dell'alta direzione (5.1), politica di sicurezza (5.2), ruoli, responsabilità e
  autorità (5.3) & punto spesso critico: ``tutto parte dall'alto''\\
6 Pianificazione & azioni su rischi e opportunità, con processo di valutazione e di trattamento del rischio
  (6.1), obiettivi di sicurezza (6.2), pianificazione delle modifiche (6.3) & qui nascono metodo di
  valutazione, criteri di accettazione e SoA\\
7 Supporto & risorse (7.1), competenze (7.2), consapevolezza (7.3), comunicazione (7.4), informazioni
  documentate (7.5) & la \textbf{consapevolezza} è molto verificata (gli auditor intervistano anche la
  reception) ed è richiesta anche da GDPR e AI Act\\
8 Attività operative & pianificazione e controllo dei processi, compresi quelli esternalizzati (8.1);
  esecuzione della valutazione del rischio (8.2) e del piano di trattamento (8.3), con evidenze & la
  clausola~6 definisce il metodo, la~8 lo applica davvero\\
9 Valutazione delle prestazioni & monitoraggio e misurazione (9.1), audit interno (9.2), riesame della
  direzione (9.3) & KPI: ``non si controlla ciò che non si misura''; serve un \textbf{programma pluriennale}
  di audit interni; riesame almeno annuale (meglio trimestrale)\\
10 Miglioramento & miglioramento continuo (10.1), non conformità e azioni correttive (10.2) & senza
  miglioramenti visibili l'auditor solleva rilievi\\
\bottomrule
\end{tabularx}
\end{center}

\begin{integrazione}[Clausole 0--3 e informazioni documentate obbligatorie]
Le clausole 0--3 (introduzione, scopo, riferimenti normativi, termini) non contengono requisiti. Le
informazioni documentate che la 27001:2022 richiede espressamente, e che l'auditor chiede per prime, sono:
perimetro (4.3), politica (5.2), processo di valutazione (6.1.2) e di trattamento del rischio (6.1.3),
\textbf{Statement of Applicability} (6.1.3), obiettivi (6.2), evidenze di competenza (7.2), evidenze che i
processi operano come pianificato (8.1), risultati della valutazione (8.2) e del trattamento (8.3),
risultati del monitoraggio (9.1), programma e risultati degli audit interni (9.2), risultati del riesame
(9.3), non conformità e azioni correttive (10.2). La 6.1.2 chiede anche di individuare i \textbf{risk
owner}; la 6.1.3 di ottenere la loro approvazione del piano di trattamento e l'accettazione dei rischi
residui.
\end{integrazione}

% =====================================================================
\section{Certificazione e audit}
\label{sec:iso-certificazione}
% =====================================================================
\paragraph{Ciclo triennale.} Il certificato dura tre anni: audit di certificazione il primo anno,
\textbf{due audit di mantenimento} negli anni successivi (su parti della norma), poi
\textbf{ricertificazione} completa. Più a lungo un'organizzazione è certificata, più gli auditor sono
esigenti: il primo anno guardano soprattutto i documenti, poi testano i processi e intervistano il
personale, chiedendo anche le \textbf{dipendenze} tra processi e unità organizzative.

\paragraph{Audit interno e riesame.} L'audit interno non è un timbro: in un caso reale l'auditor ha
contestato l'assenza di un programma pluriennale di audit interni con perimetro definito in anticipo. La
direzione riesamina l'ISMS almeno una volta l'anno e deve lasciare traccia del proprio coinvolgimento.

\paragraph{Costi.} La consulenza per arrivare alla certificazione va da poche migliaia a decine di migliaia
di euro; il costo vero è far funzionare i processi e controllare i controlli. Per l'azienda la compliance è
un costo che non produce valore diretto, ma protegge clienti e pubblico.

\paragraph{Rilievi.} Quando qualcosa non va, l'auditor solleva un rilievo, di tre gravità crescenti:
\textbf{opportunità di miglioramento} (OFI), \textbf{non conformità minore}, \textbf{non conformità
maggiore}.

\begin{integrazione}[Come funziona l'audit di certificazione]
\begin{itemize}
  \item La certificazione è rilasciata da un \textbf{organismo di certificazione} accreditato (in Italia da
    Accredia); i requisiti per questi organismi sono nella ISO/IEC~27006.
  \item L'audit iniziale ha due fasi: \textbf{stage~1} (documentazione e prontezza: perimetro, valutazione del
    rischio, SoA, politica) e \textbf{stage~2} (verifica che il sistema funzioni davvero, con campioni,
    evidenze e interviste).
  \item \textbf{Non conformità maggiore}: un requisito manca del tutto o il sistema fallisce in modo
    sistematico; blocca il rilascio (o il mantenimento) del certificato finché non è chiusa con azioni
    correttive verificate. \textbf{Minore}: mancanza isolata che non compromette il sistema; richiede un piano
    di azioni correttive, verificato all'audit successivo. \textbf{OFI} (o osservazione): non è una non
    conformità, ma un suggerimento che l'auditor si aspetta di vedere considerato.
  \item Le certificazioni sulla versione 2013 sono scadute al termine della transizione (31 ottobre 2025):
    oggi si certifica solo sulla 27001:2022.
\end{itemize}
\end{integrazione}

% =====================================================================
\section{Allegato A: i 93 controlli}
\label{sec:iso-allegato-a}
% =====================================================================
Un \textbf{controllo} è una misura che mantiene o modifica il rischio. La revisione 2022 (che sostituisce
quella del 2013) conta \textbf{93 controlli}, di cui \textbf{11 nuovi}, in quattro temi:

\begin{center}
\small
\begin{tabularx}{\linewidth}{l l c X}
\toprule
\textbf{Sezione} & \textbf{Tema} & \textbf{N.} & \textbf{Riguardano}\\
\midrule
5 & organizzativi & 37 & tutto ciò che non rientra negli altri temi (politiche, ruoli, fornitori, incidenti)\\
6 & persone & 8 & le singole persone (selezione, formazione, lavoro da remoto)\\
7 & fisici & 14 & oggetti e luoghi fisici (perimetri, scrivanie sgombre, supporti)\\
8 & tecnologici & 34 & la tecnologia\\
\bottomrule
\end{tabularx}
\end{center}

\noindent I controlli fisici sono spesso trascurati (armadi chiusi a chiave, scrivanie sgombre). La
ISO~27002:2022 assegna a ogni controllo \textbf{cinque attributi}, utili per filtrarli e ordinarli:
\begin{enumerate}
  \item \textbf{tipo}: preventivo (evita l'incidente), detective (lo rileva), correttivo (limita il danno e
    ripristina);
  \item \textbf{proprietà} di sicurezza tutelate: C, I, A;
  \item \textbf{concetti di cybersicurezza}: identificare, proteggere, rilevare, rispondere, ripristinare;
  \item \textbf{capacità operative};
  \item \textbf{domini di sicurezza}.
\end{enumerate}

Non tutti i controlli valgono per ogni organizzazione: uno studio di consulenza che non sviluppa software non
ha codice sorgente da proteggere (8.4) né ambienti di sviluppo da separare (8.31). Lo standard evolve: un
controllo oggi sufficiente può non bastare alla prossima revisione. Il focus del corso sono i
\textbf{controlli tecnologici}.

\begin{integrazione}[Le novità del 2022]
La versione 2013 aveva 114 controlli in 14 domini; nel 2022 sono stati accorpati e ridotti a 93. Gli 11
nuovi sono: \textbf{5.7} threat intelligence, \textbf{5.23} sicurezza nell'uso dei servizi cloud,
\textbf{5.30} prontezza ICT per la continuità operativa, \textbf{7.4} monitoraggio della sicurezza fisica,
\textbf{8.9} gestione della configurazione, \textbf{8.10} cancellazione delle informazioni, \textbf{8.11}
data masking, \textbf{8.12} data leakage prevention, \textbf{8.16} attività di monitoraggio, \textbf{8.23}
web filtering, \textbf{8.28} secure coding. La \emph{threat intelligence} citata a lezione come requisito
``dell'ultima revisione'' è quindi il controllo 5.7, organizzativo.
\end{integrazione}

\subsection{I controlli tecnologici 8.1--8.34}
Titoli dalle slide; idea in sintesi; nell'ultima colonna gli esempi e i collegamenti fatti a lezione
(Capitoli~\ref{lez:2026-10-19} e~\ref{lez:2026-10-23}). L'asterisco indica i controlli introdotti nel 2022
(integrazione).

{\small
\begin{longtable}{>{\raggedright\arraybackslash}p{0.8cm} >{\raggedright\arraybackslash}p{3.3cm}
  >{\raggedright\arraybackslash}p{5.4cm} >{\raggedright\arraybackslash}p{4.6cm}}
\toprule
\textbf{N.} & \textbf{Controllo} & \textbf{Idea chiave} & \textbf{Dal corso}\\
\midrule
\endfirsthead
\toprule
\textbf{N.} & \textbf{Controllo} & \textbf{Idea chiave} & \textbf{Dal corso}\\
\midrule
\endhead
\midrule
\multicolumn{4}{r}{\emph{continua}}\\
\endfoot
\bottomrule
\endlastfoot
8.1 & User endpoint devices & gestire e proteggere portatili, smartphone e postazioni: hardening,
  aggiornamenti, MDM, cifratura dei dati a riposo & protezione proporzionata ai dati trattati (il portatile
  del CEO più di quello della reception)\\
8.2 & Privileged access rights & limitare e tracciare gli account amministrativi; minimo privilegio,
  accesso temporaneo, niente credenziali condivise & amministratore locale concesso a tutti ``otto volte su
  dieci''; caso dell'account pubblicitario compromesso\\
8.3 & Information access restriction & accesso a dati e funzioni solo quanto il ruolo richiede & coerenza
  tra classificazione del dato e diritti concessi\\
8.4 & Access to source code & proteggere repository e artefatti; accessi limitati e registrati & sviluppatori
  esterni con accesso completo al repository\\
8.5 & Secure authentication & autenticazione forte e proporzionata (MFA); segreti protetti per tutto il ciclo
  di vita & MFA ovunque; mai credenziali in chiaro al neoassunto\\
8.6 & Capacity management & monitorare e pianificare risorse perché l'esaurimento non causi fermi &
  disponibilità; nel cloud è demandata al fornitore\\
8.7 & Protection against malware & difesa a strati, ripristino e consapevolezza; presumere la compromissione
  & azienda ferma un mese e mezzo senza protezione né backup\\
8.8 & Management of technical vulnerabilities & inventario degli asset, scansione, priorità secondo il
  rischio, processo di patching & server dimenticati mai aggiornati; vulnerabilità PHP di ottobre 2026\\
8.9* & Configuration management & baseline sicure, modifiche controllate e registrate & sapere in ogni
  momento quanto si è lontani dalla baseline\\
8.10* & Information deletion & cancellare in modo sicuro su ogni supporto quando il dato non serve più &
  conservazione limitata (GDPR); backup inclusi; nel cloud evidenze difficili, documentare la gestione\\
8.11* & Data masking & sostituire dati sensibili con dati fittizi realistici fuori dalla produzione & nasce
  per il GDPR ma vale per ogni dato riservato; dati sintetici\\
8.12* & Data leakage prevention & rilevare e bloccare l'esfiltrazione (email, web, supporti, rete) & richiede
  di monitorare le comunicazioni: tensione con GDPR e art.~4 Statuto dei lavoratori\\
8.13 & Information backup & backup regolari, protetti, con copie offline o immutabili e obiettivi di
  ripristino & vale solo se si può ripristinare: \textbf{test di ripristino} periodici\\
8.14 & Redundancy of information processing facilities & ridondanza per rispettare i livelli di
  disponibilità & disaster recovery, sito \emph{hot}; equilibrio con il budget\\
8.15 & Logging & registrare gli eventi rilevanti in log protetti da manomissioni & ``traccia documentale'';
  limite dei 21 giorni sui metadati email dei dipendenti\\
8.16* & Monitoring activities & sorvegliare reti e sistemi e gestire gli alert (SIEM, SOC) & alert del
  venerdì sera mai letto: terabyte esfiltrati nel weekend\\
8.17 & Clock synchronization & orologi allineati a una fonte affidabile & senza, nessuna timeline
  attendibile degli eventi\\
8.18 & Use of privileged utility programs & limitare i programmi che possono scavalcare i controlli & niente
  software con privilegi elevati senza la sicurezza\\
8.19 & Installation of software on operational systems & installazione solo di software approvato e
  verificato & whitelist tramite MDM\\
8.20 & Networks security & proteggere e controllare le reti (segregazione, cifratura, monitoraggio del
  traffico) & IoT (telecamere, DVR esposti); \textbf{zero trust}; il Wi-Fi è un rischio in sé\\
8.21 & Security of network services & funzioni di sicurezza e livelli di servizio dei servizi di rete
  interni ed esterni & nel Lab~3: SLA e mitigazione DDoS del provider\\
8.22 & Segregation of networks & segmentare per livello di fiducia e funzione, limitare il \emph{blast
  radius} & VLAN, mai reti piatte; separare OT/SCADA dal back office\\
8.23* & Web filtering & bloccare contenuti web malevoli o non ammessi & non commentato in aula\\
8.24 & Use of cryptography & crittografia con corretta gestione delle chiavi (riservatezza, autenticità,
  integrità) & forte quanto il controllo sulle chiavi; chiavi delegate al cloud = dipendenza\\
8.25 & Secure development life cycle & sicurezza in tutto il ciclo di sviluppo & rilevante solo per chi
  sviluppa\\
8.26 & Application security requirements & requisiti di sicurezza fin da specifica e progettazione & anche
  se lo sviluppo è esternalizzato; privacy by design (banca spagnola sanzionata)\\
8.27 & Secure system architecture and engineering principles & principi di architettura sicura & sistema di
  sviluppo ``ombra'' fuori perimetro, con credenziali in chiaro\\
8.28* & Secure coding & standard di codifica, revisione del codice, strumenti & OWASP come riferimento\\
8.29 & Security testing in development and acceptance & test di sicurezza prima del rilascio &
  vulnerability assessment e penetration test, non solo test funzionali\\
8.30 & Outsourced development & supervisionare la sicurezza del software sviluppato da terzi & controlli e
  audit sul fornitore\\
8.31 & Separation of development, test and production environments & ambienti separati con accessi
  diversi & assicurazione che usava dati di produzione nei test\\
8.32 & Change management & processo formale: valutare, documentare, testare, approvare & non solo IT: anche
  traslochi e riorganizzazioni; piano di rollback\\
8.33 & Test information & scegliere e proteggere i dati di test; preferire dati sintetici & i dati sintetici
  conservano le proprietà statistiche; produzione solo in casi documentati\\
8.34 & Protection of information systems during audit testing & pianificare e concordare audit e test per
  non disturbare i sistemi & penetration test su un SaaS in produzione: carico concordato prima\\
\end{longtable}
}

\begin{nota}[Buon senso, ma obbligatorio]
Molti controlli sembrano ovvi (backup, monitoraggio, aggiornamenti). Sono requisiti della norma e buone
pratiche che chiunque dovrebbe adottare, certificato o no. Il buon senso, osserva Gianluca, non è
distribuito in modo uniforme. Le \emph{quick win} esistono: il \textbf{backup} è la prima (un cliente senza
backup è rimasto bloccato da un ransomware per oltre due mesi).
\end{nota}

\subsection{Sicurezza contro privacy: DLP, logging e monitoraggio}
\label{sec:iso-privacy}
I controlli 8.12, 8.15 e 8.16 raccolgono dati sulle attività delle persone. Prima di attivarli bisogna
chiedersi se quei dati sono \textbf{dati personali}: se sì, si applica tutto il GDPR, e per i dipendenti
anche il diritto del lavoro.
\begin{itemize}
  \item \textbf{DLP (8.12)}: per bloccare l'esfiltrazione bisogna analizzare email e uso dei dispositivi,
    cioè controllare l'attività dei lavoratori. Chi lo esclude dalla SoA (lo stesso vale per l'8.8) deve avere
    una motivazione documentata molto solida.
  \item \textbf{Logging (8.15)}: secondo un documento del Garante citato in aula, i metadati della posta
    elettronica dei dipendenti si possono conservare per non più di \textbf{21 giorni}, salvo accordo
    sindacale o autorizzazione. Ma molti servizi (Microsoft~365, Google Workspace) hanno conservazioni più
    lunghe non modificabili, e la sicurezza richiede log molto più lunghi (in Italia un incidente viene
    scoperto in media dopo circa sei mesi, secondo il docente). Un documento di accountability che
    giustificava una conservazione più lunga non ha evitato una sanzione a una Regione.
  \item \textbf{Monitoraggio (8.16)}: raccogliere senza presidiare non serve; ma presidiare significa
    trattare dati personali, con le stesse cautele.
\end{itemize}

\begin{integrazione}[Art.~4 dello Statuto dei lavoratori e documento del Garante]
L'art.~4 della L.~300/1970 ammette gli strumenti da cui derivi anche la possibilità di \textbf{controllo a
distanza} dei lavoratori solo per esigenze organizzative e produttive, di sicurezza del lavoro o di tutela
del patrimonio aziendale, e previo \textbf{accordo sindacale} o, in mancanza, autorizzazione
dell'Ispettorato del lavoro; non servono per gli strumenti usati per rendere la prestazione e per registrare
accessi e presenze. Il documento di indirizzo del Garante sui programmi di gestione della posta elettronica e
sui metadati (pubblicato nel 2024) indicava in origine 7 giorni di conservazione, portati a 21 dopo la
consultazione pubblica; oltre quel termine si ricade nelle garanzie dell'art.~4.
\end{integrazione}

% =====================================================================
\section{La gestione del rischio secondo la ISO}
\label{sec:iso-rischio}
% =====================================================================

\subsection{Perché e con quali mattoni}
Secondo le slide il risk management protegge gli asset, assicura continuità e resilienza, soddisfa obblighi
legali e contrattuali, orienta gli investimenti in sicurezza e crea fiducia. Il \textbf{rischio} è
l'effetto dell'incertezza sugli obiettivi, di solito espresso come combinazione delle conseguenze di un
evento e della sua probabilità; il \textbf{rischio informativo} è la possibilità che una minaccia sfrutti
una vulnerabilità di un asset e causi un danno all'\textbf{organizzazione}.

\begin{center}
\small
\begin{tabularx}{\linewidth}{>{\raggedright\arraybackslash}p{2.4cm} >{\raggedright\arraybackslash}p{4.2cm} X}
\toprule
\textbf{Concetto} & \textbf{Definizione} & \textbf{Esempi e osservazioni}\\
\midrule
Asset & ciò che ha valore e va protetto & informazioni, hardware, software, \textbf{persone}, beni intangibili
  (reputazione, marchio, avviamento). Le persone sono l'asset più dimenticato: chi detiene conoscenza non
  documentata (\emph{key people}) è critico. Primo passo: inventario e valore degli asset\\
Minaccia & evento o attore che potrebbe danneggiare un asset & dolose (malware, phishing, DDoS, insider),
  accidentali (errore umano), ambientali (incendi, alluvioni, blackout), tecniche (bug, guasti)\\
Vulnerabilità & debolezza che una minaccia può sfruttare & patch mancanti, password deboli o di default,
  scarsa formazione, configurazioni insicure, lacune fisiche (un cassetto senza chiave), applicazioni mal
  progettate\\
Impatto & entità del danno se la minaccia sfrutta la vulnerabilità & perdite economiche, danno reputazionale,
  sanzioni, fermo operativo\\
Probabilità & possibilità che la minaccia sfrutti la vulnerabilità & di solito qualitativa (alta, media,
  bassa)\\
\bottomrule
\end{tabularx}
\end{center}

\[
\text{Rischio} = \text{Impatto} \times \text{Probabilità}
\]

\paragraph{Probabilità che cambia.} Esempio della settimana del 23/10: una vulnerabilità PHP di esecuzione
remota di codice aveva subito impatto altissimo ma probabilità bassa; appena è uscito un \emph{proof of
concept}, la probabilità è schizzata e tutti hanno corso a installare le patch.

\paragraph{Threat intelligence.} Le vulnerabilità tecniche si trovano con scansioni e penetration test,
quelle umane passano dall'ingegneria sociale. Non si deve conoscere da soli ogni minaccia: la 27001:2022
chiede di individuare le proprie \textbf{fonti di threat intelligence} (es.\ i bollettini dell'ACN), che
spesso stimano già la probabilità di sfruttamento delle vulnerabilità note e dimostrano all'auditor
l'attenzione continua al rischio.

\subsection{Il processo di valutazione}
La 27001 richiede un processo sistematico e \textbf{documentato}; le slide lo scandiscono così:
\begin{enumerate}
  \item \textbf{Contesto}: perimetro della valutazione (sistemi, reparti, organizzazione intera), criteri di
    accettazione e propensione al rischio; inventario e classificazione degli asset per criticità e valore,
    con le loro dipendenze.
  \item \textbf{Identificazione}: minacce interne ed esterne pertinenti agli asset; vulnerabilità tecniche,
    di processo, fisiche e umane (scansioni, penetration test, audit, analisi degli incidenti).
  \item \textbf{Analisi}: stima di probabilità e impatto sull'organizzazione (economico, operativo,
    reputazionale, legale) con scale coerenti.
  \item \textbf{Ponderazione} (\emph{evaluation}): livello di rischio, confronto con i criteri di
    accettazione, classifica delle priorità.
  \item \textbf{Trattamento}: scelta dell'opzione e dei controlli.
  \item \textbf{Monitoraggio e riesame}: il ciclo riparte.
\end{enumerate}

\paragraph{Prioritizzare.} Di solito si parte dal rischio più alto, ma conta anche il tempo di attuazione:
una \emph{quick win} può passare avanti a un rischio più alto che richiede un progetto lungo. Decide il
vertice, che chiede a chi valuta costi, tempi e proposta di priorità (rapporto costi/benefici; alcuni
controlli coprono un solo rischio, altri molti). Il risk manager prudente considera anche lo
\textbf{scenario peggiore} e i \emph{cigni neri}: il COVID-19 ha colto il mondo impreparato perché la sua
probabilità era stata giudicata troppo bassa.

\subsection{Scale e matrice}
Le matrici sono quasi sempre \textbf{qualitative} (basso, medio, alto; al massimo quattro livelli); le scale
quantitative sono rare. La ISO \textbf{non fissa} una matrice: si deve essere coerenti e saper difendere le
scelte. La più usata è la \textbf{3$\times$3}, quella del Lab~3 (Capitolo~\ref{lab:03}):

\begin{center}
\begin{tikzpicture}[x=1.7cm, y=1.0cm]
  \foreach \pr in {1,2,3}{
    \foreach \im in {1,2,3}{
      \pgfmathtruncatemacro{\scr}{\pr*\im}
      \ifnum\scr<3 \def\cellcol{green!25}\else\ifnum\scr<5 \def\cellcol{yellow!45}\else\def\cellcol{red!35}\fi\fi
      \fill[fill=\cellcol] (\pr-1,\im-1) rectangle (\pr,\im);
      \draw (\pr-1,\im-1) rectangle (\pr,\im);
      \node[font=\small] at (\pr-0.5,\im-0.5) {\scr};
    }
  }
  \node[font=\footnotesize] at (0.5,-0.3) {1 bassa};
  \node[font=\footnotesize] at (1.5,-0.3) {2 media};
  \node[font=\footnotesize] at (2.5,-0.3) {3 alta};
  \node[font=\footnotesize, anchor=east] at (-0.1,0.5) {1 basso};
  \node[font=\footnotesize, anchor=east] at (-0.1,1.5) {2 medio};
  \node[font=\footnotesize, anchor=east] at (-0.1,2.5) {3 alto};
  \node[font=\small\bfseries] at (1.5,-0.8) {Probabilità};
  \node[font=\small\bfseries, rotate=90] at (-1.45,1.5) {Impatto};
  \node[font=\footnotesize, anchor=west, align=left] at (3.3,1.5)
    {1--2: \textbf{basso}\\3--4: \textbf{medio}\\6--9: \textbf{alto}\\(il 5 non esiste)};
\end{tikzpicture}
\end{center}

\paragraph{Livelli di impatto sull'organizzazione (indicazioni del docente).} \textbf{Basso}: perdita
economica limitata, nessun fermo, recupero facile. \textbf{Medio}: un ``mal di testa'' misurabile (sanzioni,
fermo limitato) ma superabile. \textbf{Alto}: evento che può mettere a rischio l'\textbf{esistenza stessa}
dell'organizzazione (perdita totale dei dati senza backup). Le ancore precise della scala 1--3 usate nel
laboratorio sono nel Capitolo~\ref{lab:03}.

\paragraph{Tre dimensioni, un livello.} Con il metodo del corso la probabilità è unica per ogni rischio,
mentre l'impatto si valuta separatamente su \textbf{C}, \textbf{I} e \textbf{A}; si ottengono tre punteggi e
il livello del rischio è il loro \textbf{massimo}:
\[
\text{Livello} = \max\left(P \times I_C,\; P \times I_I,\; P \times I_A\right)
\]

\begin{esame}
\textbf{Errore tipico: fare la media} di C, I, A. Un rischio con punteggi 9/1/1 è \textbf{alto}, non medio:
basta una dimensione alta. E il livello da solo non basta: bisogna sapere \textbf{quale dimensione} lo rende
alto, per scegliere controlli che agiscano su quella (nel Lab~3, R-01 è alto per la riservatezza, R-17 per
la disponibilità).
\end{esame}

\subsection{Rischio inerente e residuo}
\begin{itemize}
  \item \textbf{Rischio inerente}: il rischio come si presenta prima del trattamento in esame. Nel Lab~3 è
    stimato \textbf{con i controlli che l'organizzazione ha oggi} (una baseline debole).
  \item \textbf{Rischio residuo}: il rischio che resta dopo i controlli scelti. Nel Lab~3 si
    \textbf{rivaluta da capo} (nuova probabilità, nuovi impatti), non si calcola ``inerente meno
    efficacia''; probabilità e impatti residui non possono superare gli inerenti.
  \item Anche il residuo è oggetto di decisione: va accettato formalmente, oppure trattato ancora.
\end{itemize}
Lo stesso schema compare nel GDPR: nella DPIA (Capitoli~\ref{lez:2026-09-28} e~\ref{lez:2026-10-05}) il
professore lo riassume come ``residuo = inerente meno mitigazione'', e se il residuo resta alto si passa alla
consultazione preventiva (art.~36) o si ferma il trattamento.

\begin{nota}[Due definizioni di ``inerente'']
Gli appunti 2025/26 definiscono il rischio inerente come quello \emph{in assenza di qualsiasi controllo}; il
materiale del Lab~3 lo stima \emph{con i controlli già in essere}. Sono due convenzioni diffuse: l'importante
è dichiarare quale si usa e applicarla in modo coerente. Gli stessi appunti mettono in guardia dalla
``matrice verde'': mostrare solo il residuo dà l'impressione che tutto sia sotto controllo anche quando il
piano di trattamento non è ancora attuato. Il residuo è ciò che si \emph{spera}, l'inerente ciò che
\emph{c'è}: verificare sempre piano e tempi di attuazione.
\end{nota}

\subsection{Propensione e tolleranza al rischio}
\begin{center}
\small
\begin{tabularx}{\linewidth}{>{\raggedright\arraybackslash}p{2.4cm} X X}
\toprule
 & \textbf{Risk appetite (propensione)} & \textbf{Risk tolerance (tolleranza)}\\
\midrule
Che cos'è & disponibilità complessiva ad assumere rischi per raggiungere gli obiettivi & livello di rischio
  accettato per specifiche attività o processi\\
Livello & strategico: decide il board con l'alta direzione & operativo: soglie e criteri numerici, diversi
  per funzione e tipo di asset\\
Nel Lab~3 & le soglie 1--2 / 3--4 / 6--9 (alto non accettabile) & owner e data obiettivo del singolo
  rischio\\
\bottomrule
\end{tabularx}
\end{center}
Esempio di Gianluca: su una scala 1--9 si accetta ciò che sta sotto 3, il resto va trattato. Nella sua
esperienza, però, quasi nessun rischio inerente che compare in un registro è davvero accettabile così com'è:
i rischi davvero minimi non vengono nemmeno percepiti, e quindi non vengono mappati.

\subsection{Le quattro opzioni di trattamento}
\begin{center}
\small
\begin{tabularx}{\linewidth}{>{\raggedright\arraybackslash}p{2.3cm} X X}
\toprule
\textbf{Opzione} & \textbf{Che cosa significa} & \textbf{Attenzione}\\
\midrule
Mitigare (ridurre) & controlli che riducono probabilità e/o impatto fino a un livello accettabile &
  l'opzione più comune; il residuo va comunque accettato dalla direzione\\
Accettare & nessuna ulteriore misura: il rischio rientra nella propensione o il trattamento costerebbe
  troppo & richiede \textbf{approvazione formale}; un rischio inerente \textbf{alto} non si può accettare\\
Evitare & cessare o non avviare l'attività che genera il rischio & corrisponde allo ``stop del
  trattamento'' visto per la DPIA\\
Trasferire & condividere il rischio con un terzo (assicurazione, outsourcing, SLA) & sposta la perdita
  economica, \textbf{non la responsabilità}; la scelta del terzo va motivata (``non basta una firma su un
  foglio''); la \emph{shared responsibility matrix} del cloud è un trasferimento parziale\\
\bottomrule
\end{tabularx}
\end{center}
Il trattamento è \textbf{strategico}, non solo tattico: non ha senso spendere per mitigare il rischio di un
processo che verrà dismesso a breve. Nella prima lezione Gianluca aveva presentato tre esiti (mitigare,
accettare, trasferire), definendo il trasferimento il ``Santo Graal'' del risk manager; le slide sul risk
management aggiungono \emph{evitare}.

\paragraph{Scegliere i controlli.} Si parte dalla \textbf{minaccia e dagli asset}, mai dal catalogo: quale
controllo interrompe \emph{quella} causa? Si sceglie l'insieme minimo sufficiente, pertinente, efficace e
proporzionato; non si inventa un rischio per dare una ``casa'' a un controllo. La 27002 offre il catalogo e
la guida all'attuazione.

\subsection{Statement of Applicability}
La \textbf{SoA} è un documento \textbf{obbligatorio} per la conformità alla 27001:
\begin{itemize}
  \item elenca \textbf{tutti} i 93 controlli dell'Allegato~A;
  \item per ciascuno dichiara se è applicabile;
  \item ne giustifica l'inclusione o l'esclusione;
  \item ne indica lo stato di attuazione.
\end{itemize}
Documenta le decisioni prese nel trattamento del rischio (quali controlli per quali rischi) ed è centrale
negli audit di certificazione. Per Gianluca conta ancora di più saper giustificare perché un controllo
\textbf{non} è attuato: l'auditor può contestare il legame tra controllo e rischio e chiedere di correggere
la SoA. Il registro va dal rischio al controllo; la SoA dal controllo al rischio.

\subsection{Piano di trattamento}
Per ogni rischio da trattare il piano indica: opzione scelta; azioni concrete; \textbf{owner} di ogni azione;
tempi; risorse (budget e persone: il management chiederà sempre ``quanto costa?''); giustificazione. La
giustificazione discende dalla propensione al rischio:
\begin{itemize}
  \item si \textbf{accetta} perché il rischio è sotto soglia;
  \item si \textbf{mitiga} perché l'inerente è sopra soglia e i controlli portano il residuo entro la soglia;
  \item si \textbf{trasferisce} ciò che resta fuori dal proprio controllo;
  \item si \textbf{evita} quando nessun residuo è accettabile.
\end{itemize}
La matrice RACI è fondamentale nella procedura di gestione dei data breach, ma non è richiesta nel piano di
trattamento.

\subsection{Il risk register}
\label{sec:iso-register}
Il registro dei rischi è un documento \textbf{centrale e vivo} che raccoglie tutti i rischi individuati.
Non ha un formato obbligatorio (Excel o piattaforme dedicate). Campi minimi:

\begin{center}
\small
\begin{tabularx}{\linewidth}{>{\raggedright\arraybackslash}p{4.2cm} X}
\toprule
\textbf{Campo} & \textbf{Perché serve}\\
\midrule
ID univoco & tracciare il rischio in tutti i documenti e progetti\\
Asset coinvolti & sapere che cosa si protegge\\
Minaccia e vulnerabilità & la causa da interrompere: da qui si scelgono i controlli\\
Probabilità e impatto (C, I, A) inerenti & base della valutazione\\
Livello di rischio calcolato & con la matrice; massimo delle tre dimensioni\\
Controlli e misure concrete & azioni, non nomi di controlli (``MFA su tutti i servizi cloud'', non ``8.5'')\\
Giustificazione ed efficacia & perché quelle misure, se agiscono su probabilità o impatto, perché il residuo
  è accettabile\\
Probabilità, impatti e livello residui & dopo il trattamento\\
Decisione di trattamento & una delle quattro opzioni\\
Risk owner & chi decide, firma il trattamento e ne risponde\\
Data obiettivo e ultima revisione & tolleranza del singolo rischio; monitoraggio\\
\bottomrule
\end{tabularx}
\end{center}
\noindent I campi delle slide sono ID, asset, minaccia, vulnerabilità, probabilità, impatto, livello, owner e
trattamento proposto; a lezione Gianluca ha aggiunto la distinzione tra valori inerenti e residui, e il Lab~3
la giustificazione, l'efficacia e le date. Il registro garantisce trasparenza e accountability ed è la base
del trattamento.

\subsection{Monitoraggio continuo e altri framework}
I rischi cambiano, nascono e muoiono: il registro non si chiude in un cassetto. Vanno riesaminati
periodicamente la valutazione, l'efficacia dei controlli, l'avanzamento del piano e i rischi emergenti, in
una vera \textbf{cultura del rischio}. Non serve reinventare la ruota: oltre alla ISO esistono
\textbf{NIST CSF} (funzioni identificare, proteggere, rilevare, rispondere, ripristinare; diffuso negli USA),
\textbf{FAIR} (quantifica il rischio in termini economici), \textbf{OCTAVE} (rischi operativi,
autovalutazione) e le metodologie \textbf{ENISA}, in particolare per i dati personali.

\begin{integrazione}[ISO 31000 e ISO/IEC 27005]
La definizione ``effetto dell'incertezza sugli obiettivi'' viene dalla ISO~31000, la norma generale sulla
gestione del rischio, da cui la 27005 deriva per la sicurezza delle informazioni. La 27005:2022 ammette due
modi di identificare i rischi: \emph{basato sugli asset} (asset, minacce, vulnerabilità: quello del corso) e
\emph{basato sugli eventi} (scenari di alto livello e loro conseguenze, senza partire dall'inventario).
\end{integrazione}

% =====================================================================
\section{Rischio ISO e rischio GDPR/ENISA}
\label{sec:iso-enisa}
% =====================================================================
Il cambio di prospettiva è il punto centrale del risk management del corso. Nel GDPR il rischio si misura
sui \textbf{diritti e le libertà degli interessati}; nella ISO sull'\textbf{organizzazione}. Esempio di
Gianluca: il DPO di un'azienda preoccupato per le conseguenze di un incidente su operatività e processi HR
non ha, solo per questo, un problema di notifica al Garante: sono valutazioni distinte.

\subsection{La metodologia ENISA per i dati personali}
ENISA ha pubblicato, oltre alla metodologia di severità dei breach, linee guida (pensate anche per le PMI)
per valutare il rischio dei trattamenti di dati personali, compatibili con la 27001 e mappate sui controlli
della 27002. Quattro passi:
\begin{enumerate}
  \item \textbf{Contesto del trattamento}: quali dati, finalità, mezzi, luogo (on premise o cloud), categorie
    di interessati, destinatari; tutte le fasi (raccolta, conservazione, uso, trasferimento, cancellazione).
    Risultato: il perimetro della valutazione.
  \item \textbf{Impatto sugli individui}, su quattro livelli: \emph{basso} (piccoli disagi superabili),
    \emph{medio} (disagi significativi: costi, stress, servizi negati), \emph{alto} (conseguenze serie: perdite
    economiche, perdita del lavoro), \emph{molto alto} (conseguenze irreversibili). Si valuta separatamente
    per riservatezza, integrità e disponibilità e \textbf{si prende il più alto}.
  \item \textbf{Probabilità della minaccia} su quattro aree, ciascuna con punteggio 1--3: risorse di rete e
    tecniche; processi e procedure; persone e terze parti; settore e dimensione. Totale 4--5 bassa, 6--8
    media, 9--12 alta.
  \item \textbf{Valutazione del rischio} con la matrice:
\end{enumerate}
\begin{center}
\small
\begin{tabular}{llll}
\toprule
\textbf{Probabilità} & \textbf{Impatto basso} & \textbf{Impatto medio} & \textbf{Impatto alto/molto alto}\\
\midrule
Bassa & basso & medio & alto\\
Media & medio & alto & alto\\
Alta & alto & alto & alto\\
\bottomrule
\end{tabular}
\end{center}
\noindent La matrice ENISA è volutamente prudente: molte combinazioni finiscono in ``alto''. Il livello di
rischio orienta la scelta delle misure ai sensi dell'art.~32.

\subsection{Il confronto}
\begin{center}
\small
\begin{tabularx}{\linewidth}{>{\raggedright\arraybackslash}p{2.5cm} X X}
\toprule
\textbf{Aspetto} & \textbf{ISO/IEC 27001} & \textbf{ENISA (dati personali)}\\
\midrule
Ambito & sicurezza di tutte le informazioni & trattamenti di dati personali\\
Chi subisce l'impatto & l'\textbf{organizzazione} (C, I, A dei suoi asset) & gli
  \textbf{interessati}: diritti e libertà\\
Metodo & libero, purché coerente e documentato & strutturato in quattro passi\\
Probabilità & stimata dall'organizzazione (spesso con threat intelligence) & questionario su quattro aree,
  punteggio 4--12\\
Combinazione di C, I, A & massimo dei tre punteggi (metodo del corso) & massimo dei tre impatti\\
Controlli & 93 controlli dell'Allegato~A, SoA & misure proporzionate al livello di rischio, mappabili sulla
  27002\\
Riferimento & certificazione ISO (volontaria) & art.~32 GDPR (obbligo di legge)\\
Esito tipico & piano di trattamento e registro dei rischi & misure tecniche e organizzative adeguate;
  eventuale DPIA (art.~35)\\
\bottomrule
\end{tabularx}
\end{center}
I due approcci si usano spesso \textbf{insieme}: ENISA per i trattamenti di dati personali, ISO per l'ISMS,
le misure ENISA mappate sui controlli 27002 e un \textbf{unico registro con doppia valutazione}. Il punto di
contatto tra i due piani: un impatto C~=~3 sull'organizzazione, nel Lab~3, è ancorato a un \textbf{breach
da notificare ex art.~33}, cioè a un rischio per gli interessati.

\begin{nota}[Tre strumenti da non confondere]
\begin{itemize}
  \item \textbf{Risk assessment ISO}: rischio per l'organizzazione, \emph{prima} dell'incidente.
  \item \textbf{Valutazione del rischio ENISA per i dati personali} (e DPIA): rischio per gli interessati,
    \emph{prima} del trattamento o dell'incidente.
  \item \textbf{Metodologia ENISA di severità dei breach} (DPC, EI, CB): gravità di una violazione
    \emph{già avvenuta}, per decidere su notifica e comunicazione (Capitolo~\ref{cap:data-breach}).
\end{itemize}
Nel Lab~1 la lezione era la stessa: \textbf{la severità non è il rischio}.
\end{nota}

\begin{integrazione}[Un esempio sui due piani]
Un ransomware cifra il gestionale di uno studio professionale per due giorni, senza esfiltrazione. Lato
ISO: impatto alto sulla \textbf{disponibilità} (scadenze dei clienti mancate, fermo, reputazione), quindi
rischio alto per l'organizzazione. Lato GDPR: è un breach di disponibilità; se i dati si ripristinano dai
backup in tempi brevi e nessun interessato subisce conseguenze, la severità può risultare bassa e non
richiedere notifica, ma la violazione va comunque registrata (art.~33(5)). Stesso evento, due valutazioni,
due registri di riferimento.
\end{integrazione}

% =====================================================================
\section{Errori tipici e messaggi chiave}
\label{sec:iso-errori}
% =====================================================================
\begin{esame}
\textbf{Messaggi chiave del docente}
\begin{itemize}
  \item L'\textbf{ISMS è la base} della 27001: un sistema basato sul rischio e in miglioramento continuo,
    non un timbro. \textbf{I processi contano più dei documenti}.
  \item Numeri: \textbf{93 controlli} in 4 temi (37 / 8 / 14 / 34); clausole \textbf{4--10} non escludibili;
    ciclo di certificazione \textbf{triennale}.
  \item La certificazione riguarda l'\textbf{organizzazione}, è volontaria e non garantisce né sicurezza né
    conformità legale.
  \item Rischio ISO = impatto sull'\textbf{organizzazione}; rischio GDPR/ENISA = impatto sugli
    \textbf{interessati}.
  \item Il risk management è il fondamento della cybersicurezza: asset, minacce, vulnerabilità, impatto,
    probabilità; processo strutturato; trattamento strategico; decisioni documentate nella \textbf{SoA};
    monitoraggio continuo.
  \item Non fermarsi al rischio inerente: la decisione si prende sul \textbf{residuo}.
  \item Ogni scelta (anche l'esclusione di un controllo) va \textbf{motivata e documentata}: conta il
    ragionamento, non il numero.
\end{itemize}
\end{esame}

\begin{esame}
\textbf{Errori tipici}
\begin{enumerate}
  \item Fare la \textbf{media} di C, I, A invece del massimo.
  \item Usare lo stesso impatto per le tre dimensioni senza motivo.
  \item \textbf{Accettare} un rischio inerente alto; lasciare un residuo alto.
  \item Credere che \textbf{trasferire} elimini la responsabilità: si sposta la perdita economica, non
    l'accountability.
  \item Confondere il rischio per l'organizzazione con il rischio per gli interessati (e il risk assessment
    con la severità del breach).
  \item Sottostimare la probabilità di un'organizzazione senza controlli.
  \item Scrivere il nome del controllo al posto della misura concreta; partire dal catalogo invece che dalla
    minaccia.
  \item Residuo uguale all'inerente senza spiegazione, o addirittura superiore.
  \item Pensare che la certificazione 27001 dimostri la conformità a GDPR, NIS2 o CRA.
  \item Escludere controlli scomodi (8.8, 8.12) dalla SoA senza una motivazione solida.
\end{enumerate}
\end{esame}

% =====================================================================
\section{Punti da verificare}
\label{sec:iso-verifica}
% =====================================================================
\begin{verifica}
\begin{itemize}
  \item \textbf{ISO/IEC 27090} (sicurezza dei sistemi di IA): citata a voce come appena approvata; numero e
    stato di pubblicazione da confermare.
  \item \textbf{27017 e 27018}: compaiono nel diagramma della famiglia nelle slide (secondo gli appunti di
    recupero) e negli appunti 2025/26, ma non sono state commentate a lezione quest'anno.
  \item \textbf{PDCA}: secondo gli appunti di recupero il ciclo compare in un'immagine della slide sull'ISMS;
    non è nel testo delle slide né nella trascrizione.
  \item \textbf{Limite dei 21 giorni} e \textbf{sanzione alla Regione} (circa \texteuro{}50.000): dal parlato;
    estremi del documento del Garante e del provvedimento da verificare. Anche la sanzione di
    \texteuro{}500.000 alla banca spagnola per mancata privacy by design.
  \item \textbf{Sovrapposizione 75--80\% con SOC~2} e \textbf{sei mesi} di tempo medio di rilevazione in
    Italia: stime del docente.
  \item \textbf{Linee guida ENISA per le PMI}: titolo e anno del documento non indicati in aula.
  \item \textbf{Tipi di controllo}: a lezione vulnerability assessment e penetration test sono stati
    portati come esempi di controlli \emph{preventivi}; sono anche (e forse soprattutto) attività di
    individuazione delle vulnerabilità. Classificazione da confermare.
\end{itemize}
\end{verifica}

\chapter{Cyber Resilience Act (CRA)}
\label{cap:cra}

\dacompletare{Fonti: slide \emph{05.~PoliTO\_CRA}; lezione del 02/11/2026
(Capitolo~\ref{lez:2026-11-02}). Scaletta prevista:
\begin{itemize}
  \item finalità e natura (regolamento di prodotto, orizzontale);
  \item ambito: prodotti con elementi digitali, elaborazione dati da remoto, software;
  \item esclusioni e regimi verticali;
  \item ruoli: fabbricante, importatore, distributore;
  \item requisiti essenziali (Allegato I, Parti I e II);
  \item prodotti importanti (Classi I e II) e critici; valutazione della conformità;
  \item obblighi di segnalazione (art.~14): vulnerabilità attivamente sfruttate e incidenti gravi;
  \item rapporto con AI Act, CSA, Data Act;
  \item date di applicazione e regime transitorio (unità immesse prima e dopo l'11/12/2027).
\end{itemize}}

\chapter{NIS2}
\label{cap:nis2}

\begin{nota}[Stato e fonti]
\textbf{Stato: non ancora trattata a lezione} (lezione frontale prevista venerdì 13/11/2026, online:
Capitolo~\ref{lez:2026-11-09}). Questo capitolo è un \textbf{inquadramento provvisorio}. Fonti usate:
slide introduttive (\emph{01.~Introductory class}), glossario del laboratorio~1, diario delle lezioni
del 21/09, 05/10, 26/10, 02/11 e 09/11, tema d'esame precedente (Capitolo~\ref{cap:tema-precedente});
come fonte secondaria gli \textbf{appunti di studenti dell'A.A.~2025/26}, sempre indicati come tali.
Tutto il resto è nei riquadri viola di integrazione.
\end{nota}

% =====================================================================
\section{Che cosa protegge e perché conta}
\label{sec:nis2-inquadramento}

La NIS2 (Direttiva (UE) 2022/2555) protegge la \textbf{continuità dei servizi essenziali}: nello
scenario della lezione introduttiva (ransomware sulla piattaforma di prenotazione di una rete
ospedaliera) si applica perché si è fermato un servizio su cui la società fa affidamento. Il bene
protetto non è il dato personale (GDPR) né il prodotto (CRA), ma il \textbf{servizio} erogato
attraverso reti e sistemi informativi. Nella ``linea del tempo'' delle slide la generazione 2022
(NIS2 e DORA) è quella che porta in primo piano resilienza, supply chain e responsabilità degli organi
di amministrazione.

Punti emersi a lezione:
\begin{itemize}
  \item per il docente la NIS2 è ``il vero futuro'': un livello di cybersicurezza molto più spinto
    dell'art.~32 GDPR, che ha dominato la generazione di norme 2016--2026 (lezione del 21/09);
  \item è una \textbf{direttiva}: va recepita dagli Stati e lascia ``27 perimetri nazionali'' diversi
    per ambito, sanzioni e vigilanza. In Italia è recepita con il \textbf{D.lgs.~138/2024}; autorità
    competente è l'\textbf{ACN} (Agenzia per la Cybersicurezza Nazionale), che ospita il
    \textbf{CSIRT Italia};
  \item riguarda le reti e le \textbf{organizzazioni}, mentre il CRA riguarda il \textbf{prodotto}; in
    Italia i soggetti obbligati sarebbero, secondo una stima a voce del docente, circa 20.000
    (lezione del 26/10);
  \item si applica solo ad \textbf{alcune categorie} di soggetti, legate a infrastrutture e servizi
    critici, con requisiti di sicurezza specifici e numerosi (lezione del 02/11). A differenza della
    DPIA del GDPR, che spesso è una zona grigia, nella NIS2 ``si entra o non si entra in una
    categoria'' (lezione del 05/10);
  \item con la NIS2 molte aziende hanno bisogno della certificazione ISO/IEC~27001 per stare sul
    mercato, perché i grandi clienti scelgono i fornitori anche in base a essa (lezione del 21/09).
\end{itemize}

\paragraph{Dalla NIS alla NIS2 (appunti A.A.~2025/26).} La prima direttiva NIS (2016) si applicava a un
gruppo ristretto di operatori e fu recepita in modo disomogeneo; la NIS2 la sostituisce allargando
l'ambito a 18 settori, distinguendo soggetti essenziali e importanti, rafforzando vigilanza e obblighi
di segnalazione e inserendo la gestione del rischio della catena di fornitura.

\begin{integrazione}[dati formali]
\begin{itemize}
  \item Direttiva (UE) 2022/2555 del 14 dicembre 2022, in vigore dal 16 gennaio 2023; termine di
    recepimento il 17 ottobre 2024. La direttiva NIS (UE) 2016/1148 è abrogata dal 18 ottobre 2024.
  \item Il limite della NIS1 era anche di metodo: erano gli Stati a individuare uno per uno gli
    ``operatori di servizi essenziali'', con criteri molto diversi tra Paesi. La NIS2 sostituisce
    l'identificazione caso per caso con un criterio generale (settore + dimensione).
  \item Il D.lgs.~4 settembre 2024, n.~138 è in vigore dal 16 ottobre 2024.
\end{itemize}
\end{integrazione}

% =====================================================================
\section{Chi è soggetto: soggetti essenziali e importanti}
\label{sec:nis2-ambito}

La NIS2 distingue \textbf{soggetti essenziali} e \textbf{soggetti importanti}. Dal glossario del corso:
le due classi hanno una \textbf{vigilanza diversa} (\emph{ex ante} per gli essenziali, \emph{ex post}
per gli importanti) e \textbf{sanzioni diverse} (Sezione~\ref{sec:nis2-vigilanza}).

Nel tema d'esame precedente (Capitolo~\ref{cap:tema-precedente}) la qualificazione è il primo passo
della risposta: la sanità è un settore dell'Allegato~I, quindi un ospedale è \textbf{essenziale} se
grande impresa e \textbf{importante} se media; i suoi fornitori di cloud e di servizi di sicurezza
gestiti sono a loro volta soggetti NIS2 (infrastrutture digitali e gestione dei servizi ICT B2B).

\paragraph{Appunti A.A.~2025/26.} Agli ambiti già coperti dalla NIS1 (energia, trasporti, sanità,
finanza, acqua, infrastrutture digitali) si aggiungono, tra gli altri, alimenti, servizi postali,
rifiuti, comunicazioni elettroniche e piattaforme digitali, gestione dei servizi ICT e alcune
manifatture strategiche. Gli Stati possono designare come essenziali ulteriori soggetti per ragioni
strategiche nazionali.

\begin{integrazione}[criterio dimensionale (art.~2) e classificazione (art.~3)]
\textbf{Regola generale (\emph{size-cap}).} La NIS2 si applica ai soggetti dei settori degli Allegati~I
e~II che sono almeno \textbf{medie imprese} secondo la Raccomandazione 2003/361/CE, cioè che superano i
limiti della piccola impresa: almeno \textbf{50 dipendenti} oppure fatturato annuo \emph{e} totale di
bilancio superiori a \textbf{10 milioni di euro}.

\textbf{Indipendentemente dalla dimensione} rientrano, tra gli altri: fornitori di reti o servizi di
comunicazione elettronica accessibili al pubblico, prestatori di servizi fiduciari, gestori di registri
dei nomi di dominio di primo livello (TLD) e fornitori di servizi DNS; il fornitore unico in uno Stato
di un servizio essenziale; i soggetti la cui perturbazione avrebbe impatto significativo su sicurezza o
salute pubblica; le \textbf{entità critiche} individuate ai sensi della Direttiva CER (UE) 2022/2557;
gli enti della pubblica amministrazione centrale.

\textbf{Essenziali (art.~3(1))}: in sintesi, i soggetti dell'Allegato~I che superano le soglie della
media impresa (cioè \textbf{grandi imprese}: almeno 250 dipendenti oppure fatturato superiore a
50 milioni e bilancio superiore a 43 milioni di euro); prestatori di servizi fiduciari qualificati,
registri TLD e fornitori DNS a prescindere dalla dimensione; fornitori di comunicazioni elettroniche
pubbliche almeno medi; PA centrale; entità critiche CER; altri soggetti designati dallo Stato.

\textbf{Importanti (art.~3(2))}: tutti gli altri soggetti degli Allegati~I e~II nell'ambito, quindi le
medie imprese dell'Allegato~I e le medie e grandi imprese dell'Allegato~II.

\medskip
\textbf{Allegato~I -- settori ad alta criticità (11)}: energia (elettricità, teleriscaldamento,
petrolio, gas, idrogeno); trasporti (aereo, ferroviario, per vie d'acqua, su strada); settore bancario;
infrastrutture dei mercati finanziari; sanità; acqua potabile; acque reflue; infrastrutture digitali
(IXP, DNS, TLD, cloud, data center, CDN, servizi fiduciari, comunicazioni elettroniche); gestione dei
servizi ICT B2B (fornitori di servizi gestiti e di servizi di sicurezza gestiti); pubblica
amministrazione; spazio.

\textbf{Allegato~II -- altri settori critici (7)}: servizi postali e di corriere; gestione dei rifiuti;
fabbricazione, produzione e distribuzione di sostanze chimiche; produzione, trasformazione e
distribuzione di alimenti; fabbricazione (dispositivi medici, computer ed elettronica, apparecchiature
elettriche, macchinari, autoveicoli, altri mezzi di trasporto); fornitori di servizi digitali
(mercati online, motori di ricerca, social network); ricerca.

\medskip
\textbf{Italia.} Il D.lgs.~138/2024 aggiunge propri allegati (pubbliche amministrazioni e ulteriori
categorie, ad esempio trasporto pubblico locale, istituti di istruzione che svolgono ricerca, soggetti
culturali) e prevede un meccanismo di \textbf{auto-registrazione} dei soggetti sulla piattaforma
dell'ACN: è il soggetto a valutare se rientra nell'ambito; l'ACN redige poi l'elenco dei soggetti
essenziali e importanti e comunica loro l'inclusione.
\end{integrazione}

\begin{verifica}
\begin{itemize}
  \item Gli appunti A.A.~2025/26 parlano in un punto di ``18 settori'' e in un altro di ``25 settori'':
    il dato corretto per la direttiva è 18 (11 + 7); il numero maggiore potrebbe contare anche i
    sottosettori o gli allegati italiani. Verificare sulle slide NIS2.
  \item Allegati italiani del D.lgs.~138/2024, finestre di registrazione sulla piattaforma ACN e
    calendario degli obblighi: dettagli da confermare sulle slide e sul sito dell'ACN.
\end{itemize}
\end{verifica}

% =====================================================================
\section{Governance: l'organo di gestione (art.~20)}
\label{sec:nis2-governance}

Il glossario del corso indica tra i tratti distintivi della NIS2 la \textbf{responsabilità degli
organi di gestione} (\emph{management liability}); nelle schede di ripasso è riassunta come
responsabilità \textbf{non delegabile}. Nella risposta modello del tema precedente
(Capitolo~\ref{cap:tema-precedente}) la gestione della supply chain va supervisionata e approvata
dall'organo di gestione (art.~20).

\paragraph{Appunti A.A.~2025/26.} Amministratori e rappresentanti legali restano responsabili della
conformità anche se affidano la cybersicurezza a specialisti interni, a consulenti esterni o al punto
di contatto NIS. La gestione del rischio è però condivisa nei fatti (CISO, responsabile dell'ISMS,
amministratori di sistema, red e blue team, punto di contatto NIS): serve molta comunicazione interna.
Risultanze, piani di trattamento e decisioni di accettazione del rischio vanno riportati all'organo di
gestione: la cybersicurezza diventa materia da consiglio di amministrazione. Il riesame della
direzione previsto dalla ISO/IEC~27001 può essere riutilizzato per dimostrare questa supervisione. Nel
laboratorio NIS2 dello scorso anno, dopo l'incidente il consiglio doveva prendere atto formalmente
dell'evento, approvare il budget di remediation e documentare le decisioni; tra le misure raccomandate,
un CISO a riporto diretto del consiglio e formazione obbligatoria per gli amministratori.

\begin{integrazione}[art.~20 e responsabilità personale]
\begin{itemize}
  \item \textbf{Art.~20(1)}: gli organi di gestione dei soggetti essenziali e importanti
    \textbf{approvano} le misure di gestione del rischio dell'art.~21, ne \textbf{sovrintendono
    l'attuazione} e possono essere ritenuti \textbf{responsabili} delle violazioni.
  \item \textbf{Art.~20(2)}: i membri degli organi di gestione devono seguire una
    \textbf{formazione} specifica; gli Stati incoraggiano formazione periodica anche per i dipendenti.
  \item Per i soggetti essenziali, in caso di inosservanza persistente, l'autorità può chiedere di
    vietare temporaneamente a chi esercita funzioni dirigenziali (fino al livello di amministratore
    delegato o rappresentante legale) di svolgerle (art.~32(5)); le persone fisiche responsabili
    possono essere chiamate a rispondere secondo il diritto nazionale (artt.~32(6) e~33).
\end{itemize}
\end{integrazione}

% =====================================================================
\section{Misure di gestione del rischio (art.~21)}
\label{sec:nis2-misure}

Nella ``grammatica comune'' delle slide introduttive la NIS2 compare con l'\textbf{art.~21}: misure di
gestione del rischio \textbf{multi-rischio} (\emph{all-hazards}), \textbf{proporzionate} a esposizione,
dimensione e impatto. È lo stesso metodo delle altre norme del corso (contesto, identificazione,
analisi, valutazione, trattamento, monitoraggio): la norma dice di essere sicuri, non quali controlli
adottare, e ``abbiamo seguito la checklist'' non è mai una difesa.

\paragraph{Supply chain.} L'\textbf{art.~21(2)(d)} rende la sicurezza della catena di
approvvigionamento una misura obbligatoria, ``compresa la qualità dei fornitori'' (slide sul
perimetro). Nella risposta modello del tema precedente: l'\textbf{art.~21(3)} chiede di considerare le
vulnerabilità specifiche di ciascun fornitore, la qualità dei loro prodotti e delle loro pratiche di
sicurezza, comprese le procedure di sviluppo sicuro; in pratica, inventario dei fornitori e dei
servizi critici, mappatura delle dipendenze, requisiti minimi nei contratti, SLA e tempi di
ripristino, rivalutazione periodica.

\paragraph{Certificazioni.} Le slide ricordano che l'\textbf{art.~24} NIS2, come gli artt.~42--43 GDPR e
il Titolo~III del Cybersecurity Act, tratta la certificazione come \textbf{elemento di prova} da
valutare, mai come immunità (Capitolo~\ref{cap:csa}).

\paragraph{Appunti A.A.~2025/26: le linee guida ACN.} L'ACN ha pubblicato linee guida e una matrice di
controlli che, pur non vincolanti, sono il riferimento nazionale più concreto. Punti riportati:
\begin{itemize}
  \item la conformità non è una checklist: il controllo cardine (indicato come \emph{GVRM-03}) chiede
    un ciclo continuo di identificazione, analisi, valutazione, trattamento e monitoraggio del rischio;
    l'ACN chiede una valutazione basata su probabilità e impatto ma non impone una metodologia;
  \item inventario aggiornato degli asset (hardware, software incluse librerie e componenti open
    source, servizi gestiti da terzi);
  \item valutazione dei fornitori e delle società del gruppo (privilegi di accesso, dipendenza,
    resilienza), da cui discendono requisiti contrattuali, SLA e analisi dei rischi a cascata;
  \item la valutazione deve essere documentata e riproducibile (minaccia, asset esposto, impatto,
    probabilità), non un'impressione;
  \item la matrice ACN suggerirebbe di rivedere ruoli e risk assessment ogni due anni, intervallo
    giudicato insufficiente a lezione (almeno annuale);
  \item tensioni da gestire con l'accountability: conservazione dei log (lunga per la sicurezza,
    limitata dalla protezione dei dati) e politiche crittografiche;
  \item lo Statement of Applicability della ISO/IEC~27001 può essere esteso per mappare i requisiti
    ACN/NIS2 sui controlli esistenti.
\end{itemize}

\begin{nota}[Segnalato l'anno scorso: logica NIS2 e logica DORA]
Negli appunti A.A.~2025/26 il confronto è indicato come sottolineato più volte dal docente e
``probabile domanda d'esame'': la \textbf{NIS2} presidia soprattutto la \textbf{postura preventiva}
(misure adeguate prima dell'incidente), il \textbf{DORA} la \textbf{resilienza e la continuità}
(rendere la perturbazione recuperabile entro parametri accettabili). Da confermare nelle lezioni di
quest'anno.
\end{nota}

\begin{integrazione}[le dieci misure minime dell'art.~21(2)]
Le misure devono comprendere almeno:
\begin{enumerate}[label=(\alph*)]
  \item politiche di analisi dei rischi e di sicurezza dei sistemi informativi;
  \item gestione degli incidenti;
  \item continuità operativa (gestione dei backup, disaster recovery) e gestione delle crisi;
  \item sicurezza della catena di approvvigionamento, compresi i rapporti con i fornitori diretti;
  \item sicurezza nell'acquisizione, nello sviluppo e nella manutenzione di reti e sistemi, compresa
    la gestione e divulgazione delle vulnerabilità;
  \item politiche e procedure per valutare l'efficacia delle misure;
  \item pratiche di igiene informatica di base e formazione;
  \item politiche e procedure sull'uso della crittografia e, se opportuno, della cifratura;
  \item sicurezza delle risorse umane, controllo degli accessi e gestione degli asset;
  \item autenticazione a più fattori o continua, comunicazioni vocali, video e testuali protette e
    sistemi di comunicazione di emergenza protetti, ove opportuno.
\end{enumerate}
Altri elementi utili: se un soggetto scopre di non rispettare le misure deve adottare senza indebito
ritardo le correzioni necessarie (art.~21(4)); per i fornitori digitali (cloud, data center, DNS,
servizi gestiti, piattaforme, ecc.) i requisiti tecnici sono dettagliati dal Regolamento di esecuzione
(UE) 2024/2690; l'art.~22 prevede valutazioni coordinate a livello UE dei rischi di catene di
approvvigionamento critiche (per esempio il 5G).
\end{integrazione}

% =====================================================================
\section{Incidenti e obblighi di notifica (art.~23)}
\label{sec:nis2-notifica}

\paragraph{Incidente e incidente significativo.} L'\textbf{incidente} (art.~6) è un evento che
compromette la \textbf{disponibilità, l'autenticità, l'integrità o la riservatezza} di dati conservati,
trasmessi o elaborati, o dei servizi offerti da reti e sistemi informativi o accessibili attraverso di
essi. L'obbligo di notifica scatta solo per l'\textbf{incidente significativo} (art.~23(3)): quello che
ha causato o \emph{può causare} una grave perturbazione operativa dei servizi o perdite finanziarie per
il soggetto, oppure si è ripercosso o può ripercuotersi su altre persone fisiche o giuridiche causando
danni materiali o immateriali considerevoli. Il ``può causare'' è giuridicamente decisivo.

\begin{esame}
\begin{itemize}
  \item La \textbf{definizione di incidente NIS2} è uscita come domanda chiusa in un appello
    precedente (Capitolo~\ref{cap:tema-precedente}). Attenzione: la NIS2 include
    l'\emph{autenticità}, la definizione GDPR di data breach no.
  \item Il docente ha ricordato che l'\textbf{early warning} va fatto entro 24 ore \textbf{anche di
    sabato o domenica} (lezione del 21/09).
  \item Per l'esame vanno conosciuti l'ambito della norma e le tempistiche di notifica, che sono ``il
    cuore dell'esame'' (lezione del 09/11).
\end{itemize}
\end{esame}

\paragraph{Le tempistiche.} L'orologio parte dalla \textbf{consapevolezza} di un incidente
significativo. Destinatario è il \textbf{CSIRT} (in Italia il CSIRT Italia, presso l'ACN). Le slide
precisano che \textbf{in Italia, da gennaio 2026}, gli incidenti significativi vanno segnalati al CSIRT
Italia.

\begin{center}
\small
\begin{tabularx}{\linewidth}{>{\raggedright\arraybackslash}p{3.2cm} >{\raggedright\arraybackslash}p{3.6cm} X}
\toprule
\textbf{Fase} & \textbf{Termine} & \textbf{Contenuto} \\ \midrule
Preallarme (\emph{early warning}) & \textbf{24 h} dalla consapevolezza
  & Segnala l'incidente; indica se si sospetta un atto illegittimo o malevolo e se può avere impatto
    transfrontaliero \\
Notifica dell'incidente & \textbf{72 h} dalla consapevolezza
  & Aggiorna il preallarme: valutazione iniziale di gravità e impatto, indicatori di compromissione
    (IoC) \\
Relazione intermedia & su richiesta del CSIRT & Aggiornamenti pertinenti \\
Relazione finale & \textbf{1 mese} dalla notifica dell'incidente
  & Descrizione dettagliata, causa di fondo, misure di mitigazione, impatto transfrontaliero \\
\bottomrule
\end{tabularx}
\end{center}

Nelle slide la tabella dei cinque ``orologi'' riassume la logica: la NIS2 conta la
\textbf{perturbazione di un servizio}, il GDPR il danno alle persone, il DORA l'impatto su una funzione.
Il confronto completo con le altre norme è nel Capitolo~\ref{cap:scadenze}: in particolare le 72 ore
NIS2 decorrono dalla consapevolezza, quelle DORA dalla notifica iniziale.

\begin{esame}
Per l'esame va studiato lo \textbf{schema delle notifiche oggi vigente}: con il pacchetto
\emph{Digital Omnibus} si andrà verso un \textbf{punto unico di segnalazione presso ENISA}, che per il
docente sarà ``la rivoluzione più importante'', ma non è ancora operativo (lezione del 21/09).
\end{esame}

\paragraph{Appunti A.A.~2025/26.} Notifica tramite il punto di contatto NIS del soggetto. Il CSIRT
risponde al preallarme entro 24 ore con un primo riscontro e, se possibile, indicazioni operative: la
logica è di \textbf{supporto} più che sanzionatoria, a differenza del GDPR. Se l'incidente è ancora in
corso dopo un mese si invia una relazione sui progressi e la finale alla chiusura. In caso di impatto
transfrontaliero il CSIRT informa gli altri CSIRT ed ENISA; per le crisi su larga scala interviene
EU-CyCLONe. Manca ancora una metodologia armonizzata di valutazione della gravità comparabile a quella
ENISA per i data breach. In pratica servono percorsi di escalation interni, personale formato, buone
capacità forensi e di logging (per rispettare le 72 ore e il mese) e dialogo con i fornitori; la
tendenza è verso una \textbf{procedura unica di gestione degli incidenti} per GDPR, NIS2, DORA e norme
di settore. Nel laboratorio NIS2 dello scorso anno (fornitore di energia, soggetto essenziale) si
aggiungevano: comunicazione al pubblico per evitare allarme, autorità di settore (ARERA), denuncia alla
Polizia Postale (non obbligatoria, ma la sua assenza può essere valutata negativamente) e
allineamento tra registro degli incidenti e registro dei data breach.

\begin{integrazione}[dettagli dell'art.~23]
\begin{itemize}
  \item \textbf{Destinatari dei servizi}: il soggetto informa senza indebito ritardo i destinatari dei
    propri servizi dell'incidente significativo che può pregiudicare la fornitura (art.~23(1)) e, in
    caso di minaccia informatica significativa, delle misure che possono adottare (art.~23(2)).
  \item \textbf{Prestatori di servizi fiduciari}: la notifica dell'incidente va fatta entro
    \textbf{24 ore} invece di 72 (art.~23(4)(b)).
  \item Il CSIRT o l'autorità rispondono al preallarme \textbf{entro 24 ore} (art.~23(5)); se
    l'incidente riguarda più Stati, informano gli altri Stati ed ENISA (art.~23(6)); possono informare
    il pubblico se necessario (art.~23(7)).
  \item Per i fornitori digitali i casi in cui un incidente è ``significativo'' sono precisati dal
    Regolamento di esecuzione (UE) 2024/2690 (art.~23(11)).
  \item \textbf{Notifiche volontarie} (art.~30): incidenti non significativi, minacce informatiche e
    quasi-incidenti, anche da parte di soggetti fuori ambito.
  \item In Italia anche la \textbf{L.~90/2024} prevede obblighi di notifica all'ACN (24 ore e
    72 ore) per pubbliche amministrazioni e altri soggetti indicati dalla legge.
\end{itemize}
\end{integrazione}

\begin{verifica}
\begin{itemize}
  \item ``Da gennaio 2026'' (slide): coerente con il calendario del D.lgs.~138/2024, che fa decorrere
    l'obbligo di notifica dal 1 gennaio 2026 per i soggetti inclusi nell'elenco ACN. Confermare sulle
    slide NIS2.
  \item Le soglie di ``incidente significativo'' fissate dall'ACN per i soggetti italiani (determinazioni
    dell'Agenzia) non sono state ancora viste a lezione.
\end{itemize}
\end{verifica}

% =====================================================================
\section{Istituzioni: CSIRT, ACN, ENISA, NIS Cooperation Group, EU-CyCLONe}
\label{sec:nis2-istituzioni}

Dalla slide sulle istituzioni (``chi chiamare alle 8:14''):
\begin{itemize}
  \item \textbf{risposta operativa}: il CSIRT Italia, dentro l'ACN (``il numero da chiamare''), la rete
    europea dei CSIRT, i CERT di settore;
  \item \textbf{vigilanza NIS2}: l'\textbf{ACN} come autorità nazionale competente, affiancata da
    \textbf{autorità di settore} per soggetti specifici;
  \item \textbf{livello UE}: \textbf{ENISA} (certificazione, \emph{threat landscape}, piattaforma unica
    di segnalazione del CRA), il \textbf{NIS Cooperation Group} (strategia e linee guida) ed
    \textbf{EU-CyCLONe}, la rete per la gestione delle crisi informatiche su larga scala.
\end{itemize}
L'ACN è anche autorità nazionale di certificazione della cybersicurezza (lezione del 09/11). Secondo il
docente l'ACN può arrivare a \textbf{fermare un fornitore} ritenuto pericoloso per un'infrastruttura
critica, anche se il fornitore non è direttamente soggetto a NIS2 (lezione del 21/09). Due conseguenze
indicate dalle slide: non si sceglie il proprio interlocutore, e i difetti di coordinamento tra queste
autorità sono un rischio che gestisce il soggetto, non l'autorità.

\paragraph{ENISA.} Per il ruolo di ENISA (mandato permanente dal Cybersecurity Act, cooperazione
operativa, verso la centralizzazione delle notifiche) si veda il Capitolo~\ref{cap:csa} e la domanda
del tema precedente (Capitolo~\ref{cap:tema-precedente}): ENISA \textbf{non} irroga sanzioni e
\textbf{non} è un CSIRT nazionale.

\paragraph{Appunti A.A.~2025/26.} Ogni Stato adotta una \textbf{strategia nazionale di cybersicurezza}
(obiettivi, governance, individuazione degli asset critici, preparazione e ripristino, politica sulla
supply chain, formazione) e designa una o più autorità competenti e almeno un CSIRT. In Italia queste
funzioni sono accentrate nell'ACN.

\begin{integrazione}[articoli di riferimento]
\begin{itemize}
  \item Art.~7 strategia nazionale; art.~8 autorità competenti e \textbf{punto di contatto unico}
    (in Italia l'ACN); art.~9 autorità di gestione delle crisi informatiche e piano nazionale.
  \item Artt.~10--11 CSIRT; art.~12 divulgazione coordinata delle vulnerabilità e \textbf{banca dati
    europea delle vulnerabilità} gestita da ENISA.
  \item Art.~14 \textbf{Gruppo di cooperazione NIS} (Stati, Commissione, ENISA): livello strategico;
    art.~15 \textbf{rete dei CSIRT}: livello tecnico-operativo; art.~16 \textbf{EU-CyCLONe}: gestione
    coordinata, a livello operativo, di incidenti e crisi su vasta scala; art.~19 revisioni tra pari.
  \item Art.~27: ENISA tiene un registro dei fornitori digitali (DNS, TLD, cloud, data center, CDN,
    servizi gestiti, piattaforme).
\end{itemize}
\end{integrazione}

\begin{verifica}
Il potere dell'ACN di ``fermare un fornitore'' non soggetto a NIS2 non è stato ricondotto a lezione a
un articolo preciso: potrebbe riferirsi alle valutazioni coordinate dei rischi della supply chain
(art.~22), ai poteri nazionali sul perimetro di sicurezza nazionale cibernetica o ad altri strumenti
nazionali. Da chiarire nella lezione del 13/11.
\end{verifica}

% =====================================================================
\section{Vigilanza ed esecuzione: \emph{ex ante} ed \emph{ex post}}
\label{sec:nis2-vigilanza}

\begin{center}
\small
\begin{tabularx}{\linewidth}{>{\raggedright\arraybackslash}p{3cm} X X}
\toprule
 & \textbf{Soggetti essenziali} & \textbf{Soggetti importanti} \\ \midrule
Vigilanza & \emph{ex ante}: controlli anche in assenza di indizi di violazione
  & \emph{ex post}: si interviene a seguito di elementi che indicano una violazione \\
Sanzione massima & fino a \textbf{10 milioni di euro} o al \textbf{2\%} del fatturato mondiale annuo
  & fino a \textbf{7 milioni di euro} o all'\textbf{1,4\%} del fatturato mondiale annuo \\
\bottomrule
\end{tabularx}
\end{center}
\noindent Fonti: glossario del laboratorio~1 (classi, vigilanza, sanzioni) e slide introduttive
(sanzione per gli essenziali). Gli appunti A.A.~2025/26 aggiungono che le sanzioni tengono conto di
proporzionalità, gravità e cooperazione del soggetto, e che nel laboratorio dello scorso anno il
rischio sanzionatorio NIS2 era giudicato basso grazie alla sostanziale conformità, nonostante lacune
di formazione: la sanzione dipende dall'adeguatezza delle misure e dalla tempestività delle notifiche,
non dal solo fatto che l'incidente sia avvenuto.

\begin{integrazione}[artt.~32--34]
\begin{itemize}
  \item \textbf{Essenziali (art.~32)}: ispezioni in loco e a distanza, audit periodici e mirati anche
    da organismi indipendenti, audit \emph{ad hoc} dopo un incidente significativo, scansioni di
    sicurezza, richieste di informazioni e di prove. Misure di esecuzione: avvertimenti, istruzioni
    vincolanti, ordine di informare i destinatari dei servizi, nomina di un responsabile del
    monitoraggio e, come \emph{extrema ratio}, sospensione di certificazioni o autorizzazioni e
    divieto temporaneo di funzioni dirigenziali.
  \item \textbf{Importanti (art.~33)}: stessi strumenti, ma attivati solo a fronte di prove o indizi
    di non conformità; niente sospensioni né divieti per i dirigenti.
  \item \textbf{Sanzioni (art.~34)}: il massimale è ``almeno'' 10 milioni o 2\% (essenziali) e
    ``almeno'' 7 milioni o 1,4\% (importanti), \textbf{se superiore}: gli Stati possono prevedere
    importi più alti. Sono possibili anche penalità di mora.
\end{itemize}
\end{integrazione}

% =====================================================================
\section{Rapporti con DORA, GDPR e CRA}
\label{sec:nis2-rapporti}

\paragraph{DORA: \emph{lex specialis}.} Il settore bancario e le infrastrutture dei mercati finanziari
sono nell'Allegato~I della NIS2, ma per le entità finanziarie il \textbf{DORA} prevale dove le norme
si sovrappongono (gestione del rischio ICT, notifica degli incidenti): la norma speciale prevale su
quella generale (glossario del corso; Simulazione~2 nel Capitolo~\ref{cap:simulazioni}). Per il resto
si veda la Sezione~\ref{sec:dora-lexspecialis}. Viceversa, un ospedale non è soggetto al DORA perché
non è un'entità finanziaria (tema precedente).

\paragraph{GDPR.} Lo stesso evento può essere incidente significativo NIS2 e data breach GDPR: le
notifiche al CSIRT e al Garante sono \textbf{separate e parallele} e una non adempie l'altra
(Capitolo~\ref{cap:scadenze}). Le due definizioni non coincidono (autenticità; dati personali sì o no) e
gli orologi partono da eventi diversi. Nel tema precedente la risposta modello suggerisce una procedura
integrata GDPR--NIS2 e registri allineati. Gli appunti A.A.~2025/26 aggiungono che le informazioni
raccolte per la notifica NIS2 possono essere in gran parte riutilizzate per quella GDPR.

\paragraph{CRA.} NIS2 = organizzazione, CRA = prodotto (lezione del 26/10): il considerando~3 del CRA
spiega che né il Cybersecurity Act né la NIS2 impongono requisiti obbligatori di sicurezza ai prodotti.
Un'applicazione può ricadere insieme sotto CRA, NIS2 (che include il settore digitale) e GDPR (lezione
del 02/11). Nel tema precedente, se l'incidente nasce da una vulnerabilità attivamente sfruttata di un
prodotto, l'obbligo CRA è del fabbricante, non dell'ospedale.

\begin{integrazione}[coordinamento con le altre norme]
\begin{itemize}
  \item \textbf{Art.~4 NIS2}: se un atto settoriale dell'Unione impone misure di gestione del rischio o
    obblighi di notifica almeno equivalenti, le corrispondenti disposizioni NIS2 non si applicano; il
    considerando~28 indica espressamente il DORA come tale atto per le entità finanziarie.
  \item \textbf{Art.~35 NIS2}: se una violazione NIS2 può comportare un data breach, l'autorità NIS
    informa l'autorità di controllo GDPR; se quest'ultima infligge una sanzione pecuniaria ai sensi
    dell'art.~58(2)(i) GDPR, l'autorità NIS \textbf{non} può infliggere una sanzione pecuniaria per la
    stessa condotta (possono restare le altre misure).
  \item Un fornitore cloud o un MSSP può essere \textbf{contemporaneamente} soggetto NIS2 (infrastrutture
    digitali o servizi ICT gestiti), responsabile del trattamento GDPR del cliente e fornitore terzo
    di servizi ICT ai sensi del DORA se il cliente è una banca: tre cornici, tre insiemi di obblighi.
\end{itemize}
\end{integrazione}

\dacompletare{Da integrare dopo la lezione del 13/11/2026 con slide e diario: ambito e allegati
(versione del corso), recepimento italiano e calendario ACN, art.~20, art.~21 e linee guida ACN,
art.~23 e soglie di significatività, sanzioni, rapporti con DORA e GDPR. Aggiornare poi le schede di
ripasso (date NIS2) e verificare la coerenza con il Capitolo~\ref{cap:scadenze}.}

\chapter{DORA}
\label{cap:dora}

\dacompletare{Fonte: slide ancora da ricevere. Scaletta prevista:
\begin{itemize}
  \item ambito: entità finanziarie e fornitori terzi di servizi ICT;
  \item quadro di gestione del rischio ICT e strategia di resilienza operativa digitale;
  \item business continuity, disaster recovery, backup;
  \item ruolo dell'organo di gestione; principio di proporzionalità;
  \item incidenti ICT gravi: classificazione e tempistiche di notifica;
  \item test di resilienza;
  \item rischio da terze parti ICT e clausole contrattuali.
\end{itemize}}

\chapter{AI Act}
\label{cap:aiact}

\dacompletare{Fonte: slide ancora da ricevere. Scaletta prevista:
\begin{itemize}
  \item definizione di sistema di IA; ruoli (provider, deployer, importatore, distributore);
  \item classificazione per rischio: pratiche vietate, alto rischio (Allegati I e III), rischio
    limitato, rischio minimo;
  \item obblighi per i sistemi ad alto rischio; FRIA;
  \item art.~25: riqualificazione come provider;
  \item art.~15: accuratezza, robustezza, cybersicurezza; presunzione tramite CRA;
  \item sanzioni e calendario di applicazione (Digital Omnibus).
\end{itemize}}

\chapter{Cybersecurity Act (CSA)}
\label{cap:csa}

\begin{nota}[Fonti]
Lezione del 09/11/2026 (Capitolo~\ref{lez:2026-11-09}). Le slide sul CSA non sono ancora disponibili:
questa è una prima stesura dal parlato; articoli e date da completare sono segnalati.
\end{nota}

\section{Natura e finalità}
Il Cybersecurity Act è un \textbf{regolamento} dell'Unione europea che introduce la \textbf{soft law
europea} della cybersicurezza. Ha due pilastri:
\begin{enumerate}
  \item il \textbf{rafforzamento di ENISA}, con un mandato e compiti propri;
  \item un \textbf{quadro europeo di certificazione} della cybersicurezza per prodotti, servizi e processi ICT.
\end{enumerate}
Nasce per superare quattro problemi:
\begin{itemize}
  \item l'\textbf{interdipendenza}: un fornitore può colpire molti clienti, in più paesi;
  \item l'\textbf{asimmetria informativa} tra venditore e acquirente;
  \item la \textbf{frammentazione} degli schemi nazionali;
  \item la mancanza di \textbf{fiducia verificabile}: servono requisiti ed evidenze.
\end{itemize}
Obiettivo: la \textbf{cyber assurance}, cioè evidenze che certi requisiti sono soddisfatti (da non
confondere con la cyber insurance). Le certificazioni del CSA sono \textbf{volontarie}: oggi in Europa
non esistono certificazioni di cybersicurezza obbligatorie.

\dacompletare{Riferimento normativo (Reg.~UE 2019/881), struttura del regolamento, articoli principali e
mandato permanente di ENISA, dalle slide.}

\section{Hard law, soft law e posizione del CSA}
\begin{center}
\small
\begin{tabularx}{\linewidth}{>{\raggedright\arraybackslash}p{3cm} X X}
\toprule
 & \textbf{Hard law} & \textbf{Soft law}\\
\midrule
Esempi & limiti di velocità; GDPR, CRA, NIS2, DORA, AI Act & policy aziendali (es.\ MFA), linee guida ENISA,
  standard ISO, schemi del CSA\\
Conseguenze & sanzioni predeterminate dalla legge, discrezionalità delimitata & non previste dalla legge,
  ampio margine (es.\ provvedimenti disciplinari per violazione di policy)\\
Tempi & 1--3 anni (o più) per una legge & adattabile in fretta\\
\bottomrule
\end{tabularx}
\end{center}
Il CSA è soft law che aiuta a rispettare la hard law: in particolare, uno schema del CSA può fornire i
requisiti concreti con cui dimostrare la conformità al \textbf{CRA} (Capitolo~\ref{cap:cra}). Offre
\textbf{la possibilità, non la certezza}, di dimostrarla.

\section{ENISA e governance}
\begin{itemize}
  \item \textbf{ENISA}, Agenzia dell'UE per la cybersicurezza: linee guida, formazione ed esercitazioni;
    cooperazione operativa tra CSIRT; \textbf{preparazione degli schemi candidati}; analisi delle minacce e
    buone pratiche.
  \item \textbf{ECCG} (\emph{European Cybersecurity Certification Group}): rappresentanti delle autorità
    nazionali; coopera alla preparazione degli schemi ed esprime un parere sulle bozze.
  \item \textbf{Commissione europea}: individua il bisogno e adotta lo schema.
  \item \textbf{Autorità nazionali di certificazione}: vigilano sulla certificazione; in Italia l'\textbf{ACN}.
  \item \textbf{Organismi di valutazione e laboratori}: eseguono i test e rilasciano i certificati.
\end{itemize}

\section{Processo di adozione di uno schema}
\begin{center}
\begin{tikzpicture}[node distance=4mm, every node/.style={draw, rounded corners, align=center,
    font=\small, minimum height=1.1cm, text width=2.9cm}]
  \node (r) {\textbf{Richiesta}\\Commissione $\rightarrow$ ENISA};
  \node (d) [right=of r] {\textbf{Bozza}\\ENISA + stakeholder};
  \node (o) [right=of d] {\textbf{Parere}\\ECCG};
  \node (a) [right=of o] {\textbf{Adozione}\\Commissione};
  \draw[->] (r) -- (d); \draw[->] (d) -- (o); \draw[->] (o) -- (a);
\end{tikzpicture}
\end{center}
\dacompletare{Atto di adozione (atto di esecuzione), programma di lavoro dell'Unione (URWP), durata e
revisione degli schemi.}

\section{Framework, schema e certificato}
\begin{itemize}
  \item \textbf{Framework}: il CSA, che fissa ruoli e regole comuni. Non contiene schemi.
  \item \textbf{Schema}: il \emph{rule book} della valutazione per una categoria: ambito, requisiti di sicurezza,
    test, livelli di affidabilità.
  \item \textbf{Certificato}: attesta che uno specifico prodotto o servizio rispetta lo schema.
\end{itemize}
Oggetti possibili di uno schema: \textbf{prodotti ICT} (firewall, chip, software), \textbf{servizi ICT} (es.\ cloud
storage), \textbf{processi ICT} (progettazione e manutenzione; ancora nessuno schema), \textbf{servizi di
sicurezza gestiti} (incident response, audit, penetration test). ISO/IEC~27001 certifica invece
l'\textbf{organizzazione}.

\section{Livelli di affidabilità e valutazione della conformità}
\begin{center}
\small
\begin{tabular}{lll}
\toprule
\textbf{Livello} & \textbf{Rischio} & \textbf{Valutazione}\\
\midrule
Base & rischi noti di base & anche autovalutazione del fabbricante\\
Sostanziale & attaccanti con risorse limitate & terza parte\\
Elevato & attaccanti con risorse significative & terza parte\\
\bottomrule
\end{tabular}
\end{center}
In pratica la certificazione è quasi sempre di terza parte. Un certificato:
\begin{itemize}
  \item dà evidenze su un oggetto \textbf{identificato con precisione}, su requisiti e test specifici, a un
    livello definito;
  \item \textbf{non} dimostra l'assenza di vulnerabilità;
  \item non esonera in automatico da responsabilità: dopo un incidente conta se i requisiti erano davvero
    rispettati (\emph{accountability});
  \item richiede \textbf{monitoraggio e manutenzione}: prodotto e minacce cambiano; una patch non invalida
    automaticamente il certificato, ma il suo impatto va valutato; una nuova versione va ricertificata.
\end{itemize}
\dacompletare{Articoli su livelli di affidabilità, autovalutazione della conformità, durata e revoca dei
certificati, sanzioni nazionali.}

\section{EUCC}
Primo schema adottato: \textbf{EUCC} (\emph{EU Common Criteria-based scheme}), basato sui Common Criteria;
riguarda \textbf{prodotti e componenti ICT} (niente servizi o processi), è applicabile dal
\textbf{27 febbraio 2025} e prevede i livelli \textbf{sostanziale ed elevato}.
Concetti: \textbf{target of evaluation} (il prodotto esatto), \textbf{security target} (il documento che lo
descrive con i requisiti dichiarati), \textbf{protection profile} (requisiti riutilizzabili per un tipo di
prodotto), metodi e analisi delle vulnerabilità, ipotesi e limiti operativi.
\dacompletare{Riferimento al regolamento di esecuzione EUCC; altri schemi in preparazione (EUCS per il
cloud, EU5G).}

\section{Rapporto con le altre norme e strumenti}
\begin{center}
\small
\begin{tabularx}{\linewidth}{>{\raggedright\arraybackslash}p{3.4cm} X}
\toprule
\textbf{Norma / strumento} & \textbf{Domanda a cui risponde}\\
\midrule
CSA & come si valutano le dichiarazioni di sicurezza: quadro comune di certificazione\\
CRA & obblighi di sicurezza e sulle vulnerabilità dei prodotti con elementi digitali\\
NIS2 & sicurezza di reti e sistemi informativi dei soggetti essenziali e importanti\\
GDPR & protezione dei dati personali (la norma più generale)\\
DORA & resilienza operativa digitale del settore finanziario\\
ISO/IEC~27001 & sistema di gestione della sicurezza dell'organizzazione\\
Marcatura CE & conformità del prodotto alle norme UE che la richiedono (possibile sovrapposizione con l'EUCC)\\
Penetration test & attacco simulato autorizzato; soft law con possibili conseguenze legali\\
\bottomrule
\end{tabularx}
\end{center}

\section{Security by design e supply chain}
Protezione e manutenzione si pianificano in progettazione: ridurre la superficie d'attacco, conoscere i
componenti (\textbf{SBOM}), gestire le segnalazioni di vulnerabilità, distribuire gli aggiornamenti. Per ogni
fornitore: il componente può essere attaccato? Possiamo operare senza? Chi controlla servizi, decisioni e
dati (governance, non solo \emph{data residency})?

\section{CSA~2}
Proposta di revisione: più poteri a ENISA, certificazione semplificata, più attenzione alla supply chain.
\dacompletare{Stato dell'iter legislativo e contenuti della proposta.}

\begin{esame}
\begin{itemize}
  \item I due pilastri: \textbf{ENISA} e \textbf{schemi di certificazione}.
  \item CSA = soft law sui prodotti; ISO/IEC~27001 = organizzazione; CRA = hard law sui prodotti.
  \item Framework $\neq$ schema $\neq$ certificato; un certificato non prova l'assenza di vulnerabilità.
  \item Livelli base/sostanziale/elevato; autovalutazione solo per il livello base.
  \item EUCC: primo schema, prodotti ICT, Common Criteria, livelli sostanziale ed elevato.
\end{itemize}
\end{esame}


% =====================================================================
\part{Laboratori}
% Un capitolo per laboratorio: caso, consegna, metodo, soluzione ragionata.
% =====================================================================
% =====================================================================
\chapter{Lab 1 -- Arvenza Industrie: la mailbox delle paghe}
\label{lab:01}
% =====================================================================

\begin{nota}[Fonti e date]
Caso presentato il 05/10/2026 (Capitolo~\ref{lez:2026-10-05}), svolto venerdì 09/10/2026
(Capitolo~\ref{lez:2026-10-09}), discusso il 12/10/2026 (Capitolo~\ref{lez:2026-10-12}). Materiale:
\emph{Arvenza Breach Case Study} (testo del caso), \emph{Arvenza Case Solution} (soluzione modello,
Parte~A), metodologia ENISA v1.0 (2013, slide \emph{03}). Caso di fantasia.
\end{nota}

\section{Il caso in sintesi}

\paragraph{Contesto.} Arvenza Industrie (Italia) gestisce le paghe tramite una mailbox condivisa
usata per scambiare documenti con un \emph{payroll bureau} esterno. Gli studenti sono il
\textbf{DPO}, che giovedì 1 ottobre alle 09:00 deve consigliare CEO, HR Director e IT Manager.

\paragraph{Cronologia (2026).}
\begin{center}
\small
\begin{tabularx}{\linewidth}{p{2.3cm} X}
\toprule
\textbf{Quando} & \textbf{Evento} \\ \midrule
Mar 15/9 & Falsa email DocuSign: un proxy di phishing (\emph{adversary-in-the-middle}) ruba password e
  cookie di sessione, aggirando l'MFA. \\
15--20/9 & L'attaccante legge la mailbox; una regola nasconde tutta la posta del payroll bureau. \\
Gio 17/9 & Richieste contraffatte di cambio IBAN dello stipendio per tre dipendenti. \\
Ven 18/9 & Il bureau telefona ai tre e respinge le richieste; la sua risposta email è nascosta
  dalla regola. \\
18--19/9 & Il log di audit smette di registrare le letture per 24 ore; l'attaccante è attivo 14
  minuti. \\
Dom 20/9 & Alert di sicurezza: sessioni revocate, ticket chiuso. Nessuno guarda dentro la mailbox. \\
25--28/9 & Stipendi pagati regolarmente. Lun 28/9 ore 15:00: il bureau cita le richieste respinte;
  il DPO viene allertato. \\
Gio 1/10, 09:00 & Riunione: il DPO dà il suo parere. \\
\bottomrule
\end{tabularx}
\end{center}

\paragraph{I dati nella mailbox.} I log registrano \textbf{214 elementi letti}, trovati soprattutto con
ricerche su IBAN, bonifici, stipendi e buste paga; nessun export, sincronizzazione o inoltro.
\begin{center}
\small
\begin{tabularx}{\linewidth}{l l X l}
\toprule
\textbf{Gruppo} & \textbf{Interessati} & \textbf{Contenuto} & \textbf{Stato} \\ \midrule
G1 & 50 dipendenti & nome, codice fiscale, IBAN (email di cambio IBAN, moduli di onboarding) & letto \\
G2 & 3 dei 50 & usati per le richieste IBAN contraffatte & letto \\
G3 & $\approx$640 & buste paga 2024--26 in ZIP AES-256, password solo via SMS; mostrano quote
  sindacali, giorni di malattia, permessi Legge~104 & letto, cifrato \\
G4 & $\approx$600 & buste paga 2023, stesso contenuto; password di gennaio--aprile inviate via email alla
  stessa mailbox & non nei log \\
G5 & $\approx$520 + 53 & attestati di malattia senza diagnosi; 31 diagnosi; 22 pratiche Legge~104 &
  non nei log \\
G6 & $\approx$900 candidati & CV, di cui 60 per posizioni riservate a categorie protette (disabilità);
  molti oltre il periodo di conservazione & non nei log \\
\bottomrule
\end{tabularx}
\end{center}

\paragraph{Le tre posizioni.}
\begin{itemize}
  \item \textbf{CEO}: nessuna notifica, solo registro interno (``frode fallita, nessun euro perso,
    volevano i nostri soldi, non i dati'').
  \item \textbf{IT Manager}: notificare solo il Garante (``i log dicono cosa è stato letto, le buste
    paga sono cifrate, non allarmiamo 640 persone'').
  \item \textbf{HR Director}: Garante e tutti i dipendenti (``non possiamo provare cosa è successo in
    quelle 24 ore; diagnosi e quote sindacali sono lì; i sindacati lo verranno a sapere'').
\end{itemize}
Il DPO può adottarne una, combinarle o respingerle tutte: conta che l'opzione poggi sui fatti e su un
uso coerente del metodo.

\section{Consegna}
Parte~A (giovedì 1/10, ore 09:00):
\begin{enumerate}
  \item due perimetri: solo ciò che è stato letto o l'intera mailbox;
  \item punteggio ENISA di G1--G6 (DPC, EI, CB, flag);
  \item la severità non è il rischio: pesare anche la probabilità;
  \item scegliere un'opzione e dire chi viene informato;
  \item scadenza: da quando decorrono le 72 ore?
\end{enumerate}
In aula: formato libero (nessun modello fornito), consegna via email entro la fine del lab.

\section{Metodologia: ENISA 2013}
\[
SE = DPC \times EI + CB
\]
\begin{itemize}
  \item \textbf{DPC} (\emph{Data Processing Context}): valore di base per categoria -- dati semplici
    1, comportamentali 2, finanziari 3, sensibili 4 -- corretto in su o in giù (da 1 a 4) per fattori
    di contesto: volume, caratteristiche del titolare o degli interessati, inesattezza, disponibilità
    pubblica, natura del dato. Se ci sono più categorie vale il \textbf{valore più alto}. Ogni
    scostamento dal valore di base va motivato.
  \item \textbf{EI} (\emph{Ease of Identification}): 0,25 trascurabile, 0,5 limitata, 0,75
    significativa, 1 massima (identificazione diretta).
  \item \textbf{CB} (\emph{Circumstances of the Breach}): \textbf{somma} di perdita di riservatezza
    (0 / +0,25 destinatari noti / +0,5 destinatari indeterminati), integrità (0 / +0,25 recuperabile /
    +0,5 non recuperabile), disponibilità (0 / +0,25 temporanea / +0,5 totale), intento malevolo
    (+0,5).
  \item \textbf{Livelli}: $SE<2$ basso; $2 \le SE < 3$ medio; $3 \le SE < 4$ alto; $SE \ge 4$ molto
    alto.
  \item \textbf{Flag}: più di 100 interessati; dati resi incomprensibili (es.\ cifratura forte con
    chiave non compromessa).
\end{itemize}

\paragraph{Mappatura di lavoro (dal caso).} $SE<2$: solo registro (art.~33(5)); $2 \le SE <3$:
notifica al Garante (art.~33); $SE \ge 3$: Garante e interessati (art.~34, salvo 34(3)). È un punto
di partenza, non una regola: ci si può discostare spiegando perché.

\section{Svolgimento (soluzione modello, Parte~A)}

\subsection{Qualificazione}
\begin{itemize}
  \item È un \textbf{data breach} ai sensi dell'art.~4(12): accesso non autorizzato a dati personali
    (riservatezza), con un tentativo di violazione dell'integrità (cambio IBAN) non riuscito. ``Nessun
    euro perso'' è irrilevante: la violazione è l'accesso illecito.
  \item \textbf{Ruoli}: Arvenza è titolare; il payroll bureau è responsabile (art.~28) e si è comportato
    bene (la verifica telefonica ha impedito la frode); anche il provider della posta è responsabile e
    va interpellato per la telemetria sulle 24 ore mancanti.
  \item \textbf{Interessati}: dipendenti, ex dipendenti (buste 2023), familiari (pratiche Legge~104),
    candidati mai assunti.
  \item \textbf{Dati art.~9}: appartenenza sindacale (quote), salute (malattia, diagnosi), disabilità
    (Legge~104, categorie protette).
\end{itemize}

\subsection{Due perimetri}
\begin{itemize}
  \item \textbf{Perimetro A} (ciò che risulta letto): G1--G3, accesso certo. Lo sostiene l'IT Manager.
  \item \textbf{Perimetro B} (l'intera mailbox): G1--G6, accesso possibile ma non provato. Lo sostiene
    l'HR Director.
\end{itemize}
Entrambi vanno nell'analisi e nella notifica. Per l'\emph{accountability} l'onere della prova è del
titolare: l'assenza di prove causata dal proprio buco nei log non è prova di assenza. Il perimetro~B
però non va trattato come provato: per questo serve il passaggio sulla probabilità. A favore di una
probabilità bassa: obiettivo chiaramente finanziario, solo 14 minuti di attività nel buco, nessun
export. A sfavore: la ricerca ``payslips'' avrebbe fatto emergere anche le buste 2023 e le email con
le password.

\subsection{Punteggio}
Regole applicate in modo coerente:
\begin{itemize}
  \item \textbf{EI = 1} ovunque (nome + codice fiscale + datore di lavoro). Gli identificativi
    contano solo nell'EI: \textbf{niente doppio conteggio} nel DPC.
  \item Riservatezza \textbf{+0,5} per i gruppi letti (attore criminale, divulgazione non
    delimitabile); \textbf{0} per quelli non nei log (esposti senza prova di trattamento illecito,
    come nell'esempio ENISA).
  \item Intento malevolo \textbf{+0,5} per tutti: caratterizza la violazione, non il singolo dato.
  \item Integrità \textbf{+0,25} solo per G2 (identità usata, alterazione recuperabile). Disponibilità 0.
  \item Gruppi misti divisi in sottogruppi; per ogni gruppo vale il DPC più alto.
\end{itemize}

\begin{center}
\small
\begin{tabularx}{\linewidth}{l c X c c c l}
\toprule
\textbf{Gruppo} & \textbf{N.} & \textbf{DPC (motivo)} & \textbf{EI} & \textbf{CB} & \textbf{SE} &
  \textbf{Livello} \\ \midrule
G1 & 50 & 2: finanziari (3) ridotti di 1, solo IBAN senza saldi & 1 & 1,0 & 3,0 & alto \\
G2 & 3 & 2: come G1 & 1 & 1,25 & 3,25 & alto \\
G3 & $\approx$640 & 4: quote sindacali, malattia, L.~104 & 1 & 1,0 & 5,0 & molto alto$^{*}$ \\
G4 & $\approx$600 & 4: come G3 & 1 & 0,5 & 4,5 & molto alto \\
G5a & $\approx$520 & 2: sensibili ridotti, solo ``malato nei giorni X'' & 1 & 0,5 & 2,5 & medio \\
G5b & 31 & 4: diagnosi & 1 & 0,5 & 4,5 & molto alto \\
G5c & 22 & 4: disabilità, spesso di un familiare & 1 & 0,5 & 4,5 & molto alto \\
G6a & $\approx$840 & 1: CV ordinari & 1 & 0,5 & 1,5 & basso \\
G6b & 60 & 3: la posizione riservata rivela una disabilità (4 sostenibile) & 1 & 0,5 & 3,5 & alto \\
\bottomrule
\end{tabularx}
\end{center}
\noindent{\small $^{*}$ punteggio grezzo, prima del flag ``dati incomprensibili''. Flag ``più di 100''
per G3, G4, G5a, G6a.}

\paragraph{Scelte al limite.}
\begin{itemize}
  \item \textbf{G1 cade esattamente su 3,0}: con riservatezza +0,25 scenderebbe a 2,75 (medio). Si
    tiene +0,5 perché il destinatario è un criminale non identificato e i danni possibili
    (sottrazione di fondi, frodi su addebiti SEPA, vishing mirato) sono quelli del livello ``alto''. La
    decisione non cambia comunque, perché la probabilità è certa.
  \item \textbf{G3}: cifratura forte con chiave mai transitata nella mailbox $\rightarrow$ flag. Due
    verifiche: robustezza delle password (uno ZIP si può attaccare offline se la password è corta o
    derivata dal codice fiscale) e metadati in chiaro (nomi dei file, oggetti delle email).
  \item \textbf{G4 gennaio--aprile} è il punto debole: file cifrati e password nella stessa mailbox,
    quindi il flag non regge per quei mesi.
  \item \textbf{G6}: CV oltre la conservazione $\rightarrow$ violazione autonoma dell'art.~5(1)(e); non
    cambia SE ma aggrava la posizione del titolare.
\end{itemize}

\subsection{La severità non è il rischio}
\begin{center}
\small
\begin{tabularx}{\linewidth}{l c X l}
\toprule
\textbf{Gruppo} & \textbf{SE} & \textbf{Probabilità di accesso/abuso} & \textbf{Rischio / obbligo} \\ \midrule
G1 & 3,0 & certa: letti e usati per la frode & alto: art.~33 + 34 \\
G2 & 3,25 & già materializzata & alto: art.~33 + 34 (priorità, al telefono) \\
G3 & 5,0* & letti ma incomprensibili & basso se password robuste: 33 + esenzione 34(3)(a) \\
G4 & 4,5 & possibile per gen.--apr., remota per il resto & medio: 33, rivalutare \\
G5a & 2,5 & bassa & basso: 33 \\
G5b/c & 4,5 & bassa (non nei log, 14 minuti, nessun export) & medio: 33, rivalutare \\
G6a & 1,5 & bassa & basso: registro \\
G6b & 3,5 & bassa & medio: 33, rivalutare \\
\bottomrule
\end{tabularx}
\end{center}
L'art.~34 richiede un rischio \emph{elevato}, cioè severità \emph{e} probabilità. Mandare centinaia di
comunicazioni per qualcosa che probabilmente non è avvenuto genererebbe paura e stress (danni che
la stessa ENISA elenca) e diluirebbe il messaggio ai 50 davvero esposti. Il Garante, informato
dell'intero perimetro, può comunque ordinare la comunicazione (art.~34(4)).

\subsection{Decisione}
\begin{soluzione}[Parere del DPO (soluzione modello)]
\begin{itemize}
  \item \textbf{Garante}: notifica \textbf{preliminare oggi entro le 15:00} (art.~33(4)),
    sull'intero perimetro A+B, con tabella di severità, ragionamento sulla probabilità e motivi del
    ritardo; notifica completa dopo la forensics.
  \item \textbf{G2}: telefonata oggi (HR + DPO), poi comunicazione scritta.
  \item \textbf{G1}: comunicazione scritta individuale oggi o domani, in linguaggio semplice e con
    consigli pratici (art.~34(2)).
  \item \textbf{G3}: nessuna comunicazione individuale, una volta verificata la robustezza delle
    password; esenzione documentata (art.~34(3)(a)).
  \item \textbf{G4--G6}: nessuna comunicazione per ora, con \textbf{trigger di rivalutazione}
    documentati.
  \item \textbf{Tutto il personale e RSU}: breve nota interna fattuale (trasparenza, non notifica
    ex art.~34).
  \item \textbf{Polizia Postale}: denuncia (accesso abusivo, sostituzione di persona, tentata truffa):
    non obbligatoria ma apprezzata dal Garante.
  \item \textbf{Responsabili}: dichiarazione scritta del bureau; log e telemetria dal provider di
    posta.
  \item \textbf{Registro} (art.~33(5)): voce completa con la motivazione di ogni scelta.
\end{itemize}
Posizione: si combina l'IT Manager (Garante) con parte dell'HR Director (interessati, ma solo dove il
rischio è alto) e si \textbf{respinge il CEO}.
\end{soluzione}

\paragraph{Perché non le tre posizioni così come sono.}
\begin{itemize}
  \item \textbf{CEO}: l'art.~33(1) esonera dalla notifica (lasciando solo la documentazione nel registro
    dell'art.~33(5)) se è \emph{improbabile} che la violazione presenti un rischio. Qui un attaccante ha letto 214 elementi e usato i dati per impersonare
    dipendenti: ``volevano i nostri soldi'' è proprio la ragione per cui il rischio è alto. Non
    notificare trasformerebbe una notifica tardiva in una violazione occultata.
  \item \textbf{IT Manager}: giusto sul Garante, sbagliato su due punti: i log con un buco di 24 ore
    non possono definire il perimetro; la cifratura protegge G3, non G1.
  \item \textbf{HR Director}: giusta sul buco nei log e sulle categorie particolari, ma troppo ampia:
    per G4--G6 la probabilità è bassa e per G3 c'è un'esenzione esplicita. La preoccupazione per i
    sindacati si gestisce con la nota interna.
\end{itemize}

\subsection{Da quando decorrono le 72 ore}
\begin{center}
\small
\begin{tabularx}{\linewidth}{l X l l}
\toprule
\textbf{Candidato} & \textbf{Argomento} & \textbf{Scadenza} & \textbf{Esito} \\ \midrule
Mar 15/9 & solo l'attaccante sapeva; conta la consapevolezza del titolare & -- & no \\
Dom 20/9 & Arvenza sapeva che un attaccante aveva avuto una sessione nella mailbox paghe; chiudere il
  ticket senza guardare non sposta la consapevolezza (EDPB 9/2022) & Mer 23/9 & la più credibile \\
Lun 28/9 15:00 & primo collegamento esplicito tra alert e dati personali & Gio 1/10 15:00 &
  difendibile al massimo \\
\bottomrule
\end{tabularx}
\end{center}
Si notifica \textbf{oggi prima delle 15:00}, così da rispettare anche la lettura più favorevole, e
nella notifica si dichiara l'alert del 20/9 spiegando il ritardo (art.~33(1)): l'alert è stato gestito
come ticket IT e mai portato al DPO. Ammetterlo è meglio che farlo scoprire al Garante; la mancata
escalation diventa una voce del remediation plan. L'art.~34 non ha un termine in ore, ma dopo 11
giorni di ritardo ``senza ingiustificato ritardo'' significa: i tre oggi, i 50 entro 24 ore.

\subsection{Contenuto delle comunicazioni}
\begin{itemize}
  \item \textbf{Notifica preliminare al Garante} (art.~33(3)): natura (violazione di riservatezza via
    phishing AiTM, accesso 15--20/9, tre richieste IBAN respinte), interessati e registrazioni (letti
    ed esposti, distinti), categorie (incluse quelle dell'art.~9), cronologia e motivi del ritardo,
    conseguenze probabili (frodi sugli stipendi o SEPA, phishing e vishing mirati, discriminazione
    se toccate le categorie particolari), severità e probabilità, misure, contatti del DPO, fase
    preliminare.
  \item \textbf{Comunicazione ai 50} (art.~34(2)): cosa è successo in parole semplici (lo stipendio è
    stato pagato correttamente); cosa può accadere; cosa ha fatto l'azienda; cosa può fare
    l'interessato (non confermare mai dati bancari via email, controllare gli addebiti); contatto del
    DPO. Per i tre impersonati, una frase esplicita sulle richieste inviate a loro nome.
\end{itemize}

\subsection{Remediation (lacune ex artt.~32 e 5)}
\begin{center}
\small
\begin{tabularx}{\linewidth}{X X}
\toprule
\textbf{Lacuna emersa} & \textbf{Misura} \\ \midrule
MFA aggirata con furto del token di sessione & MFA resistente al phishing (FIDO2/passkey), accesso
  condizionale, sessioni brevi \\
Regola malevola non rilevata & alert sulle nuove regole di inoltro/occultamento, revisione periodica \\
Alert chiuso senza valutare l'impatto sui dati & procedura: ogni compromissione di un sistema con
  dati personali va al DPO entro poche ore \\
Buco di 24 ore nei log & log estesi, conservazione $\ge$ 6 mesi, SIEM, impegno contrattuale del
  provider \\
Dati art.~9 e password nella mailbox condivisa & spostare i documenti sanitari nel sistema HR con
  accessi per ruolo, eliminare le email con password, minimizzazione (artt.~5(1)(c), 25) \\
CV oltre la conservazione & applicare il piano di conservazione (art.~5(1)(e)) \\
Cambi IBAN richiesti via email & rendere obbligatoria la verifica telefonica del bureau \\
Fattore umano & formazione mirata anti-phishing per HR e paghe \\
\bottomrule
\end{tabularx}
\end{center}

\paragraph{Fatti che cambierebbero la decisione (preparazione alla Parte~B).} Letture recuperate nel
buco su G5b, G5c o G6b $\rightarrow$ art.~34 per quelle persone; letture delle buste 2023 e delle
email con password $\rightarrow$ art.~34 per G4; password di G3 deboli $\rightarrow$ cade l'esenzione;
dati su un leak site o richiesta di estorsione $\rightarrow$ riservatezza +0,5 per tutti e art.~34
ampio; un dipendente segnala un addebito sconosciuto $\rightarrow$ conferma dell'abuso; log che
escludono l'accesso nel buco $\rightarrow$ il perimetro~B si riduce, ma resta nel registro.

\section{Errori tipici emersi nel debrief}
\begin{itemize}
  \item \textbf{Abbassare il DPC dei dati sanitari senza motivare}: gli attestati di malattia sono
    dati relativi alla salute (art.~9) e partono da DPC~4; scendere a 2 è sostenibile solo
    argomentando (Gianluca, 12/10).
  \item \textbf{Data di consapevolezza troppo tardiva} senza evidenze: un ispettore la contesterebbe;
    chi la sostiene deve spiegare perché la conoscenza effettiva è arrivata dopo.
  \item \textbf{Confondere severità e rischio}: il punteggio ENISA misura l'impatto, non la
    probabilità.
  \item \textbf{Sovrastimare il CB}: aggiungere integrità e disponibilità ``per prudenza'' quando
    nessun dato è stato alterato o perso.
  \item \textbf{Doppio conteggio} degli identificativi (codice fiscale nell'EI e di nuovo nel DPC).
  \item Trattare il perimetro~B come provato, oppure ignorarlo perché ``non è nei log''.
  \item Dimenticare categorie di interessati (es.\ gli $\approx$900 candidati, i familiari nelle pratiche
    Legge~104).
  \item Errori formali: lavoro consegnato senza allegato o senza nomi dei componenti.
\end{itemize}

\section{Collegamenti}
\begin{itemize}
  \item GDPR: artt.~4(12), 5(1)(c),(e),(f), 5(2), 9, 25, 28, 32, 33, 34 (Capitolo~\ref{cap:gdpr} e
    Capitolo~\ref{cap:data-breach}).
  \item Lezione del 05/10: consapevolezza del titolare, registro, carenza organizzativa
    (Capitolo~\ref{lez:2026-10-05}).
  \item Lab~2: stessa metodologia, ma con violazione comunicata dal responsabile e doppio binario
    GDPR/DORA (Capitolo~\ref{lab:02}).
\end{itemize}

\begin{verifica}
\begin{itemize}
  \item La soluzione modello richiama le \textbf{Linee guida EDPB 9/2022} per la nozione di
    consapevolezza: riferimento contenuto nella soluzione, non commentato in aula.
  \item La \textbf{Parte~B} (aggiornamento dell'8 ottobre) annunciata nel caso non è nel materiale
    disponibile.
  \item Gli articoli del codice penale citati dalla soluzione per la denuncia (615-ter, 494, 56 e
    640~c.p.) non sono stati discussi in aula.
\end{itemize}
\end{verifica}

% =====================================================================
\chapter{Lab 2 -- Sylvane Retraite: il canale di scambio file}
\label{lab:02}
% =====================================================================

\begin{nota}[Fonti e date]
Caso presentato da Gianluca il 12/10/2026 (Capitolo~\ref{lez:2026-10-12}), svolto venerdì 16/10/2026
(Capitolo~\ref{lez:2026-10-16}), discusso nel debrief di lunedì 19/10/2026. Materiale: README del caso,
\emph{Sylvane Retraite Breach Case Study} (PDF), \emph{Severity Assessment Worksheet} (obbligatorio),
\emph{Notification Form} (Excel, condizionale), immagine della \emph{curveball~1}. Caso di fantasia.
\textbf{Non esiste una soluzione ufficiale}: lo svolgimento ragionato è dei redattori ed è segnalato
come tale.
\end{nota}

\section{Il caso in sintesi}

\paragraph{Contesto.} Sylvane Retraite è un fondo pensione professionale francese (IRPS), titolare
dei dati dei propri iscritti. Un fornitore esterno (responsabile del trattamento) carica file
operativi su un \textbf{server SFTP nella DMZ del fondo}, che funge da ``cassetta delle lettere'':
il fondo sposta poi i file nei propri ambienti interni. L'accesso al server è con autenticazione a
\textbf{certificato} (niente password) e solo dall'IP pubblico del fornitore. Gli studenti sono il
\textbf{DPO}, che venerdì 11 aprile 2025 pomeriggio deve consigliare CEO, CISO e Head of Member
Services.

\paragraph{Cronologia (2025).}
\begin{center}
\small
\begin{tabularx}{\linewidth}{p{2.6cm} X}
\toprule
\textbf{Quando} & \textbf{Evento} \\ \midrule
Mer 2/4, 08:30 & Inizia l'accesso non autorizzato al server SFTP, attraverso i sistemi del fornitore e
  con le sue credenziali legittime. \\
2--4/4 & Circa 53 ore di permanenza; file esfiltrati. \\
Ven 4/4, 13:37 & Finestra chiusa: accesso bloccato, arrivano i forensi. Il fornitore viene a
  conoscenza dell'incidente e lo riconduce ai propri sistemi. \\
Sab 5 -- Gio 10/4 & ``Settimana silenziosa'': il fornitore completa i controlli; il fondo non sa nulla. \\
Ven 11/4, 10:44 & Comunicazione formale del fornitore via email certificata: ``qualcosa'' è stato
  esfiltrato, senza dettagli. Il fornitore ha già notificato la propria agenzia nazionale di
  cybersicurezza. \\
Ven 11/4, 11:12 & Il fondo classifica l'incidente come \textbf{grave ai sensi del DORA}; nessun
  impatto sulla continuità. \\
Ven 11/4, pomeriggio & Punto di decisione (Parte~A, ore 14:00). \\
\bottomrule
\end{tabularx}
\end{center}

\paragraph{Che cosa passava nel canale.}
\begin{center}
\small
\begin{tabularx}{\linewidth}{l X l}
\toprule
\textbf{Gruppo} & \textbf{Contenuto} & \textbf{Stato} \\ \midrule
G1 & nome e cognome dell'iscritto + valore totale della posizione pensionistica maturata (dati comuni +
  finanziari); oltre 100 iscritti, numero non ancora quantificato & confermato \\
G2 & file operativi caricati dal fornitore; categorie e volume da determinare & confermato \\
G3 & il server riceveva un \textbf{export giornaliero completo} del database: dati storici e recenti
  non separabili & indeterminato \\
G4 & quattro ambienti interni alimentati dalla DMZ (antiriciclaggio, BI, risk, middleware): non
  raggiunti & non coinvolto \\
G5 & categorie particolari (art.~9): nessuna individuata finora & nessuna \\
G6 & pubblicazione o abuso: nessuna evidenza (``l'assenza di prova non è prova di assenza'') & non
  accertato \\
\bottomrule
\end{tabularx}
\end{center}

\paragraph{Le tre posizioni.}
\begin{itemize}
  \item \textbf{CEO}: nessuna notifica, solo registro. L'accesso è passato dai sistemi del fornitore,
    il server è protetto da certificati, nulla è stato pubblicato; informare gli iscritti porterebbe
    solo un'ondata di telefonate. Si rivaluterà con nuove informazioni.
  \item \textbf{CISO}: solo CNIL. Nome, cognome e valore della posizione sono stati esfiltrati: c'è un
    rischio ex art.~33; notifica preliminare e integrazioni successive.
  \item \textbf{Head of Member Services}: CNIL e tutti gli iscritti coinvolti. Se un truffatore chiama
    un iscritto citando l'importo esatto della sua posizione (social engineering) e noi non abbiamo
    detto nulla, non abbiamo difese.
\end{itemize}

\section{Consegna}
\begin{itemize}
  \item \textbf{Tempi}: 105 minuti di \emph{wall clock} = tre giorni di \emph{exercise clock}. Le nuove
    informazioni (\emph{curveball}) arrivano in chat durante il lab.
  \item \textbf{Worksheet obbligatorio} per tutti, qualunque sia la conclusione (anche ``nessuna
    notifica''). Se la notifica è dovuta, si compila anche il \textbf{modulo di notifica} e si carica
    un unico zip con entrambi i file, non protetto da password.
  \item Nome del gruppo e dei componenti nell'intestazione e nel modulo di caricamento; file nominati
    con il nome del gruppo. Lavori non attribuibili: scartati.
  \item Compiti della Parte~A: (1) due perimetri (solo ciò che è stato preso o tutto ciò che il server
    conteneva); (2) punteggio dei gruppi di dati; (3) la severità non è il rischio; (4) scegliere chi
    informare: iscritti, CNIL, ACPR o nessuno; (5) da quando decorrono le 72 ore, e dalla conoscenza
    di chi.
  \item Indicazione data durante il lab: serve una valutazione \textbf{complessiva} della violazione,
    non un punteggio per ogni voce.
\end{itemize}

\subsection{Il worksheet}
Struttura (10 sezioni): sintesi dei fatti; DPC (tabella delle quattro categorie, con motivazione di
ogni scostamento); EI (0,25 / 0,50 / 0,75 / 1,00); CB (somma di riservatezza, integrità, disponibilità,
intento malevolo); calcolo di SE; fascia di rischio; decisione (solo registro / CNIL / CNIL e iscritti
/ segnalazione DORA all'ACPR); perimetro; \textbf{giustificazione} (la risposta da dare all'autorità
se chiede ``cosa avete fatto e perché''); \textbf{strategia} (il ragionamento interno, ``ciò che non
si può dire ad alta voce'').

\begin{nota}[Honne e tatemae]
Gianluca ha presentato le ultime due sezioni con due concetti giapponesi: \textbf{tatemae} è la
facciata, ciò che si mostra all'autorità (sezione~9); \textbf{honne} è la verità, il motivo reale
della scelta, che può includere anche il timore di essere licenziati (sezione~10). Nella realtà sono
due cose molto diverse.
\end{nota}

\subsection{Il modulo di notifica}
È la trasposizione, in Excel, del modulo online dell'autorità (art.~33): sezioni A (chi notifica: il
legale rappresentante o un delegato, con dichiarazione sulle conseguenze di dichiarazioni false),
B (notifica preliminare o completa), C (titolare), D (DPO), E (altri soggetti: il fornitore come
responsabile), F (informazioni sulla violazione: date, modalità di conoscenza, motivi del ritardo,
natura, causa, sistemi, misure in essere, categorie e numero di interessati e registrazioni),
H (misure adottate e per il futuro), L (comunicazione agli interessati, con il \textbf{testo effettivo}
della comunicazione), M (altre autorità: già indicata la notifica all'ACPR ai sensi del DORA),
N (trattamenti transfrontalieri). Alcuni campi sono precompilati perché lo studente non può
conoscerli.

\subsection{La curveball~1 (venerdì 11/4, 16:30)}
Il fornitore ha inviato una comunicazione identica a un'altra società del gruppo, che fa
\textbf{wealth management} per clienti facoltosi e di alto profilo. Stesso fornitore, stesso canale,
stessa finestra. Domande guida: il proprio perimetro e punteggio; chi è titolare dei dati della
società sorella; se il gruppo ha ora un problema che nessuna singola entità può risolvere da sola.

\section{Metodologia}
Stessa formula del Lab~1 (scale complete nel Capitolo~\ref{lab:01}):
\[
SE = DPC \times EI + CB
\]
con la mappatura di lavoro $SE<2$ solo registro, $2 \le SE <3$ CNIL, $SE \ge 3$ CNIL e iscritti
(salvo art.~34(3)). In parallelo corre un \textbf{secondo orologio}: secondo il caso, ai sensi del
DORA il fondo deve inviare all'ACPR una notifica iniziale entro 24 ore, intermedia entro 72 ore e
finale entro un mese. È una questione distinta dalla notifica GDPR.

\section{Svolgimento ragionato}

\begin{integrazione}[Svolgimento dei redattori -- non è una soluzione ufficiale]
\paragraph{Qualificazione.} È un data breach ai sensi dell'art.~4(12): esfiltrazione confermata di
dati personali degli iscritti (riservatezza) da parte di un terzo non autorizzato. Il fatto che il
vettore d'attacco sia nei sistemi del fornitore non toglie al fondo, in quanto titolare, gli obblighi
degli artt.~33 e 34; incide sulla responsabilità del fornitore (art.~28) e sui rapporti contrattuali.

\paragraph{Perimetri.} Perimetro~A: ciò che risulta preso (G1, più G2 per quanto si saprà). Perimetro~B:
tutto ciò che il server riceveva, incluso l'export giornaliero completo (G3), dove dati storici e
recenti non si separano. G4 è escluso (accesso fermo alla DMZ), G5 vuoto allo stato.

\paragraph{Punteggio di G1.}
\begin{itemize}
  \item \textbf{DPC}: dati finanziari, base 3. Il valore della posizione dà un'indicazione concreta
    della situazione patrimoniale, quindi non è un ``semplice numero di conto'' (che giustificherebbe
    2). Abbassarlo a 2 è sostenibile solo argomentando che un singolo importo, senza movimenti, dà un
    quadro limitato. Potrebbe salire a 4 se l'export storico permette di ricostruire l'andamento nel
    tempo (volume).
  \item \textbf{EI}: nome e cognome, nel contesto di un fondo specifico e associati a un importo:
    identificazione significativa o diretta (0,75--1).
  \item \textbf{CB}: riservatezza +0,5 (esfiltrazione verso un soggetto criminale, destinatari non
    delimitabili; +0,25 sostenibile se si ritiene il destinatario unico); intento malevolo +0,5;
    integrità e disponibilità 0 (nessuna alterazione, continuità intatta). CB = 1,0.
\end{itemize}
\begin{center}
\small
\begin{tabular}{cccll}
\toprule
\textbf{DPC} & \textbf{EI} & \textbf{CB} & \textbf{SE} & \textbf{Livello} \\ \midrule
2 & 1 & 1,0 & 3,0 & alto (al limite) \\
3 & 0,75 & 1,0 & 3,25 & alto \\
3 & 1 & 1,0 & 4,0 & molto alto \\
\bottomrule
\end{tabular}
\end{center}
Flag: più di 100 interessati. Il caso è costruito per stare \textbf{sul confine} tra medio e alto.

\paragraph{Severità e probabilità.} L'accesso a G1 è certo; l'abuso non è accertato, ma la
combinazione nome + importo è esattamente ciò che serve per un attacco di social engineering
credibile (``la chiamo dal fondo, la sua posizione è di X euro''). La curveball~1 aumenta la
probabilità: più società finanziarie colpite attraverso lo stesso fornitore suggeriscono un'operazione
organizzata e a scopo di lucro.

\paragraph{Decisione.}
\begin{itemize}
  \item \textbf{CNIL}: notifica \textbf{preliminare} (art.~33(4)), senza attendere il quadro completo,
    con impegno a integrarla. La posizione del CEO va respinta: protezione con certificati e colpa del
    fornitore non rilevano, perché l'attaccante ha usato le credenziali legittime del fornitore e i
    dati sono comunque usciti; ``nessun abuso rilevato'' non equivale a ``rischio improbabile''.
  \item \textbf{Iscritti}: la tesi dell'Head of Member Services è forte. Una comunicazione (art.~34)
    con consigli pratici, ad esempio ``il fondo non vi chiederà mai password o pagamenti al
    telefono'', riduce proprio il rischio principale. In alternativa, chi sceglie ``solo CNIL'' deve
    documentare i trigger che farebbero scattare la comunicazione.
  \item \textbf{ACPR}: notifica DORA separata, già avviata con la classificazione delle 11:12.
  \item \textbf{Fornitore}: richiesta formale di informazioni (categorie, volume, finestra, log),
    contestazione della settimana di silenzio (art.~33(2)), riserva di azioni contrattuali.
  \item \textbf{Registro} (art.~33(5)): tutte le scelte e i relativi motivi.
\end{itemize}

\paragraph{Le 72 ore.} Per il titolare decorrono da quando \emph{il titolare} è venuto a conoscenza:
venerdì 11/4 alle 10:44 (comunicazione del responsabile), quindi scadenza \textbf{lunedì 14/4 alle
10:44} (il weekend conta). Questo regge solo se il fondo può dimostrare di non aver potuto sapere
prima: nessuna visibilità sul server durante la settimana silenziosa, nessun alert. Il ritardo del
fornitore (conoscenza il 4/4) è una sua violazione dell'art.~33(2), non sposta indietro l'orologio
del titolare. È lo schema anticipato dal prof.~Vaciago il 05/10 (Capitolo~\ref{lez:2026-10-05}).

\paragraph{La curveball~1.} La società sorella è un \textbf{titolare distinto} dei dati dei propri
clienti: valuta e notifica per conto proprio, e per i suoi clienti di alto profilo il DPC sale
(caratteristiche particolari degli interessati). Il punteggio di Sylvane, sul proprio perimetro, non
cambia; cambiano la probabilità di abuso e la dimensione di gruppo: comunicazioni coerenti (se la
sorella informa i clienti e Sylvane no, il messaggio è incoerente), un'unica linea verso il fornitore
condiviso, coordinamento della risposta.
\end{integrazione}

\section{Indicazioni emerse nel debrief (19/10)}
\begin{esame}
\begin{itemize}
  \item \textbf{Riservatezza senza esfiltrazione}: la violazione di riservatezza c'è anche se i file
    sono stati solo visti e non portati via; l'impatto sugli interessati, però, è molto diverso
    (prof.~Vaciago).
  \item \textbf{Le fasce di rischio sono una guida, non una regola matematica}: nel worksheet le fasce
    ``$2 \le 3$'' e ``$3 \le 4$'' si sovrappongono proprio su 3. Con $SE = 3$ si è in zona grigia: si
    decide e ci si assume la responsabilità in base all'\emph{accountability}; l'autorità può valutare
    diversamente. Gruppi diversi, sullo stesso caso, sono arrivati a numeri molto diversi.
  \item \textbf{Non presumere nulla che non si possa provare} (Gianluca): ogni dichiarazione falsa
    all'autorità è un reato e chi la firma ne risponde personalmente. Nel dubbio ci si attiene alle
    evidenze. Il prof.~Vaciago aggiunge che serve il dolo (\emph{mens rea}): l'errore in buona fede
    non è punito, ma va poi spiegato a un pubblico ministero spesso digiuno di tecnica.
  \item \textbf{Responsabilità del fornitore}: se il vettore d'attacco è nell'infrastruttura del
    responsabile, è sua la colpa di essere stato violato, anche se i dati stavano sul server del
    titolare.
\end{itemize}
\end{esame}

\paragraph{Una delle soluzioni presentate.} Un gruppo ha diviso i dati in due insiemi (canale di
scambio; una copia di backup sul server con dati più ricchi, senza evidenza di accesso nei log),
notificando CNIL e iscritti (lettera con l'avvertimento che il fondo non chiederà mai password o
denaro) e segnalando all'ACPR. Ha fatto partire le 72 ore dall'11/4, potendo dimostrare di non aver
saputo prima, e ha proposto una risposta condivisa a livello di gruppo, documentando e firmando ogni
passaggio. Il docente ha giudicato la presentazione molto buona.

\section{Errori tipici}
\begin{itemize}
  \item Accettare l'argomento ``è colpa del fornitore'' come motivo per non notificare: gli obblighi
    di notifica restano del titolare.
  \item Far partire le 72 ore dal 4/4 (conoscenza del fornitore) oppure, all'opposto, rimandarle in
    attesa di dettagli: la notifica preliminare esiste proprio per questo.
  \item Trattare la fascia di rischio come una regola matematica, specialmente con $SE = 3$.
  \item Fondare la decisione su ipotesi non documentate (es.\ ``probabilmente sono solo dati
    recenti'').
  \item Confondere la notifica GDPR alla CNIL con la segnalazione DORA all'ACPR: due binari, due
    autorità, due orologi.
  \item Valutare voce per voce invece di dare un giudizio complessivo sulla violazione.
  \item Ignorare la curveball o, all'opposto, assorbire nel proprio perimetro i dati della società
    sorella, che ha un proprio titolare.
\end{itemize}

\section{Collegamenti}
\begin{itemize}
  \item GDPR: artt.~4(12), 28, 33(2), 33(4), 33(5), 34 (Capitolo~\ref{cap:data-breach}).
  \item DORA: notifica degli incidenti gravi all'autorità finanziaria (Capitolo~\ref{cap:dora}).
  \item Lezione del 05/10: ritardo del responsabile e decorrenza delle 72 ore; art.~168 Codice privacy
    sulle false dichiarazioni (Capitolo~\ref{lez:2026-10-05}).
  \item Lab~1: metodologia e scale ENISA (Capitolo~\ref{lab:01}).
\end{itemize}

\begin{verifica}
\begin{itemize}
  \item \textbf{Curveball successive}: nel materiale c'è solo la curveball~1. Dal debrief del 19/10
    (trascrizione molto disturbata) sembrano esserne emerse altre: un account non autorizzato presente
    da tempo, il coinvolgimento di altre istituzioni finanziarie e una \textbf{copia di backup} sul
    server con dati più ricchi (identificativi fiscali, dati KYC, transazioni) senza evidenza di
    accesso. Da ricostruire con certezza.
  \item \textbf{Scadenze DORA}: il caso indica notifica iniziale entro 24 ore dalla rilevazione,
    intermedia entro 72 ore, finale entro un mese. Le norme tecniche attuative del DORA fissano la
    notifica iniziale a 4 ore dalla classificazione (e comunque entro 24 ore dalla conoscenza):
    verificare nel capitolo DORA quale formulazione usa il corso.
  \item La data dell'esercizio è l'11 aprile \textbf{2025} (come nel materiale), non 2026.
  \item Lo svolgimento ragionato non è stato validato dai docenti: i punteggi sono alternative
    argomentabili, non la ``risposta giusta''.
\end{itemize}
\end{verifica}

% =====================================================================
\chapter{Lab 3 -- Risk register ISO/IEC 27001 per uno studio professionale}
\label{lab:03}
% =====================================================================

\begin{nota}[Fonte]
Laboratorio del 30/10/2026 (Capitolo~\ref{lez:2026-10-30}), presentato da Gianluca a fine lezione il
26/10 (Capitolo~\ref{lez:2026-10-26}). Fonti: le slide \emph{The Risk Register Exercise} (briefing), il
\emph{Risk Register Exercise Playbook} e il file Excel dell'esercitazione
(\texttt{lab/lab\_03/materiale}). Richiami teorici: slide \emph{04.~polito\_iso\_2026} (controlli 8.x) e
\emph{06.~27001\_2026\_2310} (\emph{Risk Management}). I due esempi R-01 e R-17 sono quelli svolti nel
playbook; le altre righe ragionate \textbf{non} sono soluzioni ufficiali e si trovano nei riquadri
\emph{Integrazione}.
\end{nota}

% ---------------------------------------------------------------------
\section{Il contesto}

Si è il \textbf{responsabile della sicurezza delle informazioni} di un piccolo \textbf{studio
professionale} (consulenza fiscale, legale o tecnica) che tratta dati personali dei clienti e documenti
riservati su servizi cloud, portatili e un gestionale di studio. Lo studio sta implementando un ISMS in
vista della certificazione ISO/IEC~27001:2022 e ha già individuato \textbf{18 scenari di rischio}. Il
compito è completarne la \textbf{valutazione} e il \textbf{trattamento}.

\paragraph{Perimetro.} Il registro riguarda \textbf{solo la clausola~8 dell'Allegato~A}, cioè i 34
controlli tecnologici. I rischi trattabili solo con controlli organizzativi, documentali o fisici
(clausole 5, 6, 7) stanno in registri separati, così nessun rischio viene valutato due volte. Ogni riga
deve quindi citare almeno un controllo 8.x: un rischio senza controllo tecnologico applicabile appartiene
ai registri organizzativi.

\paragraph{Il punto di partenza.} Gli scenari descrivono uno studio con una \textbf{baseline debole}:
quasi nessun controllo tecnologico in essere. È proprio per questo che i rischi sono stati sollevati, e
il rischio inerente va stimato con i controlli che lo studio ha oggi.

\begin{center}
\small
\begin{tabularx}{\linewidth}{l >{\raggedright\arraybackslash}p{3.2cm} X}
\toprule
\textbf{ID} & \textbf{Categoria} & \textbf{Scenario (in sintesi)} \\ \midrule
R-01 & Accessi e identità & credenziali di gestionale ed email compromesse (niente MFA, password
  riusate o condivise): accesso ai dati fiscali dei clienti \\
R-02 & Malware e phishing & ransomware da allegato o link malevolo: documenti cifrati, attività ferma \\
R-03 & Endpoint & accesso ai dati da dispositivi personali o non gestiti (BYOD, lavoro da remoto): dati
  portati via o contaminati, nessuna tracciabilità \\
R-04 & Backup e continuità & backup non eseguiti, non cifrati o mai verificati: perdita definitiva
  dell'archivio \\
R-05 & Logging e monitoraggio & nessun log centralizzato né monitoraggio: incidenti non rilevati, eventi
  non ricostruibili in un contenzioso \\
R-06 & Cloud e SaaS & documenti condivisi con link pubblici o cartelle aperte a tutti: dati dei clienti
  esposti \\
R-07 & Vulnerabilità e patch & sistemi operativi e applicazioni obsoleti: vulnerabilità note sfruttate \\
R-08 & Rete & configurazioni insicure (credenziali di default, servizi esposti inutili, protocolli
  obsoleti) \\
R-09 & Rete & rete unica e piatta (ospiti, IoT, postazioni): movimento laterale dell'attacco \\
R-10 & Cifratura e trasferimenti & documenti trasmessi in chiaro (email, chiavette USB, link cloud) \\
R-11 & Endpoint & furto o smarrimento di portatili e smartphone con documenti non cifrati \\
R-12 & Fornitori IT & amministratore di sistema esterno con privilegi permanenti e attività non
  registrata: modifiche non attribuibili \\
R-13 & Applicazioni & gestionale e portale clienti acquistati o commissionati senza requisiti né test di
  sicurezza \\
R-14 & Change management & aggiornamenti e modifiche non controllati: regressioni, fermi, configurazioni
  non conformi \\
R-15 & Backup e continuità & saturazione o indisponibilità di servizi critici (connettività, storage,
  cloud) senza ridondanza: scadenze dei clienti mancate \\
R-16 & Dati e supporti & dismissione di PC e supporti senza cancellazione certificata: dati recuperati
  da apparati venduti, smaltiti o in riparazione \\
R-17 & Servizi esposti & DDoS o indisponibilità di sito e portale clienti: interruzione e danno
  reputazionale \\
R-18 & Dati e supporti & metadati nascosti nei documenti inviati (autori, revisioni, percorsi,
  commenti): fuga di informazioni riservate \\ \bottomrule
\end{tabularx}
\end{center}

% ---------------------------------------------------------------------
\section{La consegna}

\paragraph{Il file.} Sette fogli: \emph{Group Name and members} (da compilare per primo),
\emph{Instructions} (il metodo), \emph{Risk Register} (il lavoro: 18 righe), \emph{Annex A clause 8}
(i 34 controlli con nota di applicabilità e attuazione tipica), \emph{3x3 Matrix}, \emph{Summary}
(indicatori e controlli di coerenza), \emph{Lists} (i menu a tendina).

\paragraph{Le colonne del registro.} Trenta colonne per riga:
\begin{itemize}
  \item \textbf{precompilate} (da non toccare): ID, categoria, rischio, minaccia/causa, asset; risk owner,
    data obiettivo, ultima revisione;
  \item \textbf{da compilare} (13): probabilità e impatti C, I, A inerenti; controlli 8.x applicabili;
    misure; \textbf{giustificazione}; efficacia; probabilità e impatti residui; decisione di trattamento;
  \item \textbf{calcolate} (mai scriverci dentro): punteggi per dimensione, livello inerente e residuo,
    controllo ``entro la propensione?''.
\end{itemize}

\paragraph{Regole pratiche.} Lavoro in gruppo, 90 minuti, consegna del file completato con nome del
gruppo come prefisso del nome del file, caricamento nella cartella condivisa indicata. Un lavoro non
attribuibile (gruppo o partecipanti mancanti) non viene valutato. La valutazione premia il
\textbf{ragionamento}, non la coincidenza con un numero.

% ---------------------------------------------------------------------
\section{Metodologia}

\subsection{Tre dimensioni, due passaggi}
\begin{itemize}
  \item Ogni rischio si valuta \textbf{tre volte}, una per proprietà: \textbf{C} riservatezza (accesso o
    divulgazione non autorizzati), \textbf{I} integrità (alterazione, cancellazione, inesattezza, anche
    per errore), \textbf{A} disponibilità (indisponibilità oltre il fermo tollerabile).
  \item Per ogni dimensione: $\text{punteggio} = \text{Probabilità} \times \text{Impatto}$, entrambi su
    scala \textbf{1--3}. La probabilità è unica per la riga, gli impatti sono tre.
  \item Il livello complessivo è il \textbf{massimo} dei tre punteggi, \textbf{mai la media}.
  \item Si valuta due volte: \textbf{inerente} (con i controlli di oggi) e \textbf{residuo} (dopo le
    misure scelte). Il residuo si \textbf{rivaluta da capo}: non è ``inerente meno efficacia''.
\end{itemize}

\begin{center}
\small
\textbf{Ancore della scala 1--3 (dal foglio \emph{Instructions})}\par\smallskip
\begin{tabularx}{\linewidth}{c X X}
\toprule
\textbf{Valore} & \textbf{Probabilità} & \textbf{Impatto (C, I, A)} \\ \midrule
1 & Bassa: servono condizioni insolite o un attore motivato; nessun precedente nel settore negli ultimi
  due anni & Basso: limitato e reversibile; nessun obbligo di notifica; attività ripristinata in poche ore
  senza conseguenze per i clienti \\
2 & Media: plausibile entro l'anno (configurazioni diffuse, errore umano ricorrente, minaccia
  opportunistica) & Medio: uno--due giorni di interruzione, perdita parziale di dati recuperabili,
  possibile reclamo o segnalazione all'autorità \\
3 & Alta: evento atteso o già accaduto (minaccia automatizzata e diffusa, nessun controllo, punto
  ricorrente nei verbali) & Alto: perdita definitiva di dati dei clienti, \textbf{breach da notificare ex
  art.~33 GDPR}, danno reputazionale o contenzioso \\ \bottomrule
\end{tabularx}
\end{center}

\subsection{Propensione al rischio}
\begin{center}
\small
\begin{tabularx}{\linewidth}{l l X}
\toprule
\textbf{Punteggio} & \textbf{Livello} & \textbf{Regola} \\ \midrule
1--2 & Basso & accettabile; nessuna azione obbligatoria; monitoraggio annuale al riesame della direzione \\
3--4 & Medio & accettabile solo con trattamento documentato o accettazione formale del residuo da parte
  della direzione \\
6--9 & Alto & \textbf{non accettabile}: piano di trattamento con owner e scadenza; il residuo deve scendere
  a medio o meno \\ \bottomrule
\end{tabularx}
\end{center}
\noindent (Con due fattori 1--3 i prodotti possibili sono 1, 2, 3, 4, 6, 9: il valore 5 non esiste.) Il
playbook collega i concetti della lezione: la \emph{propensione} corrisponde alle soglie, la
\emph{tolleranza} all'owner e alla scadenza del singolo rischio.

\subsection{La decisione di trattamento}
La decisione dice \textbf{che cosa fa lo studio del rischio residuo}; non misura quanto è alto il
residuo. Un residuo medio può portare una qualsiasi delle quattro decisioni.
\begin{center}
\small
\begin{tabularx}{\linewidth}{l X}
\toprule
\textbf{Decisione} & \textbf{Significato} \\ \midrule
Ridotto (\emph{Reduced}) & uno o più controlli 8.x mitigano il rischio; il residuo è accettato dalla
  direzione. Corrisponde a ``mitigare'' nelle slide \\
Accettato (\emph{Accepted}) & \textbf{nessun ulteriore controllo}: il residuo coincide con l'inerente,
  già abbastanza basso, e altri controlli non sono giustificati dai costi. Accettazione formale e motivata \\
Trasferito (\emph{Transferred}) & la perdita \textbf{economica} passa a un terzo (assicurazione cyber,
  clausole contrattuali, SLA del fornitore); la \textbf{responsabilità} del trattamento resta allo studio \\
Evitato (\emph{Avoided}) & si abbandona (o non si avvia) l'attività che genera il rischio \\ \bottomrule
\end{tabularx}
\end{center}
Due conseguenze: un rischio \textbf{inerente alto non può mai essere ``Accettato''} (senza misure resterebbe
alto); una riga ``Trasferita'' può essere inerente alta, ma deve comunque citare controlli 8.x e avere
residuo medio o basso: il trasferimento copre ciò che i controlli non possono coprire, non li sostituisce.
Nel lavoro reale ogni accettazione di un residuo medio va firmata dalla direzione.

\subsection{Scegliere i controlli}
\begin{enumerate}
  \item Partire dalla \textbf{minaccia/causa} della riga e dagli asset, \textbf{mai dal catalogo}.
  \item Chiedersi quale dei 34 controlli \textbf{interrompe quella causa}: di solito due--quattro.
  \item Tenere l'\textbf{insieme minimo} sufficiente. Non si inventa mai un rischio per dare una ``casa''
    a un controllo: il registro va dal rischio al controllo, non viceversa (il documento guidato dai
    controlli è la SoA, fuori dall'esercizio).
\end{enumerate}
Nella colonna delle misure va un'\textbf{azione concreta} (``TLS obbligatorio su tutti gli scambi con i
clienti, cifratura completa dei dischi, backup cifrati con gestione documentata delle chiavi''), non il
nome del controllo (``8.24 uso della crittografia''). Nella giustificazione: minaccia affrontata,
controllo applicato, proporzionalità, perché il residuo è accettabile; e se la misura agisce sulla
\textbf{probabilità}, sull'\textbf{impatto} o su entrambi.

\paragraph{Applicabilità.} Dal foglio \emph{Annex A clause 8}:
\begin{itemize}
  \item \textbf{26 controlli} applicabili senza riserve;
  \item \textbf{8.4, 8.25, 8.28} (codice sorgente, sviluppo sicuro, secure coding) \textbf{non applicabili}:
    lo studio non sviluppa software. Citarli è un errore;
  \item \textbf{8.11, 8.14, 8.27, 8.30, 8.34} da valutare caso per caso.
\end{itemize}
Restano \textbf{31 controlli} di cui rendere conto: ciascuno deve essere citato in qualche riga oppure
registrato come ``non implicato'', con la motivazione, nel blocco in fondo al foglio dei controlli. Il
playbook segnala tre controlli che non derivano dai titoli degli scenari e vanno cercati nel testo:
\textbf{8.30} (sviluppo esternalizzato), \textbf{8.33} (dati di test), \textbf{8.34} (protezione dei
sistemi durante audit e test).

\subsection{Controlli di coerenza prima della consegna}
\begin{enumerate}[label=(\alph*)]
  \item nessun rischio residuo è alto;
  \item ogni rischio inerente alto ha una decisione con controlli, owner e scadenza (Ridotto o
    Trasferito; mai Accettato);
  \item tutti i 31 controlli applicabili sono citati o dichiarati non implicati (l'unica regola senza
    contatore automatico: va verificata a mano);
  \item ogni accettazione di un residuo medio è motivata e firmata.
\end{enumerate}
Inoltre: probabilità e impatti residui mai superiori agli inerenti; giustificazione presente su ogni riga;
intestazione del registro compilata. L'efficacia (bassa, media, alta) non entra in alcuna formula: è il
giudizio che rende credibile il residuo.

\paragraph{Taratura.} La soluzione di riferimento usata in aula valuta \textbf{17 scenari su 18 alti}
come rischio inerente e \textbf{nessuno alto} dopo il trattamento. Se il proprio profilo inerente è
quasi tutto medio, probabilmente si sta sottostimando la probabilità di uno studio senza controlli.

% ---------------------------------------------------------------------
\section{Svolgimento ragionato}

\subsection{R-01 -- Credenziali compromesse (esempio del playbook)}
\paragraph{Inerente.} Probabilità \textbf{3}: minaccia automatizzata e diffusa (credential stuffing,
phishing), niente MFA e credenziali condivise. Impatto C \textbf{3} (accesso ai dati fiscali dei clienti:
breach notificabile e contenzioso), I \textbf{2} (dati alterabili ma recuperabili da backup), A \textbf{1}
(lo studio continua a lavorare). Punteggi C/I/A = 9/6/3 $\Rightarrow$ \textbf{alto}.

\paragraph{Controlli e misure.} \textbf{8.2, 8.3, 8.5}: MFA su tutti i servizi cloud e sul gestionale,
eliminazione degli account condivisi, permessi per ruolo con minimo privilegio, password manager, revisione
trimestrale degli accessi. Efficacia alta.

\paragraph{Residuo.} Probabilità 1, impatti 1/1/1 $\Rightarrow$ \textbf{basso}. Il ragionamento copre
entrambi gli assi: l'MFA abbatte la \textbf{probabilità}; minimo privilegio e niente account condivisi
riducono l'\textbf{impatto}, perché limitano ciò che un singolo account compromesso può raggiungere. Chi
ritiene che l'accesso a dati dei clienti resti comunque notificabile può tenere C~=~2: con probabilità 1 il
livello resta basso (1~$\times$~2~=~2). Il numero cambia, la conclusione no.
\textbf{Decisione: Ridotto.}

\subsection{R-17 -- Portale indisponibile (esempio del playbook)}
\paragraph{Inerente.} Probabilità \textbf{3} (attacco volumetrico automatizzato, un solo hosting
provider, nessuna mitigazione). Impatto A \textbf{3} (un portale che i clienti non raggiungono è il caso
tipico di danno reputazionale), C \textbf{1} (un DDoS non divulga nulla), I \textbf{1}. Punteggi 3/3/9
$\Rightarrow$ \textbf{alto}, come R-01 ma per un motivo opposto.

\paragraph{Controlli e misure.} \textbf{8.21, 8.14, 8.6}: SLA con il provider con impegni di disponibilità
e penali, mitigazione DDoS lato provider, piano di ripiego documentato (pagina di stato, canale
alternativo verso i clienti), monitoraggio della disponibilità con alert. Efficacia media: le misure
riducono frequenza e durata dei fermi, ma non eliminano la minaccia, che sta fuori dal controllo dello
studio.

\paragraph{Residuo.} Probabilità 2, impatti 1/1/2 $\Rightarrow$ punteggi 2/2/4 $\Rightarrow$
\textbf{medio}. Una parte dell'esposizione è irriducibile e viene \textbf{spostata}: la perdita economica
residua va al provider (penali) e a una polizza cyber per interruzione dell'attività; allo studio restano
il danno reputazionale e la franchigia. \textbf{Decisione: Trasferito}, non perché il residuo sia alto ma
perché il trasferimento è ciò che chiude la riga.

\begin{center}
\small
\begin{tabular}{lcc}
\toprule
& \textbf{R-01} & \textbf{R-17} \\ \midrule
Punteggi inerenti C / I / A & 9 / 6 / 3 & 3 / 3 / 9 \\
Livello inerente & alto & alto \\
Controlli & 8.2, 8.3, 8.5 & 8.21, 8.14, 8.6 \\
Efficacia & alta & media \\
Livello residuo & basso & medio \\
Decisione & Ridotto & Trasferito \\ \bottomrule
\end{tabular}
\end{center}

\begin{esame}
Stesso livello inerente, ragioni diverse: il livello da solo non dice come trattare un rischio. Bisogna
sapere \textbf{quale dimensione} lo rende alto (qui riservatezza in un caso, disponibilità nell'altro) e
scegliere controlli che agiscono su quella.
\end{esame}

\begin{integrazione}[Altre righe ragionate (proposta non ufficiale)]
Valori coerenti con le ancore del file; sono una delle risposte difendibili, non ``la'' risposta.
\begin{itemize}
  \item \textbf{R-02 Ransomware.} Inerente: P 3; C 2 (l'attaccante può aver letto dati prima di
    cifrarli), I 2, A 3 (studio fermo) $\Rightarrow$ 6/6/9, alto. Controlli \textbf{8.7, 8.23, 8.13}: EDR
    con protezione comportamentale, blocco delle macro nei file da Internet, filtro web e DNS, backup
    immutabile fuori sede con test di ripristino. I primi due agiscono sulla probabilità (tre strati sulla
    catena di consegna), il backup è il controllo \emph{correttivo} che riduce l'impatto sulla
    disponibilità. Residuo: P 1; 2/2/2 $\Rightarrow$ basso. Ridotto.
  \item \textbf{R-06 Link pubblici.} Inerente: P 3 (impostazioni predefinite del cloud ed errore umano);
    C 3 (breach notificabile), I 1, A 1 $\Rightarrow$ 9, alto. Controlli \textbf{8.3, 8.9, 8.12}: un
    controllo per ciascuna causa (restrizione degli accessi con condivisioni solo ad ospiti autenticati e a
    scadenza; baseline sicura del tenant senza link ``chiunque''; etichette di classificazione con regole DLP).
    Residuo: P 1, C resta 3 (un documento condiviso con la persona sbagliata resta un breach): 3,
    \textbf{medio}. Ridotto, con accettazione formale del residuo medio.
  \item \textbf{R-12 Amministratore esterno.} Inerente: P 2; C 3, I 3, A 2 $\Rightarrow$ alto. Controlli
    \textbf{8.2, 8.15, 8.18}: account nominativi, elevazione temporanea su richiesta, strumenti di
    amministrazione remota limitati e tracciati, log centralizzati non modificabili dal fornitore. Le
    clausole contrattuali e gli audit sul fornitore sono controlli organizzativi: vanno nell'altro registro.
    Residuo medio (C resta 3: un amministratore che aggira i controlli vede comunque tutto). Ridotto.
  \item \textbf{R-18 Metadati.} È l'unico scenario plausibilmente \emph{medio} già come inerente: P 2
    (errore ricorrente, ma i metadati rivelano qualcosa di rilevante solo a volte), C 2 (contesto interno,
    non una divulgazione massiva di dati personali), I 1, A 1 $\Rightarrow$ 4. È il caso tipico di
    \textbf{Accettato}: nessuna misura, residuo = inerente, accettazione motivata e firmata; si cita
    comunque \textbf{8.12} come controllo che lo tratterebbe (rimozione automatica dei metadati in uscita)
    se la valutazione cambiasse.
\end{itemize}
\end{integrazione}

\begin{integrazione}[Una mappatura possibile rischi--controlli]
\small
\begin{tabularx}{\linewidth}{l X l X}
\toprule
\textbf{ID} & \textbf{Controlli} & \textbf{ID} & \textbf{Controlli} \\ \midrule
R-01 & 8.2, 8.3, 8.5 & R-10 & 8.24, 8.12 \\
R-02 & 8.7, 8.23, 8.13 & R-11 & 8.1, 8.24 \\
R-03 & 8.1, 8.15 & R-12 & 8.2, 8.15, 8.18 \\
R-04 & 8.13, 8.24 & R-13 & 8.26, 8.27, 8.29, 8.30 \\
R-05 & 8.15, 8.16, 8.17 & R-14 & 8.32, 8.31, 8.19, 8.9 \\
R-06 & 8.3, 8.9, 8.12 & R-15 & 8.6, 8.14 \\
R-07 & 8.8 & R-16 & 8.10, 8.24 \\
R-08 & 8.9, 8.21 & R-17 & 8.21, 8.14, 8.6 \\
R-09 & 8.20, 8.22 & R-18 & 8.12 \\ \bottomrule
\end{tabularx}
\medskip

\noindent Con questa mappatura restano da dichiarare \emph{non implicati}, con motivazione: \textbf{8.11}
(nessuno scenario usa dati reali in test o in report verso terzi), \textbf{8.33} (nessuna minaccia riguarda
la scelta o protezione dei dati di test: in R-14 l'ambiente di test non esiste ancora) e \textbf{8.34}
(nessuno scenario riguarda audit o test su sistemi operativi; le scansioni di R-07 e i penetration test di
R-13 sono misure, da pianificare fuori orario o in staging). L'8.30 trova invece posto in R-13, dove il
portale è ``commissionato'' a un fornitore.
\end{integrazione}

% ---------------------------------------------------------------------
\section{Errori tipici}
\begin{enumerate}
  \item \textbf{Fare la media} di C, I, A invece di prendere il massimo (C~=~9, I~=~1, A~=~1 è alto).
  \item Usare \textbf{lo stesso impatto per tutte e tre le dimensioni}: quasi sempre sbagliato, rispondono
    a domande diverse.
  \item Scrivere il \textbf{nome del controllo} al posto della misura concreta.
  \item Scrivere nelle colonne con formule o fuori dai menu a tendina.
  \item Lasciare un \textbf{residuo alto} non trattato.
  \item Citare controlli \textbf{non applicabili} (8.4, 8.25, 8.28).
  \item \textbf{Sottostimare la probabilità} inerente di uno studio senza controlli (troppi rischi medi).
  \item Residuo uguale all'inerente \textbf{senza spiegazione}; o residuo \emph{superiore} all'inerente.
    Se un controllo agisce solo sulla probabilità, dirlo: i controlli detective rilevano, non prevengono.
  \item ``Accettare'' un rischio inerente alto.
  \item Dimenticare l'\textbf{accountability}: il trasferimento sposta il rischio economico, non l'obbligo.
  \item Ignorare la dimensione dei \textbf{dati personali}: un impatto C~=~3 va ricondotto a un breach
    notificabile ex art.~33 GDPR nella giustificazione.
  \item Dimenticare di rendere conto dei controlli che nessuna riga cita (regola~(c)).
\end{enumerate}

% ---------------------------------------------------------------------
\section{Collegamenti}
\begin{itemize}
  \item \textbf{ISO/IEC 27001}: il registro attua le clausole \textbf{6.1.2} (valutazione del rischio),
    \textbf{6.1.3} (trattamento) e \textbf{8.2/8.3} (esecuzione e documentazione) ed è la fonte principale
    della \textbf{Statement of Applicability}. Teoria nel Capitolo~\ref{cap:iso}; clausole e controlli
    nelle lezioni del 19/10 e del 23/10 (Capitoli~\ref{lez:2026-10-19} e~\ref{lez:2026-10-23}); risk
    register, opzioni di trattamento e SoA nella lezione del 26/10 (Capitolo~\ref{lez:2026-10-26}).
  \item \textbf{ENISA}: la metodologia per i dati personali valuta l'impatto dal punto di vista
    dell'interessato e, come qui, prende il valore più alto tra C, I, A; utile per controllare ogni impatto
    C~=~3. Da non confondere con la metodologia di severità dei breach (Capitolo~\ref{cap:data-breach}).
  \item \textbf{GDPR}: art.~32 (misure adeguate al rischio) e art.~33 (notifica) come ancora dell'impatto
    massimo.
  \item \textbf{Esame}: il docente ha indicato tra le ``regole del gioco'' il saper distinguere
    \textbf{rischio inerente e misure di mitigazione} (lezione del 09/11, Capitolo~\ref{lez:2026-11-09}).
\end{itemize}

% ---------------------------------------------------------------------
\section{Punti da verificare}
\begin{verifica}
\begin{itemize}
  \item La soluzione di riferimento del docente non è stata distribuita: si conoscono solo i due esempi del
    playbook (R-01, R-17) e la taratura ``17 alti su 18, nessuno alto dopo il trattamento''.
  \item Il file indica come scadenza di R-17 ``Q3 2026'', anteriore alla data di revisione del registro
    (19/10/2026): probabilmente un refuso del materiale.
  \item Gli esiti e i criteri di valutazione dei gruppi non sono riportati qui per scelta.
\end{itemize}
\end{verifica}

% =====================================================================
\chapter{Lab 4 -- Conformity audit CRA--GDPR: Sestante Mobility}
\label{lab:04}
% =====================================================================

\begin{nota}[Fonte]
Laboratorio del 06/11/2026 (Capitolo~\ref{lez:2026-11-06}), presentato a fine lezione il 02/11
(Capitolo~\ref{lez:2026-11-02}). Fonti: il \emph{Conformity Audit Playbook} (guida dell'auditor
esterno, con evidence room, interviste, checklist e template) e le slide \emph{Conformity Audit
Exercise -- Briefing}. Per i richiami teorici: slide \emph{05.~PoliTO\_CRA}. Il caso è inventato,
mentre il quadro normativo e le date sono reali; la data fittizia dell'audit è gennaio 2028.
Lo svolgimento ragionato non è una soluzione ufficiale: è nei riquadri \emph{Integrazione}.
\end{nota}

% ---------------------------------------------------------------------
\section{Dall'incidente all'audit}

Nella simulazione Pulse K (Sezione~\ref{sec:pulsek}) gli studenti erano il team di crisi interno al
fabbricante e correvano contro gli orologi dell'art.~14 CRA e degli artt.~33--34 GDPR. In questo
laboratorio il ruolo si rovescia: si è un \textbf{auditor esterno indipendente}, nominato dal consiglio
di amministrazione. Non c'è un orologio che corre: gli incidenti del 2027 sono chiusi (male) e il
compito è ricostruire cosa dicono dell'organizzazione e se questa può legittimamente
\textbf{dichiarare la conformità} nel 2028.

\begin{center}
\small
\begin{tabularx}{\linewidth}{l X X}
\toprule
 & \textbf{Simulazione di incidente} & \textbf{Audit esterno (questo lab)} \\ \midrule
Posizione   & team di crisi del fabbricante & auditor indipendente nominato dal CdA \\
Materiale   & \emph{inject} che arrivano in sequenza & 16 documenti e due interviste, quattro anni di storia \\
Nucleo giuridico & art.~14 CRA, artt.~33--34 GDPR & Allegati I, VII, VIII CRA; artt.~25, 28, 32, 33(5), 35, 44 GDPR \\
Prodotto finale  & decision sheet con le scadenze & registro dei finding, opinione, management letter \\
Errore tipico & mancare una scadenza & finding non supportato dalle evidenze, o opinione contraddetta dal registro \\
\bottomrule
\end{tabularx}
\end{center}

% ---------------------------------------------------------------------
\section{Il caso Sestante Mobility}

\subsection{L'azienda e il prodotto}
\begin{itemize}
  \item \textbf{Sestante Mobility S.p.A.}, Bologna (stabilimento principale in Italia): 120 dipendenti,
    fatturato 2027 di 31 milioni di euro. È una \textbf{media impresa}, quindi non beneficia delle
    agevolazioni sanzionatorie previste per le piccole imprese (art.~64(10) CRA).
  \item Circa 1.100 clienti \emph{fleet operator} (PMI) in Italia, Spagna e Polonia; un distributore,
    Iberia Telematics (Madrid), che rivende, installa e gestisce le sostituzioni in garanzia.
  \item \textbf{46.000 box installati}: Italia 27.000, Spagna 12.000 (di cui circa 1.500 immessi dal
    12/12/2027), Polonia 7.000.
\end{itemize}

\paragraph{Tre componenti, un solo prodotto (il B2).}
\begin{itemize}
  \item \textbf{box B2}: si collega alla porta OBD-II del veicolo, ha GNSS e modem propri, firmware
    aggiornato \emph{over the air} (OTA), invia la posizione ogni 30 secondi;
  \item \textbf{Sestante Cloud}: gestito da Sestante su server UE in affitto (Francoforte);
  \item \textbf{Fleet app}: mappa dei veicoli, geofence, report e \emph{driver score}.
\end{itemize}
Dati in gioco: posizione e storico dei viaggi (90 giorni), eventi di guida e orari di lavoro
(24 mesi), driver score condiviso di default con l'HR del cliente, contatti aziendali e fatturazione.
Il conducente è un dipendente con nome e cognome: si tratta di \textbf{dati personali}, non di
``dati del veicolo'', qualunque cosa sostenga il CTO nell'intervista.

\paragraph{Fornitori.}
\begin{itemize}
  \item \textbf{Vistula EMS} (Katowice) costruisce il box su progetto Sestante, senza clausole di
    sicurezza, elenco dei componenti né diritto di audit;
  \item \textbf{Updato 2.3.1}, libreria di aggiornamento di terze parti, bloccata a quella versione dal
    2022: supporta la verifica della firma, ma la funzione è disattivata nella configurazione di build;
  \item \textbf{MetricLoop Inc.} (USA) svolge analytics, scoring dei conducenti e un dataset di
    \emph{benchmarking} venduto in forma aggregata ad assicuratori: clausole contrattuali standard
    firmate, nessuna \emph{transfer impact assessment}.
\end{itemize}
Sestante appone il proprio marchio sul box: è il \textbf{fabbricante}, anche per i componenti
acquistati e per il codice open source che distribuisce.

\subsection{L'incarico e le domande del CdA}
Il comitato di audit nomina gli auditor il 12/01/2028; settimana di fieldwork 24--28/01/2028; report
al CdA il 22/02/2028. Tre eventi del 2027 hanno spinto il comitato ad agire:
\begin{itemize}
  \item \textbf{agosto 2027}: un ricercatore segnala che i box accettano un pacchetto firmware non
    firmato; il PSIRT chiude la segnalazione in quattro giorni (``serve l'accesso fisico, non è una
    vulnerabilità''), senza fix né voce di backlog;
  \item \textbf{settembre 2027}: un box reso in garanzia viene ricondizionato e rivenduto, con i dati
    della flotta precedente ancora raggiungibili;
  \item \textbf{novembre 2027}: un conducente scopre di essere tracciato anche fuori orario, con un
    punteggio di guida che arriva alla dashboard HR del datore di lavoro.
\end{itemize}
Nessuno dei tre è un incidente aperto. La domanda non è cosa si doveva fare nel 2027, ma
\textbf{che cosa l'azienda ha diritto di dichiarare di sé oggi}. Le domande sono tre:
\begin{enumerate}
  \item il fascicolo consente di tenere il B2 sul mercato, in particolare le circa 8.000 unità
    immesse dal 12/12/2027?
  \item che cosa si può confermare a \textbf{Iberia Lease}, il principale cliente spagnolo, che
    rinnoverà 12.000 box entro il \textbf{24/04/2028} solo se c'è evidenza di conformità al CRA?
  \item il \textbf{T3}, successore del B2, sarà immesso sul mercato il \textbf{15/06/2028}, ma il suo
    fascicolo non è nella evidence room: va dichiarata la \textbf{limitazione di scope} e va detto che
    cosa deve esistere prima che qualcuno firmi una dichiarazione.
\end{enumerate}

\subsection{Perimetro e regole del gioco}
\begin{itemize}
  \item Nel perimetro: Sestante come fabbricante e, per alcuni flussi, come titolare o responsabile
    del trattamento; box, cloud e app; i fatti del 2027.
  \item Fuori perimetro: la condotta dei fleet operator verso i propri conducenti e le norme nazionali
    di diritto del lavoro; il trattamento proprio del distributore; il T3 come prodotto.
  \item La evidence room (E-01\ldots E-16) è l'unica fonte di fatti: niente visite, accessi ai sistemi
    o contatti esterni. \textbf{Senza riferimento a un'evidenza, nessun finding.}
  \item I finding si scrivono contro l'azienda e il suo prodotto, mai contro i clienti.
  \item L'auditor tratta a sua volta dati personali: campiona, non copia, chiede il minimo.
  \item Dove la legge ammette due risposte difendibili si sceglie e si motiva per iscritto
    (\emph{decide and defend}).
\end{itemize}

\subsection{La evidence room}
{\small
\begin{longtable}{@{}l p{4.2cm} p{8.3cm}@{}}
\toprule
\textbf{Rif.} & \textbf{Documento} & \textbf{Che cosa contiene (in sintesi)} \\ \midrule
\endhead
E-01 & Indice della documentazione tecnica e bozza di dichiarazione & 11 dei 18 documenti del
  fascicolo; la bozza di dichiarazione per il T3 indica il modulo~A e non cita norme, specifiche o
  report di test. \\
E-02 & Risk assessment (rev.~2, ottobre 2025) & Afferma che gli aggiornamenti sono firmati e
  verificati dal boot loader; raccomanda un interruttore ``viaggio privato'' in una release futura. \\
E-03 & Periodo di supporto, sito e app store & La policy non indica alcun numero; il sito promette
  ``cinque anni di aggiornamenti di sicurezza''; l'app store non riporta né periodo di supporto né
  contatto per le vulnerabilità. \\
E-04 & Thread con il ricercatore (agosto 2027) & Firmware arbitrario installato, nessun controllo di
  firma né protezione dal rollback; il PSIRT risponde ``CVSS 4.2, non è una vulnerabilità''. \\
E-05 & Pagina ``Security'' del sito & Unico contatto pubblicato: l'indirizzo del supporto; l'indirizzo
  PSIRT esiste ma non è pubblicato. \\
E-06 & Wiki di supporto (2025) & Credenziale di manutenzione remota identica su tutti i box,
  riportata in chiaro. \\
E-07 & Configurazione di build, firmware 4.x & \texttt{OTA\_VERIFY\_SIGNATURE=false},
  \texttt{OTA\_ALLOW\_ROLLBACK=true}, \texttt{MAINT\_SHELL\_ENABLED=true},
  \texttt{LOG\_LEVEL=debug}. \\
E-08 & Foglio componenti e nota sulle dipendenze & 8 moduli su 11, senza versione, licenza né
  colonna CVE; Updato 2.3.1 ha una CVE corretta nella release successiva, aggiornamento rinviato. \\
E-09 & Penetration test (giugno 2026) & 27 finding, 9 ancora aperti: console admin senza MFA,
  endpoint di debug in produzione, token di sessione validi 30 giorni, due API key nel codice
  dell'app. \\
E-10 & Impostazioni di default e pack di onboarding & Posizione ogni 30~s anche fuori orario;
  modalità viaggio privato assente; score condiviso con l'HR di default; nessun modello di
  informativa per i conducenti. \\
E-11 & Informativa privacy e DPIA (2025) & DPIA centrata su app e cloud, telemetria del box liquidata
  come rischio ``basso''; nessun parere del DPO; l'informativa non cita USA né benchmarking. \\
E-12 & Template di accordo con i (sub)responsabili e relativo elenco & Ruoli invertiti (Sestante
  ``titolare'', l'operatore ``responsabile''); breach comunicato all'operatore ``entro 30 giorni'';
  elenco fornitori non aggiornato. \\
E-13 & Configurazione cloud e accessi & AES-256 a riposo con KMS, TLS~1.3 e HSTS; ma MFA disattivata
  sulla console admin, 14 account con accesso al DB di produzione, sessioni di 30 giorni, nessuna
  registrazione delle sessioni privilegiate; allarmi disattivati per un cliente. \\
E-14 & Procedura RMA e scambio successivo (settembre 2027) & Il distributore azzera il box ma
  l'account cloud precedente resta attivo sei giorni: 340 conducenti, 90 giorni di viaggi. Due
  ingegneri chiudono il caso: ``rischio basso, veicoli non persone'', nessuna notifica. \\
E-15 & Contratto MetricLoop & Clausole standard (moduli 2021) senza transfer impact assessment;
  scoring dei conducenti e dataset di benchmarking venduto ad assicuratori. \\
E-16 & Certificato ISO/IEC 27001 & Valido 2026--2029, ma lo scope copre solo la piattaforma cloud:
  sviluppo firmware e supporto sul campo sono esclusi. \\
\bottomrule
\end{longtable}}

\noindent Il playbook suggerisce di partire da \textbf{E-04, E-06, E-07, E-10, E-12, E-14}: contengono
i finding attesi, mentre il resto della room spiega perché nessuno se n'è accorto. Le due interviste
(CTO e DPO esterno) vanno lette come \textbf{evidenza debole}, cioè come prova di ciò che una persona
crede, da confrontare con i documenti. Esempi di contraddizioni:
\begin{itemize}
  \item il CTO dice che le password di manutenzione ``ruotano ogni trimestre'', ma la wiki (E-06) è
    ferma al 2025 e la credenziale è unica;
  \item il CTO dice che ``il distributore resetta il box'', ma la procedura (E-14) azzera il
    dispositivo e non l'account cloud;
  \item il DPO dichiara di aver chiesto nel 2026 di rifare la DPIA dopo l'aggiunta della dashboard HR,
    richiesta chiusa dal product manager; di essere stato informato del caso RMA dall'operatore e non
    dall'azienda; di lavorare quattro ore al mese e di non essere nel percorso di escalation degli
    incidenti.
\end{itemize}

\begin{esame}
\textbf{Il brief è incompleto.} La lettera d'incarico cita tre eventi, ma la room ne contiene di più:
ciò che conta spesso non è mai stato segnalato né verbalizzato (per esempio la credenziale di
manutenzione condivisa di E-06). Per ogni documento la domanda è: \emph{contraddice qualcosa che ho già
letto?} Una fonte è un'asserzione, due fonti concordanti sono un fatto.
\end{esame}

% ---------------------------------------------------------------------
\section{Consegna}

Un file di testo per gruppo, prodotto dal workbench online e organizzato in tre blocchi.
\begin{description}[style=nextline,leftmargin=1.5em]
  \item[A -- Risposte di scope] Quali obblighi CRA vincolano ciascuna delle tre popolazioni (B2 fino
    all'11/12/2027, circa 38.000 unità; B2 dal 12/12/2027, circa 8.000 unità; T3) e che ruolo GDPR ha
    Sestante per tre flussi: (a)~posizione e viaggi dei veicoli del cliente; (b)~supporto, garanzia e
    RMA; (c)~account e fatturazione dei clienti.
  \item[B -- Registro dei finding] Almeno quattro finding di livello Medium o superiore, ciascuno con
    condizione, criterio, riferimento all'evidenza, causa, effetto, severità, raccomandazione,
    responsabile e data. In più: \textbf{un controllo che funziona} e almeno un criterio esaminato e
    scartato (\emph{Tested, no finding}) con il motivo.
  \item[C -- Opinione e management letter] Tipo di opinione con le limitazioni, decisione sul breach
    di settembre 2027, tre azioni prioritarie con responsabile e data, almeno un'azione formativa
    misurabile per ogni errore umano.
  \item[Bonus] Per ogni finding, il controllo più vicino dell'Annex~A.8 della ISO/IEC~27001, oppure
    ``nessun controllo A.8 è applicabile'' con il motivo.
\end{description}

\noindent\textbf{Criteri di valutazione} (slide di briefing): accuratezza giuridica 30\%; evidenze e
finding 25\%; severità 15\%; opinione e redazione 15\%; remediation 15\%. Si valuta il ragionamento e
l'evidenza citata, non la corrispondenza con un numero.

% ---------------------------------------------------------------------
\section{Metodologia dell'audit esterno}

\subsection{Le cinque fasi}
\begin{enumerate}
  \item \textbf{Plan}: scope, criteri, periodo, popolazione, rischi; che cosa potrebbe rendere
    sbagliata la risposta?
  \item \textbf{Request}: chiedere le evidenze che proverebbero o smentirebbero ogni rischio (qui la
    room è fissa: si valuta se basterebbe a una richiesta reale).
  \item \textbf{Walk through}: seguire una transazione dall'inizio alla fine (qui: un box reso in
    garanzia, E-14 contro E-01).
  \item \textbf{Test}: rieseguire, ispezionare, osservare; confrontare le configurazioni (E-07, E-10,
    E-13) con quanto affermato in E-02 ed E-03.
  \item \textbf{Report}: registro, opinione con limitazioni, management letter.
\end{enumerate}
Come riferimenti di metodo il playbook cita ISO~19011 e ISAE~3000. L'auditor non certifica nulla: dà
al CdA una \emph{ragionevole garanzia} rispetto a criteri definiti.

\subsection{I criteri: tre famiglie da non mescolare}
\begin{itemize}
  \item \textbf{diritto UE, lato prodotto}: CRA (ambito, obblighi del fabbricante, Allegato~I,
    documentazione tecnica, valutazione della conformità, segnalazione, transizione, sanzioni). Si
    cita prima l'articolo e poi l'allegato, dicendo a quali unità e a quale versione si riferisce;
  \item \textbf{diritto UE, lato dati}: GDPR. Si cita il paragrafo quando conta: l'art.~33(5)
    (registro delle violazioni) non è l'art.~33(1) (notifica);
  \item \textbf{standard volontari e policy interne}: ISO/IEC~27001 e procedure dell'azienda. Uno
    standard non è legge: va usato come benchmark o come criterio interno a cui l'azienda si è
    impegnata, dicendo quale dei due.
\end{itemize}

\begin{esame}
\textbf{Tre avvertenze del playbook.}
\begin{enumerate}
  \item \textbf{Un certificato non è conformità.} Il certificato ISO/IEC~27001 di Sestante (E-16)
    copre solo il cloud, non il firmware, e nessuna norma del CRA lo richiede: è un'affermazione su
    un sistema di gestione, non un criterio giuridico.
  \item \textbf{L'assenza di evidenza è un finding.} ``Non evidenziato'' è una condizione legittima,
    da valutare.
  \item Dove la risposta dipende dallo stato del diritto (per esempio la voce dell'Allegato~III o le
    norme armonizzate effettivamente pubblicate), si decide, si motiva e si indica la verifica che si
    farebbe in un incarico reale.
\end{enumerate}
\end{esame}

\subsection{Le evidenze}
\begin{center}
\small
\begin{tabularx}{\linewidth}{l l X}
\toprule
\textbf{Forza} & \textbf{Tipo} & \textbf{Esempi e limiti} \\ \midrule
Massima & Configurazioni e record di sistema & E-07, E-10, E-13: mostrano lo stato in un dato
  momento, non l'intenzione \\
Alta & Documenti scritti contemporanei ai fatti & E-04, E-14: mostrano ciò che si è scritto, non ciò
  che si è fatto \\
Media & Policy, procedure, contratti & E-03, E-12, procedura RMA: riguardano il disegno, non
  l'efficacia operativa \\
Bassa & Dichiarazioni e interviste & CTO e DPO: prova di ciò che si crede, da corroborare \\
\bottomrule
\end{tabularx}
\end{center}
\begin{itemize}
  \item Un conflitto non si risolve mai in silenzio: si mostrano i due documenti e si dice quale
    descrive lo stato del sistema (un flag di build prevale su una policy; un'email dell'epoca prevale
    su un ricordo).
  \item \textbf{Si dichiara sempre la popolazione}: 46.000 box non sono lo stesso finding di 8.000
    box immessi dopo l'11/12/2027. Un template firmato da tutti i clienti (E-12) è un difetto
    sistemico anche se lo si esamina una sola volta.
  \item Ciò che non si è potuto testare va scritto come limitazione.
\end{itemize}

\subsection{Il finding e la sua severità}
Un finding ha cinque parti: \textbf{condizione} (i fatti, con riferimento E-xx), \textbf{criterio}
(uno solo, citato con precisione), \textbf{causa} (il controllo mancante o fallito, non un generico
``errore umano''), \textbf{effetto} (esposizione per persone, clienti, autorità, accesso al mercato),
\textbf{raccomandazione} (azione, responsabile, data).

La severità di audit non si sceglie: si calcola come \textbf{esposizione $\times$ ampiezza}.
\begin{center}
\small
\begin{tabularx}{\linewidth}{c X X}
\toprule
\textbf{Valore} & \textbf{Esposizione} & \textbf{Ampiezza} \\ \midrule
1 & lacuna documentale: il requisito è rispettato in pratica, ma la prova manca &
  popolazione limitata (un'unità, una persona, un fornitore, un processo) \\
2 & requisito non rispettato, senza misure compensative & popolazione significativa (un reparto, un
  segmento di clienti, i conducenti di un operatore) \\
3 & norma violata, o dichiarazione di conformità non firmabile & tutta la flotta o tutti i clienti \\
\bottomrule
\end{tabularx}
\end{center}
Prodotto: \textbf{1--2 Low; 3 Medium; 4--6 High; 8--9 Critical}. Un Critical aperto rende impossibile
un'opinione senza rilievi; i High portano a un'opinione con rilievi e condizioni.

\begin{esame}
\textbf{Due scale, mai mescolate.} La severità di audit valuta un finding per il CdA. La severità
ENISA ($SE = DPC \times EI + CB$, Capitolo~\ref{cap:schede}) misura il rischio di un data breach per
le persone e serve alla decisione di notifica. Nel caso Sestante l'ENISA si usa solo per verificare
se la conclusione ``rischio basso'' sul breach di settembre 2027 era difendibile; il finding sul
processo che l'ha prodotta si valuta con la scala di audit. Mescolarle è ``l'errore classico'' del
laboratorio.
\end{esame}

\subsection{Il controllo che funziona}
Un controllo non è il requisito: è la misura con cui l'organizzazione lo soddisfa (tecnica,
organizzativa o contrattuale). Per dichiararlo funzionante il playbook propone sei domande: esiste
(è progettato e scritto)? è implementato dove sta il rischio? opera (ci sono evidenze nel periodo)?
è efficace e non contraddetto da altre evidenze? ha un responsabile e un'escalation che funziona?
copre tutta la popolazione in scope? Se una delle sei fallisce si ha un finding, e la condizione
mancata ne è la causa. Un controllo che funziona non riceve un punteggio: va riportato perché
un'opinione che elenca solo carenze non è credibile.

\subsection{L'opinione}
\begin{center}
\small
\begin{tabularx}{\linewidth}{l X X}
\toprule
\textbf{Opinione} & \textbf{Quando} & \textbf{Nel caso Sestante (playbook)} \\ \midrule
\emph{Unqualified} (senza rilievi) & criteri soddisfatti; finding Low o Medium con piani accettati &
  solo se il registro non ha Critical o High aperti \\
\emph{Qualified} (con rilievi) & criteri soddisfatti salvo questioni specifiche e identificate &
  probabile approdo: il fascicolo è riparabile, ma non così com'è \\
\emph{Adverse} (negativa) & criteri non soddisfatti; scostamento rilevante e pervasivo & difendibile
  sul lato prodotto, se il registro è pesante \\
\emph{Disclaimer} (impossibilità di esprimersi) & evidenze insufficienti & possibile per le unità
  immesse prima dell'11/12/2027, dove il perimetro è più sottile \\
\bottomrule
\end{tabularx}
\end{center}
\noindent ``Dipende'' non è un'opinione, e non lo è nemmeno un'opinione contraddetta dal proprio
registro. L'opinione deve contenere: la risposta alle due domande del CdA, criteri e fonti, perimetro
e limitazioni, il T3 come limitazione dichiarata, periodo e popolazione per categoria di unità, i
finding che la motivano, le azioni richieste al CdA e la decisione sul breach di settembre 2027.

% ---------------------------------------------------------------------
\section{Il CRA applicato all'audit}

\subsection{Ambito e ruoli}
\begin{itemize}
  \item Il B2 è un \textbf{prodotto con elementi digitali}: hardware con software, connesso alla rete
    cellulare e al veicolo. Il cloud rientra come \textbf{soluzione di elaborazione dati da remoto}
    del prodotto (artt.~2 e~3(1)--(2)), come già visto per Pulse Cloud.
  \item \textbf{Fabbricante} è Sestante, perché commercializza il box con il proprio marchio
    (art.~3(13); il playbook richiama anche l'art.~21). Vistula EMS produce su commessa; Iberia
    Telematics è un \textbf{distributore} (artt.~19--20), con obblighi propri di verifica, ma i
    finding si scrivono contro il fabbricante.
  \item Il fabbricante risponde anche dei componenti di terze parti (dovere di diligenza
    sull'integrazione dei componenti, slide CRA ``Security by design \& by default'').
\end{itemize}

\subsection{Tempistica: la transizione vale per unità}
Il CRA applica gli obblighi di segnalazione (art.~14) dall'\textbf{11/09/2026} a tutti i prodotti,
anche quelli già sul mercato; tutto il resto (Allegato~I, art.~13, artt.~19--21) dall'\textbf{11/12/2027}
(artt.~69(2)--(3) e~71(2), secondo il playbook). La data si applica alla \textbf{singola unità}
immessa sul mercato, non al modello.

\begin{center}
\small
\begin{tabularx}{\linewidth}{l X}
\toprule
\textbf{Popolazione} & \textbf{Obblighi CRA (lettura del materiale)} \\ \midrule
B2 immessi fino all'11/12/2027 (circa 38.000) & solo segnalazione ex art.~14 (dall'11/09/2026);
  l'Allegato~I non si applica a queste unità \\
B2 immessi dal 12/12/2027 (circa 8.000) & regime completo: Allegato~I Parti I e II, art.~13,
  documentazione tecnica, valutazione della conformità, dichiarazione UE, marcatura CE, art.~14 \\
T3 (dal 15/06/2028) & regime completo; nessuna opinione possibile perché il fascicolo non è nella
  room: è una limitazione di scope \\
\bottomrule
\end{tabularx}
\end{center}
\noindent Conseguenza pratica: la flotta di Iberia Lease è \textbf{mista} (circa 1.500 box spagnoli
sono stati immessi dopo il 12/12/2027). Un finding sul fascicolo deve dire a quale popolazione si
riferisce: ``l'azienda non è conforme al CRA'', senza indicare le unità, è la frase che fa contestare
un audit.

\begin{integrazione}[Prodotti già sul mercato e modifiche sostanziali]
Nel testo del regolamento i prodotti immessi prima dell'11/12/2027 ricadono negli obblighi del CRA
solo se, da quella data, subiscono una \emph{modifica sostanziale}. Un aggiornamento firmware che
cambia in modo rilevante la funzione o il profilo di rischio potrebbe quindi trascinare anche le
unità ``vecchie'' nel regime completo. È un punto da tenere presente quando si raccomanda di
correggere il firmware via OTA su tutta la flotta.
\end{integrazione}

\subsection{Requisiti essenziali: Allegato~I, Parte~I}
Le slide del corso elencano i requisiti di prodotto; ecco quelli che il caso permette di testare.
\begin{center}
\small
\begin{tabularx}{\linewidth}{p{5.2cm} X}
\toprule
\textbf{Requisito (Allegato I, Parte I)} & \textbf{Che cosa mostra la room} \\ \midrule
nessuna vulnerabilità sfruttabile nota all'immissione & firmware senza verifica della firma, già
  dimostrato da un ricercatore (E-04, E-07); CVE nota in Updato (E-08) \\
configurazione \emph{secure by default} & shell di manutenzione attiva e log in debug (E-07);
  tracciamento fuori orario e score all'HR attivi di default (E-10) \\
vulnerabilità correggibili con aggiornamenti di sicurezza & l'OTA esiste, ma accetta pacchetti non
  firmati e consente il rollback (E-07) \\
protezione da accessi non autorizzati & credenziale identica su tutti i box (E-06); console admin
  senza MFA (E-09, E-13) \\
riservatezza e integrità dei dati & cifratura a riposo e in transito nel cloud (E-13): è il
  controllo che funziona \\
minimizzazione dei dati & posizione ogni 30~s anche fuori orario, nessuna modalità privata (E-10) \\
riduzione della superficie d'attacco & endpoint di debug in produzione (E-09), shell di manutenzione
  (E-07) \\
registrazione e monitoraggio delle attività & log di 90 giorni, nessuna registrazione delle sessioni
  privilegiate, allarmi disattivati per un cliente (E-13) \\
cancellazione sicura e permanente di dati e impostazioni & il ricondizionamento azzera il box ma non
  l'account cloud (E-14) \\
\bottomrule
\end{tabularx}
\end{center}

\subsection{Gestione delle vulnerabilità: Allegato~I, Parte~II}
\begin{itemize}
  \item \textbf{identificare e documentare componenti e vulnerabilità}: il foglio componenti (E-08)
    copre 8 moduli su 11, senza versioni né CVE; non c'è una SBOM;
  \item \textbf{correggere senza ritardo}: l'aggiornamento di Updato è rinviato ``a quando si
    toccherà la pipeline'' (E-08); 9 finding del pentest restano aperti (E-09);
  \item \textbf{test regolari}: il pentest esiste, ma senza seguito;
  \item \textbf{policy di divulgazione coordinata} e \textbf{canale di segnalazione}: il PSIRT ha
    respinto la segnalazione senza triage motivato (E-04) e il suo indirizzo non è pubblicato (E-05);
    il playbook collega questo punto anche al punto di contatto unico dell'art.~13(17);
  \item \textbf{distribuzione sicura degli aggiornamenti}: firma disattivata nella build (E-07),
    contro quanto afferma il risk assessment (E-02);
  \item \textbf{aggiornamenti per il periodo di supporto}: la policy non fissa la durata, il sito
    promette cinque anni, l'app store tace (E-03).
\end{itemize}

\begin{integrazione}[SBOM e periodo di supporto nel testo del CRA]
L'Allegato~I, Parte~II chiede di documentare i componenti anche redigendo una \emph{software bill of
materials} in formato leggibile da macchina, almeno per le dipendenze di primo livello. Il periodo di
supporto (art.~13) deve riflettere il tempo d'uso atteso del prodotto ed essere, di regola, di almeno
cinque anni, salvo prodotti destinati a un uso più breve; va indicato al momento dell'acquisto. Nel
caso Sestante la promessa commerciale di cinque anni (E-03) non è sostenuta da alcuna decisione
documentata.
\end{integrazione}

\subsection{Documentazione tecnica e valutazione della conformità}
\begin{itemize}
  \item La \textbf{documentazione tecnica} (Allegato~VII) è l'oggetto stesso dell'audit: l'indice
    (E-01) elenca 11 dei 18 documenti previsti e il risk assessment (E-02) è fermo al 2025, prima
    dei cambiamenti di architettura. Il risk assessment va tenuto aggiornato per la versione immessa
    sul mercato.
  \item \textbf{Classificazione}: dalle slide, i \textbf{prodotti importanti} sono le categorie
    dell'Allegato~III (Classi I e II) e richiedono procedure di valutazione più rigorose; i
    \textbf{prodotti critici} (Allegato~IV) richiedono un certificato europeo di cybersicurezza almeno
    di livello ``sostanziale''. Tutti gli altri sono prodotti ``di default''.
  \item Se il B2 rientri in una voce dell'Allegato~III decide se l'\textbf{autovalutazione}
    (modulo~A, controllo interno) fosse disponibile: è una questione \emph{decide and defend}
    (art.~32, Allegato~VIII). La bozza di dichiarazione (E-01, riferita al T3) usa il modulo~A ma
    non cita né norme né report di test.
\end{itemize}

\begin{integrazione}[I moduli di valutazione della conformità]
Il CRA riprende i moduli del ``nuovo quadro legislativo'' UE (Allegato~VIII):
\begin{itemize}
  \item \textbf{modulo A} -- controllo interno: il fabbricante valuta da sé, redige la documentazione
    tecnica, firma la dichiarazione UE e appone la marcatura CE;
  \item \textbf{modulo B + C} -- esame UE del tipo da parte di un organismo notificato, seguito dalla
    conformità al tipo basata sul controllo interno della produzione;
  \item \textbf{modulo H} -- garanzia di qualità totale, valutata da un organismo notificato;
  \item in alternativa, un \textbf{certificato europeo di cybersicurezza} rilasciato nell'ambito di uno
    schema del Cybersecurity Act (Capitolo~\ref{lez:2026-11-09}), dove disponibile.
\end{itemize}
Schema generale: prodotti di default $\to$ modulo~A; Classe~I $\to$ modulo~A solo se si applicano
integralmente norme armonizzate, specifiche comuni o certificazioni pertinenti, altrimenti B+C o H;
Classe~II $\to$ B+C, H o certificazione; prodotti critici $\to$ certificazione, se richiesta da atto
della Commissione. Nel caso Sestante l'Allegato~III non contiene una voce per la ``telematica di
flotta'': la scelta più difendibile è trattare il B2 come prodotto di default (modulo~A disponibile),
segnalando la verifica da fare sugli atti di esecuzione che descrivono le categorie. Anche in questa
lettura il vero problema della dichiarazione è il \emph{contenuto} (fascicolo incompleto, requisiti
non rispettati), non il percorso scelto.
\end{integrazione}

\subsection{Obblighi di segnalazione (art.~14)}
Il playbook pone la domanda in modo preciso: per la segnalazione di agosto 2027 e per la credenziale
condivisa nota dal 2025 non conta se fosse cattiva prassi, ma \textbf{se fosse da segnalare, e da
quando}. Ricordiamo i due binari (Capitolo~\ref{lez:2026-11-02}): vulnerabilità attivamente sfruttata
(24~ore, 72~ore, final report entro 14 giorni dalla disponibilità della correzione) e incidente grave
(24~ore, 72~ore, final report entro un mese). L'art.~14 si applica dall'11/09/2026 a tutte le unità,
comprese le 38.000 ``vecchie''.

\begin{integrazione}[Era da notificare?]
Una prova di concetto mostrata da un ricercatore nell'ambito di una divulgazione responsabile non è,
di per sé, evidenza affidabile che un \emph{attore malevolo} abbia sfruttato la vulnerabilità: sui
fatti della room la segnalazione di agosto 2027 non fa scattare l'obbligo dell'art.~14. Il difetto va
quindi valutato soprattutto rispetto all'Allegato~I (Parte~II per la gestione della segnalazione,
Parte~I per la vulnerabilità sfruttabile nota). La conclusione cambierebbe se emergessero indizi di
sfruttamento reale, per esempio box flashati sul campo o la credenziale condivisa usata da terzi: in
quel caso decorrerebbero subito le 24~ore. È una scelta \emph{decide and defend}: va scritta e
motivata, e la lettura opposta (prudente) è legittima se argomentata.
\end{integrazione}

\subsection{Azioni correttive e sanzioni}
Il fabbricante che sa che un prodotto non è conforme deve adottare misure correttive, fino al ritiro o
al richiamo (art.~13(21)), e cooperare con le autorità di vigilanza del mercato, che possono ordinare
le stesse misure. Le sanzioni per la violazione dei requisiti essenziali e degli obblighi del
fabbricante arrivano fino a 15 milioni di euro o al 2,5\% del fatturato mondiale annuo
(art.~64(2)). È questa l'esposizione che interessa davvero al CdA: un'autorità può ordinare di
togliere dal mercato i box installati?

% ---------------------------------------------------------------------
\section{Il lato GDPR in sintesi}

\begin{itemize}
  \item \textbf{Ruoli per flusso}: il ruolo decide chi doveva notificare cosa a settembre 2027, e il
    template di E-12 dice il contrario di ciò che emerge dai fatti.
  \item \textbf{Art.~25} (by design e by default): il box non permette di sospendere il tracciamento
    fuori orario; la conformità del titolare (l'operatore) dipende da come il fabbricante ha costruito
    il prodotto. Ciò che Sestante deve all'operatore è un prodotto che gli consenta di restare
    proporzionato.
  \item \textbf{Art.~28} e \textbf{artt.~44--46}: template con ruoli invertiti e un termine di 30
    giorni per comunicare i breach; trasferimento verso gli USA senza transfer impact assessment.
  \item \textbf{Art.~32}: cifratura presente; controllo degli accessi carente (MFA, 14 account,
    sessioni di 30 giorni).
  \item \textbf{Art.~33(1) e~33(5)}, \textbf{art.~34}: breach di settembre 2027 non notificato e
    registro delle violazioni vuoto.
  \item \textbf{Art.~35} e \textbf{artt.~12--14}: DPIA del 2025 non rifatta dopo l'arrivo della
    dashboard HR; il conducente non sa a che cosa serve il suo punteggio.
  \item \textbf{Artt.~37--39}: DPO con risorse minime e fuori dal percorso di escalation.
\end{itemize}

\begin{integrazione}[Ruolo di Sestante per i tre flussi]
\begin{itemize}
  \item \textbf{Flusso (a) -- posizione e viaggi}: per il servizio di tracciamento reso all'operatore
    Sestante è \textbf{responsabile} (art.~28), perché l'operatore decide le finalità verso i propri
    conducenti. Per lo scoring e il benchmarking venduto ad assicuratori tramite MetricLoop (E-15) le
    finalità le decide Sestante, che diventa \textbf{titolare} per quella parte: quindi
    \emph{entrambi}, distinguendo le finalità.
  \item \textbf{Flusso (b) -- supporto, garanzia, RMA}: \textbf{titolare} dei propri registri e del
    processo RMA che ha causato l'esposizione. Per i dati di viaggio trattati per conto dell'operatore,
    però, la violazione andava comunicata al titolare senza ingiustificato ritardo (art.~33(2)),
    non ``entro 30 giorni'' come nel template.
  \item \textbf{Flusso (c) -- account e fatturazione}: \textbf{titolare}.
\end{itemize}
\end{integrazione}

% ---------------------------------------------------------------------
\section{Svolgimento ragionato}

\begin{integrazione}[Registro dei finding candidati]
Una proposta, con severità calcolata come esposizione $\times$ ampiezza. Bastano quattro finding
solidi; gli altri possono finire tra le limitazioni o nella management letter.
\begin{center}
\small
\begin{tabularx}{\linewidth}{l X l l c}
\toprule
\textbf{\#} & \textbf{Condizione} & \textbf{Criterio} & \textbf{Evidenze} & \textbf{E$\times$A} \\ \midrule
1 & Aggiornamenti OTA accettati senza verifica della firma e con rollback, contro quanto dichiara il
  risk assessment; segnalazione respinta; libreria con CVE nota; nessuna SBOM &
  CRA All.~I, Parte~II & E-02, 04, 07, 08 & 3$\times$3 Critical \\
2 & Credenziale di manutenzione identica su tutti i box e shell attiva in produzione &
  CRA All.~I, Parte~I & E-06, 07, 09 & 3$\times$3 Critical \\
3 & Breach di settembre 2027 non notificato né registrato; RMA che non disattiva l'account cloud &
  GDPR art.~33(1)--(5) & E-14, E-12 & 3$\times$2 High \\
4 & Tracciamento fuori orario e score all'HR attivi di default; modalità privata mai realizzata &
  GDPR art.~25 & E-10, 02, 11 & 2$\times$3 High \\
5 & Template di accordo con ruoli invertiti e termine di 30 giorni, firmato da tutti i clienti &
  GDPR art.~28(3)--(4) & E-12 & 2$\times$3 High \\
6 & Fascicolo tecnico incompleto (11/18) e risk assessment non aggiornato & CRA All.~VII &
  E-01, 02 & 2$\times$3 High \\
7 & Trasferimento negli USA senza transfer impact assessment; informativa che non lo cita &
  GDPR artt.~44--46 & E-15, 11 & 2$\times$3 High \\
8 & Periodo di supporto non definito ma promesso (5 anni) sul sito; contatto PSIRT non pubblicato &
  CRA art.~13 / All.~I, Parte~II & E-03, 05 & 2$\times$3 High \\
\bottomrule
\end{tabularx}
\end{center}
Per l'esposizione dei finding~1 e~2 conta che, con questi difetti, la dichiarazione per le 8.000
unità post-12/12/2027 \textbf{non è firmabile}; l'ampiezza è massima perché la stessa build e la
stessa credenziale sono su ogni unità. Il finding~3 riguarda i conducenti di un solo operatore
(ampiezza~2), ma il difetto RMA è sistemico e può ripetersi a ogni reso.

\medskip\noindent\textbf{Mappatura A.8 (bonus).} 1 $\to$ A.8.24 (crittografia: firme) o A.8.8
(vulnerabilità); 2 $\to$ A.8.5 (autenticazione sicura); 3 $\to$ nessun controllo A.8 (notifica e
registro sono obblighi giuridici e organizzativi; la gestione degli incidenti sta nei controlli
organizzativi A.5), ma il difetto RMA è vicino ad A.8.10 (cancellazione delle informazioni); 4, 5,
7 $\to$ nessun controllo A.8 (scelte di finalità, di disegno o contrattuali); 6 $\to$ A.8.32 (gestione
dei cambiamenti) in senso lato.
\end{integrazione}

\begin{integrazione}[Tested, no finding e controllo che funziona]
\begin{itemize}
  \item \textbf{Art.~14 CRA sulla segnalazione di agosto 2027}: esaminato e scartato sui fatti
    disponibili (nessuna evidenza di sfruttamento malevolo), con il difetto valutato nel finding~1.
  \item \textbf{Certificato ISO/IEC~27001 (E-16)}: esaminato come prova di conformità e scartato:
    scope limitato al cloud, nessun valore come criterio CRA.
  \item \textbf{Allegato~III e modulo~A (E-01)}: esaminato; scelta motivata di considerare il B2
    prodotto di default, quindi il modulo~A non è di per sé un finding.
  \item \textbf{Distributore (E-14)}: il reset e la rivendita li fa Iberia Telematics, ma la procedura
    è di Sestante: il finding si scrive contro il fabbricante.
  \item \textbf{Controllo che funziona}: cifratura a riposo (AES-256 con KMS) e in transito (TLS~1.3,
    HSTS) nel cloud (E-13). Esiste, è implementata nella configurazione, era attiva all'esportazione
    e copre tutti gli account. Il limite (controllo degli accessi carente) è un tema diverso.
\end{itemize}
\end{integrazione}

\begin{integrazione}[Opinione e management letter]
\begin{itemize}
  \item \textbf{Opinione}: con due Critical aperti l'opinione senza rilievi è esclusa. La più
    coerente con il playbook è \textbf{con rilievi} (\emph{qualified}): il fascicolo è riparabile, a
    condizione di chiudere i finding~1 e~2 prima di firmare la dichiarazione. Un'opinione
    \textbf{negativa} (\emph{adverse}) limitata alle 8.000 unità post-12/12/2027 è difendibile se il
    registro è più pesante; va comunque distinta la posizione delle 38.000 unità precedenti, vincolate
    solo all'art.~14.
  \item \textbf{CdA}: oggi il fascicolo non sostiene la dichiarazione per le unità post-12/12/2027;
    le azioni correttive (fino a ritiro o richiamo) sono un rischio concreto.
  \item \textbf{Iberia Lease}: entro il 24/04/2028 si può confermare la conformità solo per unità
    coperte da evidenze nuove (build con firma attiva, credenziali per dispositivo, SBOM, periodo di
    supporto deciso, dichiarazione firmata), distinguendo i box della flotta mista per data di
    immissione.
  \item \textbf{T3}: nessuna opinione possibile. Prima di firmare servono il fascicolo completo
    (Allegato~VII), un risk assessment sulla nuova architettura, la SBOM, la verifica della firma
    attiva, le credenziali per dispositivo, il periodo di supporto e la scelta motivata del percorso
    di valutazione (Allegato~III o modulo~A).
  \item \textbf{Breach di settembre 2027}: rivalutarlo con la metodologia ENISA. A titolo indicativo:
    90 giorni di viaggi di conducenti identificati sono dati comportamentali (DPC~2, o~3 se consentono
    un profilo della vita quotidiana), con identificabilità piena (EI~1) e destinatario noto (CB tra
    0 e 0,25): $SE \geq 2$, cioè almeno ``medio''. La conclusione ``rischio basso, veicoli non persone''
    non regge. Quindi: notifica tardiva al Garante con i motivi del ritardo (art.~33(1)), iscrizione
    nel registro (art.~33(5)), valutazione dell'art.~34, comunicazione all'operatore. Se si decide di
    non notificare, la scelta va scritta e motivata con una valutazione, non con una frase in un'email.
  \item \textbf{Tre azioni prioritarie}: firma OTA e anti-rollback, poi revoca della credenziale
    condivisa via OTA (CTO); procedura breach con il DPO nell'escalation e RMA che disattiva l'account
    (DPO e direzione); fascicolo, SBOM e periodo di supporto (responsabile prodotto). Formazione
    misurabile: triage delle vulnerabilità per PSIRT e firmware; escalation dei breach per supporto e
    ingegneria, verificata con un'esercitazione su un caso RMA.
\end{itemize}
\end{integrazione}

% ---------------------------------------------------------------------
\section{Errori tipici}

Dal playbook e dalle slide di briefing:
\begin{itemize}
  \item frasi generiche come ``Sestante non è conforme'' o ``l'azienda non rispetta il GDPR'': conforme
    a che cosa, per quali unità, da quando, per quali dati e quali persone?
  \item ``Non c'è una SBOM'' o ``l'OTA è insicuro'' senza evidenza. Meglio: ``la verifica della firma
    è presente nella libreria e disattivata nella build (E-07), mentre il risk assessment afferma che
    gli aggiornamenti sono firmati (E-02)'';
  \item mescolare la severità di audit con quella ENISA;
  \item scrivere finding contro i clienti (gli operatori sono titolari dei dati dei propri
    conducenti) o valutare il diritto del lavoro, che è fuori perimetro;
  \item trattare il certificato ISO/IEC~27001 come prova di conformità;
  \item prendere un'intervista per un fatto senza un documento che la confermi;
  \item opinione senza rilievi accanto a un Critical aperto, o opinione ``dipende'';
  \item più criteri nello stesso finding, o citazioni imprecise (un paragrafo dell'art.~33 non vale
    l'altro);
  \item dimenticare la limitazione sul T3 o ciò che non si è potuto testare;
  \item forzare una mappatura A.8: ``nessun controllo A.8 è applicabile'' è spesso la risposta
    corretta, perché metà dei problemi non sono controlli tecnici.
\end{itemize}

% ---------------------------------------------------------------------
\section{Collegamenti}
\begin{itemize}
  \item \textbf{Teoria CRA}: lezione del 02/11 (Capitolo~\ref{lez:2026-11-02}) e capitolo tematico
    (Capitolo~\ref{cap:cra}): Allegato~I, prodotti importanti e critici, obblighi lungo la catena,
    art.~14.
  \item \textbf{Simulazione Pulse K} (Sezione~\ref{sec:pulsek}): stessi binari CRA--GDPR visti
    dall'interno durante l'incidente; anche lì una DPIA mai aggiornata e una libreria obsoleta.
  \item \textbf{Metodologia ENISA}: scheda di ripasso (Capitolo~\ref{cap:schede}).
  \item \textbf{Cybersecurity Act} (Capitolo~\ref{lez:2026-11-09}): i certificati europei come via per
    dimostrare la conformità dei prodotti critici e, in parte, di quelli importanti.
  \item \textbf{Lab~3} (risk register ISO/IEC~27001): i controlli dell'Annex~A, qui usati solo per il
    bonus.
  \item \textbf{Scadenze}: tabella riassuntiva (Capitolo~\ref{cap:scadenze}).
\end{itemize}

% ---------------------------------------------------------------------
\section{Punti da verificare}
\begin{verifica}
\begin{itemize}
  \item \textbf{Riferimenti agli allegati nel playbook}: in un punto l'Allegato~VIII è indicato come
    ``dichiarazione di conformità'' e il modulo~A come ``Allegato~II''; altrove l'Allegato~VIII è
    correttamente quello dei moduli. Nel testo del CRA il modello di dichiarazione UE dovrebbe essere
    l'Allegato~V e l'Allegato~II riguardare le informazioni per l'utente: verificare sul testo
    ufficiale.
  \item \textbf{Numerazione degli articoli} citati dal playbook: art.~13(17) (punto di contatto),
    art.~13(21) (azioni correttive), art.~21 (casi in cui gli obblighi del fabbricante si applicano ad
    altri operatori), artt.~53--54 (poteri di vigilanza), art.~64(10). Verificare sul Reg.~(UE)
    2024/2847.
  \item \textbf{Modifica sostanziale} delle unità già immesse e \textbf{moduli di valutazione} per
    classe di prodotto: contenuti di integrazione, non spiegati in aula.
  \item \textbf{Calcolo ENISA} del breach di settembre 2027: valori indicativi, non discussi in aula.
  \item Nel playbook un esempio sulla SBOM rimanda a E-09, mentre il foglio componenti è E-08:
    probabile refuso.
  \item La griglia di severità non può produrre il valore 8 (i fattori vanno da 1 a 3): la fascia
    ``8--9 Critical'' coincide in pratica con 9.
  \item Debrief del laboratorio (cinque domande delle slide): non ancora svolto in aula.
\end{itemize}
\end{verifica}


% =====================================================================
\part{Schemi di ripasso}
% =====================================================================
\chapter{Tabella comparativa delle notifiche}
\label{cap:scadenze}

Il professore ha avvertito che all'esame va identificata prima la norma applicabile e poi la
tempistica corretta, senza confondere le fasi delle diverse norme. La tabella riassume le
tempistiche viste finora.

\begin{center}
\small
\begin{tabularx}{\linewidth}{>{\raggedright\arraybackslash}p{1.7cm}
  >{\raggedright\arraybackslash}X >{\raggedright\arraybackslash}X
  >{\raggedright\arraybackslash}X >{\raggedright\arraybackslash}X}
\toprule
\textbf{Norma} & \textbf{Evento e destinatario} & \textbf{Fase 1} & \textbf{Fase 2} & \textbf{Fase 3} \\
\midrule
\textbf{GDPR} (artt.~33--34)
  & Data breach $\rightarrow$ autorità di controllo (Garante; capofila se transfrontaliero)
  & Notifica entro \textbf{72 h dalla consapevolezza} (anche preliminare)
  & Completamento: \textbf{nessun termine} (prassi $\approx$ 1 mese)
  & Interessati \textbf{senza ingiustificato ritardo}, se rischio elevato (art.~34) \\
\midrule
\textbf{NIS2} (art.~23)
  & Incidente significativo $\rightarrow$ CSIRT (Italia: CSIRT Italia, ACN)
  & Early warning \textbf{24 h dalla consapevolezza}
  & Notifica \textbf{72 h dalla consapevolezza}; relazione intermedia su richiesta
  & Final report \textbf{1 mese dalla notifica} \\
\midrule
\textbf{DORA}
  & Incidente ICT grave $\rightarrow$ autorità competente di settore
  & Notifica iniziale \textbf{4 h dalla classificazione} come grave, max \textbf{24 h dalla conoscenza}
  & Relazione intermedia \textbf{72 h dalla notifica iniziale}
  & Relazione finale \textbf{1 mese} \\
\midrule
\textbf{CRA} (art.~14(2))
  & Vulnerabilità attivamente sfruttata $\rightarrow$ CSIRT coordinatore + ENISA (Single Reporting Platform)
  & Early warning \textbf{24 h dalla conoscenza}
  & Notifica \textbf{72 h dalla conoscenza}
  & Final report \textbf{14 giorni dalla disponibilità della misura correttiva} \\
\midrule
\textbf{CRA} (art.~14(4))
  & Incidente grave sulla sicurezza del prodotto $\rightarrow$ come sopra
  & Early warning \textbf{24 h}
  & Notifica \textbf{72 h}
  & Final report \textbf{1 mese dalla notifica} \\
\midrule
\textbf{AI Act} (art.~73)
  & Incidente grave causato da un sistema di IA ad alto rischio $\rightarrow$ autorità di vigilanza del mercato
  & Notifica \textbf{15 giorni} dalla conoscenza (stabilito il nesso causale)
  & \textbf{2 giorni} se infrastruttura critica o violazione diffusa
  & \textbf{10 giorni} in caso di morte \\
\bottomrule
\end{tabularx}
\end{center}

\begin{esame}
\textbf{Trappole tipiche:}
\begin{itemize}
  \item le \textbf{72 ore di NIS2} decorrono dalla \emph{consapevolezza}; le \textbf{72 ore di DORA}
    decorrono dalla \emph{notifica iniziale};
  \item il \textbf{final report CRA sulle vulnerabilità} decorre dalla \emph{disponibilità del fix},
    non dalla consapevolezza;
  \item il \textbf{GDPR non ha final report} con scadenza;
  \item l'\textbf{early warning CRA} non è la \textbf{notifica preliminare GDPR}: sono obblighi
    distinti, verso autorità distinte;
  \item una notifica fatta a un'autorità \textbf{non adempie} gli obblighi verso un'altra: se lo stesso
    evento rientra in più norme, le notifiche corrono in parallelo;
  \item le scadenze \textbf{corrono anche nel weekend}.
\end{itemize}
\end{esame}

\begin{verifica}
\textbf{DORA, relazione finale}: gli appunti dell'anno precedente (\emph{main.pdf}) riportano in un
punto ``1 mese dalla relazione intermedia o dal suo ultimo aggiornamento'' e in un altro ``1 mese
dall'incidente''. Verificare sulle slide DORA di quest'anno.
\end{verifica}

\begin{integrazione}[DORA, relazione finale]
Le norme tecniche di attuazione sulla notifica degli incidenti (Reg.~delegato (UE) 2025/301) fissano la
relazione finale \textbf{entro un mese dall'ultima relazione intermedia aggiornata}. Le stesse norme
consentono ad alcune entità, se la scadenza cade nel fine settimana o in un festivo, di inviare entro
mezzogiorno del primo giorno lavorativo: per DORA la regola ``le scadenze corrono anche nel weekend'' va
quindi presa con cautela. Da confermare con le slide DORA.
\end{integrazione}

\chapter{Mappe concettuali}
\label{cap:mappe}

% ---------------------------------------------------------------------
\section{Quale norma si applica?}

Prima delle due competenze richieste dal corso (lezione del 02/11). Le domande \textbf{non sono
alternative}: vanno poste tutte, perché le norme si cumulano.

\begin{center}
\begin{tikzpicture}[
  node distance=6mm and 8mm,
  q/.style={draw, rounded corners, fill=blue!6, text width=6.6cm, align=left,
    font=\small, inner sep=5pt, minimum height=1.7cm},
  r/.style={draw, thick, fill=green!8, text width=5.4cm, align=left,
    font=\small, inner sep=5pt, minimum height=1.7cm},
  >=Stealth]
\node[q] (q1) {Sono coinvolti \textbf{dati personali} di persone fisiche?};
\node[r, right=of q1] (r1) {\textbf{GDPR}\\ se data breach: artt.~33--34};

\node[q, below=of q1] (q2) {Il soggetto è un'\textbf{entità finanziaria} (banca, istituto di
  pagamento, assicurazione, \ldots) o un suo fornitore ICT?};
\node[r, right=of q2] (r2) {\textbf{DORA}\\ lex specialis rispetto a NIS2; ai fornitori ICT si
  applica per via contrattuale};

\node[q, below=of q2] (q3) {Il soggetto è \textbf{essenziale o importante} in un settore degli
  Allegati I--II di NIS2?};
\node[r, right=of q3] (r3) {\textbf{NIS2}\\ in Italia D.Lgs.~138/2024, autorità ACN};

\node[q, below=of q3] (q4) {È un \textbf{prodotto con elementi digitali} (hardware o software
  installabile, connesso) messo a disposizione sul mercato UE e non coperto da un regime verticale?};
\node[r, right=of q4] (r4) {\textbf{CRA}\\ obblighi di segnalazione dall'11/9/2026; tutto dall'11/12/2027};

\node[q, below=of q4] (q5) {È o contiene un \textbf{sistema di IA}?};
\node[r, right=of q5] (r5) {\textbf{AI Act}\\ alto rischio: presunzione di conformità all'art.~15 tramite CRA};

\node[q, below=of q5] (q6) {Si vuole dimostrare la sicurezza con una \textbf{certificazione UE}?};
\node[r, right=of q6] (r6) {\textbf{CSA}\\ schemi europei di certificazione, ENISA};

\foreach \i in {1,...,6} {\draw[->] (q\i) -- node[above, font=\scriptsize]{sì} (r\i);}
\foreach \i/\j in {1/2,2/3,3/4,4/5,5/6} {\draw[->] (q\i) -- node[left, font=\scriptsize]{in ogni caso} (q\j);}
\end{tikzpicture}
\end{center}

\begin{nota}
\begin{itemize}
  \item Ospedale, società energetica, telco: NIS2 e, se ci sono dati personali, GDPR.
  \item Banca o istituto di pagamento: DORA e GDPR (NIS2 cede a DORA dove le norme si sovrappongono).
  \item Fabbricante di dispositivi IoT o di app: CRA e, quasi sempre, GDPR.
  \item Un certificato ISO 27001 o CSA \textbf{supporta} l'accountability ma \textbf{non prova} la
    conformità legale.
\end{itemize}
\end{nota}

% ---------------------------------------------------------------------
\section{Il percorso in caso di incidente}

Seconda competenza richiesta dal corso: individuare il percorso corretto quando si verifica un incidente.

\begin{center}
\begin{tikzpicture}[
  node distance=4mm,
  s/.style={draw, rounded corners, fill=orange!8, text width=12.5cm, align=left,
    font=\small, inner sep=5pt},
  >=Stealth]
\node[s] (a) {\textbf{1. Rilevazione ed escalation}: chi riceve il segnale (SOC, reperibile, call
  center) lo porta subito a chi decide, compresi DPO e legale.};
\node[s, below=of a] (b) {\textbf{2. Contenimento e preservazione delle evidenze}: isolamento, immagini
  forensi, log, catena di custodia.};
\node[s, below=of b] (c) {\textbf{3. Qualificazione}: quali norme sono coinvolte (mappa precedente)?
  Data breach? Incidente significativo/grave? Vulnerabilità attivamente sfruttata?};
\node[s, below=of c] (d) {\textbf{4. Consapevolezza}: fissare e \textbf{documentare} il momento di
  avvio di ciascun orologio, che può essere diverso per ciascuna norma.};
\node[s, below=of d] (e) {\textbf{5. Valutazione della severità}: ENISA per il GDPR; criteri di
  significatività/gravità per NIS2, DORA e CRA.};
\node[s, below=of e] (f) {\textbf{6. Notifiche in parallelo} verso ciascuna autorità, rispettando
  ogni scadenza (Capitolo~\ref{cap:scadenze}).};
\node[s, below=of f] (g) {\textbf{7. Comunicazioni}: interessati (art.~34 GDPR), utenti (CRA),
  destinatari dei servizi e pubblico (NIS2), retailer, stampa. Trasparenza: niente smentite false.};
\node[s, below=of g] (h) {\textbf{8. Registri e chiusura}: registro degli incidenti e registro dei
  data breach allineati; report finali; lezioni apprese e aggiornamento di risk assessment e DPIA.};
\foreach \x/\y in {a/b,b/c,c/d,d/e,e/f,f/g,g/h} {\draw[->] (\x) -- (\y);}
\end{tikzpicture}
\end{center}

\dacompletare{Altre mappe da aggiungere man mano: ruoli GDPR (titolare, responsabile,
sub-responsabile); classificazione dei prodotti CRA; classificazione per rischio dell'AI Act;
soggetti NIS2.}

\chapter{Schede di ripasso}
\label{cap:schede}

% ---------------------------------------------------------------------
\section{Severità di un data breach (metodologia ENISA 2013)}

\[
SE = DPC \times EI + CB
\]

\begin{center}
\small
\begin{tabularx}{\linewidth}{l X}
\toprule
\textbf{Fattore} & \textbf{Valori} \\ \midrule
\textbf{DPC} -- Data Processing Context & 1 semplici $\cdot$ 2 comportamentali $\cdot$ 3 finanziari
  $\cdot$ 4 sensibili; aumentabile o riducibile in base al contesto (minori, sicurezza, volume),
  motivando la scelta \\
\textbf{EI} -- Ease of Identification & 0,25 trascurabile $\cdot$ 0,5 limitata $\cdot$
  0,75 significativa $\cdot$ 1 massima \\
\textbf{CB} -- Circumstances of the Breach & riservatezza 0--0,5 $\cdot$ integrità 0--0,5 $\cdot$
  disponibilità 0--0,5 $\cdot$ intento malevolo +0,5 \\
\midrule
\textbf{Scala} & $SE<2$ bassa $\cdot$ $2 \le SE < 3$ media $\cdot$ $3 \le SE < 4$ alta $\cdot$
  $SE \ge 4$ molto alta \\ \bottomrule
\end{tabularx}
\end{center}

\begin{esame}
\begin{itemize}
  \item Non conta ``indovinare il numero'' ma \textbf{motivare} ogni valore (accountability).
  \item La metodologia non distingue minori e adulti: se ci sono soggetti vulnerabili (minori,
    pazienti, lavoratori), conviene alzare il DPC motivando e comunicare comunque agli interessati.
  \item Per i minori la comunicazione ex art.~34 va ai \textbf{genitori}.
  \item Anche l'indisponibilità temporanea dei dati è un data breach.
\end{itemize}
\end{esame}

\begin{nota}
Scala tratta dalle slide del caso Pulse K (A.A.~2026/27). Gli appunti dell'anno precedente usano
soglie leggermente diverse (es.\ 2--2,5 media, 2,5--3 medio-alta): fa fede la versione di quest'anno.
\end{nota}

% ---------------------------------------------------------------------
\section{Ruoli nelle diverse norme}

\begin{center}
\small
\begin{tabularx}{\linewidth}{l X}
\toprule
\textbf{Norma} & \textbf{Ruoli principali} \\ \midrule
GDPR    & Titolare (decide finalità e mezzi; notifica i breach) $\cdot$ responsabile (tratta per conto
  del titolare; avvisa il titolare senza ingiustificato ritardo) $\cdot$ sub-responsabile $\cdot$
  contitolari $\cdot$ interessato $\cdot$ DPO \\
NIS2    & Soggetti essenziali e importanti $\cdot$ organi di gestione (responsabilità non delegabile)
  $\cdot$ punto di contatto NIS $\cdot$ CSIRT e autorità competente \\
DORA    & Entità finanziarie $\cdot$ fornitori terzi di servizi ICT $\cdot$ organo di gestione
  $\cdot$ funzione di controllo del rischio ICT \\
CRA     & Fabbricante $\cdot$ importatore $\cdot$ distributore \\
AI Act  & Provider $\cdot$ deployer $\cdot$ importatore $\cdot$ distributore $\cdot$ fabbricante del
  prodotto (art.~25: riqualificazione come provider) \\ \bottomrule
\end{tabularx}
\end{center}

% ---------------------------------------------------------------------
\section{Date da ricordare}

\begin{center}
\small
\begin{tabularx}{\linewidth}{l X}
\toprule
\textbf{Data} & \textbf{Evento} \\ \midrule
25/05/2018 & Il GDPR diventa applicabile \\
17/01/2025 & DORA diventa applicabile \\
11/09/2026 & CRA: si applicano gli obblighi di segnalazione (art.~14) \\
11/12/2027 & CRA: applicazione integrale (requisiti dell'Allegato I) \\
02/12/2027 & AI Act: obblighi per i sistemi ad alto rischio dell'Allegato III (dopo l'AI Omnibus) \\
02/08/2028 & AI Act: obblighi per i sistemi ad alto rischio dell'Allegato I (dopo l'AI Omnibus) \\ \bottomrule
\end{tabularx}
\end{center}

\begin{integrazione}[Altre date]
NIS2: D.lgs.~138/2024 in vigore dal 16/10/2024 (recepimento italiano). AI Act: in vigore dal 01/08/2024;
divieti dell'art.~5 dal 02/02/2025; obblighi GPAI dal 02/08/2025. CSA: in vigore dal 27/06/2019.
Le date dell'AI Omnibus sono dalle slide introduttive (Capitolo~\ref{cap:aiact}).
\end{integrazione}

\dacompletare{Rivedere le date NIS2, DORA e AI Act dopo le relative lezioni.}


% =====================================================================
\part{Quiz ed esame}
% =====================================================================
\chapter{Formato dell'esame: cosa sappiamo}
\label{cap:formato}

\section{Struttura (da fonte non ufficiale)}

L'unica traccia disponibile è un messaggio nel gruppo Telegram del corso, presentato come elenco
delle domande di un appello precedente (Capitolo~\ref{cap:tema-precedente}). Da quel messaggio
risulta questa struttura:
\begin{itemize}
  \item \textbf{4 domande chiuse} a risposta multipla, alcune nella forma ``\emph{select the WRONG
    answer}'' (individuare l'affermazione errata);
  \item \textbf{1 domanda aperta} su un caso pratico, divisa in sotto-domande: quale norma si applica,
    cosa va fatto secondo la norma, quali obblighi valgono a prescindere dall'incidente;
  \item domande in \textbf{inglese}.
\end{itemize}

\begin{verifica}
\begin{itemize}
  \item È \textbf{una sola fonte}, non verificata. Per la stessa regola dell'audit spiegata in aula,
    una fonte sola è un'asserzione, non un fatto: va confermata con il docente o con altri studenti.
  \item L'appello si riferisce con ogni probabilità a un anno accademico precedente. Allora il CRA non
    era ancora operativo (obblighi di segnalazione dall'11/9/2026), mentre quest'anno è stato centrale:
    è ragionevole aspettarsi domande sul CRA e su casi CRA--GDPR.
\end{itemize}
\end{verifica}

\section{Indicazioni date dai docenti durante il corso}

\begin{itemize}
  \item L'esercizio di \textbf{notifica alle autorità} è ``molto probabile'' (lezione del 02/11).
  \item Si valuta la capacità di \textbf{qualificare il caso} (quale norma) e di indicare il
    \textbf{percorso corretto} in caso di incidente; non serve ricordare a memoria gli allegati.
  \item Le \textbf{scadenze} vanno ricordate con precisione, senza confondere le fasi delle diverse
    norme (Capitolo~\ref{cap:scadenze}). Il professore ha parlato di una tabella riassuntiva
    ``prodotta per voi per l'esame''.
  \item I \textbf{dettagli del caso} contano (es.\ una DPIA mai aggiornata): vanno letti e usati nella
    risposta.
  \item Nelle domande aperte si premia il \textbf{ragionamento motivato} (necessità, proporzionalità,
    accountability), non l'elenco di articoli.
\end{itemize}

\begin{verifica}
Non è chiaro se la tabella riassuntiva delle scadenze sarà consultabile durante l'esame o se è solo
materiale di studio. Chiedere al docente.
\end{verifica}

\section{Metodo per la domanda aperta}

\begin{enumerate}
  \item \textbf{Qualifica}: elenca le norme applicabili \emph{e} spiega in una riga perché le altre
    non si applicano (es.\ ``DORA no: non è un'entità finanziaria'').
  \item \textbf{Ruoli}: chi è titolare, responsabile, soggetto essenziale/importante, fabbricante.
  \item \textbf{Orologi}: per ogni norma, da quando decorre e verso quale autorità.
  \item \textbf{Severità}: valuta il rischio per gli interessati (ENISA) e decidi su art.~33 e 34.
  \item \textbf{Misure}: contenimento, comunicazioni, registri, misure organizzative e contrattuali.
  \item \textbf{Usa i dettagli del caso} per motivare ogni scelta.
\end{enumerate}

\chapter{Tema d'esame di un appello precedente}
\label{cap:tema-precedente}

\begin{nota}[Fonte]
Messaggio nel gruppo Telegram del corso (``per i posteri''), riportato testualmente. Il testo delle
opzioni delle domande chiuse non è disponibile: per ciascuna si indica cosa va saputo. Le domande
chiuse ricostruite si trovano nel quiz (Capitolo~\ref{cap:quiz}).
\end{nota}

\begin{enumerate}
  \item \emph{GDPR legal basis (select the WRONG answer)}
  \item \emph{NIS2 incident definition}
  \item \emph{ENISA: what is it, what role does it have}
  \item \emph{DORA strategies (select the WRONG answer)}
  \item \emph{Open question: a hospital with some IT service providers for storage (cloud) and security.
    There is a data breach/incident.}
    \begin{enumerate}[label=\Alph*)]
      \item \emph{What regulation is involved?}
      \item \emph{According to the regulation, what does the hospital have to do?}
      \item \emph{Regardless of the incident, what does the hospital have to do regarding the
        management of the supply chain?}
    \end{enumerate}
\end{enumerate}

% ---------------------------------------------------------------------
\section{Domande chiuse: cosa sapere}

\paragraph{1. Basi giuridiche GDPR (art.~6).} Sono sei, tutte di pari valore: consenso, esecuzione di
un contratto, obbligo legale, interessi vitali, interesse pubblico o pubblici poteri, legittimo
interesse. Affermazioni errate tipiche:
\begin{itemize}
  \item ``il consenso è sempre necessario'' o ``prevale sulle altre basi'';
  \item ``un'autorità pubblica può fondarsi sul legittimo interesse per i propri compiti'' (escluso
    dall'art.~6(1));
  \item ``la base contrattuale copre anche il marketing'';
  \item ``si può condividere i dati con terzi per il loro marketing in base al legittimo interesse''
    (laboratorio GDPR: serve il consenso).
\end{itemize}

\paragraph{2. Definizione di incidente NIS2 (art.~6).} Evento che compromette la
\textbf{disponibilità, l'autenticità, l'integrità o la riservatezza} di dati conservati, trasmessi o
elaborati, o dei servizi offerti dai sistemi informativi e di rete o accessibili attraverso di essi.
Va distinto dall'\textbf{incidente significativo} (art.~23(3)): quello che ha causato o può causare
una grave perturbazione operativa dei servizi o perdite finanziarie, oppure danni materiali o
immateriali considerevoli ad altre persone. Il ``può causare'' è giuridicamente decisivo.
Attenzione: NIS2 include l'\emph{autenticità}, la definizione GDPR di data breach no.

\paragraph{3. ENISA.} Agenzia dell'UE per la cybersicurezza. Con il Cybersecurity Act ha ottenuto un
\textbf{mandato permanente} e un ruolo attivo: supporto alle politiche UE e agli Stati membri,
cooperazione operativa (rete dei CSIRT), situational awareness, preparazione degli schemi europei di
certificazione su richiesta della Commissione, sensibilizzazione e formazione, gestione della
\textbf{Single Reporting Platform} del CRA. \textbf{Non} irroga sanzioni, \textbf{non} è un CSIRT
nazionale e \textbf{non} rilascia certificati alle imprese.

\paragraph{4. Strategie DORA.} La strategia di resilienza operativa digitale comprende: tolleranza al
rischio ICT, obiettivi, architettura, rilevazione, test; \textbf{strategia multi-vendor}; gestione e
conformità dell'outsourcing ICT; \textbf{strategia di comunicazione} per gli incidenti ICT. È definita
e approvata dall'\textbf{organo di gestione} dell'entità finanziaria. Affermazioni errate tipiche: che
la approvi il fornitore ICT, che sia identica al risk assessment GDPR, che DORA imponga una checklist
fissa (vale invece la proporzionalità).

% ---------------------------------------------------------------------
\section{Domanda aperta: risposta modello}

\begin{soluzione}[A) Quali norme sono coinvolte]
\begin{itemize}
  \item \textbf{GDPR}. L'ospedale tratta dati personali, anche \textbf{sanitari} (categoria
    particolare, art.~9), ed è \textbf{titolare}. Il fornitore cloud e il fornitore di sicurezza sono
    \textbf{responsabili del trattamento} (art.~28), ed eventuali loro fornitori sono sub-responsabili.
    Se sono coinvolti dati personali, l'incidente è un \textbf{data breach} (art.~4(12)).
  \item \textbf{NIS2}. La sanità è un settore dell'Allegato~I: l'ospedale è \textbf{soggetto
    essenziale} se grande impresa, \textbf{importante} se media (art.~3). In Italia la norma è il
    D.Lgs.~138/2024, con l'ACN come autorità e il CSIRT Italia per le notifiche.
  \item \textbf{I fornitori} sono a loro volta soggetti NIS2: i servizi di cloud computing rientrano
    nelle infrastrutture digitali e i servizi di sicurezza gestiti nella gestione dei servizi ICT
    (B2B), entrambi nell'Allegato~I. Hanno obblighi di notifica propri, distinti da quelli
    dell'ospedale.
  \item \textbf{Perché non le altre}: DORA no (l'ospedale non è un'entità finanziaria); CRA solo se
    l'incidente nasce da una vulnerabilità attivamente sfruttata in un prodotto con elementi digitali,
    e in quel caso l'obbligo è del fabbricante, non dell'ospedale; i dispositivi medici sono soggetti
    all'MDR ed esclusi dal CRA.
\end{itemize}
\end{soluzione}

\begin{soluzione}[B) Cosa deve fare l'ospedale secondo le norme]
\textbf{Gestione immediata}
\begin{itemize}
  \item Attivare la procedura (idealmente integrata GDPR--NIS2), contenere l'incidente insieme ai
    fornitori, preservare le evidenze (log, immagini forensi, catena di custodia) e informare l'organo
    di gestione.
  \item Ottenere dal fornitore le informazioni necessarie: il responsabile deve informare il titolare
    \textbf{senza ingiustificato ritardo} (art.~33(2)).
  \item Fissare e documentare il momento di consapevolezza per ciascuna norma.
\end{itemize}
\textbf{GDPR}
\begin{itemize}
  \item Valutare la severità per gli interessati (ENISA). Con dati sanitari il DPC è 4, quindi la
    severità è quasi certamente alta o molto alta.
  \item \textbf{Art.~33}: notificare il Garante entro \textbf{72 ore} dalla consapevolezza, anche in
    forma preliminare, completando in seguito.
  \item \textbf{Art.~34}: comunicare ai pazienti, in linguaggio chiaro, se il rischio è elevato (con
    dati sanitari tipicamente sì). La comunicazione deve contenere natura della violazione, contatti
    del DPO, conseguenze probabili, misure adottate e consigli agli interessati.
  \item Annotare la violazione nel \textbf{registro dei data breach} (art.~33(5)), anche se non si notifica.
\end{itemize}
\textbf{NIS2} (se l'incidente è significativo)
\begin{itemize}
  \item \textbf{Early warning} al CSIRT entro \textbf{24 ore}, indicando se si sospetta un atto
    illecito o un impatto transfrontaliero.
  \item \textbf{Notifica dell'incidente} entro \textbf{72 ore}, con valutazione iniziale di gravità,
    impatto e indicatori di compromissione; \textbf{relazione intermedia} su richiesta.
  \item \textbf{Relazione finale} entro \textbf{un mese} dalla notifica: descrizione, causa, misure,
    impatti transfrontalieri.
  \item Informare, se opportuno, i destinatari dei servizi impattati.
\end{itemize}
\textbf{Inoltre}
\begin{itemize}
  \item Le notifiche al Garante e al CSIRT sono \textbf{separate e parallele}: l'una non sostituisce l'altra.
  \item Valutare la denuncia alla Polizia Postale (non obbligatoria, ma la sua assenza può essere
    valutata negativamente).
  \item Allineare registro degli incidenti e registro dei data breach; far approvare le misure di
    remediation dall'organo di gestione; aggiornare risk assessment e DPIA.
\end{itemize}
\end{soluzione}

\begin{soluzione}[C) Supply chain: obblighi a prescindere dall'incidente]
\begin{itemize}
  \item \textbf{NIS2, art.~21(2)(d)}: la sicurezza della catena di approvvigionamento fa parte delle
    misure obbligatorie di gestione del rischio. Secondo l'\textbf{art.~21(3)} vanno considerate le
    vulnerabilità specifiche di ciascun fornitore, la qualità dei loro prodotti e delle loro pratiche di
    sicurezza e le procedure di sviluppo sicuro.
  \item In pratica (linee guida ACN): inventario dei fornitori e dei servizi critici; mappatura delle
    dipendenze e dei rischi a cascata; requisiti minimi di sicurezza nei contratti; SLA e tempi di
    ripristino; rivalutazione periodica; supervisione e approvazione da parte dell'organo di gestione
    (art.~20, responsabilità non delegabile).
  \item \textbf{GDPR, art.~28}: scegliere solo responsabili che offrano \textbf{garanzie sufficienti};
    stipulare un contratto (DPA) con il contenuto dell'art.~28(3), cioè istruzioni documentate,
    riservatezza, misure dell'art.~32, regole sui sub-responsabili, assistenza al titolare (anche per
    gli artt.~33--36), cancellazione o restituzione dei dati, diritto di audit.
  \item Prevedere una \textbf{clausola di notifica dei breach} con tempi stretti (es.\ 24 ore), così da
    poter rispettare le 72 ore.
  \item \textbf{Localizzazione dei dati}: se il cloud tratta dati fuori dall'UE servono gli strumenti
    del Capo~V (decisione di adeguatezza oppure SCC con TIA).
  \item Clausole su cooperazione negli incidenti, continuità operativa, backup, \textbf{exit strategy}
    e reversibilità.
  \item Le certificazioni dei fornitori (ISO 27001, 27017, 27018) sono utili come evidenza ma
    \textbf{non bastano} da sole: le verifiche vanno documentate (accountability).
\end{itemize}
\end{soluzione}

\chapter{Simulazioni d'esame}
\label{cap:simulazioni}

\section{Simulazione 1 -- Pulse K (CRA e GDPR)}
Caso svolto in aula il 02/11/2026: vedi Sezione~\ref{sec:pulsek}. Per esercitarsi, rifare i quattro
task senza guardare le soluzioni e confrontare il risultato con la lezione.

% ---------------------------------------------------------------------
\section{Simulazione 2 -- Istituto di pagamento (DORA e GDPR)}

\begin{nota}
Caso costruito sul modello della domanda aperta dell'appello precedente; non proviene dai docenti.
\end{nota}

\paragraph{Scenario.} \emph{PayNest S.p.A.} è un istituto di pagamento italiano con un'app usata da
600.000 clienti in Italia e Spagna. L'hosting è affidato a un cloud provider statunitense certificato
nell'EU--US Data Privacy Framework; il monitoraggio di sicurezza (SOC) a un MSSP italiano.
\begin{itemize}
  \item \textbf{Venerdì, 22:00}: il SOC rileva un ransomware che cifra il database delle transazioni.
    Il servizio di pagamento resta fermo per 9 ore.
  \item \textbf{Sabato, 02:00}: il comitato di crisi classifica l'incidente come grave.
  \item \textbf{Sabato, 10:00}: i log confermano l'esfiltrazione di nome, IBAN e storico delle
    transazioni di 120.000 clienti. Il vettore di accesso sono le credenziali VPN di un tecnico
    dell'MSSP, prive di MFA.
  \item I dati vengono ripristinati dai backup entro domenica.
\end{itemize}

\paragraph{Domande.}
\begin{enumerate}[label=\Alph*)]
  \item Quali norme si applicano? Perché NIS2 non è la norma di riferimento per la notifica?
  \item Quali notifiche deve fare PayNest, a chi ed entro quando?
  \item Calcolare la severità ENISA e decidere sugli artt.~33 e 34 GDPR.
  \item Quali misure contrattuali e organizzative vanno adottate verso i fornitori?
\end{enumerate}

\begin{soluzione}[A) Norme applicabili]
\begin{itemize}
  \item \textbf{DORA}: PayNest è un'entità finanziaria (istituto di pagamento). Il settore bancario e
    finanziario è anche nell'Allegato~I di NIS2, ma DORA è \textbf{lex specialis}: dove le norme si
    sovrappongono (gestione del rischio ICT, notifica degli incidenti) si applica DORA.
  \item \textbf{GDPR}: data breach con perdita di riservatezza (esfiltrazione) e di disponibilità
    (cifratura). PayNest è titolare; cloud provider e MSSP sono responsabili.
  \item \textbf{Trasferimento extra UE}: il cloud provider è certificato nel Data Privacy Framework,
    quindi il trasferimento si fonda su una \textbf{decisione di adeguatezza}. Restano gli obblighi di
    trasparenza e di monitoraggio.
  \item I due fornitori sono \textbf{fornitori terzi di servizi ICT} ai sensi di DORA, e rispondono
    verso PayNest per via contrattuale.
\end{itemize}
\end{soluzione}

\begin{soluzione}[B) Notifiche e tempistiche]
\begin{center}
\small
\begin{tabularx}{\linewidth}{l X l}
\toprule
\textbf{Obbligo} & \textbf{Decorrenza} & \textbf{Scadenza} \\ \midrule
DORA, notifica iniziale  & 4 h dalla classificazione (Sab 02:00), max 24 h dalla conoscenza (Ven 22:00)
  & \textbf{Sab 06:00} (vale la prima) \\
DORA, relazione intermedia & 72 h dalla notifica iniziale & Mar 06:00 se inviata Sab 06:00 \\
DORA, relazione finale & 1 mese & vedi Capitolo~\ref{cap:scadenze} \\
GDPR, art.~33 (Garante) & Consapevolezza della compromissione di dati personali (Sab 10:00) & Mar 10:00 \\
GDPR, art.~34 (clienti) & Rischio elevato & Senza ingiustificato ritardo \\ \bottomrule
\end{tabularx}
\end{center}
\begin{itemize}
  \item Destinatari: per DORA l'autorità competente di settore (per un istituto di pagamento italiano
    la Banca d'Italia); per il GDPR il Garante, capofila perché lo stabilimento principale è in Italia
    e i clienti sono anche in Spagna.
  \item L'MSSP e il cloud provider devono informare PayNest senza ingiustificato ritardo (art.~33(2) GDPR).
  \item Le scadenze cadono nel weekend e corrono comunque.
\end{itemize}
\end{soluzione}

\begin{soluzione}[C) Severità e decisione]
\begin{itemize}
  \item DPC = 3 (dati finanziari: IBAN e transazioni, che permettono frodi);
  \item EI = 1 (nome e IBAN identificano direttamente);
  \item CB: riservatezza 0,5 (esfiltrazione verso destinatari ignoti e malevoli) + disponibilità 0,25
    (temporanea, dati ripristinati dai backup) + intento malevolo 0,5 = 1,25.
\end{itemize}
\[
SE = 3 \times 1 + 1{,}25 = 4{,}25 \quad\Rightarrow\quad \text{molto alta}
\]
Quindi notifica al Garante (art.~33) \textbf{e} comunicazione ai 120.000 clienti (art.~34). Ai clienti
vanno indicati il rischio di phishing mirato e di frodi sull'IBAN, le misure adottate e le precauzioni
da prendere. I valori scelti, in particolare 0,25 per la disponibilità, vanno motivati.
\end{soluzione}

\begin{soluzione}[D) Misure verso i fornitori]
\begin{itemize}
  \item Gestione del rischio ICT da terze parti (DORA): \textbf{registro delle informazioni} sugli
    accordi con i fornitori ICT, valutazione del rischio prima e durante il rapporto, strategia
    multi-vendor ed \textbf{exit strategy}.
  \item Clausole contrattuali obbligatorie (DORA): requisiti di sicurezza, assistenza negli incidenti,
    diritti di audit e ispezione, limiti al subappalto, diritti di recesso, livelli di servizio.
  \item GDPR art.~28: DPA con entrambi i fornitori, clausola di notifica rapida dei breach, misure
    dell'art.~32. L'MFA obbligatoria sugli accessi remoti dei fornitori è la misura che avrebbe impedito
    l'incidente.
  \item Follow-up: verifica della chiusura delle carenze emerse e riesame da parte dell'organo di gestione.
\end{itemize}
\end{soluzione}

\begin{verifica}
\begin{itemize}
  \item Banca d'Italia come autorità DORA competente per gli istituti di pagamento italiani: verificare
    sulle slide DORA.
  \item Articoli DORA sulle terze parti (in genere artt.~28 e 30): gli appunti dell'anno precedente
    citano anche un ``art.~13''; verificare.
  \item L'app di PayNest potrebbe rientrare nel CRA come software installabile, ma l'incidente nasce da
    credenziali compromesse e non da una vulnerabilità del prodotto: non scatta l'art.~14 CRA.
    Lettura da confermare.
\end{itemize}
\end{verifica}

\dacompletare{Aggiungere una simulazione per ogni nuovo laboratorio (es.\ audit Sestante Mobility,
simulazione NIS2).}

\chapter{Quiz}
\label{cap:quiz}

\begin{nota}
Domande in inglese, come all'esame; spiegazioni in italiano. Le domande sono raggruppate per argomento
e ognuna rimanda alla lezione in cui il tema è stato trattato (link cliccabile). Le domande su ENISA,
NIS2 e DORA (strategia) ricostruiscono le domande chiuse dell'appello precedente
(Capitolo~\ref{cap:tema-precedente}); le opzioni sono nostre. Tutte le domande sono scritte dai
redattori, non dal docente. Soluzioni in fondo al capitolo.
\end{nota}

% =====================================================================
\section{Fondamenti}

\begin{domanda}[Fonti del diritto UE -- \rif{2026-09-21}{21/09}]\label{q:regolamento}
An EU \emph{regulation}, unlike a directive:
\begin{opzioni}
  \item must be transposed into national law before it applies;
  \item is directly applicable in all Member States;
  \item binds only EU institutions;
  \item is a form of soft law.
\end{opzioni}
\end{domanda}

\begin{domanda}[Hard e soft law -- \rif{2026-11-09}{09/11}]\label{q:softlaw}
Which of the following is \textbf{soft law}?
\begin{opzioni}
  \item the Cyber Resilience Act;
  \item the GDPR;
  \item ISO/IEC 27001;
  \item the NIS2 Directive as transposed by D.lgs.~138/2024.
\end{opzioni}
\end{domanda}

\begin{domanda}[Acronimi -- \rif{2026-09-21}{21/09}]\label{q:dpa}
In a contract between a controller and a cloud provider, the acronym ``DPA'' most likely refers to:
\begin{opzioni}
  \item the Data Protection Authority;
  \item the Data Processing Agreement under Art.~28 GDPR;
  \item the Data Protection Act of the provider's country;
  \item the Digital Product Assessment under the CRA.
\end{opzioni}
\end{domanda}

\begin{domanda}[Dati personali -- \rif{2026-09-21}{21/09}]\label{q:anonimi}
Which statement about anonymisation is \textbf{correct}?
\begin{opzioni}
  \item Removing name and surname is enough to make data anonymous.
  \item Pseudonymised data are no longer personal data.
  \item Data that are truly anonymous fall outside the scope of the GDPR.
  \item Anonymisation is a legal basis under Art.~6 GDPR.
\end{opzioni}
\end{domanda}

% =====================================================================
\section{GDPR}

\begin{domanda}[GDPR, basi giuridiche -- \rif{2026-09-25}{25/09}]\label{q:gdpr-basi}
Legal bases for processing under Art.~6 GDPR. \textbf{Select the WRONG statement.}
\begin{opzioni}
  \item Consent must be freely given, specific, informed and unambiguous.
  \item Performance of a contract covers only the processing strictly necessary for that contract.
  \item Public authorities may rely on legitimate interest for processing carried out in the
    performance of their tasks.
  \item Different legal bases may apply to different purposes within the same service.
\end{opzioni}
\end{domanda}

\begin{domanda}[GDPR, ruoli -- \rif{2026-09-25}{25/09}]\label{q:ruoli-fatto}
Two companies sign an agreement stating that each is an independent controller, but in practice they
jointly decide why and how the data are processed. Under the GDPR they are:
\begin{opzioni}
  \item independent controllers, because the contract says so;
  \item joint controllers, because the factual situation prevails over the contract;
  \item controller and processor, depending on who pays;
  \item outside the GDPR, since the roles are disputed.
\end{opzioni}
\end{domanda}

\begin{domanda}[GDPR, ruoli -- \rif{2026-11-09}{09/11}]\label{q:autorizzato}
An employee of the controller who processes personal data following the controller's instructions is:
\begin{opzioni}
  \item a processor (Art.~28);
  \item a sub-processor;
  \item a person authorised to process the data under the authority of the controller;
  \item a joint controller.
\end{opzioni}
\end{domanda}

\begin{domanda}[GDPR, legittimo interesse -- \rif{2026-09-28}{28/09}]\label{q:legittimo}
About legitimate interest (Art.~6(1)(f)). \textbf{Select the WRONG statement.}
\begin{opzioni}
  \item It requires a balancing test against the data subject's interests and fundamental rights.
  \item The legitimate interest invoked is the one of the data subject.
  \item It can be relied upon, for example, for some forms of video surveillance.
  \item The data subject has the right to object to processing based on it.
\end{opzioni}
\end{domanda}

\begin{domanda}[GDPR, principi -- \rif{2026-09-25}{25/09}]\label{q:principi}
Which of the following is \textbf{not} a principle listed in Art.~5 GDPR?
\begin{opzioni}
  \item purpose limitation;
  \item data minimisation;
  \item storage limitation;
  \item freedom of contract.
\end{opzioni}
\end{domanda}

\begin{domanda}[GDPR, ambito -- \rif{2026-09-25}{25/09}]\label{q:domestica}
A private individual publishes on a public website information about named colleagues of a
volunteer group. The household exemption:
\begin{opzioni}
  \item applies, because the individual acts for non-commercial purposes;
  \item does not apply, because the data are made accessible to an indefinite number of people;
  \item applies only if the data are not sensitive;
  \item applies only outside working hours.
\end{opzioni}
\end{domanda}

\begin{domanda}[GDPR, sanzioni -- \rif{2026-09-28}{28/09}]\label{q:sanzioni}
For infringements of the security obligations (including data breach notification duties), the GDPR
administrative fines can reach:
\begin{opzioni}
  \item EUR 20 million or 4\% of total worldwide annual turnover;
  \item EUR 10 million or 2\% of total worldwide annual turnover;
  \item EUR 35 million or 7\% of total worldwide annual turnover;
  \item EUR 15 million or 2.5\% of total worldwide annual turnover.
\end{opzioni}
\end{domanda}

\begin{domanda}[GDPR, categorie di dati -- \rif{2026-10-12}{12/10}]\label{q:malattia}
Sick-leave certificates of employees are:
\begin{opzioni}
  \item common personal data, because they do not state the diagnosis;
  \item data concerning health, a special category under Art.~9;
  \item judicial data under Art.~10;
  \item not personal data, because they are produced by a doctor.
\end{opzioni}
\end{domanda}

\begin{domanda}[GDPR, DPIA -- \rif{2026-10-05}{05/10}]\label{q:dpia}
After a DPIA, the residual risk of a new processing remains high despite the planned measures. The
controller must:
\begin{opzioni}
  \item start the processing and record the decision;
  \item consult the supervisory authority before starting the processing (Art.~36);
  \item notify a data breach;
  \item ask the data subjects for their consent, which removes the risk.
\end{opzioni}
\end{domanda}

% =====================================================================
\section{Data breach e severità ENISA}

\begin{domanda}[Data breach, definizione -- \rif{2026-10-05}{05/10}]\label{q:def-breach}
Which of the following is \textbf{not} a personal data breach?
\begin{opzioni}
  \item an unencrypted laptop with customer data is lost on a train;
  \item ransomware makes the HR database unavailable for a day;
  \item a phishing email reaches an employee, who deletes it without clicking;
  \item an employee browses customer records without any business need.
\end{opzioni}
\end{domanda}

\begin{domanda}[Data breach, riservatezza -- \rif{2026-10-19}{19/10}]\label{q:accesso}
An attacker accessed a file server containing personal data; logs show no evidence of exfiltration.
Under the GDPR:
\begin{opzioni}
  \item there is no breach until exfiltration is proven;
  \item there is a confidentiality breach, because unauthorised access is enough;
  \item there is only an availability breach;
  \item the event is relevant only under NIS2.
\end{opzioni}
\end{domanda}

\begin{domanda}[Data breach, 72 ore -- \rif{2026-10-05}{05/10}]\label{q:decorrenza}
The 72-hour deadline of Art.~33 GDPR runs from:
\begin{opzioni}
  \item the moment the breach occurred;
  \item the moment the controller became aware of the breach;
  \item the end of the forensic investigation;
  \item the first business day after the breach.
\end{opzioni}
\end{domanda}

\begin{domanda}[Data breach, responsabile -- \rif{2026-10-02}{02/10}]\label{q:responsabile}
A cloud provider acting as processor suffers a personal data breach. Who notifies the supervisory
authority?
\begin{opzioni}
  \item The processor, directly and instead of the controller.
  \item The controller, after the processor has informed it without undue delay.
  \item Both, independently and with the same deadline.
  \item ENISA, through the Single Reporting Platform.
\end{opzioni}
\end{domanda}

\begin{domanda}[Data breach, registro -- \rif{2026-10-05}{05/10}]\label{q:registro}
About the internal breach register (Art.~33(5)). \textbf{Select the WRONG statement.}
\begin{opzioni}
  \item It must document the facts, the effects and the remedial action taken.
  \item Breaches assessed as low risk and not notified need not be recorded.
  \item It allows the supervisory authority to verify compliance with Art.~33.
  \item A temporary loss of availability is recorded even if no notification is required.
\end{opzioni}
\end{domanda}

\begin{domanda}[Data breach, disponibilità -- \rif{2026-10-12}{12/10}]\label{q:disponibilita}
A ransomware attack makes personal data unavailable for six hours; the data are fully restored from
backup and there is no exfiltration. Under the GDPR this is:
\begin{opzioni}
  \item not a personal data breach, because the data were restored;
  \item a personal data breach (loss of availability), to be at least recorded in the breach register;
  \item a breach only if the ransom was paid;
  \item a NIS2 matter only.
\end{opzioni}
\end{domanda}

\begin{domanda}[Severità ENISA -- \rif{2026-10-02}{02/10}]\label{q:formula}
In the ENISA methodology the severity of a personal data breach is computed as:
\begin{opzioni}
  \item $SE = DPC + EI + CB$
  \item $SE = DPC \times EI + CB$
  \item $SE = (DPC + CB) \times EI$
  \item $SE = \max(DPC, EI, CB)$
\end{opzioni}
\end{domanda}

\begin{domanda}[Severità ENISA -- \rif{2026-11-02}{02/11}]\label{q:severita}
A breach exposes live location data (behavioural data, DPC~=~2) of fully identifiable users (EI~=~1),
with loss of confidentiality (0.5) and malicious intent (+0.5). Severity?
\begin{opzioni}
  \item 2, medium
  \item 2.5, medium
  \item 3, high
  \item 4, very high
\end{opzioni}
\end{domanda}

\begin{domanda}[Severità ENISA -- \rif{2026-10-12}{12/10}]\label{q:severita-salute}
Medical certificates (DPC~=~4) of employees are sent by mistake to the wrong internal mailing list;
identification is significant (EI~=~0.75); only confidentiality is affected (0.5), no malicious intent.
Severity?
\begin{opzioni}
  \item 2.5, medium
  \item 3, high
  \item 3.5, high
  \item 4.5, very high
\end{opzioni}
\end{domanda}

\begin{domanda}[GDPR, art.~34 -- \rif{2026-11-02}{02/11}]\label{q:minori}
A high-risk breach concerns personal data of children aged 5--12. The Art.~34 communication must be
addressed to:
\begin{opzioni}
  \item the children directly;
  \item the holders of parental responsibility (the parents);
  \item the children's schools;
  \item no one: communication is not required for minors, only notification to the authority.
\end{opzioni}
\end{domanda}

% =====================================================================
\section{ISO/IEC 27001 e risk management}

\begin{domanda}[ISO 27001 -- \rif{2026-10-19}{19/10}]\label{q:annexa}
ISO/IEC 27001:2022 Annex~A contains:
\begin{opzioni}
  \item 114 controls in 14 domains;
  \item 93 controls in 4 themes (organisational, people, physical, technological);
  \item 10 mandatory clauses and no controls;
  \item 34 technological controls only.
\end{opzioni}
\end{domanda}

\begin{domanda}[ISO 27001 -- \rif{2026-10-19}{19/10}]\label{q:certificazione}
About ISO/IEC 27001 certification. \textbf{Select the WRONG statement.}
\begin{opzioni}
  \item Clauses 4--10 are mandatory requirements of the ISMS.
  \item The certificate follows a three-year cycle with surveillance audits.
  \item The certificate guarantees that the organisation will not suffer security incidents.
  \item Certification concerns the organisation and its ISMS, not a single product.
\end{opzioni}
\end{domanda}

\begin{domanda}[Risk management -- \rif{2026-10-26}{26/10}]\label{q:max-cia}
In the course's risk register a risk is scored C~=~1, I~=~1, A~=~3 (scale 1--3). Its level is:
\begin{opzioni}
  \item 1, the minimum;
  \item about 1.7, the average;
  \item 3, the maximum of the three;
  \item 5, the sum.
\end{opzioni}
\end{domanda}

\begin{domanda}[Risk management -- \rif{2026-10-26}{26/10}]\label{q:trattamento}
Which of the following is \textbf{not} a risk treatment option?
\begin{opzioni}
  \item mitigate (modify) the risk with controls;
  \item transfer or share the risk;
  \item avoid the risk by stopping the activity;
  \item ignore the risk without documenting it.
\end{opzioni}
\end{domanda}

\begin{domanda}[Risk management -- \rif{2026-10-26}{26/10}]\label{q:soa}
The Statement of Applicability (SoA):
\begin{opzioni}
  \item lists the Annex~A controls, states whether each is applied and justifies inclusions and
    exclusions;
  \item is the certificate issued by the certification body;
  \item is the GDPR record of processing activities;
  \item is required only for organisations subject to NIS2.
\end{opzioni}
\end{domanda}

\begin{domanda}[ISO vs GDPR -- \rif{2026-10-23}{23/10}]\label{q:impatto}
Which statement about \emph{impact} is \textbf{correct}?
\begin{opzioni}
  \item ISO/IEC 27001 and the ENISA severity methodology measure impact on the same object.
  \item In ISO/IEC 27001 impact is assessed on the organisation; in the GDPR/ENISA approach on the
    rights and freedoms of data subjects.
  \item In ISO/IEC 27001 impact is assessed only on data subjects.
  \item The GDPR does not consider impact.
\end{opzioni}
\end{domanda}

\begin{domanda}[Rischio -- \rif{2026-10-05}{05/10}]\label{q:residuo}
Residual risk is:
\begin{opzioni}
  \item the risk before any control is applied;
  \item the risk remaining after the mitigation measures have been applied;
  \item the risk transferred to an insurer;
  \item the risk that the authority accepts on behalf of the controller.
\end{opzioni}
\end{domanda}

\begin{domanda}[ISO e conformità -- \rif{2026-11-02}{02/11}]\label{q:iso}
A company holds an ISO/IEC 27001 certificate. \textbf{Select the WRONG statement.}
\begin{opzioni}
  \item The certificate can support the demonstration of accountability.
  \item The certificate proves that the company is GDPR and CRA compliant.
  \item The certificate does not exempt the company from data breach notification duties.
  \item Auditors refer to ISO/IEC 27002 to assess how the controls are implemented.
\end{opzioni}
\end{domanda}

% =====================================================================
\section{Cyber Resilience Act}

\begin{domanda}[CRA, ambito -- \rif{2026-10-26}{26/10}]\label{q:cra-ambito}
Which of the following is \textbf{not} within the scope of the CRA?
\begin{opzioni}
  \item a GPS smartwatch for children connected to a cloud backend;
  \item an installable mobile app connected to a remote server;
  \item a medical device certified under the MDR;
  \item a home Wi-Fi router.
\end{opzioni}
\end{domanda}

\begin{domanda}[CRA, transizione -- \rif{2026-11-02}{02/11}]\label{q:cra-unita}
For CRA applicability, the cut-off date of 11 December 2027 is assessed with reference to:
\begin{opzioni}
  \item the date the product model was first launched;
  \item each individual unit placed on the market;
  \item the manufacturer's date of incorporation;
  \item the date of the last firmware update.
\end{opzioni}
\end{domanda}

\begin{domanda}[CRA, segnalazioni -- \rif{2026-11-02}{02/11}]\label{q:cra-final}
Under the CRA, the final report on an actively exploited vulnerability is due:
\begin{opzioni}
  \item within 72 hours of becoming aware of the vulnerability;
  \item within one month of becoming aware of the vulnerability;
  \item no later than 14 days after a corrective or mitigating measure is available;
  \item there is no deadline for the final report.
\end{opzioni}
\end{domanda}

\begin{domanda}[CRA e GDPR -- \rif{2026-11-02}{02/11}]\label{q:cra-gdpr}
The same incident is both a severe incident affecting a product (CRA) and a personal data breach
(GDPR). \textbf{Select the WRONG statement.}
\begin{opzioni}
  \item The CRA early warning is due within 24 hours of awareness.
  \item The GDPR notification is due within 72 hours of awareness.
  \item Sending the CRA early warning to the CSIRT also fulfils the GDPR notification duty.
  \item The two notifications run in parallel, towards different authorities.
\end{opzioni}
\end{domanda}

\begin{domanda}[CRA e AI Act -- \rif{2026-11-02}{02/11}]\label{q:cra-ai}
A product with digital elements that is also a high-risk AI system:
\begin{opzioni}
  \item is outside the CRA, since the AI Act prevails;
  \item is presumed compliant with the AI Act cybersecurity requirements (Art.~15) if it meets CRA
    Annex~I (Parts I and II) and this is shown in the CRA EU declaration of conformity;
  \item must obtain two separate CE markings, one per regulation;
  \item is automatically classified as a critical product under the CRA.
\end{opzioni}
\end{domanda}

\begin{domanda}[Metodo di audit -- \rif{2026-11-06}{06/11}]\label{q:fonti}
In the conformity audit method used in the course, a statement confirmed by two independent and
concordant sources is treated as:
\begin{opzioni}
  \item an assertion, to be taken with caution;
  \item an established fact (evidence);
  \item a finding of non-conformity;
  \item irrelevant, unless signed by the CEO.
\end{opzioni}
\end{domanda}

% =====================================================================
\section{Cybersecurity Act}

\begin{domanda}[ENISA -- \rif{2026-11-09}{09/11}]\label{q:enisa}
Which statement about ENISA is \textbf{correct}?
\begin{opzioni}
  \item ENISA imposes administrative fines on essential entities that fail to report incidents.
  \item ENISA received a permanent mandate with the Cybersecurity Act and supports Member States and EU
    institutions, e.g.\ through the CRA Single Reporting Platform.
  \item ENISA directly issues ISO/IEC 27001 certificates to companies.
  \item ENISA is the national CSIRT of Italy.
\end{opzioni}
\end{domanda}

\begin{domanda}[CSA, certificazione -- \rif{2026-11-09}{09/11}]\label{q:livelli}
The assurance levels of European cybersecurity certificates under the CSA are:
\begin{opzioni}
  \item low, medium, high, very high;
  \item basic, substantial, high;
  \item bronze, silver, gold;
  \item self-assessed, audited, certified.
\end{opzioni}
\end{domanda}

\begin{domanda}[CSA -- \rif{2026-11-09}{09/11}]\label{q:csa-cra}
About the relationship between CSA and CRA. \textbf{Select the WRONG statement.}
\begin{opzioni}
  \item The CSA creates a European framework of certification schemes.
  \item A CSA certificate can help demonstrate conformity with the CRA.
  \item The CSA replaces the CRA for certified products.
  \item EUCC is the first scheme adopted under the CSA and is based on the Common Criteria.
\end{opzioni}
\end{domanda}

% =====================================================================
\section{NIS2, DORA e AI Act}

\begin{nota}
Norme non ancora trattate a lezione (al 11/11/2026): domande basate sull'appello precedente e
sull'inquadramento dei Capitoli~\ref{cap:nis2}--\ref{cap:aiact}.
\end{nota}

\begin{domanda}[NIS2]\label{q:nis2-incidente}
Under NIS2, an \emph{incident} is:
\begin{opzioni}
  \item an event compromising the availability, authenticity, integrity or confidentiality of stored,
    transmitted or processed data or of the services offered by, or accessible via, network and
    information systems;
  \item a breach of security leading to the accidental or unlawful destruction, loss, alteration,
    unauthorised disclosure of, or access to, personal data;
  \item a vulnerability for which there is reliable evidence that a malicious actor has exploited it;
  \item only an event that has actually caused a financial loss above a fixed threshold.
\end{opzioni}
\end{domanda}

\begin{domanda}[DORA]\label{q:dora-strategie}
Digital operational resilience strategy under DORA. \textbf{Select the WRONG statement.}
\begin{opzioni}
  \item It includes the entity's ICT risk tolerance and resilience objectives.
  \item It may include a multi-vendor strategy for ICT third-party service providers.
  \item It includes a communication strategy for ICT-related incidents.
  \item It must be approved by the main ICT third-party service provider before adoption.
\end{opzioni}
\end{domanda}

\begin{domanda}[DORA]\label{q:dora-iniziale}
Under DORA, the initial notification of a major ICT-related incident must be submitted:
\begin{opzioni}
  \item within 72 hours of awareness;
  \item within 4 hours from the classification as major, and no later than 24 hours from awareness;
  \item within 24 hours from the classification as major, with no other limit;
  \item within one month, together with the final report.
\end{opzioni}
\end{domanda}

\begin{domanda}[NIS2 e DORA]\label{q:lex-specialis}
A bank is in a sector covered by both NIS2 and DORA. For ICT incident reporting:
\begin{opzioni}
  \item NIS2 prevails, being a directive;
  \item the bank must report the same incident twice, under both regimes;
  \item DORA applies as lex specialis;
  \item neither applies: the GDPR covers it.
\end{opzioni}
\end{domanda}

% =====================================================================
\section{Tempistiche trasversali}

\begin{domanda}[Tempistiche -- \rif{2026-11-02}{02/11}]\label{q:72h}
About ``72 hours'' in EU law. \textbf{Select the WRONG statement.}
\begin{opzioni}
  \item GDPR: notification to the supervisory authority within 72 hours of awareness of the breach.
  \item NIS2: incident notification within 72 hours of awareness of the significant incident.
  \item DORA: intermediate report within 72 hours of awareness of the incident.
  \item CRA: vulnerability notification within 72 hours of awareness.
\end{opzioni}
\end{domanda}

% =====================================================================
\section{Soluzioni}

\newcommand{\sol}[2]{\item[\ref{#1}.] \textbf{(#2)}.}
\begin{description}[leftmargin=2.5em,labelwidth=2em,font=\normalfont\bfseries]
  \sol{q:regolamento}{b} Il regolamento si applica direttamente; la direttiva vincola sul risultato e va
    recepita (es.\ NIS2 con il D.lgs.~138/2024).
  \sol{q:softlaw}{c} Gli standard ISO sono volontari: diventano rilevanti perché aiutano a dimostrare
    il rispetto della hard law.
  \sol{q:dpa}{b} Trappola segnalata dal docente: nei contratti ``DPA'' è il \emph{Data Processing
    Agreement}; l'autorità è la \emph{data protection authority}.
  \sol{q:anonimi}{c} Togliere nome e cognome non basta e i dati pseudonimizzati restano personali; solo
    i dati davvero anonimi escono dal GDPR.
  \sol{q:gdpr-basi}{c} L'art.~6(1) esclude il legittimo interesse per i trattamenti delle autorità
    pubbliche nell'esecuzione dei loro compiti: per quelli si usa l'interesse pubblico o l'obbligo legale.
  \sol{q:ruoli-fatto}{b} Conta la situazione di fatto, non ciò che dice il contratto: punto su cui il
    docente ha detto che sarà severo.
  \sol{q:autorizzato}{c} Il dipendente è una persona autorizzata che agisce sotto l'autorità del
    titolare; il responsabile è un soggetto esterno.
  \sol{q:legittimo}{b} Il legittimo interesse è del titolare (o di terzi); l'interessato ha diritti
    fondamentali.
  \sol{q:principi}{d} I principi dell'art.~5 sono liceità, correttezza e trasparenza; limitazione della
    finalità; minimizzazione; esattezza; limitazione della conservazione; integrità e riservatezza (più
    l'accountability).
  \sol{q:domestica}{b} È il principio del caso Lindqvist: la pubblicazione su Internet rende i dati
    accessibili a un numero indefinito di persone ed esclude l'esenzione domestica.
  \sol{q:sanzioni}{b} ``Domanda insidiosa'': per sicurezza e data breach il tetto è la metà di quello
    generale (art.~83(4) contro art.~83(5)).
  \sol{q:malattia}{b} Ogni dato relativo alla salute ricade nell'art.~9: nel modello ENISA è DPC~4,
    che si può abbassare solo motivando.
  \sol{q:dpia}{b} Se il rischio residuo resta elevato si consulta preventivamente l'autorità
    (art.~36). Conta il residuo, non l'inerente.
  \sol{q:def-breach}{c} Senza compromissione dei dati personali non c'è violazione: l'email di
    phishing cancellata è al più un evento di sicurezza.
  \sol{q:accesso}{b} La riservatezza è violata con il solo accesso non autorizzato; non presumere
    l'assenza di esfiltrazione se non la si può provare.
  \sol{q:decorrenza}{b} Dalla consapevolezza del titolare, che va però dimostrata: una data di
    consapevolezza tardiva va giustificata con evidenze.
  \sol{q:responsabile}{b} Solo il titolare notifica all'autorità; il responsabile informa il titolare
    senza ingiustificato ritardo (art.~33(2)).
  \sol{q:registro}{b} Il registro documenta \emph{ogni} violazione, anche se non notificata: è la
    prova dell'accountability.
  \sol{q:disponibilita}{b} Anche l'indisponibilità temporanea è un data breach. Se il rischio è basso
    può bastare la registrazione interna.
  \sol{q:formula}{b} È la ``regola d'oro'' da ricordare a memoria.
  \sol{q:severita}{c} $SE = 2 \times 1 + 0{,}5 + 0{,}5 = 3$: alta (è il Gruppo~B del caso Pulse K).
  \sol{q:severita-salute}{c} $SE = 4 \times 0{,}75 + 0{,}5 = 3{,}5$: alta ($3 \le SE < 4$).
  \sol{q:minori}{b} Giuridicamente gli interlocutori sono i titolari della responsabilità genitoriale.
  \sol{q:annexa}{b} 93 controlli in 4 temi (37 organizzativi, 8 persone, 14 fisici, 34 tecnologici).
    La (a) descrive la vecchia versione 2013.
  \sol{q:certificazione}{c} La certificazione attesta il sistema di gestione, non l'assenza di
    incidenti.
  \sol{q:max-cia}{c} Il livello è il massimo di C, I, A: fare la media è un errore tipico.
  \sol{q:trattamento}{d} Accettare un rischio è legittimo, ma va deciso e documentato; ignorarlo no.
  \sol{q:soa}{a} La SoA documenta le decisioni sui controlli ed è il punto in cui si rende conto delle
    esclusioni.
  \sol{q:impatto}{b} Non confondere i due piani: organizzazione (ISO) contro interessati (GDPR/ENISA).
  \sol{q:residuo}{b} Inerente meno mitigazioni: è il residuo che guida la decisione.
  \sol{q:iso}{b} La ISO è uno standard tecnico e non prova la conformità legale: è il suggerimento
    dato da Gianluca per il laboratorio di audit.
  \sol{q:cra-ambito}{c} I dispositivi medici sono soggetti all'MDR, regime verticale che esclude il CRA.
    Le app connesse e commercializzate rientrano: il software è un prodotto.
  \sol{q:cra-unita}{b} Le date di cut-off si valutano per singola unità: unità dello stesso modello
    possono avere obblighi diversi.
  \sol{q:cra-final}{c} Il termine di 14 giorni decorre dalla disponibilità del fix, non dalla
    consapevolezza. Per gli incidenti gravi, invece, il final report è dovuto entro un mese.
  \sol{q:cra-gdpr}{c} Una notifica a un'autorità non adempie gli obblighi verso un'altra: l'early
    warning CRA non è la notifica preliminare GDPR.
  \sol{q:cra-ai}{b} È la presunzione di conformità vista il 02/11, che collega l'art.~15 AI Act e
    l'Allegato~I del CRA.
  \sol{q:fonti}{b} Una fonte è un'asserzione; due fonti concordanti sono un fatto accertato.
  \sol{q:enisa}{b} Le sanzioni NIS2 sono irrogate dalle autorità nazionali; i certificati sono
    rilasciati da organismi di valutazione o autorità nazionali (e ISO non è uno schema UE); il CSIRT
    italiano è presso l'ACN.
  \sol{q:livelli}{b} Base, sostanziale, elevato (\emph{basic, substantial, high}).
  \sol{q:csa-cra}{c} Il CSA è soft law e aiuta a dimostrare la conformità al CRA (hard law), ma non lo
    sostituisce.
  \sol{q:nis2-incidente}{a} (b) è la definizione GDPR di data breach; (c) è la vulnerabilità
    attivamente sfruttata del CRA; (d) è falsa: rileva anche il danno potenziale e nella definizione non
    c'è alcuna soglia.
  \sol{q:dora-strategie}{d} La strategia è definita e approvata dall'organo di gestione dell'entità
    finanziaria; ai fornitori ICT gli obblighi arrivano per via contrattuale.
  \sol{q:dora-iniziale}{b} Doppio limite: 4 ore dalla classificazione come grave e comunque non oltre
    24 ore dalla conoscenza. Le 72 ore di DORA riguardano invece la relazione intermedia.
  \sol{q:lex-specialis}{c} DORA è lex specialis rispetto a NIS2 per le entità finanziarie.
  \sol{q:72h}{c} Le 72 ore di DORA decorrono dalla notifica iniziale, non dalla consapevolezza.
\end{description}


% =====================================================================
\appendix
% =====================================================================
% =====================================================================
% Appendice -- Glossario e acronimi
% =====================================================================
\chapter{Glossario e acronimi}
\label{app:glossario}

\begin{nota}[Fonte]
Rielaborato in italiano dal glossario distribuito con il primo laboratorio (\emph{Glossary -- acronyms,
laws \& regulations}), più alcune voci emerse a lezione. I riferimenti normativi completi sono nei
capitoli del manuale; qui una riga per voce.
\end{nota}

\section{Acronimi}
\begin{longtable}{>{\bfseries}p{0.17\textwidth}p{0.76\textwidth}}
\toprule
\normalfont\textbf{Sigla} & \textbf{Significato} \\
\midrule
\endhead
ACN & Agenzia per la Cybersicurezza Nazionale: autorità italiana competente per la NIS2 e autorità
  nazionale di certificazione; ospita lo CSIRT Italia. \\
BCR & \emph{Binding Corporate Rules}: regole vincolanti d'impresa per i trasferimenti infragruppo
  (art.~47 GDPR). \\
CB & \emph{Circumstances of the Breach}: fattore additivo della metodologia ENISA (perdita di
  riservatezza, integrità, disponibilità, intento malevolo). \\
CIA & \emph{Confidentiality, Integrity, Availability}: le tre proprietà di sicurezza; i tipi di
  violazione le rispecchiano. \\
CJEU & Corte di giustizia dell'UE (Lussemburgo): interpretazione vincolante del diritto UE. Da non
  confondere con la CEDU. \\
CRA & \emph{Cyber Resilience Act}, Reg.~(UE) 2024/2847. \\
CSA & \emph{Cybersecurity Act}, Reg.~(UE) 2019/881. \\
CSIRT & \emph{Computer Security Incident Response Team}: riceve le segnalazioni NIS2 e CRA (in Italia
  CSIRT Italia, presso l'ACN). \\
DORA & \emph{Digital Operational Resilience Act}, Reg.~(UE) 2022/2554 (settore finanziario). \\
DPA & \emph{Data Protection Authority} (autorità di controllo). Attenzione: anche \emph{Data
  Processing Agreement}, il contratto dell'art.~28 GDPR. \\
DPC & \emph{Data Processing Context}: punteggio base ENISA (da 1 a 4) che dipende dal tipo di dati. \\
DPIA & Valutazione d'impatto sulla protezione dei dati (art.~35 GDPR), obbligatoria per i trattamenti
  probabilmente ad alto rischio. \\
DPO & \emph{Data Protection Officer}, responsabile della protezione dei dati (artt.~37--39 GDPR). \\
ECtHR / CEDU & Corte europea dei diritti dell'uomo (Strasburgo, Consiglio d'Europa): non è
  un'istituzione UE. \\
EDPB & Comitato europeo per la protezione dei dati: riunisce le autorità di controllo e adotta linee
  guida (es. 9/2022 sulla notifica dei data breach). \\
EDPS & Garante europeo della protezione dei dati, competente per le istituzioni UE. \\
EI & \emph{Ease of Identification}: moltiplicatore ENISA (0,25; 0,5; 0,75; 1). \\
ENISA & Agenzia dell'UE per la cybersicurezza; mandato permanente con il CSA. \\
ESAs & Autorità europee di vigilanza (EBA, ESMA, EIOPA): sorveglianza ai sensi di DORA. \\
EUCC & Primo schema europeo di certificazione adottato con il CSA, basato sui Common Criteria. \\
FRIA & Valutazione d'impatto sui diritti fondamentali (AI Act), per alcuni \emph{deployer} di sistemi
  ad alto rischio. \\
GPAI & Modelli di IA per finalità generali (AI Act). \\
ISMS & Sistema di gestione della sicurezza delle informazioni (ISO/IEC~27001). \\
LED & Direttiva (UE) 2016/680 sul trattamento dei dati in ambito penale e di polizia: fuori dal GDPR. \\
PDCA & \emph{Plan--Do--Check--Act}: ciclo di miglioramento continuo dei sistemi di gestione ISO. \\
PDE & \emph{Product with Digital Elements}: prodotto con elementi digitali, ambito del CRA. \\
RTO / RPO & \emph{Recovery Time / Recovery Point Objective}: obiettivi di ripristino e di perdita
  massima di dati. \\
SBOM & \emph{Software Bill of Materials}: distinta dei componenti software (CRA). \\
SCCs & Clausole contrattuali tipo approvate dalla Commissione per i trasferimenti (art.~46 GDPR). \\
SE & \emph{Severity}: risultato ENISA, $SE = DPC \times EI + CB$. \\
SoA & \emph{Statement of Applicability}: elenco dei controlli dell'Annex~A ISO/IEC~27001 applicati
  o esclusi, con giustificazione. \\
TIA & \emph{Transfer Impact Assessment}: valutazione richiesta con SCC e BCR dopo Schrems~II. \\
TOMs & Misure tecniche e organizzative (artt.~24, 25, 32 GDPR). \\
\bottomrule
\end{longtable}

\section{Norme e strumenti}
\begin{longtable}{>{\bfseries}p{0.27\textwidth}p{0.66\textwidth}}
\toprule
\normalfont\textbf{Strumento} & \textbf{In una riga} \\
\midrule
\endhead
Carta dei diritti fondamentali UE, artt.~7--8 & Rispetto della vita privata e protezione dei dati
  personali come diritto fondamentale autonomo. \\
Direttiva 95/46/CE & Direttiva madre sulla protezione dei dati, abrogata dal GDPR. \\
Direttiva ePrivacy 2002/58/CE & Riservatezza nelle comunicazioni elettroniche; primo obbligo di
  notifica dei breach (operatori telecom); cookie. \\
GDPR, Reg.~(UE) 2016/679 & Protezione dei dati personali; applicabile dal 25~maggio~2018. \\
NIS2, Dir.~(UE) 2022/2555 & Soggetti essenziali e importanti, misure di gestione del rischio,
  notifica degli incidenti significativi. \\
DORA, Reg.~(UE) 2022/2554 & Resilienza operativa digitale del settore finanziario; \emph{lex
  specialis} rispetto alla NIS2. \\
CSA, Reg.~(UE) 2019/881 & Mandato permanente di ENISA; quadro europeo di certificazione. \\
CRA, Reg.~(UE) 2024/2847 & Requisiti di cybersicurezza dei prodotti con elementi digitali;
  marcatura CE. \\
AI Act, Reg.~(UE) 2024/1689 & Regolamento sull'IA basato sul rischio (pratiche vietate, alto
  rischio, trasparenza, GPAI). \\
Convenzione 108 & Primo trattato internazionale vincolante sulla protezione dei dati (Consiglio
  d'Europa, 1981). \\
Convenzione di Budapest & Convenzione sul cybercrime (2001): reati informatici armonizzati e
  cooperazione. \\
D.lgs.~196/2003 & Codice privacy italiano, adeguato al GDPR dal D.lgs.~101/2018. \\
D.lgs.~138/2024 & Recepimento italiano della NIS2; autorità competente l'ACN. \\
ISO/IEC~27001:2022 & Requisiti dell'ISMS; certificabile; Annex~A con 93 controlli. \\
ISO/IEC~27002 & Linee guida di attuazione dei controlli dell'Annex~A. \\
ISO/IEC~27701 & Estensione privacy dell'ISMS (PIMS). \\
\bottomrule
\end{longtable}


\end{document}
